\ifx\danslelivre\undefined
\RequirePackage{luatex85}
\documentclass[10pt]{../fichiers-configuration-latex/smfart}
\usepackage{polyglossia}
%\usepackage[francais]{babel}
\usepackage{fontspec}
\usepackage[math-style=ISO,bold-style=upright,sans-style=upright]{unicode-math}
%\usepackage[mathletters]{ucs}
%\usepackage[utf8]{inputenc}
%\usepackage[T1]{fontenc}
\usepackage[unicode,colorlinks=true,linkcolor=blue,citecolor=Sepia,urlcolor=magenta]{hyperref}
\setmainlanguage{french}
\setmainfont[Ligatures=TeX]{Linux Libertine O}
% variante possible : TeXGyrePagella ?
\setmathfont{XITS Math}

% maths
\usepackage{amsmath}
%\usepackage{amssymb}
%\usepackage{mathrsfs}
%\usepackage{stmaryrd}
%\usepackage{euscript}

% couleurs
\usepackage[usenames,dvipsnames]{xcolor}

% graphiques
\usepackage{graphicx}
\usepackage[all]{xy}
% Xy-pic 3.8.9's PDF driver assumes that \fontdimen8\textfont3 is a
% legacy TeX rule width.  Under LuaTeX with unicode-math/XITS Math it is
% instead 6.62pt, which turns every arrow shaft into a thick black bar.
% Override only xypdf's private cached width; do not alter the math font.
\makeatletter
\AtBeginDocument{%
  \@ifundefined{xP@preclw}{}{%
    \edef\xP@preclw{0.4pt}%
    \edef\xP@lw{\xP@dim\xP@preclw}%
    \PackageInfo{sga4-xypdf-compat}{Xy-pic PDF line width set to \xP@preclw}%
  }%
}
\makeatother
%\usepackage{tikz}
%\usetikzlibrary{matrix,arrows,calc}
%\tikzset{column sep=3em,row sep=2em,text height=1.5ex, text depth=0.25ex,>=angle 90}
%\tikzset{description/.style={fill=white,inner sep=1pt}}

% hyperliens, 'etiquettes
\usepackage{xr-hyper}
\usepackage{../fichiers-configuration-latex/smfthm}
%\ifx\cacheretiquettes\undefined
%\usepackage[inline]{showlabels} \else \fi % pour rendre visibles les 'etiquettes
%\usepackage{../fichiers-configuration-latex/makeidx} % pour faire un index
%\usepackage[pagebackref,unicode]{hyperref}  % pour voir o`u sont cit'ees les r'ef'erences
\hypersetup{unicode,colorlinks=true,linkcolor=blue,citecolor=Sepia,urlcolor=magenta,
  hypertexnames=false}
%\hypersetup{breaklinks=true} %permet le retour `a la ligne dans les liens trop longs. Incompatible avec pdflatex
%\renewcommand*{\backref}[1]{{\small [cit\'e page(s)~#1.]}}

\usepackage{xspace}
\usepackage{pifont} % pour expose XVII

\renewcommand{\thefootnote}{\roman{footnote}} % pour mettre les notes en bas de page sous la forme (ii) et non (2), qui ressemble trop (?) a un exposant.

\newcounter{pageoriginal}
\newcommand{\pageoriginale}{\refstepcounter{pageoriginal}\marginpar{\footnotesize\textbf{\thepageoriginal}}\xspace}
\let\pageoriginaled\pageoriginale

\usepackage{ifthen}
\newboolean{original}
\def\sisi{\ifthenelse{\boolean{original}}}
\setboolean{original}{false}

\newcounter{cnote}
\def\nde#1{\ifthenelse{\boolean{original}}{}{\stepcounter{cnote}\def\thefootnote{\arabic{cnote}}\protect\footnote[\value{cnote}]{N.D.E.: #1}}}
\def\ndemark{\ifthenelse{\boolean{original}}{}{\stepcounter{cnote}\def\thefootnote{\arabic{cnote}}\protect\footnotemark[\value{cnote}]}}
\def\ndetext#1{\ifthenelse{\boolean{original}}{}{\def\thefootnote{\arabic{cnote}}\protect\footnotetext[\value{cnote}]{N.D.E.: #1}}}

\setmainlanguage{french}
\setotherlanguage{english}

% unicode-math deliberately disables the legacy \Arrowvert glyph.  The
% source uses it once as a rotated equality marker; \Vert has the intended
% appearance and keeps the historical diagram buildable with LuaLaTeX.
\renewcommand{\Arrowvert}{\Vert}


%%
%% Pour afficher le numéro de version Git (exige que LuaTeX soit appelé avec --shell-escape)
%%

\catcode`\%=12
\directlua{
function command_output(cmd)
  local output_pipe = io.popen(cmd)
  local line_iter = output_pipe and output_pipe:lines()
  local line = line_iter and line_iter()
  if output_pipe then output_pipe:close() end
  return line
end
local git_command1 = [[git log --pretty=format:"%h %ai" -1]]
local git_command2 = [[if git diff --quiet HEAD ; then echo clean ; else echo dirty ; fi]]
function git_status()
  local status = command_output(git_command1)
  if status then
    tex.tprint({luatexbase.catcodetables.CatcodeTableLaTeX, string.char(0x5c).."texttt{"}, {-2,status}, {luatexbase.catcodetables.CatcodeTableLaTeX, "}"})
    if not (command_output(git_command2) == "clean") then tex.sprint(-2, " (modifications non validées)") end
  else
    tex.sprint(-2, "?")
  end
end
}
\catcode`\%=14
\def\showgitstatus{\directlua{git_status()}}


% macros pour LaTeX
\renewcommand{\mathsf}{}

% Police différente pour U+2113 SCRIPT SMALL L
\setmathfont[range="2113]{latinmodern-math.otf}

% Give every explicit SGA label its own destination and preserve the printed
% number supplied as the optional argument.  The former definition inherited
% whichever counter anchor happened to precede it and ignored #1 entirely.
\makeatletter
\def\etiquette[#1]#2{%
  \phantomsection
  \protected@edef\@currentlabel{#1}%
  \label{#2}%
}
\makeatother
\def\eqhyperref[#1]#2{(\hyperref[#1]{#2})}
\def\pagehyperref[#1]#2{(page \hyperref[#1]{#2})}
\def\hyperrefext[#1][#2]#3{\hyperref[#2]{#1 #3}}
\def\eqhyperrefext[#1][#2]#3{\hyperref[#2]{#1 #3}}     %cf. exposes 17
\def\neweqhyperrefext[#1][#2]#3{\hyperref[#2]{#1 #3}}  %cf. exposes 17

\def\shat#1{\protect\widehat{#1}}
\def\stilde#1{\protect\widetilde{#1}}
%%%%%%%%%%%%%%%%%%


\def\etcn{etc.\xspace}
\def\etco{etc...\xspace}
\def\etc{\sisi{\etco}{\etcn}}
\def\ldotsn{\sisi{...}{\etcn}}

\def\cf{{cf.}\kern.3em}
\def\Cf{{Cf.}\kern.3em}
\def\ie{{i.e.}\kern.3em}
\def\iev{{i.e.},\xspace}
\def\loccito{loc.\kern2pt cit.\xspace}
\def\loccitn{\textit{loc.\kern2pt cit.}\xspace}
\def\loccit{\sisi{\loccito}{\loccitn}}
\def\EGA{\textup{EGA}\kern.2em}
\def\SGA{\textup{SGA}\kern.3em}

\def\quoi{}

\def\quav{\quad,\quad}
\def\quaet{\quad \text{et}\quad}
\def\fibreproduct#1#2#3{#1 ×_{#3} #2}
\let\fprod\fibreproduct
\def\ccheck#1#2{\mbox{${\check{\check{\rm #1}}}^{\lower2pt\hbox{$\scriptstyle #2$}}$}}
\def\sheaf#1{\protect\mathcal{#1}}
\def\vieuxsheaf#1{{#1}} % cf. problème avec sheaf{a\Phi_i} par exemple...
\DeclareMathOperator{\Hom}{Hom}
\DeclareMathOperator{\can}{can}
\DeclareMathOperator{\Ind}{Ind}
\DeclareMathOperator{\ind}{ind}
%\DeclareMathOperator{\Hhomtop}{\mathcal{H}\!\mathit{omtop}}
\DeclareMathOperator{\Hhomtop}{\mathbf{Homtop}}
\DeclareMathOperator{\Point}{Point}
\DeclareMathOperator{\Pro}{Pro}
\DeclareMathOperator{\Top}{Top}
\DeclareMathOperator{\Mod}{Mod}
\DeclareMathOperator{\Cart}{Cart}
%\DeclareMathOperator{\Mmorsite}{\mathcal{M}\!\mathit{orsite}}
\DeclareMathOperator{\Mmorsite}{\mathbf{Morsite}}
\DeclareMathOperator{\loc}{loc}
\DeclareMathOperator{\Loc}{Loc}
\DeclareMathOperator{\Pic}{Pic}
\DeclareMathOperator{\Et}{Et}
\DeclareMathOperator{\card}{card}
\DeclareMathOperator{\Ens}{Ens}
\DeclareMathOperator{\h}{\underline{H}}
\DeclareMathOperator{\Fonct}{Fonct}
\DeclareMathOperator{\Ker}{Ker}
\DeclareMathOperator{\but}{but}
\DeclareMathOperator{\Ob}{Ob}
\DeclareMathOperator{\ob}{ob}
\DeclareMathOperator{\Lib}{Lib}
\DeclareMathOperator{\esj}{esj}
\DeclareMathOperator{\id}{id}
\DeclareMathOperator{\Id}{Id}
\DeclareMathOperator{\ad}{ad}
\DeclareMathOperator{\Iim}{Im}
\DeclareMathOperator{\pr}{pr}
\DeclareMathOperator{\Cat}{Cat}
\DeclareMathOperator{\car}{car}
\DeclareMathOperator{\SysInd}{SysInd}
\DeclareMathOperator{\Sup}{Sup}
\DeclareMathOperator{\Coker}{Coker}
\DeclareMathOperator{\Filt}{Filt}
\DeclareMathOperator{\acc}{acc}
\DeclareMathOperator{\sur}{sur}
\DeclareMathOperator{\cat}{cat}
\DeclareMathOperator{\Ord}{Ord}
\DeclareMathOperator{\Gr}{Gr}
\DeclareMathOperator{\Cov}{Cov}
\DeclareMathOperator{\Caf}{Caf}
\DeclareMathOperator{\coker}{coker}
\DeclareMathOperator{\coim}{coim}
\DeclareMathOperator{\im}{im}
\DeclareMathOperator{\Ouv}{Ouv}
\DeclareMathOperator{\diag}{diag}
\DeclareMathOperator{\Esp}{Esp}
\DeclareMathOperator{\TOP}{TOP}
\DeclareMathOperator{\Unif}{Unif}
\DeclareMathOperator{\Iind}{\mathbf{Ind}}
\DeclareMathOperator{\Hhomcart}{\mathbf{Homcart}}
\DeclareMathOperator{\Hhomtopess}{\mathbf{Homtopess}}
\DeclareMathOperator{\Ppoints}{\mathbf{Points}}
\DeclareMathOperator{\Ppoint}{\mathbf{Point}}
\DeclareMathOperator{\Ttors}{\mathbf{Tors}}
\DeclareMathOperator{\Ssex}{\mathbf{Sex}}
\DeclareMathOperator{\Ffib}{\mathbf{Fib}}
\DeclareMathOperator{\Vvois}{\mathbf{Vois}}
\DeclareMathOperator{\Ppointess}{\mathbf{Pointess}}
\DeclareMathOperator{\scrHom}{\mathbf{Hom}}
\DeclareMathOperator{\sheafPic}{\mathbf{Pic}}
\DeclareMathOperator{\sheafAut}{\mathbf{Aut}}
\DeclareMathOperator{\Hhomtoplocan}{\mathbf{Homtoplocan}}
\DeclareMathOperator{\Hhomtopan}{\mathbf{Homtopan}}
\DeclareMathOperator{\Pptgeom}{\mathbf{Ptgéom}}
\DeclareMathOperator{\Morsite}{Morsite}
\DeclareMathOperator{\Fib}{Fib}
\DeclareMathOperator{\Fais}{Fais}
\DeclareMathOperator{\kar}{kar}
\DeclareMathOperator{\Groupes}{Groupes}
\DeclareMathOperator{\Restr}{Restr}
\DeclareMathOperator{\Isom}{Isom}
\DeclareMathOperator{\Prol}{Prol}
\DeclareMathOperator{\Res}{Res}
\DeclareMathOperator{\Tors}{Tors}
\DeclareMathOperator{\esp}{esp}
\DeclareMathOperator{\es}{es}
\DeclareMathOperator{\od}{od}
\DeclareMathOperator{\gr}{gr}
\DeclareMathOperator{\sob}{sob}
\DeclareMathOperator{\dfn}{dfn}
\DeclareMathOperator{\Sch}{Sch}
\DeclareMathOperator{\ouv}{ouv}
\DeclareMathOperator{\cons}{cons}
\DeclareMathOperator{\Nerf}{Nerf}
\DeclareMathOperator{\Points}{Points}
\DeclareMathOperator{\Espsob}{Espsob}
\DeclareMathOperator{\Homtop}{Homtop}
\DeclareMathOperator{\Pt}{Pt}
\DeclareMathOperator{\Quot}{Quot}
\DeclareMathOperator{\glob}{glob}
\DeclareMathOperator{\Homfib}{Homfib}
\DeclareMathOperator{\Qouv}{Qouv}
\DeclareMathOperator{\Quotouv}{Quotouv}
\DeclareMathOperator{\ens}{ens}
\DeclareMathOperator{\ab}{ab}
\DeclareMathOperator{\an}{an}
\DeclareMathOperator{\fs}{fs}
\DeclareMathOperator{\red}{red}
\DeclareMathOperator{\Ab}{Ab}
\DeclareMathOperator{\cosq}{cosq}
\DeclareMathOperator{\Ann}{Ann}
\DeclareMathOperator{\Zar}{Zar}


%%%%%%%%%%%%
%%%%% SGA4-2
\DeclareMathOperator{\Sex}{Sex}
\DeclareMathOperator{\vois}{vois}
\DeclareMathOperator{\Vois}{Vois}
\DeclareMathOperator{\Ext}{Ext}
\DeclareMathOperator{\ext}{ext}
\DeclareMathOperator{\Fer}{Fer}
\DeclareMathOperator{\Add}{Add}
\DeclareMathOperator{\Cocart}{Cocart}
\DeclareMathOperator{\cosk}{cosk}
\DeclareMathOperator{\plat}{plat}

\DeclareMathOperator{\faisc}{faisc}
\DeclareMathOperator{\sk}{sk}
\DeclareMathOperator{\ch}{ch}

\def\HH{\mathbf{H}}
\def\qc{\mathrm{qc}}
\def\et{\mathrm{et}}
\def\coh{\mathrm{coh}}

\def\Hvr#1{%
\begin{array}[b]{l}
{\kern2pt\scriptstyle r}\\[-8pt]
{\kern1pt\scriptstyle \vee}\\[-5pt]
H^{#1}
\end{array}}

\def\Hvv#1{%
\begin{array}[b]{l}
\kern2pt\rotatebox{90}{${\scriptscriptstyle\ll}$}\\[-6pt]
H^{#1}
\end{array}}

\def\Hvi#1{%
\begin{array}[b]{l}
{\kern0.2pt\scriptscriptstyle \infty}\\[-8pt]
{\kern1pt\scriptstyle \vee}\\[-5pt]
H^{#1}
\end{array}}

\def\Hv#1{%
\begin{array}[b]{l}
{\kern1pt\scriptstyle \vee}\\[-5pt]
H^{#1}
\end{array}}

\DeclareMathOperator{\cart}{cart}
\DeclareMathOperator{\Ffaisc}{\mathbf{Faisc}}


\def\Ev#1{%
\begin{array}[b]{l}
{\kern1pt\scriptstyle \vee}\\[-5pt]
E'_{#1}
\end{array}}

\DeclareMathOperator{\Cont}{Cont}
\DeclareMathOperator{\Ccont}{\mathscr{C}\!\mathit{ont}}
\DeclareMathOperator{\Site}{Site}
\newcommand\MM{\mathbf{M}}
\def\limtop#1{{\displaystyle{\varprojlim_{#1}\!\mbox{\rm top }}}}
\def\Limtop{\mathop{\varprojlim\!\mbox{\rm top}}}
\DeclareMathOperator{\spec}{spec}
\def\ss{\scriptstyle{\mathscr{S}}}
\DeclareMathOperator{\Gal}{Gal}
\def\smallm{\scriptstyle{\mathscr{M}}}
\DeclareMathOperator{\End}{End}
\DeclareMathOperator{\ass}{ass}
\DeclareMathOperator{\zar}{zar}
\DeclareMathOperator{\Ogg}{Ogg}
\DeclareMathOperator{\res}{res}
\DeclareMathOperator{\tr}{tr}
\DeclareMathOperator{\Supp}{Supp}
\DeclareMathOperator{\codim}{codim}
\DeclareMathOperator{\rk}{rk}
\DeclareMathOperator{\cur}{cur}
\DeclareMathOperator{\cdqc}{cdqc}
\newcommand\DD{\mathbf{D}}
\DeclareMathOperator{\cl}{cl}
\def\JJ{\mathbf{J}}
\def\FF{\mathbf{F}}
\def\BB{\mathbf{B}}
\DeclareMathOperator{\Alget}{Alget}
\DeclareMathOperator{\Ass}{Ass}
\DeclareMathOperator{\supp}{supp}
\DeclareMathOperator{\rad}{rad}
\DeclareMathOperator{\prof}{prof}
\DeclareMathOperator{\cint}{int}
\DeclareMathOperator{\sic}{sic}
\DeclareMathOperator{\Tor}{Tor}
\DeclareMathOperator{\tors}{tors}
\DeclareMathOperator{\Sym}{Sym}
\DeclareMathOperator{\parf}{parf}
\DeclareMathOperator{\tor}{tor}
\DeclareMathOperator{\defn}{defn}
\DeclareMathOperator{\Div}{Div}
\DeclareMathOperator{\Pol}{Pol}
\DeclareMathOperator{\Fex}{Fex}
\DeclareMathOperator{\Lie}{Lie}
\DeclareMathOperator{\tot}{tot}
\DeclareMathOperator{\Aut}{Aut}
\DeclareMathOperator{\HOM}{HOM}
\DeclareMathOperator{\pch}{pch}
\DeclareMathOperator{\Ch}{Ch}
\DeclareMathOperator{\EXT}{EXT}
\DeclareMathOperator{\PIC}{PIC}
\DeclareMathOperator{\Rf}{Rf}
\DeclareMathOperator{\inj}{inj}
\newcommand\VV{\mathbf{V}}

\def\Xv{%
\begin{array}[b]{l}
{\kern1pt\scriptstyle \vee}\\[-5pt]
X
\end{array}}

%%%%%%%
\newcommand\TT{\mathbf{T}}
\newcommand\XX{\mathbf{X}}
\newcommand\EE{\mathbf{E}}
%%%%%%

\def\qedsymbol{\textup{C.Q.F.D.}}
\let\CQFD\qed
\let\cqfd\qed

%%%%%\bigdot
\font\bigsymb=cmsy10 at 4pt
\def\bigdot{{\kern1.2pt\raise 1.5pt\hbox{\bigsymb\char15}}}
%%%%%

\def\Lright{\mathop{\vcenter{\hbox{\hbox{$\longrightarrow$}\lower8pt \hbox{\kern-16pt\hbox{$\longrightarrow$}}}}}}

\def\Aa{{\mathcal A}}
\def\Bb{{\mathcal B}}
\def\Cc{{\mathcal C}}
\def\Dd{{\mathcal D}}
\def\Ee{{\mathcal E}}
\def\Ff{{\mathcal F}}
\def\Gg{{\mathcal G}}
\def\Hh{{\mathcal H}}
\def\Ii{{\mathcal I}}
\def\Jj{{\mathcal J}}
\def\Kk{{\mathcal K}}
\def\Ll{{\mathcal L}}
\def\Mm{{\mathcal M}}
\def\Nn{{\mathcal N}}
\def\Oo{{\mathcal O}}
\def\Pp{{\mathcal P}}
\def\Qq{{\mathcal Q}}
\def\Rr{{\mathcal R}}
\def\Ss{{\mathcal S}}
\def\Tt{{\mathcal T}}
\def\Uu{{\mathcal U}}
\def\Vv{{\mathcal V}}
\def\Ww{{\mathcal W}}
\def\Xx{{\mathcal X}}
\def\Yy{{\mathcal Y}}
\def\Zz{{\mathcal Z}}

\newcommand{\CC}{\mathbf{C}}
\newcommand{\GG}{\mathbf{G}}
\newcommand{\KK}{\mathbf{K}}
\newcommand{\LL}{\mathbf{L}}
\newcommand{\NN}{\mathbf{N}}
\newcommand{\PP}{\mathbf{P}}
\newcommand\QQ{\mathbf{Q}}
\newcommand\RR{\mathbf{R}}
\newcommand\ZZ{\mathbf{Z}}

\def\to{\mathchoice{\longrightarrow}{\rightarrow}{\rightarrow}{\rightarrow}}
\def\from{\mathchoice{\longleftarrow}{\leftarrow}{\leftarrow}{\leftarrow}}
\def\mto{\mathchoice{\longmapsto}{\mapsto}{\mapsto}{\mapsto}}
\def\mfrom{\mathrel{\from\joinrel\mapstochar\,}}
\def\hto{\mathrel{\lhook\joinrel\to}}
\def\hfrom{\mathrel{\from\joinrel\rhook}}

\def\implies{\mathchoice{\Longrightarrow}{\Rightarrow}{\Rightarrow}{\Rightarrow}}
\def\implied{\mathchoice{\Longleftarrow}{\Leftarrow}{\Leftarrow}{\Leftarrow}}
\def\iff{\mathchoice{\Longleftrightarrow}{\Leftrightarrow}{\Leftrightarrow}{\Leftrightarrow}}

\def\To#1{\mathchoice{\xrightarrow{\labelstyle~~#1~}}{\stackrel{#1}{\longrightarrow}}{}{}}
\def\Mto#1{\mathrel{\mapstochar\To{#1}}}
\def\From#1{\mathchoice{\xleftarrow{\labelstyle~~#1~}}{\stackrel{#1}{\longleftarrow}}{}{}}

\newcommand{\isomto}{\overset{\sim}{\longrightarrow}}
\newcommand{\isomfrom}{\overset{\sim}{\longleftarrow}}

\newcommand{\IFF}{\unskip~$\iff$~\ignorespaces}
\newcommand{\THEN}{\unskip~$\implies$~\ignorespaces}

\let\leq\leqslant
\let\geq\geqslant
\let\le\leqslant
\let\ge\geqslant

\def\cf{{cf.}\kern.3em}
\def\Cf{{Cf.}\kern.3em}
\def\ie{{i.e.}\kern.3em}
\def\iev{{i.e.},\xspace}
%\def\loccito{loc.\kern2pt cit.\xspace}
%\def\loccitn{\textit{loc.\kern2pt cit.}\xspace}
%\def\loccit{\sisi{\loccito}{\loccitn}}
\def\EGA{\textup{EGA}\kern.2em}
\def\SGA{\textup{SGA}\kern.3em}

\let\oldS\S
\def\S{\oldS\kern .15em}
\def\ptbl{.\kern .2em}
\def\resp{resp.\kern.3em}
\def\Exp{Exp.\kern.3em}

\let\bar\overline
\let\hat\widehat
\let\tilde\widetilde
\let\dpl\displaystyle
\let\ldots\dots
\let\emptyset\varnothing
\let\skpt\skippointrait
\let\epsilon\varepsilon

\def\sfrac#1#2{{#1}/{#2}}

\DeclareMathOperator{\cd}{cd}
%\DeclareMathOperator{\Hhom}{\mathcal{H}\!\mathit{om}}
\DeclareMathOperator{\Hhom}{\mathbf{Hom}}
%\DeclareMathOperator{\Eext}{\mathcal{E}\!\mathit{xt}}
\DeclareMathOperator{\Eext}{\mathbf{Ext}}
\DeclareMathOperator{\Spec}{Spec}
\DeclareMathOperator{\Tr}{Tr}

\let\emptyset=\varnothing

%%%%%%%%%%%%%%%
%%% expose XVII
%%%%%%%%%%%%%%%

\makeatletter
\let\old@optexpose\optexpose
\newcommand{\empty@optexpose}[1]{}
\newcommand{\paren@optexpose}[1]{\old@optexpose{#1}(}
\newcommand{\refnoexpose}[1]{\let\optexpose\empty@optexpose\ref{#1}\let\optexpose\old@optexpose}%\refnoexpose{X.cor5.2} engendre "5.2" au lieu de "X 5.2"
\newcommand{\neweqref}[1]{\let\optexpose\paren@optexpose\textup{\ref{#1})}\let\optexpose\old@optexpose}%\neweqref{XII.eq4.1} engendre "XII (4.1)" au lieu de "(XII 4.1)"
\makeatother
\newcommand{\refvb}[2][VI${}_{\text{B}}$ ]{\sisi{#1\refnoexpose{#2}}{\ref{#2}}}

\newcommand{\leftexp}[2]{{\vphantom{#2}}^{#1}{#2}}
\newcommand{\AAA}{\mathbf{A}}
\newcommand{\Sf}{\mathfrak{S}}
\DeclareMathOperator*{\quotprojlim}{\text{``$\varprojlim$''}}
\DeclareMathOperator*{\quotinjlim}{\text{``$\varinjlim$''}}
\DeclareMathOperator{\Fl}{F\sisi{\ell}{l}}
\DeclareMathOperator{\Pl}{P\sisi{\ell}{l}}
%\DeclareMathOperator{\cl}{c\sisi{\ell}{l}}
\DeclareMathOperator{\tdf}{tdf}
\DeclareMathOperator{\rg}{rg}
\newcommand{\iInd}{\sisi{I}{i}nd}
\newcommand{\iind}{ind}
\newcommand{\ppro}{pro}
\newcommand{\tTor}{\sisi{T}{t}or}
\newcommand{\ttor}{tor}
\newcommand{\chap}{Chap.}
\newcommand{\Ffex}{\Ff\mathit{ex}}
\newcommand{\sheafk}{\underline{k}}
\newcommand{\sheafn}{\underline{n}}
\newcommand{\sheafGamma}{\underline{\Gamma}}

\sisi{
\DeclareMathOperator{\simarrow}{\rule{.15mm}{1.8mm}\kern -3.4pt\rightsquigarrow}
\newcommand{\sheafell}{\underline{\ell}}
\newcommand{\sheafZZ}{\underline{\ZZ}}
}{
\newcommand{\nderef}[1]{la N.D.E.~(\ref{#1}) page~\pageref{#1}}
}

%%%%%%%%%%%%%%%%%
%%%%% expose XVII
%%%%%%%%%%%%%%%%%



\externaldocument{../01/01.tex}
\externaldocument{../02/02.tex}
\externaldocument{../03/03.tex}
\externaldocument{../04/04.tex}
\externaldocument{../05/05.tex}
\externaldocument{../05bis/05bis.tex}
\externaldocument{../06/06.tex}
\externaldocument{../07/07.tex}
\externaldocument{../08/08.tex}
\externaldocument{../09/09.tex}
%\externaldocument{../10/10.tex}
\externaldocument{../11/11.tex}
\externaldocument{../12/12.tex}
\externaldocument{../13/13.tex}
\externaldocument{../14/14.tex}
\externaldocument{../15/15.tex}
\externaldocument{../16/16.tex}
\externaldocument{../17/17.tex}
\externaldocument{../18/18.tex}
\externaldocument{../19/19.tex}

\setboolean{original}{false}

\begin{document}

\begin{center}
XVII. Cohomologie à supports propres\\
P. Deligne 
\end{center}

\selectlanguage{english}version : \showgitstatus\selectlanguage{french}
\tableofcontents
\mainmatter

\else
\chapter{Cohomologie à supports propres}\etiquette[XVII]{XVII}
%\addtocontents{toc}{P. Deligne}
\begin{center}
P. Deligne 
\end{center}
\fi

\setcounter{pageoriginal}{251}

\footnotetext[1]{Le présent exposé et le suivant, rédigés en 1968 et
  1969, reprennent et complètent les exposés oraux de A.~\textsc{Grothendieck}
  (de printemps 1964). Le rédacteur, qui n'assistait pas au séminaire
  oral, s'est partiellement inspiré des notes de A.~\textsc{Grothendieck}.}


\section*{Introduction}\pageoriginale

Dans cet exposé est développé le formalisme de la cohomologie à
support propre. Les questions de variance ont été traitées avec assez
de soin, alors qu'elles n'étaient que rendues plausibles dans le
séminaire oral. Ceci explique la longueur de l'exposé, dont les
paragraphes \hyperref[XVII.sec1]{1} à \hyperref[XVII.sec4]{4} sont consacrés aux catégories et à la
topologie générale. Il est très vivement recommandé au lecteur de ne
lire que les paragraphes \hyperref[XVII.sec5]{5} et \hyperref[XVII.sec6]{6}, où se
trouve concentrée la substance géométrique de l'exposé (\S \hyperref[XVII.sec5]{5}:
construction et variance de la cohomologie à support propre, théorèmes
de changement de base, de finitude, de \tTor-dimension finie et de
comparaison; \S \hyperref[XVII.sec6]{6}: théorie du morphisme \sisi{tracé}{trace} pour un morphisme
quasi-fini et plat).

Dans les paragraphes \hyperref[XVII.sec1]{1} et \hyperref[XVII.sec2]{2}, on « rappelle » quelques résultats sur les catégories dérivées. Le \S \hyperref[XVII.sec1]{1}
${\rm n^{\circ}}$ 1 est consacré au formalisme des signes. Dans le \S
\hyperref[XVII.sec3]{3}, on traite du problème de recoller deux formalismes de
variance. Dans le \S \hyperref[XVII.sec4]{4}, on introduit les résolutions plates et on
étudie les propriétés spéciales des résolutions flasques de
\textsc{Godement}.

Dans\pageoriginale cet exposé, les foncteurs $R^{q}f_{!}$ images
directes supérieures à supports propres, et le foncteur $Rf_{!}$ qui
leur donne naissance, ne sont définis que pour $f$ un morphisme
compactifiable (\hyperref[XVII.def3.2.1]{3.2.1}). Dans l'appendice, rédigé par
\textsc{B. Saint-Donat}, on montre comment étendre la définition aux
morphismes séparés de type fini de but quasi-compact quasi-séparé.\nde{\label{XVII.ndeConrad}le théorème de compactification de Nagata assure que tout morphisme séparé de type fini de but quasi-compact quasi-séparé est compactifiable. Voir B.~Conrad, « Deligne's notes on Nagata compactifications », \emph{J. Ramanujan Math.\ Soc.} \textbf{22} (2007), 205--257, cité [C] dans la suite, 4.1, ainsi que « Erratum for ``Deligne's notes on Nagata compactifications'' », \emph{J. Ramanujan Math.\ Soc.} \textbf{24} (2009), 427--428.}

Le \S \hyperref[XVII.sec5]{5} ${\rm n^{\circ}}$ \hyperref[XVII.sec5]{5} (cohomologie
d'un produit symétrique) ne servira plus dans ce séminaire. Il sera
utilisé dans \SGA 5 pour raffiner le théorème de rationalité des
fonctions $L$.\nde{voir \SGA 4$\frac{1}{2}$, Fonctions $L$ modulo $\ell^n$ et modulo $p$.}

Le \S \hyperref[XVII.sec6]{6} ${\rm n^{\circ}}$ \hyperref[XVII.sec3]{3} (théorie de la
trace pour coefficients \emph{continus}) ne sera utilisé dans l'exposé
XVIII que dans le cas relativement facile des groupes lisses; le
lecteur intéressé par le théorème de dualité de Poincaré (dualité
globale), et prêt à admettre un argument transcendant, pourra même se
dispenser complètement de lire ce \S \hyperref[XVII.sec6]{6} ${\rm n^{\circ}}$
\hyperref[XVII.sec3]{3}, ainsi que la plus grande partie du \S \refnoexpose{XVIII.sec1}
de XVIII.

\setcounter{section}{-1}
\section{Préliminaires terminologiques}\etiquette[0]{XVII.sec0}

\subsection{}\etiquette[0.1]{XVII.subsec0.1}
Le signe $=$ placé entre deux groupes de symboles désignant des objets
d'une catégorie signifiera parfois (par abus de notations) que ces
objets sont canoniquement isomorphes. La catégorie et l'isomorphisme
canonique devront en principe avoir été définis au préalable. Dans un
diagramme, le signe $=$ désignera alors l'isomorphisme lui-même.

\setcounter{subsection}{2}
\setcounter{subsubsection}{-1}
\subsubsection{}\etiquette[0.2.0]{XVII.subsubsec0.2.0}
Soit $f:S\rightarrow S'$ un morphisme de sites (\hyperrefext[IV][IV.subsubsec4.9.3]{4.9.3}). Si $U$ et
$U'$ sont des objets de $S$ et $S'$, un $f$-morphisme de $U$ dans $U'$
sera par définition un morphisme de $U$ dans $f^{\ast}U'$. Pour les
sites étales de schémas, on retrouve la notion usuelle.

\subsubsection{}\etiquette[0.2.1]{XVII.subsubsec0.2.1}
Si\pageoriginale $\mathcal{F}$ et $\mathcal{F}'$ sont des faisceaux
sur $S$ et $S'$, un $f$-morphisme de $\mathcal{F}$ dans
$\mathcal{F}'$ sera indifféremment
\begin{enumerate}
\renewcommand{\theenumi}{\roman{enumi}}
\renewcommand{\labelenumi}{(\theenumi)}
\item un morphisme de $f^{\ast}\mathcal{F}'$ dans $\mathcal{F}$

\item un morphisme de $\mathcal{F}'$ dans $f_{\ast}\mathcal{F}$

\item une fonction qui, à chaque $f$-morphisme $\varphi$ d'un objet
  $U$ de $S$ dans un objet $U'$ de $S'$ associe une fonction de
  $\mathcal{F}'(U')$ dans $\mathcal{F}(U)$, et ce de façon
  compatible avec la composition de $\varphi$ avec une flèche de $S$
  ou $S'$ \footnote{Ces terminologies ne sont pas compatibles à
  l'identification des objets de $S$ aux faisceaux associés.}.
\end{enumerate}

On voit sur (i), (ii) et (iii) que les faisceaux d'ensemble forment
une catégorie fibrée et cofibrée sur la catégorie des sites.

Itou pour les faisceaux de modules à gauche sur des sites
annelés. Itou pour les faisceaux étales sur des schémas; ce n'est pas
immédiatement un cas particulier de ce qui précède, car les sites
forment en fait une 2-catégorie et « site étale de $X$ »  n'est
qu'un pseudo-foncteur en $X$ (\hyperrefext[VII][VII.subsec1.4]{1.4}).

La définition précédente fait des faisceaux d'ensembles (\resp\ldots)
sur des sites variables une catégorie fibrée sur celle des sites, ayant
pour catégories fibres les \emph{catégories opposées} aux catégories
usuelles de faisceaux d'ensembles (\resp\ldots).

\subsection{}\etiquette[0.3]{XVII.subsec0.3}
Soient $I$ un ensemble fini, $\epsilon$ une fonction de $I$ à valeurs
dans $\{+1,-1\}$, $(\Aa_i)_{i\in I}$ une famille de catégories
additives graduées par des foncteurs de translation $T_i$ (\cite{SGA4-17:CD} I \sisi{1.10}{1.1.0})
et $F$ un multifoncteur multiadditif des catégories $\Aa_i$
dans une catégorie additive $\Aa$ graduée par le foncteur de
translation $T$. On suppose $F$ covariant (\resp contravariant) en les
$i$ tels que $\epsilon_i=1$ (\resp $\epsilon_i = -1$). Modifiant et
complétant\pageoriginale (\cite{SGA4-17:CD} \sisi{I 1.1}{II 2.1.1}), on dira que $F$ est un
\emph{foncteur gradué} si on s'est donné une famille
$(\varphi_i)_{i\in I}$ d'isomorphismes de foncteurs entre $T\circ F$
et $F\circ T_i^{\varepsilon_{i}}$, telle que les diagrammes suivants soient
anticommutatifs, pour $i\neq j$ dans $I$,
\[
\xymatrix{
F \circ (T_i^{\varepsilon_i}, T_j^{\varepsilon_j})  \ar[rr]^{\varphi_j
  \ast T_i^{\varepsilon_i}} \ar[dd]^{\varphi_i\ast T_j^{\varepsilon_j}}& &T\circ F \circ T_i^{\varepsilon_i} \ar[dd]^{T\ast
  \varphi_i}\\
 & - & \\
 T\circ F \circ T_j^{\varepsilon_j}  \ar[rr]^{T\ast \varphi_j} & &T\circ T
  \circ F.
}
\]
On laisse au lecteur le soin de définir le composé de deux foncteurs
gradués et de vérifier que c'est encore un foncteur gradué.

\subsection{}\etiquette[0.4]{XVII.subsec0.4}
Soient $\Aa$ une catégorie additive et $I$ un ensemble
fini. Pour $i\in I$, on désigne par $1_i$ le $i$\textsuperscript{ème}
vecteur de base de $\ZZ^{I}$. Un \emph{complexe} $I$\textsuperscript{uple} (\resp \emph{un complexe} $I$\textsuperscript{uple} \sisi{naïf}{\emph{naïf}}) de
$\Aa$ consiste en
\begin{enumerate}
\renewcommand{\theenumi}{\alph{enumi}}
\renewcommand{\labelenumi}{(\theenumi)}
\item une famille $(K^{k})_{k\in \ZZ^{I}}$ d'objets de
  $\Aa$;

\item pour chaque $i\in I$ et chaque $k\in \ZZ^{I}$, une flèche
  $d^{k}_i:K^{k}\rightarrow K^{k+1_i}$, ces flèches vérifiant
  $d^{k+1_i}_i d^{k}_i=0$ et, pour $i\neq j$
  $d^{k+1_j}_i d^{k}_j+d^{k+1_i}_j d^{k}_i=0$ (\resp
  $d^{k+1_j}_i d^{k}_j=d^{k+1_i}_j d^{k}_i$).
\end{enumerate}

\subsection{}\etiquette[0.5]{XVII.subsec0.5}
On\pageoriginale appellera \emph{\ppro-objet} d'une catégorie $\sisi{C}{\Cc}$ un
foncteur de $\sisi{C}{\Cc}$
dans $(\Ens)$ qui soit limite inductive filtrante (selon une petite
catégorie filtrante) de foncteurs représentables (\cf \hyperrefext[I][I.sec8.10]{8.10}). Tout
\ppro-objet est limite inductive selon un petit ensemble ordonné
filtrant de foncteurs représentables. Si $X_{i}$ est un système
projectif d'objets de $\sisi{C}{\Cc}$, indexé par une petite catégorie filtrante,
on désigne par $\quotprojlim X_{i}$ le \ppro-objet $\varinjlim
h^{X_i}$.

\subsection{}\etiquette[0.6]{XVII.subsec0.6}
Le lecteur dualisera \hyperref[XVII.subsec0.5]{0.5} au cas des \iind-objets (\Cf \hyperrefext[I][I.sec8.2]{8.2}).

\subsection{}\etiquette[0.7]{XVII.subsec0.7}
Conformément à la nouvelle terminologie, on appelle \emph{schéma} ce
qui s'appelait autrefois préschéma, et on appelle \emph{schéma séparé}
ce qui s'appelait autrefois schéma.

\subsection{}\etiquette[0.8]{XVII.subsec0.8}
La catégorie des \emph{schémas annelés} est la catégorie dont les
objets sont les schémas dont le site étale est muni d'un faisceau
d'anneaux, une flèche de $(S,\Aa)$ dans $(T,\Bb)$
étant un couple formé d'un morphisme de schémas $f$ de $S$ dans $T$,
et d'un homomorphisme $\varphi$ de faisceaux d'anneaux de
$f^{\ast}\Bb$ dans $\Aa$. Le morphisme de schémas $f$
est dit \emph{induit par} $(f,\varphi)$.

\subsection{}\etiquette[0.9]{XVII.subsec0.9}
Si $X$ est un schéma sur $Y$, on désigne par $(X/Y)^n$ le produit
fibré $n$-uple de $X$ sur $Y$.

\subsection{}\etiquette[0.10]{XVII.subsec0.10}
Soit $S$ un schéma. On appellera \emph{grand site étale} (\resp
\emph{grand site} fppf, \resp \emph{grand site} fpqc) de $S$ le site
$\Sch/{S}$, muni de la topologie étale (\resp fppf, \resp fpqc)
(\SGA 3 IV 6.3). On fera attention que ce n'est pas un $\sisi{U}{\Uu}$-site ($\sisi{U}{\Uu}$
étant l'univers \sisi{fini}{fixé}, $\Sch/{S}$ consistant en les schémas $\in
\sisi{U}{\Uu}$).

On\pageoriginale appelle \emph{petit site} fpqc (\resp \emph{petit
  site} fppf) de $S$ le site des schémas plats sur $S$ (\resp plats de
  présentation finie) muni de la topologie fpqc (\resp
  fppf). Lorsqu'il faudra éviter une confusion, on appellera
  \emph{petit site étale} de $S$ le site $S_{\et}$ (\hyperrefext[VII][VII.subsec1.2]{1.2}).

\subsection{}\etiquette[0.11]{XVII.subsec0.11}
Dans bien des cas, le rédacteur s'est permis de parler de diagrammes
commutatifs de morphismes de sites là où il eût fallu parler de
diagrammes essentiellement commutatifs, \ie commutatifs à isomorphisme
près (\cf IV 3.2.2). Le lecteur pourra vérifier que les arguments que
nous donnons s'appliquent aussi à la situation générale.

\subsection{}\etiquette[0.12]{XVII.subsec0.12}
La terminologie « \emph{schéma cohérent} »  pour « schéma
quasi-compact quasi-séparé »  a été subrepticement introduite par
endroits par A.~Grothendieck.

\subsection{}\etiquette[0.13]{XVII.subsec0.13}
Un faisceau de torsion $F$ sur un schéma $S$ sera dit \emph{premier
  aux caractéristiques résiduelles de} $S$ s'il est limite inductive
  de ses sous-faisceaux annulés par des entiers $n$ inversibles sur
  $S$.

\section{Les catégories dérivées}\etiquette[1]{XVII.sec1}

\subsection{Foncteurs exacts (les règles de signe)}\etiquette[1.1]{XVII.subsec1.1}

\subsubsection{}\etiquette[1.1.1]{XVII.subsubsec1.1.1}
Pour les théorèmes fondamentaux relatifs aux catégories dérivées, je
renvoie à \textsc{Verdier} \cite{SGA4-17:CD} et \cite{SGA4-17:V2}. Dans ce nᵒ, une
attention toute spéciale a été accordée aux problèmes de signes.

Rappelons que si $f:X\rightarrow Y$ est un morphisme de complexes
(dans une catégorie additive, sous-entendue par la suite), son
\emph{cône} $C(f)$\pageoriginale est défini par
$$\etiquette[1.1.1.1]{XVII.eq1.1.1.1}
C(f)^{n}=X^{n+1}\oplus Y^{n}\quad d^{n}=-d^{n+1}_X+f^{n+1}+d^{n}_Y{\quoi}.\leqno(1.1.1.1)
$$
Lorsque $X=0$ (\resp $Y=0$), on a $C(f)=Y$ (\resp $C(f)=X[1]$), de
sorte que le diagramme commutatif
\[
\xymatrix{
0        \ar@{=}[r] \ar[d]&   0   \sisi{\ar@{=}[r]}{\ar[r]} \ar[d]&  X
\ar@{=}[r]\ar[d]^{f}  & X \ar[d]\\
X        \ar[r]^{f}   &  Y    \ar@{=}[r] & Y     \ar[r]        & 0
}
\]
définit un « triangle »
\[
X\xrightarrow{f}Y\rightarrow C(f)\xrightarrow{j}X[1]{\quoi}.
\]

Dans la catégorie des complexes à homotopie près, on appelle
\emph{distingué} un triangle isomorphe à un triangle de ce type, et
\emph{antidistingué} un triangle qui devient distingué quand on change
le signe de ses flèches. On vérifie que le triangle défini par une
suite exacte de complexes scindée en chaque degré (\cite{SGA4-17:CD} p.~10, \cf la
démonstration de \hyperref[XVII.lem1.1.5.4]{1.1.5.4}) est antidistingué et que, à isomorphisme
près, tout triangle antidistingué est obtenu ainsi.

\subsubsection*{${}^{\ast}$1.1.2.}\etiquette[*1.1.2]{XVII.subsubsec*1.1.2}
Les calculs de signes s'effectuent le plus aisément à l'aide des « règles formelles »  suivantes. On écrit un élément de $C(f)^{n}$
sous la forme $1\otimes x+y$ $(x\in X^{n+1}, y\in Y^{n})$, $1$ étant une
« cellule »  de dimension $1$, et on pose $d(1)\otimes x=f(x)$. On
a donc
\[
d(1\otimes x+y)=d1\otimes x-1\otimes dx+dy=-1\otimes
dx+f(x)+dy{\quoi}._{\ast}
\]

\begin{enonce*}{Définition 1.1.3}\etiquette[1.1.3]{XVII.def1.1.3}
\begin{enumerate}
\renewcommand{\theenumi}{\roman{enumi}}
\renewcommand{\labelenumi}{(\theenumi)}
\item Un\pageoriginale foncteur gradué (\hyperref[XVII.subsec0.3]{0.3}) d'une
  catégorie triangulée $\Aa$ dans une catégorie triangulée $\Bb$ est
  dit exact s'il transforme triangles distingués en triangles
  distingués.

\item Un foncteur gradué contravariant d'une catégorie triangulée
  $\Aa$ dans une catégorie triangulée $\Bb$ est dit
  \emph{exact} si pour tout triangle distingué $(X,Y,Z,u,v,w)$ de
  $\Aa$ le triangle $(F(Z),F(Y), F(X), F(v),F(u),TF(w))$ est
  distingué (dans cette définition, on identifie $FX$ à $TFTX$ grâce à
  la graduation de $F$).

\item Soient $(\Aa_i)_{i\in I}$ une famille finie de
  catégories triangulées, $\epsilon$ une fonction de $I$ dans
  $\{+1,-1\}$ et $\Aa$ une catégorie triangulée. Un foncteur
  gradué $F$ du produit des $\Aa_i$ dans $\Aa$,
  covariant en les $i$ tels que $\epsilon(i)=1$ et contravariant en
  les $i$ tels que $\epsilon(i)=-1$ est dit \emph{exact} si quel que
  soit $i\in I$, les foncteurs de $\Aa_i$ dans $\Aa$
  déduits de $F$ en fixant toutes les variables, sauf la $i$\textsuperscript{ème}, sont des foncteurs exacts.
\end{enumerate}
\end{enonce*}

Un composé de foncteurs exacts est encore un foncteur exact.

\setcounter{subsubsection}{3}
\subsubsection{}\etiquette[1.1.4]{XVII.subsubsec1.1.4}
Si $(K^{n_1\ldots n_p}, d_1,\ldots,d_p)$ est un complexe multiple
(\hyperref[XVII.subsec0.4]{0.4}), le complexe simple associé est défini par
$$\etiquette[1.1.4.0]{XVII.eq1.1.4.0}
K^{n}=\prod_{\sum n_i=n}K^{(n_i)}\quav d=\sum d_i{\quoi}.\leqno(1.1.4.0)
$$

Quand, exceptionnellement, on le définira par une somme plutôt que par
un produit, ce fait sera signalé explicitement.

Soit $I$ un ensemble fini. On se propose, d'après \textsc{Cartan-Eilenberg} et
\textsc{J.P. Serre}, d'expliquer comment à tout complexe naïf $I$\textsuperscript{uple} $K$ est canoniquement\pageoriginale associé un complexe
$I$\textsuperscript{uple} (\hyperref[XVII.subsec0.4]{0.4}).

Pour tout ordre total $<$ sur $I$, soit $K(<)$ le complexe
$I$\textsuperscript{uple} suivant:
$$\etiquette[1.1.4.1]{XVII.eq1.1.4.1}
\begin{cases}
K(<)^{\sheafk }=K^{\sheafk }\\
d(<)_i^{\sheafk }=(-1)^{\sum\limits_{j<i} k_j} d_i{\quoi}.
\end{cases}\leqno{(1.1.4.1)}
$$

Si $<_1$ et $<_2$ sont deux ordres \sisi{totals}{totaux} sur $I$ \emph{l'isomorphisme
  canonique $\tau$ entre $K(<_1)$ et $K(<_2)$} est par définition
  celui donné par
$$\etiquette[1.1.4.2]{XVII.eq1.1.4.2}
\tau^{\sheafk }=(-1)^{\epsilon(<_1,<_2)},\quad\mbox{\rm où}\quad
\epsilon(<_1,<_2)=\sum_{\substack{i<_1j\\ \sisi{i<_2 i}{j<_2 i}}}k_ik_j{\quoi}.\leqno{(1.1.4.2)}
$$

\subsubsection*{1.1.4.2.1}\etiquette[1.1.4.2.1]{XVII.par1.1.4.2.1}
On vérifie que ces isomorphismes canoniques établissent un système
transitif d'isomorphismes entre les $K(<)$ pour $<$ ordre total sur
$I$; ces $K(<)$ forment donc un complexe $I$\textsuperscript{uple} unique
à isomorphisme unique près, qu'on désignera dans ce nᵒ
par la notation $K(\ast)$ et qu'on appelle le \emph{complexe
\emph{$I$\textsuperscript{uple}} associé ou complexe naïf} $I$\textsuperscript{uple}
$K$. On prendra garde que $K(\ast)^{\sheafk}$ \emph{n'est pas
canoniquement isomorphe à $K^{\sheafk}$: seul le choix d'un ordre
total sur $I$ permet d'identifier ces deux objets.}

Si $K$ est un complexe naïf $I$\textsuperscript{uple}, et si $i\in I$, on
désigne par $K[1_i]$ le complexe naïf suivant, de différentielles
notées $d_j[1_i]$:
$$\etiquette[1.1.4.3]{XVII.eq1.1.4.3}
\begin{cases}
K[1_i]^{\sheafk } = K^{\sheafk +1_i}\\
d_j[1_i]^{\sheafk } = d_j^{\sheafk +1_{i}}\quad{\mbox{\rm si }} i\neq
j\\
d_i[1_i]^{\sheafk }= -d_i^{\sheafk +1_i}{\quoi}.
\end{cases}\leqno{(1.1.4.3)}
$$

Si\pageoriginale $K$ est un complexe $I$\textsuperscript{uple}, on pose
$$\etiquette[1.1.4.4]{XVII.eq1.1.4.4}
\begin{cases}
K[1_i]^{\sheafk } = K^{\sheafk +1_i}\\
d_j[1_i]^{\sheafk } = -d_j[1_i]^{\sheafk +1_i}\sisi{}{.}
\end{cases}\leqno{(1.1.4.4)}
$$

Si $K$ est un complexe naïf $I$\textsuperscript{uple} et $<$ un ordre
total sur $I$, de sorte que $K(\ast)$ est donné par \eqhyperref[XVII.eq1.1.4.1]{1.1.4.1}, on
définit un isomorphisme $\sigma$ entre $K[1_i](\ast)$ et
$K(\ast)[1_i]$ par la formule
$$\etiquette[1.1.4.5]{XVII.eq1.1.4.5}
\sigma^{\sheafk }=(-1)^{\sum\limits_{j<i}k_j} \id_{K^{\sheafk +1_i}}{\quoi}.\leqno{(1.1.4.5)}
$$

\setcounter{subsubsection}{4}
\setcounter{paragraph}{5}
\paragraph{}\etiquette[1.1.4.6]{XVII.par1.1.4.6}
On vérifie que cet isomorphisme ne dépend pas du choix de l'ordre
total $<$; il définit donc un \emph{isomorphisme}, dit
\emph{canonique, entre} $K[1_i](\ast)$ et $K(\ast)[1_i]$.

\paragraph{}\etiquette[1.1.4.7]{XVII.par1.1.4.7}
Quels que soient $i$, $j$ dans $I$, si $K$ est un complexe ou un
complexe naïf $I$\textsuperscript{uple}, on identifie de façon évidente
$K[1_i][1_j]$ et $K[1_j][1_i]$.

Pour $i\neq j$ dans $I$, le diagramme suivant d'isomorphismes
canoniques du type \hyperref[XVII.par1.1.4.6]{1.1.4.6} ou \hyperref[XVII.par1.1.4.7]{1.1.4.7} est alors \emph{anticommutatif}:
$$\etiquette[1.1.4.7.1]{XVII.eq1.1.4.7.1}
\vcenter{
\xymatrix{
K[1_i][1_j](\ast) \ar@{=}[r] \ar@{<->}[d]& K[1_j][1_i](\ast)
\ar@{<->}[r] & K[1_j](\ast)[1_i]\ar@{<->}[d] \\
K[1_i](\ast)[1_j] \ar@{<->}[r] & K(\ast)[1_i][1_j]  \ar@{=}[r]& K(\ast)[1_j][1_i]\rlap{.}
}}\leqno{(1.1.4.7.1)}
$$


\subsubsection{}\etiquette[1.1.5]{XVII.subsubsec1.1.5}
Soit\pageoriginale $F:\Aa\rightarrow \Bb$ un foncteur
contravariant additif de catégories additives. On définit comme suit
l'\emph{extension de $F$ aux complexes:} $F:C(\Aa)\rightarrow
C(\Bb)$. Si $K$ est un complexe de $\Aa$, on pose
$$\etiquette[1.1.5.1]{XVII.eq1.1.5.1}
\begin{cases}
F(K)^{k}=F(K^{-k})\\
d^{k}_{F(K)}=(-1)^{k+1}F(d^{-k-1})\sisi{}{.}
\end{cases}\leqno(1.1.5.1)
$$

L'\emph{isomorphisme canonique $\sigma$ entre $F(K)[1]$ et $F(K[-1])$}
est défini par
$$\etiquette[1.1.5.2]{XVII.eq1.1.5.2}
\sigma^{\sisi{\sheafk }{k}}=(-1)^{k+1}\id_{F(K^{-k-1})}{\quoi}.\leqno(1.1.5.2)
$$

Si $F:\Aa\rightarrow \Bb$ et $G:\Bb\rightarrow
\sisi{C}{\Cc}$ sont deux foncteurs contravariants additifs, on définit
l'isomorphisme canonique $\rho$ entre $G{\circ}F(K)$ \emph{et}
$G(F(K))$ par
$$\etiquette[1.1.5.3]{XVII.eq1.1.5.3}
\rho^{k}=(-1)^{k}\sisi{}{\id_{G\circ F(K^k)}}{\quoi}.\leqno(1.1.5.3)
$$

\begin{enonce*}{Lemme 1.1.5.4}\etiquette[1.1.5.4]{XVII.lem1.1.5.4}
\begin{enumerate}
\renewcommand{\theenumi}{\roman{enumi}}
\renewcommand{\labelenumi}{(\theenumi)}
\item Pour tout foncteur contravariant additif
  $F:\Aa\rightarrow \Bb$, les formules \eqhyperref[XVII.eq1.1.5.1]{1.1.5.1},
  \eqhyperref[XVII.eq1.1.5.2]{1.1.5.2} définissent un foncteur contravariant exact de
  $K(\Aa)$ dans $K(\Bb)$.

\item Pour tout couple $(G,F)$ de foncteurs contravariants additifs
  composables, la formule \eqhyperref[XVII.eq1.1.5.3]{1.1.5.3} définit un isomorphisme de
  foncteurs gradués entre l'extension aux complexes de $GF$ et le
  composé des extensions aux complexes de $G$ et $F$.
\end{enumerate}
\end{enonce*}

\begin{enonce*}[remark]{Preuve}
Un\pageoriginale morphisme de degré $i:f:K\rightarrow L$ est au choix
\begin{enumerate}
\renewcommand{\theenumi}{\alph{enumi}}
\renewcommand{\labelenumi}{(\theenumi)}
\item un morphisme de $K$ dans $L[i]$

\item un morphisme de $K[-i]$ dans $L$

\item une famille $f^{n}:K^{n}\rightarrow L^{n+i}$ qui commute ou
  anticommute aux différentielles selon la parité de $i$.
\end{enumerate}
\end{enonce*}

Si $f:K\rightarrow L[i]$ est un morphisme de degré $i$, alors
$F(f):F(L)[-i]\rightarrow K$ est encore un morphisme de degré $i$;
pour $i=1$, on a
$$\etiquette[1.1.5.5]{XVII.eq1.1.5.5}
F(f)^{n}=(-1)^{n+1}F(f^{-n-1}){\quoi}.\leqno(1.1.5.5)
$$

Si $0\rightarrow K\xrightarrow{u} L\xrightarrow{v} M\rightarrow 0$ est
une suite exacte de complexes scindée degré par degré, le choix d'un
scindage permet d'écrire la différentielle de $L$ sous la forme
\[
d_L=d_K+d_M+f{\quoi},
\]
où $f$ est un morphisme de degré un de $M$ dans $K$. De même, la suite
exacte duale
\[
0\longrightarrow F(M)\longrightarrow F(L)\longrightarrow
F(K)\longrightarrow 0
\]
définit un morphisme de degré un $f^{\ast}$ de $F(K)$ dans $F(\sisi{L}{M})$. On
a $f^{\ast}=F(f)$. Il en résulte que l'image par $F$ d'un triangle
$(K,L,M,u,v,f)$ défini par une suite exacte courte de complexes est
encore un triangle du même type et, d'après \hyperref[XVII.subsubsec1.1.1]{1.1.1}, ceci prouve (i).

Sous les hypothèses (ii), pour $f$ de degré un, on a
$G(F(f))^{n}=(-1)^{n+1}(-1)^{-n}G\circ F(f)^{n}=-G\circ F(f)^{n}$; en
particulier, l'isomorphisme de graduation $G(F(K[1]))\simeq
G(F(K))[1]$ est l'isomorphisme ayant pour composantes $-1$.

Le\pageoriginale diagramme
\[
\xymatrix{
G\circ F(K)[1]  \ar@{<->}[r]^{\rho[1]} \ar@{=}[d] & G(F(K))[1] \ar@{<->}[d]\\
G\circ F(K[1])  \ar@{<->}[r]^{\rho} & G(F(K[1]))
}
\]
est donc commutatif, ce qui prouve (ii).

\subsubsection{}\etiquette[1.1.6]{XVII.subsubsec1.1.6}
Soit $s=+$, ou $s=-$; soient $I$ un ensemble fini,
$\epsilon:I\rightarrow \{+,-\}$ une fonction, $(\Aa_i)_{i\in
  I}$ et $\Aa$ des catégories additives, et $F$ un
multifoncteur multiadditif des catégories $\Aa_i$ dans
$\Aa$, covariant en les arguments d'indice $i$ tel que
$\epsilon(i)=+$ et contravariant en les autres. On se propose de
définir l'extension $K^{s}(F)$ (ou simplement $F$) de $F$ aux
complexes comme multifoncteur exact des catégories
$K^{s\epsilon(i)}(\Aa_i)$ ($s\epsilon(i)=\pm$, selon la règle
des signes) dans $K^{s}(\Aa)$, \footnote[1]{La construction
  fournit également une extension de $F$ en un multifoncteur des
  catégories $C^{s\epsilon(i)}(\Aa_i)$ dans
  $C^{s}(\Aa)$.} présentant la même variance que $F$. On
expliquera ensuite en quel sens cette construction est compatible à la
composition des foncteurs. D'autres conventions de \emph{degrés} sont
possibles, et seront utilisées; ce point est indépendant des questions
de \emph{signe} considérées ici.

\paragraph{}\etiquette[1.1.6.1]{XVII.par1.1.6.1}
Si $F$ est contravariant à une variable, on définit $K^{s}(F)$ par les
formules (\hyperref[XVII.subsubsec1.1.5]{1.1.5}).

\paragraph{}\etiquette[1.1.6.2]{XVII.par1.1.6.2}
Si $F$ est covariant, on définit le foncteur $K^{s}(F)$ par passage au
quotient à partir du composé du foncteur évident du produit des
catégories $C^{s}(\Aa_i)$ dans la catégorie des complexes
$I$\textsuperscript{uples} naïfs, du foncteur (\hyperref[XVII.subsubsec1.1.4]{1.1.4})
de la catégorie des complexes $I$\textsuperscript{uples} naïfs dans celle
des\pageoriginale complexes $I$\textsuperscript{uples} (\hyperref[XVII.par1.1.4.2.1]{1.1.4.2.1}), et du
foncteur complexe simple associé \eqhyperref[XVII.eq1.1.4.0]{1.1.4.0}:
\[
K^{s}(F)=C^{s}(F)(\ast)_{s.a}{\quoi}.
\]

Soient des complexes $K_j\in K^{s}(\Aa_j)$, et $i\in I$; on
pose $K_j[1_i]=K_j$ pour $j\neq i$, et $K_i[1_i]=K_{i}[1]$. On a alors
des isomorphismes évidents:
\[
C^{s}(F)(K_{j}[1_i])\simeq C^{s}(F)(K_j)[1_i]{\quoi},
\]
et
\[
C^{s}(F)(K_j)(\ast)[1_i]_{s.a}\simeq
C^{s}(F)(K_j)(\ast)_{s.a}[1]{\quoi},
\]
d'où, via \eqhyperref[XVII.eq1.1.4.5]{1.1.4.5}, un \emph{isomorphisme canonique} $\sigma$ entre
$K^{s}(F)(K_{j}[1_i])$ et $K^{s}(F)(K_{j})[1]$. Il résulte de \hyperref[XVII.par1.1.4.7]{1.1.4.7}
que ces isomorphismes définissent une graduation de $K^{s}(F)$.

\paragraph{}\etiquette[1.1.6.3]{XVII.par1.1.6.3}
Dans le cas général, $F$ est le produit de foncteur contravariants
canoniques $\Aa_i\rightarrow \Aa^{\circ}_i$ (pour
$\epsilon(i)=-$) et d'un foncteur covariant $F^{+}$; on définit le
foncteur $K^{s}(F)$ par composition à partir de \hyperref[XVII.par1.1.6.1]{1.1.6.1}
et \hyperref[XVII.par1.1.6.2]{1.1.6.2}.

\begin{enonce*}{Proposition 1.1.7}\etiquette[1.1.7]{XVII.prop1.1.7}
Sous les hypothèses précédentes, le foncteur $K^{s}(F)$ est un
foncteur multiadditif exact.
\end{enonce*}

Il suffit de vérifier \hyperref[XVII.prop1.1.7]{1.1.7} sous les hypothèses de
\hyperref[XVII.par1.1.6.1]{1.1.6.1} (où la proposition est \hyperref[XVII.lem1.1.5.4]{1.1.5.4} (i)) ou de
\hyperref[XVII.par1.1.6.2]{1.1.6.2}. Pour vérifier que les foncteurs déduits de $K^{s}(F)$ en
fixant toutes les variables, sauf la $i$\textsuperscript{ème}, sont
exacts, on peut expliciter $K^{s}(F)$ en terme d'un ordre sur $I$ pour
lequel $i$ soit le plus petit élément; un triangle distingué
type\pageoriginale \eqhyperref[XVII.eq1.1.1.1]{1.1.1.1} est alors transformé en triangle
distingué type.

\setcounter{subsubsection}{7}
\subsubsection{}\etiquette[1.1.8]{XVII.subsubsec1.1.8}
Soient $I$ et $J$ deux ensembles finis, $\Psi:I\rightarrow J$ une
application surjective, $\epsilon_{I}:I\rightarrow \{+,-\}$ et
$\epsilon_{J}:J\rightarrow \{+,-\}$ des fonctions, et des
multifoncteurs multiadditifs, covariants en les arguments pour lesquels
$\epsilon =+$ et contravariant en les autres:
\begin{align*}
G_j&: \prod_{\Psi(i)=j}\Aa_i\longrightarrow
\Bb_j\qquad (j\in J)\\
F&:  \prod_{j\in J} \Bb_j\longrightarrow \sisi{C}{\Cc.}
\end{align*}

Choisissons des ordres totaux sur $I$ et $J$, de telle sorte que
$\Psi$ soit une application croissante. On définit alors
l'\emph{isomorphisme canonique $\rho$ entre l'extension aux complexes
  de $F\circ (G_j)$ et le composé des extensions aux complexes de $F$
  et des $G_j$} par la formule
$$\etiquette[1.1.8.1]{XVII.eq1.1.8.1}
\rho^{\sheafk }=\sisi{\prod_i \eta_i^{k_i}}{(-1)^{A(\sheafk )}}:\; \mbox{\rm automorphisme de }
F\circ (G_j)(K_{i}^{k_i\epsilon\sisi{}{_I}(i)\epsilon\sisi{}{_J}(\Psi(i))}){\quoi}.\leqno(1.1.8.1)
$$%
avec \sisi{$\eta_i=-1$ si $\epsilon_I(i)=\epsilon_{J}(\Psi(i))=-$,
$\eta_i=1$ sinon.}{\[
A(\sheafk )=\sum_{\epsilon_I(i)=\epsilon_J(\Psi(i))=-} k_i + \sum_{\epsilon_J(j)=-} \sum_{\substack{\Psi(a)=\Psi(b)=j\\a<b}} k_a k_b.
\]}%Voir l'Erratum dans SGA 4 1/2, p. 312

Cette formule généralise \eqhyperref[XVII.eq1.1.5.3]{1.1.5.3}, on vérifie que $\rho$ ne dépend
pas des ordres choisis sur $I$ et $J$, et est un isomorphisme de
foncteurs gradués (\cf \hyperref[XVII.lem1.1.5.4]{1.1.5.4} (ii)).

Ces isomorphismes $\rho$ vérifient la condition évidente de cocycle
pour un composé triple.

\subsubsection{}\etiquette[1.1.9]{XVII.subsubsec1.1.9}
Soient\pageoriginale $I$ un ensemble fini, $i$ et $j$ deux éléments
distincts de $I$, $I'=I\backslash\{i,j\}$, $(\Aa_i)_{i\in I}$ et
$\Aa$ des catégories additives et $F:(\Aa_i)_{i\in
  I}\rightarrow \Aa$ et $G:(\Aa_i)_{i\in
  I'}\rightarrow \Aa$ des multifoncteurs multiadditifs. On
suppose que $\Aa_i=\Aa_j$ et que $F$ est contravariant
en la $i$\textsuperscript{ème} variable et covariant en la $j$\textsuperscript{ième}. Pour $X$, $Y\in \Ob \Aa_i$, on désignera par
$F_{X,Y}$ le multifoncteur des $(\Aa_i)_{i\in I'}$ dans
$\Aa$ obtenu en fixant les $i$\textsuperscript{ième} et $j$\textsuperscript{ième}
arguments:
\[
F_{X,Y}(\ldots)=F(\ldots\sisi{}{,} X,\ldots,Y,\ldots){\quoi}.
\]
On appelle \emph{morphisme de contraction} $c:F\rightarrow G$ la
donnée pour tout $X\in \Aa$ d'un morphisme de foncteurs
\[
c_{X}:F_{X,X}\longrightarrow G{\quoi},
\]
tel que pour toute flèche $f:X\rightarrow Y$ de $\Aa_i$, le
diagramme
\[
\xymatrix{
F_{Y,X} \ar[d] \ar[r] & F_{X,X} \ar[d]^{c_X}\\
F_{Y,Y}  \ar[r]_{c_Y}     & G
}
\]
soit commutatif.

Les conventions \hyperref[XVII.subsubsec1.1.6]{1.1.6} sont motivées par le

\begin{enonce*}{Lemme 1.1.9.1}\etiquette[1.1.9.1]{XVII.lem1.1.9.1}
Soient $c:F\rightarrow G$ un morphisme de contraction et
$(K_{\ell})_{\ell \in I}$ des complexes dans les $\Aa_{\ell}$, avec
$K_{i}=K_{j}$. Soit $<$ un ordre total sur $I$ pour lequel $i$ soit le
prédécesseur de $j$. Soit, avec les notations de
\hyperref[XVII.subsubsec1.1.6]{1.1.6}, pour $k\in \ZZ$,
\[
c^{k}:K^{s}(F)(K_{\ell})^{k}\longrightarrow K^{s}(G)(K_{\ell})^{k}
\]
le\pageoriginale morphisme dont la restriction à $F(\ldots
K_{i}^{k_i}K_{j}^{k_j}\ldots)$ vaut $0$ si $k_i\neq k_j$, et vaut
$c_{K_i^{k_i}}$ si $k_{i}=k_{j}$. Alors, $(c^{k})_{k\in \ZZ}$
est un morphisme de complexes.
\end{enonce*}

On laisse au lecteur le soin de vérifier ce lemme, et le fait que le
\emph{morphisme de contraction} $(c^{k})$ ne dépend pas de l'ordre
choisi sur $I$.

L'exemple standard de morphisme de contraction est celui de la loi de
composition
\[
\Hom(Y,Z)\otimes \Hom(X,Y)\longrightarrow \Hom(X,Z)\sisi{}{.}
\]

\begin{enonce*}[remark]{Exemples 1.1.10}\etiquette[1.1.10]{XVII.exem1.1.10}
\begin{enumerate}
\renewcommand{\theenumi}{\roman{enumi}}
\renewcommand{\labelenumi}{(\theenumi)}
\item Si on applique les définitions de \hyperref[XVII.subsubsec1.1.6]{1.1.6} au foncteur
  \emph{produit tensoriel}, on retrouve la notion usuelle de produit
  tensoriel de deux complexes, la différentielle étant donnée par
\[
d(x\otimes y)=dx\otimes y+(-1)^{\deg x}x\otimes dy{\quoi}.
\]

\item Si on applique les définitions de \hyperref[XVII.subsubsec1.1.6]{1.1.6} au
  foncteur $\Hom(X,Y)$, \emph{en considérant $Y$ comme première
  variable et $X$ comme seconde variable}, convention qu'on suivra
  toujours, cette convention permet d'identifier, pour deux complexes
  $X$, $Y$, le $k$\sisi{-ème}{-ième} composant du complexe $\Hom^{\bigdot}(X,Y)$ défini dans
  \hyperref[XVII.subsubsec1.1.6]{1.1.6} à $\prod_n\Hom(X^{n},Y^{n+k})$, et on retrouve la
  différentielle habituelle: pour $f\in \Hom(X^{n},Y^{n+k})$, on a
\[
df=d\circ f-(-1)^{k}f\circ d{\quoi}.
\]

\item Les conventions précédentes ont pour vertu que les systèmes de
  flèches suivante sont des morphismes de complexes, sans qu'on doive
  les\pageoriginale perturber par aucun signe (\hyperref[XVII.subsubsec1.1.9]{1.1.9}):
\begin{align*}
& \Hom^{\bigdot}(L,M)\otimes \Hom^{\bigdot}(K,L)\longrightarrow
  \Hom^{\bigdot}(K,M):f\otimes g\mapsto f\circ g\\
& \Hom^{\bigdot}(L,M)\otimes L\longrightarrow M: f\otimes x\mapsto f(x){\quoi}.
\end{align*}
\end{enumerate}
\end{enonce*}

\subsubsection*{${}^{\ast}$1.1.10.1}\etiquette[*1.1.10.1]{XVII.par1.1.10.1}
La convention \eqhyperref[XVII.eq1.1.4.5]{1.1.4.5}--(\hyperref[XVII.par1.1.6.2]{1.1.6.2}) se justifie comme suit, pour
$F=\otimes$: un élément de $((K\otimes L)[1])^{n}$ s'écrit sous la
forme (\hyperref[XVII.subsubsec*1.1.2]{1.1.2}) (en se rappelant que $(K\otimes L)[1]=C(K\otimes
L\rightarrow 0)$ \eqhyperref[XVII.eq1.1.1.1]{1.1.1.1}) $\sum 1\otimes(x^{p}\otimes y^{q})$, un
élément de $(K[1]\otimes L)^{n}$ sous la forme $\sum (1\otimes
x^{p})\otimes y^{q}$ et un élément de $(K\otimes L[1])^{n}$ sous la
forme $\sum x^{p}\otimes (1\otimes y^{q})$; le signe apparaît quand
on interchange $x^{p}$ et $1$. Pour $F=$ dual, la convention \eqhyperref[XVII.eq1.1.5.2]{1.1.5.2}
se justifie ainsi: un élément de $K^{\ast}[1]$ s'écrit $1\otimes
\omega$; sa valeur sur un élément $1'\otimes x$ de $K[-1]$ ($1'$
cellule de dimension $-1$) est $<1\otimes \omega, 1'\otimes
x>=(-1)^{\deg\omega}<1, 1'><\omega,x>=(-1)^{\deg\omega}<\omega,x>$;
c'est la convention adoptée.${}_{\ast}$

\begin{enonce*}[remark]{Remarque 1.1.11}\etiquette[1.1.11]{XVII.rem1.1.11}
« Rappelons »  qu'une \emph{catégorie sous-additive} est une
catégorie dont les ensembles de flèches $\Hom(X,Y)$ sont munies de
lois de groupes abéliens, la composition des morphismes étant
biadditive. Un \emph{foncteur additif} $F$ entre catégories
sous-additives et un foncteur qui vérifie la formule
\[
F(f+g)=F(f)+F(g)
\]
pour $f$, $g\in \Hom(X,Y)$. Les foncteurs additifs d'une catégorie
sous-additive $\Aa$ dans un autre $\Bb$ forment une
catégorie, notée $\Hhom_{\text{\rm add}}(\Aa,\Bb)$.
\end{enonce*}

Soit\pageoriginale $(D)$ l'unique catégorie sous-additive, d'ensemble d'objets
indexés par $\ZZ$, $(X\sisi{_n}{^n})_{n\in \ZZ}$, avec
$$\etiquette[1.1.11.1]{XVII.eq1.1.11.1}
\begin{cases}
\Hom(X^{n},X^{n})=\ZZ,\;\mbox{\rm engendré par l'identité}\\
\Hom(X^{n},X^{n+1})=\ZZ,\;\mbox{\rm engendré par une flèche
  notée $d_n$}\\
\Hom(X^{i},X^{j})=0\;\mbox{si } j-i\neq 0\;\mbox{\rm ou } 1{\quoi}.
\end{cases}\leqno(1.1.11.1)
$$

Soit $D$ le complexe de $(D)$ défini par $D^{n}=X^{n}$ et
$d^{n}=d^{n}$.

Pour toute catégorie sous-additive $\Aa$, le foncteur
$G\mapsto G(D)$ de $\Hhom_{\text{\rm add}}((D),\Aa)$
dans $C(\Aa)$ est un isomorphisme.

La construction \hyperref[XVII.subsubsec1.1.6]{1.1.6} s'étend au cas d'un foncteur multiadditif de
catégories sous-additives $\Aa_i$ dans une catégorie additive
$\Aa$.

Soit
\[
F^{\bigdot}=(F^{n},d^{n})_{n\in \ZZ}
\]
un complexe de foncteurs multiadditifs de catégories additives
$\Aa_i$ dans une catégorie additive $\Aa$. On peut
considérer $F^{\bigdot}$ comme un foncteur de $(D)$ dans une catégorie
de multifoncteurs, ou encore comme un foncteur multiadditif
\[
F_1:\prod\Aa_i\times (D)\longrightarrow \Aa{\quoi}.
\]
On définit l'\emph{extension de $F\sisi{}{^{\bigdot}}$ aux complexes} par la formule
$$\etiquette[1.1.11.2]{XVII.eq1.1.11.2}
F^{\bigdot}(K_{1}\ldots K_{n})=F_{1}(K_{1}\ldots K_{n},D){\quoi}.\leqno(1.1.11.2)
$$

\setcounter{subsubsection}{11}
\subsubsection{}\etiquette[1.1.12]{XVII.subsubsec1.1.12}
Soient $\Aa$ et $\Bb$ deux catégories additives et
\[
F^{\bigdot}=(F^{n},d^{n})_{n\in \ZZ}
\]
un complexe de foncteurs additifs de $\Aa$ dans
$\Bb$. Supposons que les foncteurs\pageoriginale $F^{i}$
admettent des adjoints à droite $F_{i}$; si
$\leftexp{t}{d}^{i}:F_{i+1}\rightarrow F_i$ est le transposé de $d^{i}$, soit
$F$. le complexe de foncteurs additifs de $\Bb$ dans
$\Aa$, de composantes les $F_{i}$, et d'opérateur différentiel
\[
(d_\bigdot)^{i}=(-1)^{i+1}\,\leftexp{t}{d}^{-i-1}:F_{-i}\longrightarrow
F_{-i-1}{\quoi}.
\]
Soit ${}^{t}$ le foncteur \emph{contravariant} de $(D)$ dans $(D)$
donné par
\[
\leftexp{t}{X}^{n}=X^{-n}\quav \leftexp{t}{d}^{n}=d^{-n-1}{\quoi}.
\]
Le diagramme suivant est commutatif à isomorphisme canonique près
$$\etiquette[1.1.12.1]{XVII.eq1.1.12.1}
\vcenter{
\xymatrix@C=6em{
\mathcal{A} \times (D) \times \mathcal{B}  \ar[d]^{(\Id,{}^t,\Id)}\ar[r]^(.55){(F^{n}(A), \Id)} &
\mathcal{B}\times \mathcal{B}  \ar[r]^{\Hom_{\mathcal{B}}}&
(\Ab)\ar@{=}[d] \\
\mathcal{A} \times (D) \times \mathcal{B}
\ar[r]^(.6){(\Id,(F_{\bigdot})^n(B))} & \mathcal{A} \times \mathcal{A}
\ar[r]^{\Hom_{\mathcal{A}}} & (\Ab) \rlap{.}
}}\leqno{(1.1.12.1)}
$$
Ceci exprime la formule d'adjonction
\begin{equation*}
\Hom(F^{n}(A),B)\overset{\sim}{\longleftrightarrow}\Hom(A,F_{n}(B)){\quoi}.
\end{equation*}
Les conventions générales \hyperref[XVII.def1.1.3]{1.1.3} fournissent donc un isomorphisme
\begin{equation*}
\Hom^{\bigdot}(F^{\bigdot}(K),L)\overset{\sim}{\longleftrightarrow}\Hom\sisi{}{^\bigdot}(K,F_{\bigdot}(L)){\quoi}.
\end{equation*}

\begin{enonce*}{Définition 1.1.13}\etiquette[1.1.13]{XVII.def1.1.13}
Si $K$ est un complexe d'une catégorie abélienne $\Aa$, on
appelle \emph{tronqué à droite (\resp à gauche)} de $K$ en dimension
$n$, et on désigne par $\tau_{\leq n}(K)$ \emph{(\resp $\tau_{\geq
    n}(K)$)} les complexes suivants
\begin{align*}
\tau_{\leq n}(K) &: \ldots K^{n-2}\longrightarrow
K^{n-1}\longrightarrow \Ker(d^n)\longrightarrow 0\\
\tau_{\geq n}(K) &: 0\longrightarrow \coker (d^{n-1})\longrightarrow
K^{n+1}\longrightarrow K^{n+2}\ldots.
\end{align*}
\end{enonce*}

\begin{enonce*}{Proposition 1.1.14}\etiquette[1.1.14]{XVII.prop1.1.14}
\begin{enumerate}
\renewcommand{\theenumi}{\roman{enumi}}
\renewcommand{\labelenumi}{(\theenumi)}
\item $H^{i}(\tau_{\leq n}(K))=0$\pageoriginale pour $i>n$ et $H^{i}(\tau_{\geq
n}(K))=0$ pour $i<n$.

\item Le morphisme canonique $\tau_{\leq n}(K)\rightarrow K$ (\resp
$K\rightarrow \tau_{\geq n}K$) induit des isomorphismes sur les
$H^{i}$ pour $i\leq n$ (\resp $i\geq n$).
\end{enumerate}
\end{enonce*}

Le foncteur $\tau_{\leq n}$ (\resp $\tau_{\geq n}$) transforme
morphismes homotopes en morphismes homotopes. Dès lors, il résulte de
\sisi{\hyperref[XVII.subsubsec1.1.4]{1.1.4}}{\hyperref[XVII.prop1.1.14]{1.1.14}} que le foncteur \hyperref[XVII.def1.1.13]{1.1.13} garde
un sens dans la catégorie dérivée.

\setcounter{subsubsection}{14}
\subsubsection{}\etiquette[1.1.15]{XVII.subsubsec1.1.15}
Si $K$ est un complexe double, de première (\resp deuxième)
différentielle $d'$ (\resp $d''$), on désigne par $\tau''_{\leq n}(K)$
le sous-complexe double de $K$ tel que
\begin{equation*}
\tau''_{\leq n}(K)^{pq}=
\begin{cases}
K^{pq}&\mbox{\rm si } q<n\\
\Ker(d'') &\mbox{\rm si } q=n\\
0 &\mbox{\rm si } q>n
\end{cases}
\end{equation*}
on définit de même $\tau''_{\geq n}(K)$, $\tau'_{\leq n}(K)$ et $\tau'_{\geq n}(K)$.

\subsubsection{}\etiquette[1.1.16]{XVII.subsubsec1.1.16}
On aura aussi à considérer les « tronqués bêtes »  d'un complexe
$K$, définis par
\begin{xalignat*}{2}
\sigma_{\leq n}(K)^{i}=
\begin{cases}
K^{i} &\mbox{\rm si } i\leq n\\
0 &\mbox{\rm si } i>n
\end{cases}
&&
\sigma_{\geq n}(K)^{i}=
\begin{cases}
K^{i} &\mbox{\rm si } i\geq n\\
0 &\mbox{\rm si } i<n{\quoi}.
\end{cases}
\end{xalignat*}

\subsection{Foncteurs dérivés}\etiquette[1.2]{XVII.subsec1.2}\pageoriginale

Soient $\Aa$ une catégorie triangulée et $S$ un système
multiplicatif saturé (\cite{SGA4-17:CD} \sisi{}{I }2.1.2). Si $A\in \Ob \Aa$, on
désignera par $A^{+}$ et $A^{-}$ les \iind\ et \ppro-objets ((\hyperref[XVII.subsec0.5]{0.5})
et (\hyperref[XVII.subsec0.6]{0.6}))
\begin{equation*}
A^{+}=\quotinjlim_{A\xrightarrow{s}A'} A' \qquad A^{-}=\quotprojlim_{A'\xrightarrow{s}A}A'{\quoi},
\end{equation*}
où $s$ parcourt la catégorie filtrante des flèches de $S$ de source ou
but $A$.

Les foncteurs $A\mapsto A^{+}$ et $A\rightarrow A^{-}$ rendent
inversibles les éléments de $S$ et définissent des foncteurs
pleinement fidèle de $\Aa(S^{-1})$ dans les catégories des Ind
et \ppro-objets de $\Aa$ respectivement.

\begin{enonce*}{Définition 1.2.1}\etiquette[1.2.1]{XVII.def1.2.1}
Soit $F$ un multifoncteur covariant exact des catégories triangulées
$(\Aa_i)_{0<i\leq n}$ dans la catégorie triangulée
$\Aa$ et soient $S_i$ et $S$ des systèmes multiplicatifs
saturés des catégories $\Aa_i$ et $\Aa$.
\begin{enumerate}
\renewcommand{\theenumi}{\roman{enumi}}
\renewcommand{\labelenumi}{(\theenumi)}
\item Le foncteur dérivé droit $RF$ de $F$ (relativement aux $S_i$ et
à $S$) est le foncteur des catégories $\Aa_i(S^{-1}_i)$ dans
la catégorie des \iInd-objets de $\Aa(S^{-1})$ qui rend
commutatif le diagramme suivant.
\[
\xymatrix@C=5em{%csy
  (\Aa_i)  \ar[d]\ar[r]^{(A_i)\sisi{\simarrow}{\mapsto}  F(A^{+}_i)} & \Ind
  \Aa\ar[d] \\
\Aa_i(S^{-1}_i)\ar[r]^{RF} & \Ind\Aa(S^{-1}).
}
\]

\item On\pageoriginale dit que $RF$ est défini en $(A_i)$ si le
\iInd-objet $RF(A_i)$ de $\Aa(S^{-1})$ est « essentiellement
constant », \ie provient d'un objet de $\Aa(S^{-1})$; cet
objet est appelé \emph{la valeur de} $RF$ en $(A_i)$.

\item On dit que $F$ est \emph{dérivable à droite si} $RF$ est partout
défini; on désignera alors encore par $RF$ le foncteur des
$\Aa_i(S^{-1}_i)$ dans $\Aa(S^{-1})$ qu'il définit.

\item Une famille $(A_i)$ d'objets des $\Aa_i$ est dite
\emph{déployée pour} $RF$ si le morphisme canonique de $F(A_i)$ dans
$RF(A_i)$ est un isomorphisme.
\end{enumerate}
\end{enonce*}

On laisse au lecteur le soin de définir par dualité le dérivé gauche
$LF$ de $F$ (à valeur dans $\Pro \Aa(S^{-1})$) et d'étendre
les définitions précédentes aux cas où $F$ est covariant en certains
arguments et contravariant en d'autres.

Lorsque $RF$ est partout défini, la définition qu'on en a donné ici
coïncide avec celle de Verdier (\cite{SGA4-17:CD} p.~39), comme on le vérifie
facilement.

\begin{enonce*}{Proposition 1.2.2}\etiquette[1.2.2]{XVII.prop1.2.2}
Soient $(A'_1, A_1,A''_1)$ un triangle distingué dans $\Aa_1$,
et\break $A_i\in \Ob \Aa_i$ $(i\neq 1)$.
\begin{enumerate}
\renewcommand{\theenumi}{\roman{enumi}}
\renewcommand{\labelenumi}{(\theenumi)}
\item Si $(A'_{1},(A_{i}))$ et $(A''_{1},(A_{i}))$ sont déployées pour
  $RF$, alors la famille $(A_{1},(A_{i}))$ l'est aussi.

\item Si $RF$ est défini en $(A'_1, (A_i))$ et en $(A''_1, (A_1))$, il
l'est aussi en $(A_1, (A_i))$ et le triangle $(RF(A'_1,(A_i)),
RF(A_1,(A_i)),RF(A''_1,(A_i)))$ est distingué dans
$\Aa(S^{-1})$.
\end{enumerate}
\end{enonce*}

Prouvons\pageoriginale que (ii) $\Rightarrow$ (i). Puisque $F$ est exact,
la flèche $$(F(A'_1,(A_i)), F(A_1,(A_i)), F(A''_1,(A_i)))\rightarrow
(RF(A'_1,(A_i)), RF(A_1,(A_i)), RF(A''_1,(A_i)))$$ est un morphisme de
triangles distingués. Par hypothèse, cette flèche est un isomorphisme
en deux sommets; c'en est donc un au troisième, comme on voulait.

\begin{enonce*}{Lemme 1.2.2.1}\etiquette[1.2.2.1]{XVII.lem1.2.2.1}
Si $\Delta=(X,Y,Z)$ est un triangle distingué d'une catégorie triangulée
$\Aa$ et $S$ un système multiplicatif saturé de $\Aa$,
alors le triangle $\Delta^{+}=(X^{+},Y^{+},Z^{+})$ de
$\Ind\Aa$ est limite de triangles distingués de $\Aa$.
\end{enonce*}

Désignons par $\mathcal{L}$ la catégorie des morphismes de triangle
distingués de source $\Delta$ dont les flèches appartiennent à $S$.
\begin{enumerate}
\renewcommand{\theenumi}{\alph{enumi}}
\renewcommand{\labelenumi}{(\theenumi)}
\item Si $\Delta=(X,Y,Z)$ et $\Delta'=(X',Y',Z')$ sont deux triangles
distingués de $\Aa$ et si $f:\Delta\rightarrow \Delta'$ est un
$\Aa(S^{-1})$-morphisme de triangles, il existe dans
$\Aa$ des morphismes de triangles
$s:\Delta'\rightarrow\Delta''$ et $f_1:\Delta\rightarrow \Delta''$
tels que les composantes de $s$ soient dans $S$ et que le diagramme
\[
\vcenter{
\xymatrix{
\Delta                \ar[rd]^{f_1}  \ar[d]_{f}  &\\
\Delta'          \ar[r]^{s}    & \Delta''
}}\tag{1}
\]
soit commutatif dans $\Aa(S^{-1})$.

\sisi{\noindent}{}%
Il existe un diagramme commutatif
\[
\vcenter{
\xymatrix{
X \ar[rd]^{f_1} \ar[d]_{f}&  &\\
X' \ar[r]^{s\sisi{}{_X}} & X''\sisi{}{_0}&
}}(s\sisi{}{_X} \in S, f_1 \in \Fl \Aa)\tag{2}
\]
et\pageoriginale un morphisme $\Delta'\xrightarrow{s\sisi{_1}{}}\Delta''\sisi{_1}{_0}$
dans $S$, admettant $s\sisi{}{_X}$ pour composante. Raisonnons de même en les
autres sommets: on obtient un diagramme du type (1), dans lequel
toutefois $f_1$ n'est pas encore un morphisme de triangle \emph{dans}
$\Aa$, les diagrammes
\[
\vcenter{
\xymatrix{
Z  \ar[r]^{f\sisi{}{_1}} \ar[d]_{w}& Z''_{0} \ar[d]_{w\sisi{'}{''}}\\
X         \ar[r]^{f\sisi{}{_1}} & X''_{0}
}}\tag{3}
\]
par exemple pouvant ne pas être commutatifs \emph{dans} $\Aa$,
seulement dans $\Aa(S^{-1})$. Il existe
$t\sisi{_0}{_X}:X''_0\rightarrow X''_1$ dans $S$ tel que $t\sisi{}{_X}w''f\sisi{}{_1}=t\sisi{}{_X}f\sisi{}{_1}w$
et $t:\Delta''_0\rightarrow \Delta''_1$ admettant $t\sisi{_0}{_X}$
pour composante. Remplaçons $\Delta''_0$ par $\Delta''_1$,
$s$ par $ts$ et $f_1$ par $tf_1$; dans le diagramme de type (1)
obtenu, (3) est cette fois commutatif; procédant de même aux autres
sommets, on obtient (1). On laisse au lecteur le soin de vérifier de
même:

\item Si $f$, $g:\Delta\rightrightarrows \Delta'$ sont deux morphismes
de triangles distingués dans $\Aa$, égaux dans
$\Aa(S^{-1})$, alors il existe un morphisme de triangles
distingués à flèches dans $S$, $s:\Delta'\rightarrow \Delta''$, tel
que $\sisi{S}{s}f=\sisi{S}{s}g$.

\item La catégorie $\Ll$ est filtrante à droite.

\item On achève la démonstration du lemme en notant que
\[
\Delta^{+}=\quotinjlim_{s\in \Ll}{\mbox{\rm (but
de $s$).}}
\]
\end{enumerate}

Prouvons\pageoriginale \hyperref[XVII.prop1.2.2]{1.2.2} (ii). Soient $X'$ et $X''$ des objets de
$\Aa(S^{-1})$ représentant $RF(A'_1,(A_i))$ et $RF(A''_1,
(A_i))$, et complétons le morphisme de degré 1 de $X''$ dans $X'$ en
un triangle distingué $(X',X,X'')$. Remplaçant les $A_i$, $A'_1$,
$A_1$ et $A''_1$ par le but d'une flèche de \sisi{$S$}{$S_i$ ou $S_1$} de source $A_i$,
$A'_1$, $A_1$ ou $A''_1$, on se ramène au cas où il existe un
diagramme commutatif
\[
\xymatrix{
F(A''_1, (A_i)) \ar[d] & X''  \ar[l]_-{a''} \ar[d]\\
F(A'_1, (A_i))    & X'  \ar[l]_-{a'}\rlap{,}
}
\]
où $a''$ et $a'$ sont des sections de flèches de $F(A''_1,(A_i))$ dans
$X''$ et de $F(A'_1,(A_i))$ dans $X'$. En vertu de l'axiome $TR$ 3 de
\cite{SGA4-17:CD} on peut compléter ce diagramme en un morphisme de triangles
\begin{align*}
(X'XX'')& \longrightarrow
(F(A'_1,(A_i)),F(A_1,(A_i)),F(A''_1,(A_i)))\\
& \longrightarrow
(RF(A'_1,(A_i)),RF(A_1,(A_i)),RF(A''_1,(A_i))){\quoi}.
\end{align*}

La flèche composée est un morphisme de triangles, de source un
triangle distingué et de but une limite (dans
$\Ind\Aa(S^{-1})$) de triangles distingués (\hyperref[XVII.lem1.2.2.1]{1.2.2.1}). Quel que
soit $Y$ dans $\Aa(S^{-1})$, le diagramme suivant sera
commutatif et ses lignes exactes
\[
\text{\tiny
$\xymatrix@C=1em{
\Hom\sisi{'}{(}Y, X''[\sisi{}{-}1]) \ar[r]\ar[d]& \Hom(Y,X') \ar[r] \ar[d]& \Hom(Y,X)
\ar[r]\ar[d]& \Hom(Y,X'') \ar[r] \ar[d]&
\Hom(Y,X'[1]) \ar[d]\\
\Hom(Y,RF(A''\sisi{}{_1},A_i)[\sisi{}{-}1])\ar[r] & \Hom(Y,RF(A'_1,A_i)\sisi{}{)}\ar[r] &  \Hom(Y,RF(A_1,A_i))\ar[r] &
\Hom(Y,RF(A''_1,A_i)\sisi{}{)}\ar[r] & \Hom(Y,RF(A'_1,A_i)[1])
}$}
\]
Par\pageoriginale hypothèse, les flèches autres que la flèche centrale
sont des isomorphismes. Le lemme des \hyperref[XVII.sec5]{5} montre que la flèche centrale
aussi est un isomorphisme, quel que soit $Y$, donc que le triangle
distingué $(X',X,X'')$ est isomorphe au triangle ($RF(A'_1,(A_i))$,
$RF(A_1,(A_i))$, $RF(A''_1,(A_i))$). Ceci achève la démonstration.

Si $\Aa$ et $\Bb$ sont deux catégories abéliennes et
$F$ un foncteur covariant additif de $\Aa$ dans $\Bb$,
$F$ définit des foncteurs exacts de $K(\Aa)$ dans
$K(\Bb)$, de $K^{+}(\Aa)$ dans $K^{+}(\Bb)$ et
de $K^{-}(\Aa)$ dans $K^{-}(\Bb)$ (\hyperref[XVII.prop1.1.7]{1.1.7}).

\begin{enonce*}{Définition 1.2.3}\etiquette[1.2.3]{XVII.def1.2.3}
\begin{enumerate}
\renewcommand{\theenumi}{\roman{enumi}}
\renewcommand{\labelenumi}{(\theenumi)}
\item On appelle \emph{foncteurs dérivés droits de $F$} les foncteurs
de $D^{+}(\Aa)$ dans $\Ind D^{+}\sisi{\Bb}{(\Bb)}$ et de
$D(\Aa)$ dans $\Ind D(\Bb)$ dérivés droits des
extensions de $F$ aux complexes $K^{+}(\Aa)\rightarrow
K^{+}(\Bb)$ et $K(\Aa)\rightarrow K(\Bb)$.

\item Un objet $A$ de $\Aa$ est dit \emph{acyclique à droite}
pour $F$, ou, par abus de langage, \emph{acyclique pour} $RF$, si le
complexe réduit à $A$ en degré $0$ est déployé pour $RF$.
\end{enumerate}
\end{enonce*}

On s'intéressera surtout au cas où $RF$ est partout défini; on peut
alors le considérer comme un foncteur de $D^{+}(\Aa)$ dans
$D^{+}(\Bb)$, ou de $D(\Aa)$ dans $D(\Bb)$, selon le cas.

\begin{enonce*}{Proposition 1.2.4}\etiquette[1.2.4]{XVII.prop1.2.4}
Le diagramme suivant est commutatif.
\[
\xymatrix{
D^{+}(\mathcal{A})  \ar[r] \ar[d]^{RF}& D(\mathcal{A})\ar[d]^{RF}        \\
 \Ind \;D^{+}(\mathcal{B}) \ar[r] & \Ind \; D(\mathcal{B}).
}
\]
\end{enonce*}

Si\pageoriginale $K$ est un complexe nul en degré $<n$ et si
$s:K\rightarrow L$ est un quasi-isomorphisme, le morphisme composé
$K\xrightarrow{s}L\rightarrow \sigma_{\geq n}(L)$
(\hyperref[XVII.subsubsec1.1.6]{1.1.6}) est encore un quasi-isomorphisme. Les
quasi-isomorphismes de $K$ dans un complexe borné inférieurement
forment donc une sous-catégorie (pleine) cofinale de la catégorie de
tous les quasi-isomorphismes de source $K$. La proposition en résulte
formellement.

En particulier, si le foncteur $RF$\sisi{,}{} est partout défini sur
$D(\Aa)$, il est partout défini sur $D^{+}(\sisi{\Bb}{\Aa})$, et
le foncteur dérivé $RF:D^{+}(\Aa)\rightarrow
D^{+}(\Bb)$ est restriction du foncteur dérivé
$RF:D(\Aa)\rightarrow D(\Bb)$.

\setcounter{subsubsection}{4}
\subsubsection{}\etiquette[1.2.5]{XVII.subsubsec1.2.5}
On laisse au lecteur le soin de:
\begin{enumerate}
\renewcommand{\theenumi}{\roman{enumi}}
\renewcommand{\labelenumi}{(\theenumi)}
\item Généraliser la définition \hyperref[XVII.def1.2.3]{1.2.3} au cas d'un complexe de
multifoncteurs, covariant en certaines variables et contravariant en
d'autres;

\item Dualiser la définition \hyperref[XVII.subsubsec1.2.5]{1.2.5} (i) pour définir les dérivés
gauches;

\item Etendre la proposition \hyperref[XVII.prop1.2.4]{1.2.4} au cas d'un complexe borné de
multifoncteurs;

\item Dualiser la proposition \hyperref[XVII.subsubsec1.2.5]{1.2.5} (iii) en
considérant des foncteurs dérivés gauches et l'inclusion de
$D^{-}(\Bb)$ dans $D(\Bb)$.
\end{enumerate}

\begin{enonce*}{Proposition 1.2.6}\etiquette[1.2.6]{XVII.prop1.2.6}
Soient $\Aa$ et $\Bb$ deux catégories abéliennes,
$𝒫$ une partie de $\Ob(\Aa)$ et $F$ un foncteur additif de
$\Aa$ dans $\Bb$. On suppose que
\begin{enumerate}
\renewcommand{\theenumi}{\alph{enumi}}
\renewcommand{\labelenumi}{(\theenumi)}
\item tout objet de $\Aa$ est quotient d'un élément de $𝒫$

\item pour toute suite exacte courte
\[
0\longrightarrow P\longrightarrow Q\longrightarrow R\longrightarrow
0{\quoi},
\]
où $Q$ et $R$ sont dans $𝒫$, $P$ appartient aussi à $𝒫$ et la
$0$-suite $0\rightarrow FP\rightarrow FQ\rightarrow FR\rightarrow
0$\pageoriginale est exacte. Alors, tout objet de $𝒫$ est
acyclique pour $LF$.
\end{enumerate}
\end{enonce*}

Si $A\in \Ob \Aa$, les résolutions gauches de $A$ par des
objets de $𝒫$ formeront un système cofinal dans la catégorie
filtrante des résolutions gauches de $A\sisi{_0}{}$. Il suffit donc de
prouver que si un complexe
\[
\ldots\longrightarrow A_2\longrightarrow A_1\longrightarrow A_0
\longrightarrow A\longrightarrow 0
\]
est acyclique et à objets dans $𝒫$, son image est acyclique. Le
coupant en suites exactes courtes
$$\etiquette[1.2.6.i]{XVII.eq1.2.6.i}
0\longrightarrow \Iim(d_{i-1})\longrightarrow A_i\longrightarrow
\Iim(d_i)\longrightarrow 0\qquad\sisi{(1.2.6. i)\quad}{}(i\geq 0)\sisi{}{\leqno(1.2.6.i)}%\tag{1.2.6.$i$}
$$
on voit par récurrence sur $i$ que $\Iim(d_i)\in \PP$ et que l'image
de \sisi{(\hyperref[XVII.prop1.2.6]{1.2.6} $i$)}{\eqhyperref[XVII.eq1.2.6.i]{1.2.6.i}} par $F$ est une suite exacte courte. L'image du complexe
est donc acyclique.

\begin{enonce*}{Proposition 1.2.7}\etiquette[1.2.7]{XVII.prop1.2.7}
Soit $F$ un foncteur additif d'une catégorie abélienne dans une
catégorie abélienne $\Bb$. Si tout objet de $\Aa$ est
quotient (\resp sous-objet) d'un objet acyclique pour $LF$ (\resp
$RF$) (\hyperref[XVII.def1.2.3]{1.2.3}), le foncteur dérivé $LF:D^{-}(\Aa)\rightarrow
D^{-}(\Bb)$ (\resp $RF:D^{+}(\Aa)\rightarrow
D^{+}(\Bb)$) est partout défini et tout complexe borné
supérieurement (\resp inférieurement) d'objets acycliques pour $LF$
(\resp $RF$) est déployé.
\end{enonce*}

Il suffit de considérer le dérivé gauche $LF$. Si $K$ est un complexe
borné supérieurement, les quasi-isomorphismes $s:L\rightarrow K$ avec
$L$ borné supérieurement et à objets acycliques forment un système
cofinal dans la catégorie des quasi-isomorphismes de but $K$. Il
suffit de montrer que pour $K$ à objets acycliques,
$F(s):F(L)\rightarrow F(K)$ est un quasi-isomorphisme;\pageoriginale
on aura alors dans le cas général $F(L)=LF(K)$. Soit $M$ le cône de
$s$. Le complexe est à objets acyclique\sisi{}{s} et $H^{\ast}(M)=0$; il faut
prouver que $H^{\ast}(F(M))=0$.

Soit $i$ un entier. Les complexes déployés forment une sous-catégorie
triangulée, de sorte que tout complexe borné d'objets acycliques,
notamment $\tau_{\geq i}(M)$, est déployé. Puisque $H^{j}(\tau_{\geq
i}(\sisi{L}{M}))=0$ pour $j>i$, on a, pour $j>i$,
$H^{j}(F(M))=H^{j}(F(\tau_{\geq i}(M)))=H^{j}LF(\tau_{\geq i}(M))=0$.

\begin{enonce*}{Proposition 1.2.8}\etiquette[1.2.8]{XVII.prop1.2.8}
Soit $F^{\bigdot}$ un complexe borné inférieurement de foncteurs de
$\Aa$ dans $\Bb$. Si pour tout $K\in
K^{+}(\Aa)$ existe un quasi-isomorphisme $L\xrightarrow{s}K$,
où $L\in K^{+}(\Aa)$ est déployé pour tous les $RF^{i}$, alors
le foncteur dérivé $RF^{\bigdot}:D^{+}(\Aa)\rightarrow
D^{+}(\Bb)$ existe et tout complexe \emph{borné
inférieurement} déployé pour tous les $RF^{i}$ est déployé pour $RF^{\bigdot}$.
\end{enonce*}

\sisi{\noindent}{}%
La preuve en est laissée au lecteur, ainsi que celle de l'énoncé dual,
ou des variantes obtenues en prenant $F$ contravariant, ou en prenant
un complexe borné de foncteurs et en travaillant dans $D(\Aa)$
et $D(\Bb)$.

\begin{enonce*}{Définition 1.2.9}\etiquette[1.2.9]{XVII.def1.2.9}
Soit $F$ un foncteur de $D^{b}(\Aa)$ dans $D(\Bb)$ et
$d$ un intervalle de $\ZZ$. On dit que $F$ est
\emph{d'amplitude cohomologique $\subset d$} si pour tout $A\in \Ob
\Aa$, $H^{i}F(A)=0$ pour $i\not\in d$.
\end{enonce*}

Très souvent, $d$ sera un intervalle de la forme $[0,x]$ ou $[-x,0]$
et on dira alors que $F$ est de dimension cohomologique \resp
homologique $\leq x$. On dira que $F$ est de dimension finie s'il
existe un intervalle fini $d$ de $\ZZ$ tel que $F$ soit
d'amplitude $\subset d$.

\begin{enonce*}{Proposition 1.2.10}\etiquette[1.2.10]{XVII.prop1.2.10}
Supposons\pageoriginale remplies les hypothèses de \hyperref[XVII.prop1.2.7]{1.2.7}, et que le
foncteur $RF:D^{+}(\Aa)\rightarrow D^{+}(\Bb)$ (\resp
$LF:D^{-}(\Aa)\rightarrow D^{-}(\Bb)$) est de
dimension finie. Alors tout complexe d'objets acycliques pour $RF$
(\resp $LF$) est déployé, et le foncteur $RF:D(\Aa)\rightarrow
D(\Bb)$ (\resp $LF:D(\Aa)\rightarrow D(\Bb)$) est partout défini.
\end{enonce*}

Pour la démonstration, on renvoie à Verdier \cite{SGA4-17:CD}. Il existe un énoncé
analogue pour les complexes bornés de foncteurs. On peut encore
remplacer l'exposant $+$ (\resp $-$) par $b$.

\section{Catégories fibrées en catégories dérivées}\etiquette[2]{XVII.sec2}

\subsection{Introduction}\etiquette[2.1]{XVII.sec2.1}

Le rédacteur insiste pour que le lecteur s'abstienne de lire ce \S. On
y donne le formalisme du théorème trivial de dualité et on y
résout la perplexité soulevée par Artin dans \hyperrefext[XII][XII.sec4]{4}.

Pour les notions de catégories sur une autre, de catégorie fibrée, de
catégorie cofibrée et de clivage, on renvoie à \SGA 1 VI. On établit
entre autres dans cet exposé une équivalence entre les notions de
catégories fibrées clivées normalisées sur une catégorie $\sisi{C}{\Cc}$, et de
pseudo-foncteur de $\sisi{C}{\Cc}^{\circ}$ dans $(\Cat)$.

\begin{enonce*}{Définition 2.1.1}\etiquette[2.1.1]{XVII.def2.1.1}
Si $X$ est un site annelé, on désigne par $D(X)$ la catégorie dérivée
de la catégorie abélienne des faisceaux de modules à gauche sur $X$.
\end{enonce*}

Quand il y a lieu d'éviter une confusion, on écrit plutôt
$D(X,\Aa)$, $\Aa$ désignant le faisceau d'anneaux.

Si\pageoriginale $f:X\rightarrow Y$ est un morphisme de sites annelés,
$f$ définit des foncteurs dérivés
\begin{itemize}
\item[(1)] $Rf_{\ast}:D^{+}(X)\rightarrow D^{+}(Y)$ et

\item[(2)] $Lf^{\ast}:D^{-}(Y)\rightarrow D^{-}(X)$,
\end{itemize}
et, par passage aux catégorie\sisi{}{s} opposées, des foncteurs
\begin{itemize}
\item[$\rm(1^{\circ})$] $Rf_{\ast}:D^{+}(X)^{\circ}\rightarrow
D^{+}(Y)^{\circ}$

\item[$\rm(2^{\circ})$] $Lf^{\ast}:D^{-}(Y)^{\circ}\rightarrow
D^{-}(X)^{\circ}$,
\end{itemize}
« fonctoriels »  en $f$. Plus précisément, $(1^{\circ})$ et
$(2^{\circ})$ définissent deux pseudo-foncteurs de la catégorie des
sites dans celle des catégories triangulées, l'un covariant et l'autre
contravariant. On trouve ainsi deux catégories sur celle des sites,
l'une cofibrée, de fibres les $D^{+}(X)^{\circ}$, et de foncteurs
image directe les $Rf_{\ast}$, l'autre fibrée, de fibres les
$D^{-}(X)^{\circ}$, et de foncteurs image réciproque les
$\sisi{\LL}{L} f^{\ast}$. Si le foncteur $f^{\ast}$ est de dimension
homologique finie (\ie de « \tTor-dimension finie ») on peut dans
(2) remplacer l'exposant par un exposant $+$. On obtient donc deux
catégories sur la catégorie $\Ss_{\tdf}$ d\sisi{a}{e}s sites annelés et
des morphismes de site annelés de \tTor-dimension finie, catégories dont
les fibres sont les catégories $D^{+}(X)^{\circ}$. Le « théorème
trivial de dualité »  de \textsc{Verdier} (\cite{SGA4-17:CD} p.~48) exprime que ces deux
catégories sont canoniquement isomorphes; elles sont donc fibrées et
cofibrées; les foncteurs « image directe »  sont les foncteurs
$Rf_{\ast}$\sisi{}{ }(sic), et les foncteurs « image réciproque »  sont les
foncteurs $Lf^{\ast}$\sisi{}{ }(sic).

\setcounter{subsection}{1}
\setcounter{subsubsection}{1}
\subsubsection{}\etiquette[2.1.2]{XVII.subsubsec2.1.2}

Le langage des catégories fibrées permet de résoudre aisément la
perplexité soulevée par \textsc{Artin} en \hyperrefext[XII][XII.sec4]{4}: il suffit d'appliquer la
proposition suivante à la catégorie fibrée et cofibrée sur celle des
sites\pageoriginale de fibres les opposées des catégories de faisceaux
sur les sites correspondants (\cf 0.2).

Soit $\Ee$ une catégorie sur $\Bb$, $f:y\rightarrow x$
une flèche de $\Bb$, et $F\in \Ob(\Ee_x)$. Rappelons
qu'un couple $(G,\varphi)$ $(G\in \Ob \Ee_y, \varphi\in
\Hom_f(G,F))$ est une \emph{image réciproque au sens strict} de $F$
par $f$ si, quels que soient $g:\sisi{Z}{z}\rightarrow y$ et $H\in \Ob
\Ee\sisi{_Z}{_z}$, on a
\[
\Hom_g(H,G)\xrightarrow[\sim]{\varphi}\Hom_{fg}(H,F){\quoi}.
\]
De même pour les images directes. Si $\Ee$ est fibrée ou cofibrée
sur $\Bb$, cette notion se réduit à celle d'image réciproque.

\begin{enonce*}{Proposition 2.1.3}\etiquette[2.1.3]{XVII.prop2.1.3}
Soient $\epsilon$ une catégorie sur $\Bb$, $y\in \Ob
\Bb$, $F\in \Ob \Ee_y$ et un diagramme commutatif dans
$\Bb$:
$$\etiquette[2.1.3.1]{XVII.eq2.1.3.1}
\vcenter{
\xymatrix{
y'  \ar[r]^{g'} \ar[d]_{f'}& y\ar[d]^{f}         \\
x'          \ar[r]    & x.
}}\leqno{(2.1.3.1)}
$$
On suppose que les images directes et réciproques $f_{\ast}F$,
$g^{\ast}f_{\ast}F$, ${g'}^{\ast}F$ et $f'_{\ast}{g'}^{\ast}F$
existent au sens strict. Il existe alors une et une seule flèche
$\varphi \in \Fl \Ee_{x'}$, dite \emph{flèche de changement de
base}, rendant commutatif le diagramme
$$\etiquette[2.1.3.2]{XVII.eq2.1.3.2}
\vcenter{
\xymatrix{
                             &  g'^{\ast}F \ar[ld]  \ar[r]    & F \ar[d]\\
f'_{\ast}g'^{\ast}F \ar[rd]_{\varphi}  &                                & f_{\ast}F\\
  & g^{\ast}f_{\ast}F \ar[ru]&.
}}\leqno{(2.1.3.2)}
$$
Si\pageoriginale de plus les images directes $g'_{\ast}{g'}^{\ast
}F$ et $g'_{\ast}f'_{\ast}{g'}^{\ast}F$ existent au sens strict, alors la flèche $\varphi$ coïncide avec
celle définie en \emph{\hyperrefext[XII][XII.sec4]{4}}. Si ${f'}^{\ast}f'_{\ast}{g'}^{\ast}F$
et ${f'}^{\ast}f'_{\ast}F$ existent au sens strict, elle coïncide
encore avec la $2$\sisi{ᵒ}{\textsuperscript{nde}} flèche définie en \emph{\hyperrefext[XII][XII.sec4]{4}}.

En particulier les deux flèches de changement de base
\emph{\hyperrefext[XII][XII.sec4]{4}} coïncident.
\end{enonce*}

Par hypothèse, les flèches du diagramme \eqhyperref[XVII.eq2.1.3.2]{2.1.3.2} induisent des
bijections
\sisi{\begin{align*}
& \Hom_y(F,F)=\Hom_{g'}({g'}^{\ast}F,F)=\Hom_{fg'}({g'}^{\ast}F,f_{\ast}F)
\quad\mbox{\rm et}\\
&
\Hom_{x'}(f'_{\ast}{g'}^{\ast}F,g^{\ast}f_{\ast}F)=\Hom_{f'}({g'}^{\ast}F,
g^{\ast}f_{\ast}F)=\Hom_{gf'}({g'}^{\ast}F, f_{\ast}F){\quoi}.
\end{align*}}{\[\Hom_{x'}(f'_{\ast}{g'}^{\ast}F,g^{\ast}f_{\ast}F)=\Hom_{f'}({g'}^{\ast}F,
g^{\ast}f_{\ast}F)=\Hom_{gf'}({g'}^{\ast}F, f_{\ast}F).\]}
Puisque $fg'=gf'$, $\Hom_{fg'}({g'}^{\ast}F,f_{\ast}F)$ est identique
à $\Hom_{gf'}({g'}^{\ast}F,f_{\ast}F)$ et les flèches du diagramme
\sisi{(2.1.2.2)}{\eqhyperref[XVII.eq2.1.3.2]{2.1.3.2}} induisent des \sisi{bijections}{applications}
\[
\Hom_y(F,F)\sisi{=}{\to}\Hom_{fg'}({g'}^{\ast}F,f_{\ast}F)=\Hom_{x'}(f'_{\ast}{g'}^{\ast}F,g^{\ast}f_{\ast}F){\quoi},
\]
et l'image de $1_F$ est la seule flèche de $\Ee_{x'}$ rendant
\eqhyperref[XVII.eq2.1.3.2]{2.1.3.2} commutatif.%Voir l'Erratum dans SGA 4 1/2, p. 312.

Par définition, dans \hyperrefext[XII][XII.sec4]{4}, la flèche de changement de base $\psi$
était celle rendant commutatif le diagramme suivant
$$\etiquette[2.1.3.3]{XVII.eq2.1.3.3}
\vcenter{
\xymatrix{
f'_{\ast}g'{}^{\ast}F \ar[d]^{\psi}  & g'{}^{\ast}F \ar[d]^{u_4} \ar[l] \ar@{=}[r] &
g'{}^{\ast}F \ar[r] \ar[d]^{u_3}   & F \ar@{=}[r] \ar[d]^{u_2}&  F \ar[d]^{u_1}\\
g^{\ast}f_{\ast}F   & {f'}^{\ast}g^{\ast}f_{\ast}F  \ar[l] \ar@{=}[r] &
g'f^{\ast}f_{\ast}F   \ar[r] & f^{\ast}f_{\ast}F \ar[r] &  f_{\ast}F.
}}\leqno{(2.1.3.3)}
$$
Les flèches horizontales et $u_1$ sont les flèches canoniques; $u_2$
est donc la flèche d'adjonction, $u_3={g\sisi{^{1}}{'}}^{\ast}(u_2)$, $u_4$ se
déduit de l'isomorphisme entre ${f'}^{\ast}g^{\ast}$ et
${g'}^{\ast}f^{\ast}$, et la flèche $\psi$ se déduit de $u_4$ par
adjonction.

L'application\pageoriginale composée de source ${g'}^{\ast}F$ et de
but $f_{\ast}F$ est celle qui apparaît dans \eqhyperref[XVII.eq2.1.3.2]{2.1.3.2} et il en est
donc de même dans le diagramme commutatif suivant, dont le carré de
gauche est extrait de \sisi{(2.1.2.3)}{\eqhyperref[XVII.eq2.1.3.3]{2.1.3.3}}
\[
\xymatrix{
g'{}^{\ast}F \ar[r]^{u_4} \ar[d]& f'{}^{\ast}g{}^{\ast}f_{\ast}F
\ar[r] \ar[d]&
f_{\ast}F \ar@{=}[d]\\
f'_{\ast}g'{}^{\ast}F  \ar[r]^{\psi} & g^{\ast}f_{\ast}F \ar[r] &
f_{\ast}F\rlap{.}
}
\]
Cela signifie que $\psi$ vérifie la propriété caractéristique de
$\varphi$.

La définition de $\varphi$ est autoduale, ce qui dispense de démontrer
la seconde partie de \hyperref[XVII.prop2.1.3]{2.1.3}.

\subsection{Catégories fibrées en catégories triangulées}\etiquette[2.2]{XVII.subsec2.2}

\begin{enonce*}{Définition 2.2.1}\etiquette[2.2.1]{XVII.def2.2.1}
On appelle « catégorie additive sur une catégorie $\Bb$ »
une catégorie $\Aa$ sur $\Bb$ dont les ensembles de
flèches $\Hom_f(X,Y)$ (pour $f\in \Fl (\Bb)$, $X$ (\resp
$Y$) au-dessus de la source (\resp du but) de $f$) ont été munis de
lois de groupes abéliens, de sorte que:
\begin{enumerate}
\renewcommand{\theenumi}{\roman{enumi}}
\renewcommand{\labelenumi}{(\theenumi)}
\item Quels que soient $x\xrightarrow{f}y\xrightarrow{g}z$ dans
$\Bb$ et $X$, $Y$, $Z$ dans $\Aa$ au-dessus de $\sisi{s}{x}$,
$y$ et $z$ respectivement, la loi de composition:
\[
\Hom_g(Y,Z)\times \Hom_f
(X,Y)\xrightarrow{\circ}\Hom_{gf}(X,Z)
\]
est biadditive.

\item Les catégories fibres sont additives.
\end{enumerate}
\end{enonce*}

On n'aura pas à utiliser ici que ces lois de groupe abélien, quand
elles veulent bien exister, sont uniquement déterminées par
$\Aa$, $\Bb$ et $\rho$.\nde{ici $\rho$ désigne le foncteur structurel $\Aa\to \Bb$.} Un $\Bb$-foncteur
$F:\Aa_1\rightarrow \Aa_2$ entre catégories additives
sur $\Bb$ est dit \emph{additif} s'il induit des
homomorphismes
\[
F:\Hom_{\Aa_1,f}(X,Y)\longrightarrow
\Hom_{\Aa_{2},f}(X,Y){\quoi}.
\]

\begin{enonce*}{Définition 2.2.2}\etiquette[2.2.2]{XVII.def2.2.2}
On\pageoriginale appelle « catégorie abélienne sur une catégorie
$\Bb$ »  une catégorie $\Aa$ additive sur
$\Bb$, dont les fibres sont abéliennes, et telle que pour
toute flèche $f:x\rightarrow y$ de $\Bb$ le bifoncteur
\[
\Hom_f(X,Y):\Aa^{\circ}_x\times \Aa_y\longrightarrow
(\Ab)
\]
soit exact à gauche en $X$ et $Y$.
\end{enonce*}

On n'aura pas à utiliser ici que si $\Aa$ est une catégorie
fibrée et cofibrée sur $\Bb$ (\SGA 1 VI 6) elle est abélienne
sur $\Bb$ dès que ses fibres sont des catégories abéliennes.

\begin{enonce*}{Définition 2.2.3}\etiquette[2.2.3]{XVII.def2.2.3}
On appelle « catégorie triangulée sur une catégorie $\Bb$ » une catégorie $\Aa$ additive sur $\Bb$, munie
d'un foncteur $T:\Aa\rightarrow \Aa$ et dont les
fibres $\Aa_x$ pour $x\in \Ob \Bb$ sont munies
d'ensembles $\Delta_x$ de triangles, ce de sorte que;
\begin{enumerate}
\renewcommand{\theenumi}{\roman{enumi}}
\renewcommand{\labelenumi}{(\theenumi)}
\item $T$ est un $\Bb$-automorphisme additif de $\Aa$

\item chaque catégorie fibre $\Aa_x$, munie du foncteur $T_x$
induit par $T$ et de l'ensemble $\Delta_x$, est une catégorie
triangulée \emph{(\cite{SGA4-17:CD} \sisi{}{I 1.}1.1)}.

\item si les triangles $(X,Y,Z,u,v,w)$ et $(X',Y',Z',u',v',w')$
appartiennent à $\Delta_x$ et $\Delta_{x'}$ respectivement, tout
diagramme
\[
\xymatrix{
            & Z   \ar[ldd]_{w}&            &  &           & Z'   \ar[ldd]|\hole_(.3){w'} & \\
            &   & Y\ar[lu]_{v} \ar[rrrr]_(.4){\sisi{j}{g}}&  &     &       &Y'\ar[lu]_{v'} \\
X   \ar[rru]^{u} \ar[rrrr]_{f}&             &          &  &   X'   \ar[rru]^{u'} &              &
}
\]
où $u'f=gu$ peut se prolonger en un morphisme de triangles
$(f,g,h):(X,Y,Z)\rightarrow (X',Y',Z')$.
\end{enumerate}
\end{enonce*}

$T$\pageoriginale sera appelé le foncteur de translation, on écrira
souvent $A[1]$ plutôt que $T(A)$, $A[n]$ désignant alors $T^{n}(A)$
$(n\in\ZZ)$.

Les triangles appartenant aux ensembles $\Delta_x$ seront dits
\emph{distingués} ou \emph{exacts}.

\begin{enonce*}{Proposition 2.2.4}\etiquette[2.2.4]{XVII.prop2.2.4}
Soient $\Aa$ une catégorie triangulée sur $\Bb$, $x$
un objet de $\Bb$, $(X,Y,Z,u,v,w)$ un triangle distingué de
$\Aa_x$. Quels que soient $f:y\rightarrow x$ et l'objet $A\in
\Ob \Aa_y$, la suite infinie suivante est exacte
\begin{align*}
\ldots  \longrightarrow \Hom_f(A,X) \longrightarrow \Hom_f(A,Y)&\longrightarrow
\Hom_f(A,Z)\\
 & \longrightarrow \Hom_f(A,X[1])\longrightarrow
\Hom_f(A,Y[1])\longrightarrow\ldots
\end{align*}
\end{enonce*}

Enoncé analogue pour $f:x\rightarrow y$ et $A\in \Ob \Aa_y$.

On laisse au lecteur le soin de démontrer cette proposition en
paraphrasant \textsc{Verdier} \cite{SGA4-17:CD} p.~4. On lui laisse aussi le soin de donner
un sens à la proposition suivante et de la vérifier:

\begin{enonce*}{Proposition 2.2.5}\etiquette[2.2.5]{XVII.prop2.2.5}
Soit $\Aa$ une catégorie additive (\resp abélienne, \resp
triangulée) sur $\Bb$. Pour tout foncteur
$\Bb'\rightarrow \Bb$ de but $\Bb$, la
catégorie $\Aa'=\fprod{\Aa}{\Bb'}{\Bb}$ est
additive (\resp abélienne, \resp triangulée) sur $\Bb$.
\end{enonce*}

\begin{enonce*}[remark]{Exemple 2.2.6}\etiquette[2.2.6]{XVII.exem2.2.6}
La catégorie des faisceaux de modules sur des sites annelés variables
est abélienne sur la catégorie des sites annelés.
\end{enonce*}

\begin{enonce*}[remark]{Exemple 2.2.7}\etiquette[2.2.7]{XVII.exem2.2.7}
Soit $\Aa$ une catégorie additive sur $\Bb$. La
catégorie $C(\Aa/\Bb)$ des complexes de $\Aa$
sur $\Bb$ a pour objets les complexes des catégories
fibres. Une flèche $f$ de $(X^n,d)$ dans $(Y^n,d)$ est un système
de\pageoriginale flèches $f^{n}:X^{n}\rightarrow Y^{n}$ tel que
$\sisi{fd}{df}^{n}=f^{n+1} d$. On pose $p(f)=p(f^{n})$ (cette flèche ne dépend
pas de $n$ car les $p(d)$ sont des
identités). $C(\Aa/\Bb)$ est une catégorie additive
sur $\Bb$. On définit un foncteur « de translation »  en posant
\[
(K[1])^{n}=K^{n+1}\quav (d[1])\sisi{_n}{^n}=-d\sisi{_{n+1}}{^{n+1}}{\quoi}.
\]
\end{enonce*}

\begin{enonce*}[remark]{Exemple 2.2.8}\etiquette[2.2.8]{XVII.exem2.2.8}
Soit $\Aa$ une catégorie additive sur $\Bb$. \emph{La
catégorie $K(\Aa/\Bb)$ des complexes de $\Aa$
sur $\Bb$ à homotopie près} a les mêmes objets que
$C(\Aa/\Bb)$. On l'obtient en décrétant nulles, dans
$C(\Aa/\Bb)$, les flèches homotopes à zéro; en
d'autres termes, pour $f:x\rightarrow y$ dans $\Bb$,
$X^{\bigdot}\in \Ob C(\Aa/\Bb)_x$ et $Y^{\bigdot}\in
\Ob C(\Aa/\Bb)_y$, on pose
\[
\Hom_f(X^{\bigdot},
Y^{\bigdot})=H^{0}\Hom^{\bigdot}_f(X^{\bigdot},Y^{\bigdot})
\]
\end{enonce*}

Le foncteur de translation passe au quotient.

Si $F:\Bb'\rightarrow \Bb$ est un foncteur de but
$\Bb$, et si
$\Aa'=\fprod{\Aa}{\Bb'}{\Bb}$, on
laisse au lecteur le soin d'identifier $C(\Aa'/\Bb')$
à $\fprod{C(\Aa/\Bb)}{\Bb'}{\Bb}$ et
$K(\Aa'/\Bb')$ à
$\fprod{K(\Aa/\Bb)}{\Bb'}{\Bb}$. En
particulier, les fibres de la catégorie $K(\Aa/\Bb)$
sont les catégories $K(\Aa_x)$ et, en tant que telles, sont
triangulées.

\begin{enonce*}{Proposition 2.2.9}\etiquette[2.2.9]{XVII.prop2.2.9}
\begin{enumerate}
\renewcommand{\theenumi}{\roman{enumi}}
\renewcommand{\labelenumi}{(\theenumi)}
\item Pour les structures additionnelles définies en \hyperref[XVII.exem2.2.8]{2.2.8}, la
catégorie $K(\Aa/\Bb)$ est triangulée sur $\Bb$.

\item La formation de la catégorie additive (\resp triangulée)
$C(\Aa/\Bb)$ (\resp $K(\Aa/\Bb)$) est
compatible à tout changement de catégorie base
$\Bb'\rightarrow \Bb$.
\end{enumerate}
\end{enonce*}

Les\pageoriginale conditions (i) et (ii) de la définition \hyperref[XVII.def2.2.3]{2.2.3} se
vérifient fibre par fibre. (iii) se vérifie en paraphrasant \textsc{Verdier}
\cite{SGA4-17:CD} p.~4.

\begin{enonce*}[remark]{Remarque 2.2.10}\etiquette[2.2.10]{XVII.rem2.2.10}
Ce qui précède s'applique aussi aux catégorie\sisi{}{s}
$K^{+}(\Aa/\Bb)$, $K^{-}(\Aa/\Bb)$ et
$K^{b}(\Aa/\Bb)$ obtenues en se limitant
respectivement aux complexes bornés inférieurement, supérieurement, ou bornés.
\end{enonce*}

\begin{enonce*}[remark]{Remarque 2.2.11}\etiquette[2.2.11]{XVII.rem2.2.11}
Prenant pour $\Bb$ la catégorie finale, on retrouve les
catégories usuelles de complexes.
\end{enonce*}

\setcounter{subsection}{2}
\setcounter{subsubsection}{11}
\subsubsection{}\etiquette[2.2.12]{XVII.subsubsec2.2.12}
Soit $F$ un pseudo-foncteur d'une catégorie $\Bb^{\circ}$ dans
$(\Cat)$. On sait (\SGA 1 VI) que $F$ définit une catégorie fibrée sur
$\Bb$, de fibres les catégories $F(x)$. Quels que soient
$f:x\rightarrow y$ dans $\Bb$ et $X\in \Ob F(\sisi{n}{x})$, $Y\in \Ob
F(y)$, on a par définition
$$\etiquette[2.2.12.1]{XVII.eq2.2.12.1}
\Hom_f(X,Y)=\Hom_{F(x)}(X,F(f)(Y)){\quoi}.\leqno(2.2.12.1)
$$
On sait (ibidem) que la construction précédente définit une
équivalence entre la 2-catégorie des pseudo-foncteurs de
$\Bb^{\circ}$ dans $(\Cat)$ et la 2-catégorie des catégories
fibrées sur $\Bb$.

Partons d'un pseudo-foncteur de $\Bb^{\circ}$ dans la
2-catégorie ayant pour objets les catégories additives (\resp
abéliennes, \resp triangulées) et pour 1-flèches les foncteurs additifs
(\resp exacts à gauche, \resp triangulés). On vérifie aussitôt que la
catégorie définie par \sisi{(2.1.12.1)}{\eqhyperref[XVII.eq2.2.12.1]{2.2.12.1}} est additive (\resp abélienne, \resp
triangulée). Réciproquement:

\begin{enonce*}{Proposition 2.2.13}\etiquette[2.2.13]{XVII.prop2.2.13}
La\pageoriginale construction précédente établit une équivalence entre
\begin{enumerate}
\renewcommand{\theenumi}{\alph{enumi}}
\renewcommand{\labelenumi}{(\theenumi)}
\item La $2$-catégorie des pseudo-foncteurs de $\Bb^{\circ}$
dans la $2$-catégorie ayant pour objets les catégories additives
(\resp abéliennes, \resp triangulées), et pour $1$-flèches les
foncteurs additifs (\resp exacts à gauche, \resp triangulés).

\item La $2$-catégorie des catégories additives (\resp abéliennes,
\resp triangulées) sur $\Bb$, qui sont fibrées sur $\Bb$.
\end{enumerate}
\end{enonce*}

La théorie développée en \SGA 1 VI permet de ne vérifier que le point
suivant:

\setcounter{subsubsection}{13}
\setcounter{paragraph}{0}
\paragraph{}\etiquette[2.2.13.1]{XVII.par2.2.13.1}
Soient $\Aa$ une catégorie additive (\resp abélienne, \resp
triangulée) sur $\Bb$, $f:x\rightarrow y$ une flèche de
$\Bb$ et supposons qu'un foncteur image réciproque
$f^{\ast}:\Aa_y\rightarrow \Aa\sisi{_n}{_x}$ existe; $f^{\ast}$
est alors additif (\resp exact à gauche, \resp triangulé).

Le foncteur $f^{\ast}$ donne lieu à des isomorphismes
\[
\Hom_f(X,Y)\simeq \Hom_x (X,f^{\ast}(Y)){\quoi}.
\]
Il faut prouver dans le cas triangulé\sisi{e}{}, que $f^{\ast}$ est triangulé
pour l'isomorphisme $Tf^{\ast}=f^{\ast}T$ ($T$ foncteur de
translation) qui rend commutatif les diagrammes suivants
$$\etiquette[2.2.13.2]{XVII.eq2.2.13.2}
\vcenter{
\xymatrix{
\Hom_f(X,Y)  \ar@{=}[r] \ar@{=}[d]& \Hom_x(X,f^{\ast}Y)\ar@{=}[rd]\\
\Hom_f(TX,TY)\ar@{=}[rd]           &      & \Hom_x(TX,Tf^{\ast}Y) \\
             &      \Hom_x(TX,f^{\ast}TY)\ar@{=}[ru] &
}}\leqno{(2.2.13.2)}
$$

\begin{itemize}
\item[\emph{Cas additif}:] Soit\pageoriginale un diagramme
$$\etiquette[2.2.13.3]{XVII.eq2.2.13.3}
\vcenter{
\xymatrix@C=2.5cm@R=1.75cm{
f^{\ast} X \ar@<-3pt>[d]_{f^{*}(u)}\ar@<3pt>^{f^{\ast}(v)}[d]        \ar[r]^{p} & X
\ar@<-3pt>_{u}[d] \ar@<3pt>^{v}[d]\\
f^{\ast} X' \ar[r]^{p'}& X'\rlap{.}
}}\leqno{(2.2.13.3)}
$$
On a
$p'f^{\ast}(u+v)=(u+v)p=up+vp=p'f^{\ast}(u)+p'f^{\ast}(v)=p'(f^{\ast}(u)+f^{\ast}(v))$
et donc $f^{\ast}(u+v)=f^{\ast}(u)+f^{\ast}(v)$.

\item[\emph{Cas abélien}:] par définition (\hyperref[XVII.subsec2.2]{2.2}), le foncteur
$\Hom(X,f^{\ast}Y)=\Hom_f(X,Y)$ est exact à gauche en $Y$, quel que
soit $X$, donc $f^{\ast}$ est exact à gauche.

\item[\emph{Cas triangulé}:] soit $(Y',Y,Y'')$ un triangle distingué dans
$\Aa_y$ et soit \sisi{$(f^{\ast}Y',f^{\ast}Y)$}{$(f^{\ast}Y',f^{\ast}Y, X'')$ un triangle distingué dans $\Aa_x$}. En vertu de \hyperref[XVII.def2.2.3]{2.2.3},
le morphisme canonique entre les bases de ces triangles peut se
prolonger en un morphisme de triangles, donnant lieu au diagramme
suivant
\[
\xymatrix@C=2em{
&  X''       \ar[rrr]^{\sisi{}{s}}\ar[ldd]   &     & &**[r]f^{\ast}Y''\ar[rrr]\ar[ldd]   &  & & Y''\ar[ldd] \\                      &     & f^{\ast}Y   \sisi{\ar[rrr]|\hole}{\ar@{=}[rrr]|\hole}\ar[lu]& &      &f^{\ast}Y \ar[rrr]|\hole\sisi{}{\ar[ul]}&&& Y\ar[lu]\\
  f^{\ast}Y' \sisi{\ar[rrr]}{\ar@{=}[rrr]}\ar[rru]         &&& f^{\ast}Y'\sisi{}{\ar[rru]}   \ar[rrr] && & Y'\ar[rru]&&.
}
\]

Prouvons que $s$ est un isomorphisme, ce qui achèvera la
démonstration. Quel que soit $X$ dans $\Ob \Aa_x$, on dispose
du diagramme commutatif:
\[
\text{\scriptsize
$\xymatrix@C=1em{
\Hom_x(X,f^{\ast}Y') \ar[r] \ar@{=}[d]& \Hom_x(X,f^{\ast}Y) \ar@{=}[d]\ar[r]& \Hom_x(X,X'')\ar[d]
\ar[r]& \Hom_x(X,f^{\ast}Y'[1])\ar[r]\ar@{=}[d]
& \Hom_x(X,f^{\ast}Y[1]) \ar@{=}[d]\\
\Hom_x(X,f^{\ast}Y') \ar[r]\ar@{=}[d]&  \Hom_x(X,f^{\ast}Y)\ar[r] \ar@{=}[d]&
\Hom_x(X,f^{\ast}Y'') \ar[r]\ar@{=}[d]&  \Hom_x(X,f^{\ast}Y'[1])\ar@{=}[d] \ar[r]&
\Hom_x(X,f^{\ast}Y[1]) \ar@{=}[d]\\
\Hom_f(X,Y)  \ar[r]& \Hom_f(X,Y)\ar[r]  &
\Hom_f(X,Y'') \ar[r]& \Hom_f(X,Y'[1])\ar[r] &
\Hom_f(X,Y[1]).
}$}
\]

Les\pageoriginale premières et dernières lignes sont exactes en vertu
de \hyperref[XVII.prop2.2.4]{2.2.4}. En vertu du lemme des cinq, l'application induite par $s$ de
$\Hom(X,X'')$ dans $\Hom(X,f^{\ast}Y'')$ est toujours bijective et $s$
est donc un isomorphisme.
\end{itemize}

On a évidemment des résultats duaux pour les catégories cofibrées.

\subsection{Formule triviale de dualité}\etiquette[2.3]{XVII.subsec2.3}

Soit $\Aa$ une catégorie sur une catégorie base~$\Bb$,
et soit $S$ un ensemble de flèches de $\Aa$, chacune se
projetant sur une identité de $\Bb$. La catégorie
$\Aa(S^{-1})$ déduite de $\Aa$ par calcul de fractions
est une catégorie sur $\Bb$ (vue la propriété universelle).

Soit $\Bb'\rightarrow \Bb$ un foncteur; on désigne par
$\Aa'$ le produit fibré
$\fprod{\Aa}{\Bb'}{\Bb}$ et par $S'$ l'image
réciproque de $S$ dans $\Aa'$. Le foncteur composé
$\Aa'\rightarrow \Aa\rightarrow \Aa(S^{-1})$
rend inversible les flèches de $S'$, donc se factorise par
$\Aa'({S'}^{-1})$. On trouve ainsi un $\Bb'$-foncteur
canonique
$$\etiquette[2.3.1]{XVII.eq2.3.1}
\Aa'(S^{-1})\longrightarrow
\fprod{\Bb'}{\Aa(\sisi{S}{{S'}}^{-1})}{\Bb}{\quoi}.\leqno(2.3.1)
$$

\begin{enonce*}{Lemme 2.3.2}\etiquette[2.3.2]{XVII.lem2.3.2}
Si $\Bb'$ est une sous-catégorie de $\Bb$, et si
chaque fois qu'une flèche composée $h=f\circ g$ de $\Bb$
appartient à $\Bb'$, $f$ et $g$ appartiennent à
$\Bb'$, alors le foncteur \eqhyperref[XVII.eq2.3.1]{2.3.1} est un isomorphisme.
\end{enonce*}

Notons\pageoriginale que \eqhyperref[XVII.eq2.3.1]{2.3.1} induit toujours une bijection sur les
objets, et que, dans le cas présent,
$\fprod{\Bb'}{\Aa(S^{-1})}{\Bb}$ est une
sous-catégorie de $\Aa(S^{-1})$.

Rappelons la construction « explicite »  de $\Aa(S^{-1})$
\cite[\sisi{Chap.}{\chap}~I]{SGA4-17:GZ}. Etant donné deux objets $X$ et $Y$ de $\Aa$, on
considère les « composés formels »  $f_1\ldots f_n$ où chaque
$f_i$ est soit une flèche de $\Aa$, soit l'inverse d'une
flèche de $S$. L'ensemble $\Hom_{\Aa(S^{-1})}(X,Y)$ sera le
quotient de l'ensemble de ces « composés formels »  par la
relation d'équivalence engendrée par les relations
\begin{alignat*}{2}
& f_1\ldots uv\ldots f_n\equiv f_1\ldots (u\circ v)\ldots f_n &\quad &
\mbox{\rm pour } u,v\;\mbox{\rm dans } \Fl(\Aa)\\
& f_1\ldots ss^{-1}\ldots f_n\equiv f_1\ldots f_n && \mbox{\rm pour }
s\in S\\
& f_1\ldots s^{-1}s\ldots f_n\equiv f_1\ldots f_n && \mbox{\rm pour }
s\in S{\quoi}.
\end{alignat*}

Si on applique ces constructions à $\Aa$ et à la
sous-catégorie $\Aa'$, on vérifie que $\Aa'(S^{-1})$
est identique à la sous-catégorie de $\Aa(S^{-1})$ image
réciproque de $\Bb'$, du fait que si un « composé formel » $f_1\ldots f_n$ a pour composé dans $\Bb$ une flèche de
$\Bb'$, l'image de chaque $f_i$ sera aussi une flèche de
$\Bb'$.

Supposons maintenant que $\Aa$ soit une catégorie triangulée
sur $\Bb$, et que les ensembles $S_x=S\cap
\Fl(\Aa_x)$ pour $x\in \Ob (\Bb)$ soient des
systèmes multiplicatifs saturés (\cite{SGA4-17:CD} \sisi{Ch.}{\chap}~I \S~2\sisi{. 1}{ nᵒ 1}).

\begin{enonce*}{Définition 2.3.3}\etiquette[2.3.3]{XVII.def2.3.3}
Sous les hypothèses précédentes la catégorie $\Aa/\Bb$
sera dite presque dérivable (\resp dérivable, \resp codérivable)
relativement à $S$ si, quel que soit le foncteur
$F:\Bb'\rightarrow \Bb$, le foncteur \eqhyperref[XVII.eq2.3.1]{2.3.1} est un
isomorphisme (\resp et si $\Aa(S^{-1})$ est fibrée, \resp
cofibrée sur $\Bb$).
\end{enonce*}

On\pageoriginale désignera par $I$ la catégorie du diagramme suivant
$$\etiquette[2.3.4]{XVII.eq2.3.4}
I:\underset{0}{\times} \xrightarrow{F} \underset{1}{\times}{\quoi}.\leqno{(2.3.4)}
$$

D'après \hyperref[XVII.prop2.2.13]{2.2.13}, il revient au même de se donner une catégorie
triangulée $\Aa$ sur $I$, fibrée (\resp cofibrée) sur
$\sisi{\Aa}{I}$, ou de se donner ses fibres $\Aa_0$ et
$\Aa_1$ et le foncteur triangulé image réciproque
$F^{\ast}:\Aa_1\rightarrow \Aa_0$ (\resp et le
foncteur triangulé image directe $F_{\ast}:\Aa_0\rightarrow
\Aa_1$). Si $A_0\in \Ob \Aa_0$ et $A_1\in \Ob
\Aa_1$, on a \eqhyperref[XVII.eq2.2.12.1]{2.2.12.1}
\[
\Hom_F(A_0,A_1)=\Hom(A_0,F^{\ast}A_1)
\]
(\resp $\Hom_F(A_0,A_1)=\Hom(F_{\ast}A_0,A_1)$).

\begin{enonce*}{Proposition 2.3.5}\etiquette[2.3.5]{XVII.prop2.3.5}
Soi\sisi{}{en}t $\Aa$ une catégorie triangulée sur $I$ et $S$ comme en
(\hyperref[XVII.def2.3.3]{2.3.3}).
\begin{enumerate}
\renewcommand{\theenumi}{\roman{enumi}}
\renewcommand{\labelenumi}{(\theenumi)}
\item $\Aa$ est presque dérivable relativement à $S$, et
$\Aa(S^{-1})$ est triangulée sur $I$.

\item Si $\Aa$ est fibrée (\resp cofibrée), le foncteur
$F^{\ast}$ (\resp $F_{\ast}$) est dérivable à droite (\resp à gauche)
relativement à $S_0$ et $S_1$ (\hyperref[XVII.def1.2.1]{1.2.1}) si et seulement si
$\Aa$ est dérivable (\resp codérivable) sur $I$.
\end{enumerate}
\end{enonce*}

En vertu de \hyperref[XVII.lem2.3.2]{2.3.2}, les catégories fibres de
$\Aa(S^{-1})$ sont les catégories $\Aa_0(S^{-1}_0)$ et
$\Aa_1(S^{-1}_1)$, où $S_i=S\cap \Fl(\Aa_i)$ $(i=0,1)$. Avec les
notations du \S \hyperref[XVII.sec1]{1} ${\rm n^{\circ}}$ \hyperref[XVII.sec2]{2}, les
foncteurs $A_0\rightarrow A^{-}_0$ et $A_1\rightarrow A^{+}_1$
identifient $\Aa_0(S^{-1}_0)$ et $\Aa_1(S^{-1}_1)$ à des catégories de
\ppro\ et \iind-objets de $\Aa_0$ et $\Aa_1$. Cela permet de définir une
catégorie sur $I$, de fibres $\Aa_0(S^{-1}_0)$ et $\Aa_1(S^{-1})$, en
posant, pour $A_0\in \Aa_0$ et $A_1\in \Aa_1$:
\[
\Hom_F(A_0,A_1)=\Hom_F(A^{-}_0,A^{+}_1)=\varinjlim_{A'_0{\xrightarrow[s]{}}A_{\sisi{p}{0}}}
\varinjlim_{A_1{\xrightarrow[t]{}}{A'_1}}\Hom(A'_0,A'_1){\quoi},
\]
où\pageoriginale $s$ et $t$ parcourent les éléments de $S_0$ et $S_1$ de but (de
source) $A_0$ ($A_1$). La catégorie $\Aa$ s'envoie dans cette
catégorie, et on vérifie aussitôt qu'elle vérifie la propriété
universelle de $\Aa(S^{-1})$, sur $I$ et après tout changement
de base.

Pour vérifier (i), il suffit de prouver que cette catégorie, soit
$\Aa(S^{-1})$, est triangulée, et il suffit de vérifier la
condition \hyperref[XVII.def2.2.3]{2.2.3} (iii). Soient donc $(X_0,Y_0,Z_0)$ et $(X_1,Y_1,Z_1)$
des triangles distingués de $\Aa_0$ et $\Aa_1$,
donnant lieu aux triangles distingués $(X^{-}_0,Y^{-}_0,Z^{-}_0)$ et
$(X^{+}_1,Y^{+}_1,Z^{+}_1)$ de $\Aa_0(S^{-1}_0)$ et
$\Aa_1(S^{-1}_1)$.
$$\etiquette[2.3.5.1]{XVII.eq2.3.5.1}
\vcenter{
\xymatrix{
X^{-}_{0}  \ar[r]^{u} \ar[d]& X^{+}_1  \ar[d]\\
Y^{-}_{0}  \ar[r]^{v} & Y^{+}_1\rlap{.}\\
}}\leqno{(2.3.5.1)}
$$

En vertu du lemme \hyperref[XVII.lem1.2.2.1]{1.2.2.1}, les triangles $(X^{-}_0,Y^{-}_0,Z^{-}_0)$
et $(X^{+}_1,Y^{+}_1,Z^{+}_1)$ sont limites de triangles distingués;
il existe donc des morphismes de triangles distingués
\begin{equation*}
\begin{split}
(X^{-}_0,Y^{-}_0,Z^{-}_0) &\longrightarrow (X'_0, Y'_0,Z'_0)\\
(X'_1, Y'_1, Z'_1) &\longrightarrow (X^{+}_1, Y^{+}_1, Z^{+}_1){\quoi},
\end{split}\tag{2.3.5.2}
\end{equation*}\etiquette[2.3.5.2]{XVII.eq2.3.5.2}%L'enlèvement de ce commentaire introduira un espace blanc non désiré.
et\pageoriginale un diagramme
\[
\xymatrix{
X'_{0}  \ar[r]^{u'} \ar[d]& X'_1 \ar[d] \\
Y'_{0}  \ar[r]^{v'} & Y'_1\rlap{,}
}
\]
compatibles avec \eqhyperref[XVII.eq2.3.5.1]{2.3.5.1}. La catégorie $\Aa$ étant
triangulée sur $I$, ce diagramme peut se prolonger en un morphisme de
triangles
\[
(X'_0,Y'_0,Z'_0)\longrightarrow (X'_1, Y'_1, Z'_1){\quoi}.
\]
Composant ce dernier avec les morphismes \eqhyperref[XVII.eq2.3.5.2]{2.3.5.2}, on voit que le
diagramme \eqhyperref[XVII.eq2.3.5.1]{2.3.5.1} peut se prolonger en un morphisme de triangles, ce
qui établit \hyperref[XVII.def2.2.3]{2.2.3} (iii).

Vérifions (ii), supposant par exemple $\Aa$ fibrée sur $I$. Si
$A_1\in \Ob (\Aa_1(S^{-1}_1))$, une image réciproque de $A_1$
dans $\Aa(S^{-1})$ est un objet $B$ de
$\Aa_0(S^{-1}_0)$ tel que, dans $\Aa_0(S^{-1}_0)$, on
ait
\begin{align*}
\Hom_{\Aa_0(S^{-1}_0)}(A_0,B) & =\Hom_F(A^{-}_0,
A^{+}_1)=\Hom(A^{-}_0,
F^{\ast}(A^{+}_1))\\
&=\Hom_{\Aa_0(S^{-1}_0)}
(A_0,F^{\ast}(A^{+}_1)){\quoi}.
\end{align*}
Il existe un tel objet $B$ si et seulement si le \iInd-objet
$F^{\ast}(A^{+}_1)$ de $\Aa_0(S^{-1})$ est essentiellement
constant (\ie représentable), et $B$ représente alors cet \iind-objet,
qui n'est autre que $RF(A_1)$ (\hyperref[XVII.def1.2.1]{1.2.1}), d'où l'assertion.

\begin{enonce*}{Corollaire 2.3.6}\etiquette[2.3.6]{XVII.cor2.3.6}
Avec les notations de \hyperref[XVII.def2.3.3]{2.3.3}, si $\Aa$ est presque dérivable
sur $\Bb$, alors $\Aa(S^{-1})$ est une catégorie
triangulée sur $\Bb$.
\end{enonce*}


Tout d'abord, la catégorie fibre $\Aa(S^{-1})_x$ peut être
identifiée à $\Aa_x(S^{-1}_x)$: on applique l'hypothèse au
changement de base $\{x\}\rightarrow \Bb$. Cette\pageoriginale
catégorie est triangulée, munie d'un ensemble $\Delta_x$ de triangles,
ce qui donne un sens à l'énoncé précédent. Pour vérifier la condition
2.2.3 (iii), il suffit de le faire après tout changement de base
$I\rightarrow \Bb$, et l'assertion résulte de \hyperref[XVII.prop2.3.5]{2.3.5} (i).

Il résulte aussitôt de la démonstration de \hyperref[XVII.prop2.3.5]{2.3.5} (ii) que:

\begin{enonce*}{Théorème 2.3.7}\etiquette[2.3.7]{XVII.thm2.3.7}
\emph{(formule triviale de dualité).}

Soit $\Aa$ une catégorie triangulée fibrée et cofibrée sur
$I$, et désignons par $F_{\ast}$ et $F^{\ast}$ les foncteurs image
directe et image réciproque. Soit $S=S_0\cup S_1$ un système
multiplicatif comme en \hyperref[XVII.def2.3.3]{2.3.3}. On suppose que $LF_{\ast}$ est défini
en $A_0\in \Ob \Aa(S^{-1}_0)$ et que $Rf^{\ast}$ est défini en
$A_1\in \Ob\Aa_1(S^{-1}_1)$. On a alors
\[
\Hom_{\Aa_0(S^{-1}_0)}(A_0,RF^{\ast}A_1)=\Hom_{\Aa\sisi{_{(S^{-1})}}{(S^{-1})}}(A_0,A_1)=\Hom_{\Aa_1(S^{-1}_1)}(LF_{\ast}A_0,A_1){\quoi}.
\]
\end{enonce*}

La position des symboles $R$, $L$, ${}^\ast$ et ${}_{\ast}$ dans \hyperref[XVII.thm2.3.7]{2.3.7}
n'est aberrante qu'en apparence; en effet, dans la catégorie fibrée et
cofibrée 0.2 des faisceaux sur des sites variables, les catégories
fibres sont les catégories opposées des catégories usuelles de
faisceaux.

\subsection{Catégories fibrées en catégories dérivées}\etiquette[2.4]{XVII.subsec2.4}

\setcounter{subsubsection}{-1}
\subsubsection{}\etiquette[2.4.0]{XVII.subsubsec2.4.0}
On désigne dans ce ${\rm n^{\circ}}$ par $\Aa$ une catégorie
triangulée cofibrée sur $\Bb$ et par $S$ un système
multiplicatif de morphismes de $\Aa$, tous au-dessus d'une
identité, et tels que $S_x=S\cap \Fl(\Aa_x)$ soit un système
multiplicatif saturé de $\Aa_x$ \cite[\sisi{}{I }1.2.1]{SGA4-17:CD} pour tout
$x\in \Ob (\Bb)$.

On\pageoriginale recherche un critère pour que $\Aa$ soit
codérivable relativement à $S$ (\hyperref[XVII.def2.3.3]{2.3.3}), ce qui implique que
$\Aa(S^{-1})$ soit triangulée sur $\Bb$ lorsqu'on la
munit de l'ensemble des triangles isomorphes à l'image d'un triangle
distingué de $\Aa$ (\hyperref[XVII.prop2.3.5]{2.3.5}).

\begin{enonce*}{Proposition 2.4.1}\etiquette[2.4.1]{XVII.prop2.4.1}
Soient $\Aa$, $\Bb$ et $S$ comme ci-dessus. Pour que
$\Aa$ soit codérivable relativement à $S$, il suffit que la
condition suivante soit vérifiée:
\begin{itemize}
\item[$(\ast)$] $\forall x \in \Ob\Bb$, $\forall A\in
\Ob(\Aa_x)$, $\exists s:A'\rightarrow A$ dans $S$\sisi{}{,} $\forall g$,
$f:x\xrightarrow{f}y\xrightarrow{g}\sisi{Z}{z}$, $f_{\ast}A'$ est déployé à
gauche pour le foncteur $g_{\ast}$.
\end{itemize}
\end{enonce*}

Désignons par $\Aa^{1}$ la sous-catégorie pleine de
$\Aa$ formée des objets dont toutes les images directes sont
déployés relativement à tous les foncteurs image directes. Alors \hyperref[XVII.prop1.2.2]{1.2.2}
montre que $\Aa^{1}$ est une sous-catégorie triangulée de
$\Aa$, évidemment cofibrée. On pose $S^{1}=S\cap \text{Fl}
(\Aa^{1})$.

\begin{enonce*}{Lemme 2.4.1.1}\etiquette[2.4.1.1]{XVII.lem2.4.1.1}
Supposons que $S$ soit stable par images directes. Alors,
$\Aa/\Bb$ est codérivable et $\Aa(S^{-1})$
peut se calculer par un calcul de fractions à droite.
\end{enonce*}

Vérifions la condition c) de \cite[I 2.2]{SGA4-17:GZ}, si possible moins triviale
que les autres. On se donne $f$ et $s$ $(s\in S)$ et on doit trouver
$$\etiquette[2.4.1.2]{XVII.eq2.4.1.2}
\vcenter{
\xymatrix{
A'        \ar@{-->}[r]      & B'          \\
A  \ar[u]^{s}   \sisi{\ar[r]}{\ar[r]_f}          & B\rlap{,}  \ar@{-->}[u]^{t}
}}\leqno{(2.4.1.2)}
$$
un\pageoriginale diagramme commutatif \eqhyperref[XVII.eq2.4.1.2]{2.4.1.2} avec $t\in S$. Soit
$F$ l'image de $f$ dans $\Bb$. Par hypothèse $F_{\ast}(s)$ est
encore dans $S$, de sorte que, d'après \textsc{Verdier}, il existe un diagramme
commutatif
\[
\xymatrix{
F_{\ast}(A')  \ar[r] & B'\\
F_{\ast}(A)    \ar[r]_{f}\sisi{\ar[u]_{s}}{\ar[u]^{F_\ast(s)}}   & B\ar[u]_{t}
}
\]
donc aussi un diagramme \eqhyperref[XVII.eq2.4.1.2]{2.4.1.2}.

Du fait que $\Aa(S^{-1})$ se calcule par calcul de fractions à
droite, il est trivial que sa formation commute au changement de base
$c:\Bb'\rightarrow\Bb$.

Ce lemme s'applique en particulier à $\Aa^{1}$ et $S^{1}$; le
lemme suivant implique donc \hyperref[XVII.prop2.4.1]{2.4.1}:

\begin{enonce*}{Lemme 2.4.1.3}\etiquette[2.4.1.3]{XVII.lem2.4.1.3}
Sur $\Bb$ et après tout changement de base
$c:\Bb'\rightarrow \Bb$, le foncteur canonique
$\Phi:\Aa\sisi{_1}{^1}(\sisi{{S^1}}{(S^1)}^{-1})\rightarrow \Aa(S^{-1})$ est
une équivalence de catégorie.
\end{enonce*}

Construisons un foncteur quasi-inverse au foncteur $\Phi$. Quel que
soit $A\in \Ob \Aa$, choisissons $s_A\in S$ de but $A$ et de
source $A^1\in \Ob \Aa^{1}$ avec $s_A=\Id_{A}$ pour $A\in
\Ob \Aa^{1}$. Soit $f$ une flèche de $\Aa$, de but
$A$, de source $B$ et d'image $F$ dans $\Bb$. Puisque
$s_{B}\in S$ et que $F_{\ast}A^{1}\in \Ob \Aa^{1}$, on peut
trouver $B\sisi{^{1}}{'}$, \sisi{$v$ et $s\in S^{1}$}{$s\in S^1$ et $u$} rendant commutatif le diagramme
suivant
$$\etiquette[2.4.1.4]{XVII.eq2.4.1.4}
\vcenter{
\xymatrix{
A      \ar[r]^{f} &  B      \ar[r]_{s_B} & B\sisi{_1}{^1}       \\
A\sisi{_1}{^1} \ar[u]_{s_A}\ar[r]^{a} & F_{\ast}A\sisi{_1}{^1}  \ar[u]& B'\rlap{.} \ar[l]_{s} \ar[u]_u
}}\leqno{(2.4.1.4)}
$$

La\pageoriginale flèche $f^{1}=us^{-1} a$ de
$\Aa^{1}(\sisi{{S^1}}{(S^1)}^{-1})$ ne dépend alors que de $f$.

Soient $f$ et $g$ deux flèches composables, d'images $F$ et $G$ dans
$\Bb$. Pour prouver que $(g\circ f)^{1}=g^{1}\circ f^{1}$,
contemplons le diagramme suivant, où figurent des diagrammes \eqhyperref[XVII.eq2.4.1.4]{2.4.1.4}
pour $f$ et $g$:
\[
\xymatrix@C=2em{
                      & A \ar[rr]^{f}&      & B\ar[rrr]^{g}     & &
                      & C  &       \\
 A^{1} \ar[ru] \ar[rr]^{a}&          &F_{\ast}A^{1}\ar[ru] \ar[rrr]^(.6){a'}& & & \sisi{GF}{(GF)}_{\ast}A^1\ar[ru] &
              & &C^{1}\ar[llu]\\
& & & B^1 \ar[rrr]^(.4){b}|(.58)\hole\ar[uu]|\hole & & & G_{\ast}B^1\ar[uu] &     \\
&  & B\sisi{}{'}\ar[uu]^{s}\ar[rrr]^{b'}\ar[ru]_{u}       &   & &G_{\ast}B'\ar[ru]\ar[uu]^(.3){G_{\ast}s}    &  &C''\ar[ll]^{r}\ar[r]_{w} & C'\ar[llu]_t\ar[uu]_{v}\rlap{.}
}
\]

Puisque $t\in S$, il existe $C''$, $r\in S^{1}$ et $w$ qui le rende\sisi{}{nt}
commutatif. Par définition,\break $g^{1}\circ f^{1}=vt^{-1} b\, u\, s^{-1}a$
et $(g\circ f)^{1}=vw(G_{\ast}(s)r)^{-1}a'a$, et on conclut par la
commutativité du diagramme.

On a donc défini un foncteur de $\Aa$ dans
$\Aa^{1}(\sisi{{S^1}}{(S^1)}^{-1})$; il se factorise en
$\Psi:\Aa(S^{-1})\rightarrow \Aa^{1}(\sisi{{S^1}}{(S^1)}^{-1})$, la propriété
universelle de $\Aa(S^{-1})$. Le composé $\Phi\Psi$ est
isomorphe à l'identité car son composé avec le foncteur
$J:\Aa\rightarrow \Aa(S^{-1})$ est isomorphe à $J$. On
vérifie de même que $\Psi\Phi$ est isomorphe à l'identité.

\begin{enonce*}{Corollaire 2.4.2}\etiquette[2.4.2]{XVII.cor2.4.2}
Avec les notations de \hyperref[XVII.prop2.4.1]{2.4.1}, $\Aa$ est codérivable
relativement à $S$ dès que, pour tout ensemble fini $B$ de flèches de
$\Bb$, la condition $(\ast)$ est vérifiée lorsqu'on se limite
à ne considérer que les images directes par les flèches de $B$.
\end{enonce*}

Quelle\pageoriginale que soit la catégorie $\Bb'$ d'ensemble
de flèches fini, la prop. \hyperref[XVII.prop2.4.1]{2.4.1} sera applicable après tout changement
de catégorie de base $\Bb'\rightarrow \Bb$. On prendra
pour $\Bb'$ les catégories des diagrammes suivants:
\[
  \begin{gathered}
    *\\
    (0)
  \end{gathered}\qquad
  \begin{gathered}
    *\to*\\
    (1)
  \end{gathered}\qquad
  \begin{gathered}
    *\to*\to*\\
    (2)
  \end{gathered}\qquad
  \begin{gathered}
    *\to*\to*\to*\\
    (3)
  \end{gathered}
\]
Considérons les catégories $\Aa_x(S^{-1}_x)$ $(x\in \Ob
\Bb)$. Etant donné $f\in \Fl(\Bb)$ et $A_1$, $A_2$
dans ces catégories au-dessus de la source et du but de $f$, on
définit $\Hom_f(A_1,A_2)$ comme étant l'ensemble analogue défini après
le changement de catégorie base de type (1) défini par $f$. On définit
la composition des morphismes à l'aide de changements de catégorie
base de type (3). On construit ainsi une catégorie
$\Aa(S^{-1})^{\ast}$ sur $\Bb$, dont la formation est
compatible à tout changement de base; on vérifie qu'elle est solution
du problème universel qui définit $\Aa(S^{-1})$ en utilisant
les changements de base de type (0), (1) et (2).

Le résultat suivant est le dual de \hyperref[XVII.cor2.4.2]{2.4.2}.

\begin{enonce*}{Proposition 2.4.3}\etiquette[2.4.3]{XVII.prop2.4.3}
Soit $\Aa$ une catégorie triangulée fibrée sur $\Bb$
et $S$ comme en \hyperref[XVII.prop2.4.1]{2.4.1}. Pour que $\Aa$ soit dérivable
relativement à $S$, il suffit que pour tout ensemble fini $B$ de
flèches de $\Bb$, on ait
\begin{itemize}
\item[$(\ast')$] $\forall x \in
\Ob(\Bb)$, $\forall A\in\Ob (\Aa_x)$, $\exists
s:A\rightarrow A'$ dans $S$,
$\forall g,f:Z\xrightarrow{g}Y\xrightarrow{f}X$ dans $B$, $f^{\ast}A'$
est déployé à droite pour le foncteur $g^{\ast}$.
\end{itemize}
\end{enonce*}

\section{Recollement de catégories fibrées ou
  cofibrées}\etiquette[3]{XVII.sec3}\pageoriginale

\subsection{Introduction}\etiquette[3.1]{XVII.subsec3.1}

Soit $f:Y\rightarrow X$ un morphisme séparé de type fini d'un schéma
$Y$ dans un schéma $X$, ou une application continue d'un espace
localement compact dans un espace localement compact. On sait alors
définir un foncteur « image directe à support propre », noté
$f_{!}$, de la catégorie des faisceaux abéliens sur $Y$ dans la
catégorie des faisceaux abéliens sur $X$ (\cf \sisi{\hyperref[XVII.subsec5.3]{5.3}}{\hyperref[XVII.subsec6.1]{6.1}} dans le cas des
schémas, et \textsc{Verdier} \cite{SGA4-17:V3} pour le cas des espaces localement compacts
séparés). Lorsque $Y$ et $X$ sont des schémas de type fini sur
$\CC$, ces deux définitions sont compatibles en un sens
évident. Malheureusement dans le cas des schémas, contrairement à ce
qui se passe dans le cas des espaces localement compacts, les
foncteurs dérivés du foncteur $f_{!}$ sont pathologique; on peut
toutefois définir un foncteur « image directe à support propre »
raisonnable, directement de la catégorie $D^{+}(Y)$ dans
$D^{+}(X)$. Pour le définir, on factorise le morphisme $f$ en une
immersion ouverte et un morphisme propre: $f=\bar{f}j$. Le foncteur
$Rf_{!}$ cherché se définit alors comme composé du « prolongement
par $0$ »  $j_{!}$ et du foncteur $R\bar{f}_{\ast}$, dérivé du
foncteur image directe par $\bar{f}$.

Ces foncteurs $Rf_{!}$, « image directe à support propre »  ne
sont utilisables que si on vérifie pour eux un « formalisme de
variance »  incluant, pour un composé $f=gh$, une formule
\[
Rf_{!}=Rg_{!} Rh_{!}
\]
qui s'exprime le mieux en terme de catégories cofibrées. On dispose
d'un tel formalisme séparément pour les « prolongements par
zéro » et pour les images\pageoriginale directes par un morphisme
propre, et le problème est de recoller ces formalismes. Pour éviter
d'être embouteillés plus tard par des calculs idiots, on se propose
dans ce \S de dresser la liste des vérifications élémentaires requises
pour mener à bien un tel recollement. On le fera dans un cadre
légèrement plus abstrait qu'il n'est indispensable, espérant que le
résultat pourra resservir.

La construction esquissée plus haut du foncteur $Rf_{!}$ exige que le
morphisme $f$ puisse se factoriser en une immersion ouverte et un
morphisme propre. Sur le modèle de \textsc{Hartshorne} \cite[III 8 p.~189]{SGA4-17:RD}, on
introduit au nᵒ \hyperref[XVII.sec2]{2} la notion adéquate de morphisme
compactifiable.

\subsection{Morphismes compactifiables}\etiquette[3.2]{XVII.subsec3.2}

\begin{enonce*}{Définition 3.2.1}\etiquette[3.2.1]{XVII.def3.2.1}
Un morphisme $f$ d'un $S$-schéma $X$ dans un $S$-schéma $Y$
quasi-compact, quasi-séparé, est dit \emph{$S$-compactifiable} s'il existe un
$S$-schéma $P$ propre sur $S$ et une factorisation de $f$ en un
morphisme quasi-fini séparé $j$ de $X$ dans $\fprod{P}{Y}{S}$ suivi de
la projection de $\fprod{P}{Y}{S}$ dans $Y$
\[
\xymatrix{
                                & &  &\fprod{P}{Y}{S} \ar[rd]\ar[ld]& \\
X \ar[rr]^(.7){f} \ar[rrru]^{j}\ar[rrrd]& &Y \ar[rd]&                & P\ar[ld]\\
                                & &  &       S        &.
}
\]
\end{enonce*}

Lorsque $S=Y$, on parlera simplement de morphisme compactifiable.

\begin{enonce*}{Proposition 3.2.2}\etiquette[3.2.2]{XVII.prop3.2.2}
\begin{enumerate}
\renewcommand{\theenumi}{\roman{enumi}}
\renewcommand{\labelenumi}{(\theenumi)}
\item Un\pageoriginale morphisme $S$-compactifiable est séparé de type
fini.

\item Si $Y$ est quasi-compact, quasi-séparé, un morphisme quasi-fini
séparé de but $Y$ est \sisi{}{$S$-}compactifiable.

\item Le composé de deux morphismes $S$-compactifiables est encore
$S$-compactifiable.

\item Soit un diagramme de $S$-schémas quasi-compacts quasi-séparés
\[
\xymatrix{
\fprod{X}{Y'}{Y} \ar[d]_{f'} \ar[r] & X \ar[d]^{f}      \\
Y' \ar[r]& Y\rlap{.}
}
\]
Si $f$ est $S$-compactifiable, alors $f'$ est $S$-compactifiable.

\item Soit $g$ un morphisme $S$-compactifiable;\nde{il suffit de supposer que $g$ est un morphisme quasi-compact quasi-séparé de $S$-schémas.} pour qu'un morphisme
composé $g\circ f$ soit $S$-compactifiable, il faut que $\sisi{h}{f}$ le soit.
\end{enumerate}
\end{enonce*}

Les assertions (i) (ii) et (iv) sont triviales. Pour prouver (iii) et
(v), considérons un composé $f=gh$, $g$ étant \sisi{}{$S$-}compactifiable.

Supposons $gh$ \sisi{}{$S$-}compactifiable, d'où un diagramme
\[
\xymatrix{
                     & \fprod{P}{Y}{S} \ar[d] \ar[r]^{v} & \fprod{P}{Z}{S}\ar[d] \\  X  \ar[ru]^{u} \ar[r]^{h}   & Y                \ar[r]^{g} &    Z\rlap{.}
}
\]
Le composé $vu$ est quasi-fini; $u$ est donc quasi-fini et (v) est
prouvé.

Si $h$ est compactifiable, on dispose d'un diagramme
\[
\xymatrix{
  &&                & &\sisi{Q\times P \times Z}{\fprod{\fprod{Q}{P}{S}}{Z}{S}} \ar[d]&\\
  &&\fprod{Q}{Y}{S} \ar[d]\ar[rru]\ar[rdd]& & \fprod{P}{Z}{S} \ar[d]\ar[rdd] & \\
X \ar[rrrdd]\ar[rru] \ar[rr]^{h}&& Y \sisi{\ar[rru]}{\ar[rru]|(.2)\hole} \ar[rdd]\sisi{\ar[rr]^{g}}{\ar[rr]^g|(.25)\hole}            & & Z\ar[ldd]                & \\
  &&                &Q\ar[d]&                  &P\ar[lld] \\
  &&                &S&                   &,
}
\]
et\pageoriginale $f=fh$ se factorise par un morphisme quasi-fini
séparé dans $\fprod{(\fprod{Q}{P}{S})}{Z}{S}$.

Nagata affirme dans \sisi{[5]}{\cite{SGA4-17:Na}} que tout morphisme de type fini entre schémas
noethériens intègres séparés est compactifiable.\nde{plus généralement, tout morphisme de schémas séparé de type fini de but quasi-compact quasi-séparé est compactifiable. Voir \nderef{XVII.ndeConrad}.} Le rédacteur avoue ne
pas avoir compris la démonstration.

Fixons maintenant un schéma $S$ et désignons par $(S)$ la
sous-catégorie de la catégorie des $S$-schémas dont les objets sont
les $S$-schémas quasi-compacts, quasi-séparés et dont les flèches sont
les $S$-morphismes $S$-compactifiables.\nde{Grâce au théorème de compactification, on peut également travailler dans la sous-catégorie de la catégorie des schémas (absolus) dont les objets sont les schémas quasi-compacts quasi-séparés et dont les flèches sont les morphismes séparés de type fini (\ie compactifiables). Notons que la stabilité des morphismes compactifiables par composition est équivalente à la forme générale du théorème de compactification.}

\begin{enonce*}{Proposition 3.2.3}\etiquette[3.2.3]{XVII.prop3.2.3}
Dans la catégorie $(S)$:
\begin{enumerate}
\renewcommand{\theenumi}{\roman{enumi}}
\renewcommand{\labelenumi}{(\theenumi)}
\item Les immersions ouvertes (\resp les morphismes propres)
définissent une sous-catégorie $(S,i)$ (\resp $(S,p)$) de $(S)$ ayant
mêmes objets que $(S)$.

\item Les produits fibrés existent dans $(S,p)$ et sont des produits
fibrés dans $(S)$.

\item Tout morphisme $f$ est le composé $f=pj$ d'un morphisme propre
$p$ et d'une immersion ouverte.
\end{enumerate}
\end{enonce*}

En\pageoriginale vertu de \hyperref[XVII.prop3.2.2]{3.2.2} (iii) (iv) et (v), les produits fibrés
existent dans $(S)$ et sont des produits fibrés dans la catégorie de
tous les schémas. Puisque le produit fibré $X=\fprod{X_1}{X_2}{Y}$ de
deux schémas propres sur $Y$ est encore propre sur $Y$, et qu'un
schéma est propre sur $X$ si et seulement si il est propre sur $X_1$
et $X_2$, la condition (ii) est vérifiée.

Si $f$ est un morphisme de $(S)$, il admet par hypothèse, dans $(S)$,
une factorisation $f=pj'$, où $p$ est propre et $j\sisi{}{'}$ quasi-fini
séparé. En vertu de \EGA IV 18.12.13, $j'$ admet une factorisation
$j'=qj$, où $q$ est fini et $j$ une immersion ouverte; de plus, en
vertu de \hyperref[XVII.prop3.2.2]{3.2.2} (ii), $q$ et $j$ sont dans $(S)$. On a alors $f=(pq)j$,
ce qui prouve (iii). L'assertion (i) résulte de \hyperref[XVII.prop3.2.2]{3.2.2} (iii).

\setcounter{subsubsection}{3}
\subsubsection{}\etiquette[3.2.4]{XVII.subsubsec3.2.4}
Dans la suite de ce \S, on suppose donnée une catégorie $(S)$, munie
de deux sous-catégories $(S,i)$ et $(S,p)$, telles que
\begin{enumerate}
\renewcommand{\theenumi}{\roman{enumi}}
\renewcommand{\labelenumi}{(\theenumi)}
\item $\Ob(S)=\Ob(S,i)=\Ob(S,p)$;

\item les produits fibrés existent dans $(S,p)$ et sont des produits
fibrés dans $(S)$;

\item $\forall f\in \Fl(S)$, $\exists p\in \Fl(S,p)$, $\exists
j\in \Fl(S,i)$, $f=pj$.
\end{enumerate}

Les flèches de $(S,i)$ (\resp de $(S,p)$) s'appelleront les immersions
ouvertes (\resp les morphismes propres).

\begin{enonce*}{Définition 3.2.5}\etiquette[3.2.5]{XVII.def3.2.5}
\begin{enumerate}
\renewcommand{\theenumi}{\roman{enumi}}
\renewcommand{\labelenumi}{(\theenumi)}
\item Une \emph{compactification} d'un morphisme $f$ de $(S)$ (\resp
d'un couple $(f,g)$, \resp d'un triple $(f,g,h)$ de morphismes
composables) est un diagramme commutatif \eqhyperref[XVII.eq3.2.5.1]{3.2.5.1} (\resp
\eqhyperref[XVII.eq3.2.5.2]{3.2.5.2}, \resp \eqhyperref[XVII.eq3.2.5.3]{3.2.5.3}) dont les flèches horizontales sont des
immersions ouvertes et dont les flèches obliques sont propres.

\item \emph{Un\pageoriginale morphisme de compactifications} est un
  morphisme de
diagrammes \eqhyperref[XVII.eq3.2.5.1]{3.2.5.1} (\resp\ldots) qui soit l'identité sur la 1\sisi{ère}{\textsuperscript{ière}}
ligne verticale et propre en tous les sommets
$$
\etiquette[3.2.5.1]{XVII.eq3.2.5.1}
\etiquette[3.2.5.2]{XVII.eq3.2.5.2}
\etiquette[3.2.5.3]{XVII.eq3.2.5.3}
\begin{gathered}[b]
\xymatrix{
Y        \ar[d]_{f} \ar@{^{(}->}[r]^{i}                & \bar{Y} \ar[ld]^{\bar{f}}\\
X      & &}\\
(3.2.5.1)
\end{gathered}\qquad
\begin{gathered}[b]
\xymatrix{
Z        \ar@{^{(}->}[r]^{j} \sisi{\ar[d]}{\ar[d]_g}& \bar{Z}\ar@{^{(}->}[r]^{j'}\sisi{\ar@{-}[ld]^{\bar{g}}}{\ar[ld]^{\bar{g}}}
&\bar{\bar{Z}} \ar[ld]^{\bar{\bar{g}}}\\
Y \ar[d]_{f}\ar@{^{(}->}[r]^{i} & {\bar{Y}} \ar[ld]^{\bar{f}}&\\
X      & &}\\
(3.2.5.2)
\end{gathered}\qquad
\begin{gathered}[b]
\xymatrix@=3em{
\bigdot \ar@{^{(}->}[r]_(.4){\sisi{h}{k}} \ar[d]_{h}& \bigdot \ar@{^{(}->}[r]_(.4){\sisi{h}{k}'}\ar[ld]^(.4){\bar{h}} & \bigdot
\ar@{^{(}->}[r]_(.4){\sisi{h}{k}''} \ar[ld]^(.4){\bar{\bar{h}}}&\bigdot \ar[ld]^(.4){\bar{\bar{\bar{h}}}}\\
\bigdot \ar@{^{(}->}[r]_(.4){j}\ar[d]_{g} & \bigdot \ar@{^{(}->}[r]_(.4){j'}
\ar[ld]^(.4){\bar{g}}& \bigdot \ar[ld]^(.4){\bar{\bar{g}}}\\
\bigdot
\ar@{^{(}->}[r]_(.4){i} \ar[d]_{f}&\bigdot \ar[ld]^(.4){\bar{f}}\\
\bigdot}\\
(3.2.5.3)
\end{gathered}
$$
\end{enumerate}
\end{enonce*}

\begin{enonce*}{Proposition 3.2.6}\etiquette[3.2.6]{XVII.prop3.2.6}
\begin{enumerate}
\renewcommand{\theenumi}{\roman{enumi}}
\renewcommand{\labelenumi}{(\theenumi)}
\item Tout morphisme (\resp tout couple $(f,g)$, \resp tout triple
$(f,g,h)$, de morphismes composables) admet une compactification.

\item La catégorie des compactifications de $f$ est filtrante à
gauche.\nde{« filtrante à gauche » signifie « cofiltrante » (\hyperrefext[I][I.def2.7]{2.7}). Pour les axiomes $L_1$ et $L_2$ dans la démonstration de \hyperref[XVII.prop3.2.6]{3.2.6}, \cf \hyperrefext[V][V.par8.1.1.1]{8.1.1.1}.}

\item Si des diagrammes $\sisi{(3.2.4.2)}{\eqhyperref[XVII.eq3.2.5.2]{3.2.5.2}}_i$ $(i=1,2)$ sont des
compactifications de $(f,g)$, il existe une compactification
$\sisi{(3.2.4.2)}{\eqhyperref[XVII.eq3.2.5.2]{3.2.5.2}}_3$ de $(f,g)$ et des morphismes de compactifications de
$\sisi{(3.2.4.2)}{\eqhyperref[XVII.eq3.2.5.2]{3.2.5.2}}_3$ dans les compactifications $\sisi{(3.2.4.2)}{\eqhyperref[XVII.eq3.2.5.2]{3.2.5.2}}_i$.
\end{enumerate}
\end{enonce*}

\begin{enumerate}
\renewcommand{\theenumi}{\roman{enumi}}
\renewcommand{\labelenumi}{(\theenumi)}
\item Le cas d'un morphisme n'est autre que \hyperref[XVII.subsubsec3.2.4]{3.2.4}
(iii). Traitons le cas d'un triple; avec les notations de \eqhyperref[XVII.eq3.2.5.3]{3.2.5.3},
il suffit de construire successivement des compactifications de $f$,
$g$ et $h$, de $i\bar{g}$ et $j\bar{h}$, et enfin de
$j'\bar{\bar{h}}$.

\item Soient\pageoriginale $\eqhyperref[XVII.eq3.2.5.1]{3.2.5.1}_1$ et $\eqhyperref[XVII.eq3.2.5.1]{3.2.5.1}_2$ deux
compactifications de $f$. La projection $\sisi{\rho}{p}$ de
$\fprod{\bar{Y}_1}{\bar{Y}_2}{S}$ dans $X$ est propre et $Y$ s'y
envoie par $(i_1,i_2)$, qui (\hyperref[XVII.subsubsec3.2.4]{3.2.4} (iii)) peut se compactifier en
$(i_1,i_2)=qj$. La compactification $f=(p\sisi{\cdot}{\circ} q)j$ domine
$\sisi{(3.2.4.1)}{\eqhyperref[XVII.eq3.2.5.1]{3.2.5.1}}_1$ et $\sisi{(3.2.4.1)}{\eqhyperref[XVII.eq3.2.5.1]{3.2.5.1}}_2$.

Si $r$ et $s$ sont deux morphismes de $\sisi{(3.2.4.1)}{\eqhyperref[XVII.eq3.2.5.1]{3.2.5.1}}_1$ dans
$\sisi{(3.2.4.1)}{\eqhyperref[XVII.eq3.2.5.1]{3.2.5.1}}_2$, le noyau $K$ de la double flèche
$\bar{Y}_{1} \overset{r}{\underset{s}\rightrightarrows} \bar{Y}_{2}$ s'envoie par
un morphisme propre dans $\bar{Y}_1$ et $i_1$ se factorise par $K$:
soit $i_1=ki$. Si $i=pi_3$ est une compactification de $i$,
l'existence de la compactification $f=(\bar{f}_1kp)i_3$ montre que
l'axiome $L_2$ est satisfait.

\item Appliquons l'axiome $L_1$ des catégories filtrantes aux
compactifications de $f$ et $g$ déduites des $\sisi{(3.2.4.2)}{\eqhyperref[XVII.eq3.2.5.2]{3.2.5.2}}_i$
$(i=1,2)$. On trouve ainsi $\bar{Y}_3$ et $\bar{Z}_3$ donnant lieu à
des diagrammes commutatifs $\eqhyperref[XVII.eq3.2.5.1]{3.2.5.1}_i$ $(i=1,2)$
\[
\xymatrix{
&         &\bar{Z}_i \ar[d] \ar@{^{(}->}[r]& \bar{\bar{Z}}_i\ar[d]\\
\bar{Z}_3 \ar[rru] \ar[d]^{\bar{g}_3} &&Y   \ar@{^{(}->}[r]      & \bar{Y}_i\\
Y \ar@{=}[rru]\ar[r]_{i_3} & \bar{Y}_3 \ar[rru]&&\sisi{(3.2.5.1)_i}{},
}
\]
et il reste à trouver une compactification de $i_3\bar{g}_3$ par
$\bar{\bar{Z}}_3$, telle que $\bar{\bar{Z}}_3$ puisse s'envoyer par un
morphisme propre dans $\bar{\bar{Z}}_i$, de façon à donner lieu à
un cube commutatif $(i=1,2)$. Si le but d'un morphisme entre
compactifications de $i_3\bar{g}_3$ a cette propriété, la source l'a
encore et en vertu de (ii), il suffira de s'occuper séparément des cas
$i=1$ et $i=2$. Soit\pageoriginale $T$ le produit fibré de $\bar{Y}_3$
et $\bar{\bar{Z}}_i$ sur $\bar{Y}_i$: alors $T$ figure dans un cube
commutatif analogue à celui qu'on cherche, sauf que la flèche $k$ de
$\bar{Z}_3$ dans $T$ pourrait ne pas être une immersion
ouverte. Posant $k=pj$ ($p$ propre, $j$ immersion ouverte) et rempla\c
cant $T$ par la source $\bar{\bar{Z}}_3$ de $p$, on trouve le cube
cherché.
\end{enumerate}

\begin{enonce*}{Remarque 3.2.7}\etiquette[3.2.7]{XVII.rem3.2.7}
La proposition \hyperref[XVII.prop3.2.3]{3.2.3} aurait encore été valable si, dans la définition
\hyperref[XVII.def3.2.1]{3.2.1} des morphismes compactifiables, on avait omis d'exiger que $Y$
soit quasi-compact quasi-séparé et demandé en compensation que $j$
soit une immersion ouverte quasi-compacte et non seulement un
morphisme quasi-fini.
\end{enonce*}

\subsection{Recollement}\etiquette[3.3]{XVII.subsec3.3}

Rappelons que les hypothèses \hyperref[XVII.subsubsec3.2.4]{3.2.4} sont satisfaites.

\subsubsection{}\etiquette[3.3.1]{XVII.subsubsec3.3.1}
Supposons donnés:
\begin{itemize}
\item[$\alpha$)] pour chaque $X\in \Ob (S)$, une catégorie $F(X)$,

\item[$\beta$)] une $(S,i)$-catégorie cofibrée $F_i$ et une $(S,p)$
catégorie cofibrée $F_{p}$,

\item[$\gamma$)] pour chaque $X\in \Ob(S)$, des isomorphismes entre
$F(X)$, $(F_{i})_X$ et $(F_{p})_X$,

\item[$\delta$)] un scindage normalisé de $F_{i}$ et un scindage
normalisé de $F_{p}$.
\end{itemize}

En vertu de (\SGA 1 VI), il « revient au même » de se donner plutôt $\alpha$) et:
\begin{itemize}
\item[(i)] Pour tout morphisme propre $p:X\rightarrow Y$, un foncteur
$p_{\ast}:F(X)\rightarrow F(Y)$,

\item[(ii)] pour tout couple composable $(p,q)$ de morphismes propres,
un\pageoriginale isomorphisme de foncteur
$c_{p,q}:p_{\ast}q_{\ast}\leftrightarrow (pq)_{\ast}$,

\item[(i\sisi{}{'}) et (ii\sisi{)'}{')},] données analogues pour $(S,i)$,
\end{itemize}
ces données vérifiant les conditions suivantes:
\begin{itemize}
\item[(a)] si $p$ est une identité, alors $p_{\ast}$, $c_{p,q}$ et
$c_{q,p}$ sont des identités,

\item[(b)] pour tout triple composable de morphismes propres, on a
\[
c_{p,qr}\circ (p_{\ast}\ast c_{q,r})=c_{pq,r}\circ (c_{p,q}\ast
r_{\ast}){\quoi},
\]

\item[(a\sisi{)'}{')} et (b\sisi{)'}{')},] conditions analogues pour $(S,p)$.
\end{itemize}

Supposons \emph{donné} de plus,
\begin{itemize}
\item[(iii)] pour tout diagramme commutatif $D$
\[
\vcenter{
\xymatrix{
X \ar[r]^{j'} \ar[d]^{p'}& Y\ar[d]^{p} \\
Z \ar[r]_{j}             &           T
}}\quad j,j' \in \Fl(S,i)\text{; }
  p,p' \in \Fl(S,p),
\]
un isomorphisme de foncteur $d(D):p_{\ast}j'_{\ast}\leftrightarrow
j_{\ast}p'_{\ast}$.
\end{itemize}

\begin{enonce*}{Proposition 3.3.2}\etiquette[3.3.2]{XVII.prop3.3.2}
Si des données {\rm (i) (ii) ${\rm (i')}$ ${\rm (ii')}$ (iii)}
vérifient les conditions {\rm (a) (b) ${\rm (a')}$ ${\rm (b')}$} ainsi
que {\rm(c)} et ${\rm (c')}$ énoncées ci-dessous, alors, il existe une
et « essentiellement »  une seule catégorie $F$ cofibrée sur $S$,
munie d'isomorphismes $\fprod{F}{(S,i)}{S}=F_i$ et
$\fprod{F}{(S,p)}{S}=F_p$, compatibles à la donnée \hyperref[XVII.subsubsec3.3.1]{3.3.1} $\gamma$)
et tels que les isomorphismes $d(D)$ de {\rm(iii)} soient les
isomorphismes composés
$j_{\ast}p'_{\ast}=(jp')_{\ast}=(pj')_{\ast}=p_{\ast}j'_{\ast}$.
\end{enonce*}

\begin{itemize}
\item[(c)] pour tout diagramme commutatif du type suivant
\[
\xymatrix{
X_1 \ar@{^{(}->}[r]^{i} \ar[d]^{p_1}& X_2 \ar@{^{(}->}[r]^{i'}
\ar[d]^{p_2}& X_3 \ar[d]^{p_3}\\
Y_1 \ar@{^{(}->}[r]^{j}& Y_2 \ar@{^{(}->}[r]^{j'}& Y_3,\\
}
\]
dont\pageoriginale les flèches horizontales sont des immersions
ouvertes et dont les flèches verticales sont propres, le triangle
suivant d'isomorphismes $d$ de (iii) est commutatif:
\[
\xymatrix@R=1cm{
p_{3{\ast}} i'_{\ast} i_{\ast} \ar[rr] \ar@{-}[rd]&  & j'_{\ast}j_{\ast}p_{1{\ast}}\ar@{-}[ld]\\
                               &j'_{\ast}p_{2{\ast}}
                                i_{\ast}&.
}
\]

\item[(c')] Pour tout diagramme commutatif du type suivant
\[
\xymatrix{
X_1  \ar[r]^{p} \ar[d]^{i_1}& X_2  \ar[r]^{p'} \ar[d]^{i_2}& X_3 \ar[d]^{i_3}          \\
Y_1         \ar[r]^{q} & Y_2  \ar[r]^{q'} & Y_3\rlap{,}
}
\]
dont les flèches horizontales sont propres et dont les flèches
verticales sont des immersions ouvertes, le triangle suivant
d'isomorphismes $d$ de (iii) est commutatif:
\[
\xymatrix{
i_{3{\ast}} p'_{\ast} p_{\ast} \ar[rr] \ar@{-}[rd]&    &q'_{\ast} q_{\ast}
i_{1{\ast}} \\
&q'_{\ast} i_{2{\ast}} p_{\ast}  \ar@{-}[ru]&.
}
\]
\end{itemize}

\setcounter{subsubsection}{2}
\paragraph{Définition de $f_{\ast}$}\etiquette[3.3.2.1]{XVII.par3.3.2.1}\pageoriginale
Soit $f$ un morphisme de $Y$ dans $X$. Chaque compactification de $f$,
de diagramme \sisi{(\hyperref[XVII.def3.2.5]{3.2.5} (i))}{\eqhyperref[XVII.eq3.2.5.1]{3.2.5.1}}, définit un foncteur composé $\bar{f}_{\ast}i_{\ast}$ de $F(Y)$ dans $F(X)$. Un morphisme de compactifications
définit un diagramme commutatif
$$\etiquette[3.3.2.2]{XVII.eq3.3.2.2}
\vcenter{
\xymatrix{
Y  \ar@{^{(}->}[r]^{i_1} \ar@{=}[d]& \bar{Y}_1  \sisi{\ar[d]}{\ar[d]^p}\\
Y  \ar@{^{(}->}[r]^{i_2} \ar[d]        & \bar{Y}_2 \ar[ld]^{\bar{f}_2}\\
X                         &.
}}\leqno{(3.3.2.2)}
$$
La donnée (iii) nous fournit donc un isomorphisme
\[
\bar{f}_{2{\ast}} i_{2{\ast}}=\bar{f}_{2\ast} \circ
(i_{2\ast}\circ 1_{Y{\ast}})=\bar{f}_{2{\ast}}\circ
(p_{\ast}\circ i_{1{\ast}})=\bar{f}_{1{\ast}}\circ
i_{1{\ast}}{\quoi}.
\]
L'axiome ${\rm(c')}$ garantit de plus que l'isomorphisme associé à un
composé de morphismes de compactification est le composé des
isomorphismes associés à chacun d'eux. La catégorie des morphismes de
compactification étant filtrante (\hyperref[XVII.prop3.2.6]{3.2.6} (ii)), on obtient ainsi un
système transitif d'isomorphismes entre les foncteurs associés aux
diverses compactification de $f$. Choisissant l'une d'elles, on
définit le foncteur $f_{\ast}$.

\setcounter{paragraph}{2}
\paragraph{Homomorphismes de transitivité}\etiquette[3.3.2.3]{XVII.par3.3.2.3}

Toute compactification d'un couple $(f,g)$ de morphismes composables
(diagramme \eqhyperref[XVII.eq3.2.5.2]{3.2.5.2}) définit une compactification de $f$, une de $g$,
une de $f\circ g=(\bar{f}\bar{\bar{g}})(j'j)$, et un isomorphisme
$c_{f,g}:(fg)_{\ast}\isomto f_{\ast}g_{\ast}$, composé des
isomorphismes
\[
(fg)_{\ast}=(\bar{f}\,\bar{\bar{g}})_{\ast}(j'j)_{\ast}=\bar{f}_{\ast}(\bar{\bar{g}}_{\ast}j'_{\ast})j_{\ast}=\bar{f}_{\ast}(i_{\ast}\bar{g}_{\ast})j_{\ast}=(\bar{f}_{\ast}i_{\ast})(\bar{g}_{\ast}j_{\ast})=f_{\ast}g_{\ast}
\]
(où\pageoriginale l'isomorphisme médian provient de (iii)).

Prouvons que cet isomorphisme ne dépend pas de la compactification
choisie de $(f,g)$. En vertu de \hyperref[XVII.prop3.2.6]{3.2.6} (iii), il suffira de comparer
les isomorphismes $c_{f,g}$ obtenus à partir de deux compactifications
$\sisi{(3.2.4.2)}{\eqhyperref[XVII.eq3.2.5.2]{3.2.5.2}}_i$ $(i=1,2)$ telles qu'existe un morphisme de
compactification de la première dans la seconde. Les flèches de ces
morphismes seront désignées par $p$.

Considérons le diagramme suivant d'isomorphismes de foncteurs. Pour
faciliter sa lecture, on a omis d'écrire les ${\ast}$.
\[
\xymatrix{
\bar{f}_2 \bar{\bar{g}}_2j'_2j_2\sisi{}{\ar@{}[rrd]|{(1)}} \ar@{-}[rr] \ar@{-}[ddd] \ar@{-}[rd]& &
\bar{f}_2\bar{\bar{g}}_2pj'_1j_1
\ar@{-}[rr] \ar@{-}[ld] \ar@{-}[rd]& &\bar{f}_1\bar{\bar{g}}_1j'_1j_1 \ar@{-}[ddd]\sisi{}{\ar@{}[lld]|{(2)}}\\
 \ar@{}[ddr]|{(3)} & \bar{f}_2\bar{\bar{g}}_2j'_2pj_1\ar@{-}[ru] \ar@{-}[dd]& &
\bar{f}_2p\bar{\bar{g}}_1j'_1j_1 \ar@{-}[dd]\ar@{-}[ru]& \ar@{}[ldd]|{(5)}\\
 &  &  (4)  & &\\
\bar{f}_2i_2\bar{g}_2j_2 \ar@{-}[r] & \bar{f}_2i_2\bar{g}_2pj_1 & &
\bar{f}_2pi_1\bar{g}_1j_1 \ar@{-}[r]&  \bar{f}_1i_1\bar{g}_1j_1\\
& &  \bar{f}_2i_21_Y\bar{g}_1j_1 \ar@{-}[ru]\ar@{-}[lu]& & .\\
}
\]

\begin{enumerate}
\renewcommand{\theenumi}{\alph{enumi}}
\renewcommand{\labelenumi}{\theenumi)}
\item L'isomorphisme\pageoriginale composé de la 1\sisi{ère}{\textsuperscript{ière}} ligne est l'isomorphisme
\hyperref[XVII.par3.3.2.1]{3.3.2.1} entre deux définitions de $(fg)_{\ast}$.

\item L'isomorphisme composé de la dernière ligne est l'isomorphisme
\hyperref[XVII.par3.3.2.1]{3.3.2.1} entre deux définitions de $f_{\ast}g_{\ast}$.

\item Les isomorphismes verticaux extrêmes sont les isomorphismes
$c_{f,g}$ entre $(fg)_{\ast}$ et $f_{\ast}g_{\ast}$ déduits de l'une
ou l'autre compactification.
\end{enumerate}
Le problème est donc de prouver que le
bord extérieur du diagramme est commutatif. Un diagramme commutatif de
foncteurs reste commutatif quand on compose chacun d'eux avec un même
foncteur. Ceci rappelé, la commutativité de (\sisi{i}{1}) résulte de (c)
appliqué au diagramme
\[
\xymatrix{
Z  \ar@{^{(}->}[r]^{j_1} \ar@{=}[d]& \bar{Z}_1  \ar@{^{(}->}[r]^{j'_1} \ar[d]& \bar{\bar{Z}}_1 \ar[d]\\
Z  \ar@{^{(}->}[r]^{j_2} & \bar{Z}_2  \ar@{^{(}->}[r]^{j'_2} & \bar{\bar{Z}}_2\rlap{.}
}
\]
La commutativité de (3) et (5) est triviale, celle de (2), qui
concerne essentiellement des morphismes propres, l'est aussi. Reste à
considérer l'hexagone (4).

Considérons le cube commutatif suivant
\[
\xymatrix{
                              &**[r]\bar{Z}_2  \ar@{^{(}->}[rrr]^{ j'_2}
                              \ar[dd]^(.35){\bar{g}_2}|\hole &    &   &
                              **[r]\bar{\bar{Z}}_2 \ar[dd]^{\bar{\bar{g}}_2}\\
  \bar{Z}_1\sisi{}{\ar@{-->}[rd]}  \ar@{^{(}->}[rrr]^(.65){ j'_1} \ar[ru] \ar[dd]^{\bar{g}_1}   &           &               &
                              \bar{\bar{Z}}_1 \ar[ru]\ar@{-->}[rd] \ar[dd]^(.35){\bar{\bar{g}}_1}|\hole&\\
& **[r] Y\ar[rrr]_{i_2}  \ar@{^{(}->}[rrr]  & && **[r]\bar{Y}_2 \\
Y  \ar@{^{(}->}[rrr]_{i_1}   \ar@{=}[ru]  &           &               &
                              \bar{Y}_1\ar[ru]  &.\\
}
\]

Il\pageoriginale y a six façons de parcourir les arêtes de ce cube
$\bar{Z}_1$ à $\bar{Y}_2$, chacune « adjacente »  à deux autres;
on dispose d'un isomorphisme entre les foncteurs composés associés à
deux chemins adjacents (\ie séparés par une seule face) et la
commutativité de (4) résultera de la commutativité de l'hexagone ainsi
obtenu.

La morale de ce genre de diagramme est qu'une arête représente un
foncteur, une face un isomorphisme entre foncteurs composés et un
volume une condition de compatibilité (entre les faces qui en sont le
bord). On décomposera le cube en deux prismes à base triangulaire (par
le « plan »  $\bar{Z}_1 \bar{\bar{Z}}_1 Y \bar{Y}_2$) pour
se ramener aux hypothèses ${\rm(c')}$.

De façon précise, on complète l'hexagone en le diagramme
d'isomorphismes suivant (pour faciliter sa lecture, les ${\ast}$
ont été omis):
\[
\xymatrix{
                   & & \bar{\bar{g}}_2 j'_2p      & &   \\
i_2\bar{g}_2 p  \ar@{-}[rru] \ar@{-}[rd]\ar@{-}[dd] &   &   &&   \bar{\bar{g}}_2 p j'_1 \ar@{-}[ld]\ar[dd]\ar@{-}[llu]\\
 & i_2(g_2p) \ar@{-}[rr]&      & (\bar{\bar{g}}_2p)j'_1 & \\
i_21_{Y}\bar{g}_1 \ar@{-}[ru] \ar@{-}[rrd]&      &&   &
  p\bar{\bar{g}}_1 j'_1 \ar@{-}[lu] \ar@{-}[lld]\\
&  &  pi_1 \bar{g}_1  & &
}
\]
Les\pageoriginale triangles sont trivialement commutatifs et les
pentagones le sont en vertu de ${\rm(c')}$.

\paragraph{Condition de cocycles}\etiquette[3.3.2.4]{XVII.par3.3.2.4}

Vérifions que les isomorphismes $c_{f,g}$ construits en \hyperref[XVII.par3.3.2.3]{3.3.2.3}
vérifient une condition de cocycle pour un triple $(f,g,h)$ de
morphismes composables. Introduisons une compactification de
$(f,g,h)$, de diagramme \sisi{(3.2.4.3)}{\eqhyperref[XVII.eq3.2.5.3]{3.2.5.3}} et considérons le diagramme
d'isomorphismes de foncteurs suivant:
\[
\sisi{\xymatrix{
\bar{f} i \bar{g} j \bar{h} k \ar@{-}[rr] \ar@{-}[dd]&
&\bar{f}\bar{\bar{g}} j' j\bar{h} k \ar@{-}[dd]\\
 & \bar{f}\bar{\bar{g}}j' \bar{\bar{h}}
k' k \ar@{-}[rd]&\\
\bar{f} i \bar{g}\bar{\bar{h}}k' k \ar@{-}[ru] \ar@{-}[rr]&    &
\bar{f}\bar{\bar{g}}\bar{\bar{h}}k'' k' k.
}}{\xymatrix{
\bar{f} i \bar{g} j \bar{h} k \ar@{-}[rr]^{c_{f,g}*h_*} \ar@{-}[dd]_{f_**c_{g,h}}&&
\bar{f}\bar{\bar{g}} j' j\bar{h} k \ar@{-}[dd]\ar@{-}[rddd]^{c_{fg,h}}\\
&\\
\bar{f} i \bar{g}\bar{\bar{h}}k' k \ar@{-}[rrrd]_{c_{f,gh}} \ar@{-}[rr] && \bar{f}\bar{\bar{g}}j' \bar{\bar{h}}
k' k \ar@{-}[rd]\\
& &   & \bar{f}\bar{\bar{g}}\bar{\bar{\bar{h}}}k'' k' k.
}}
\]
La commutativité du carré est triviale, celle des triangles résulte de
(c) et ${\rm(c')}$. La commutativité du bord extérieur étant ce qu'il
fallait démontrer, ceci achève la démonstration de \hyperref[XVII.prop3.3.2]{3.3.2}.

\subsubsection{}\etiquette[3.3.3]{XVII.subsubsec3.3.3}

Supposons données sur $(S,i)$ et $(S,p)$ respectivement des catégories
fibrées munies d'un scindage normalisé $F_i$ et $F_p$, et supposons
que $F_i$ et $F_p$ ont même fibre $F(X)$ en tout $X\in \Ob S$.

D'après \SGA 1 VI, il « revient au même »  de se donner plutôt:
\begin{itemize}
\item[(i bis)] pour tout morphisme propre $p:X\rightarrow Y$, un
foncteur $p^{\ast}:F(Y)\rightarrow F(X)$,

\item[(ii bis)] pour tout couple composable $(p,q)$ de morphismes
propres, un isomorphisme de foncteur
$c_{p,q}:q^{\ast}p^{\ast}\leftrightarrow (pq)^{\ast}$,

\item[(${\rm i'}$ bis)] et\pageoriginale (${\rm ii'}$ bis) données analogues pour
$(S,i)$, ces données vérifiant les conditions (a bis) (b bis) (${\rm a'}$ bis)
(${\rm b'}$ bis) duales de \hyperref[XVII.subsubsec3.3.1]{3.3.1} (a) (b) ${\rm (a')}$ ${\rm (b')}$.

Supposons \emph{donné} de plus,

\item[(iii bis)] pour tout diagramme commutatif $D$
\[
\vcenter{
\xymatrix{
X \ar[r]^{j'} \ar[d]^{p'}&  Y\ar[d]^{p} \\
Z \ar[r]^{j}& T
}}\quad j,j' \in \Fl(S,i)\text{; }p\sisi{}{,} p' \in \Fl(S,p)
\]
un isomorphisme de foncteurs $d(D):{j'}^{\ast}p^{\ast}\leftrightarrow
{p'}^{\ast}j^{\ast}$.
\end{itemize}

On laisse au lecteur d'énoncer les conditions (c bis) et (${\rm c'}$
bis) duales de \hyperref[XVII.prop3.3.2]{3.3.2} (c) et ${\rm (c')}$.

Des raisonnements parallèles à ceux qui précèdent prouvent alors la
proposition suivante.

\begin{enonce*}{Proposition 3.3.4}\etiquette[3.3.4]{XVII.prop3.3.4}
Si des données {\rm (i bis) (ii bis) (${\rm i'}$ bis) (${\rm ii'}$
bis) (iii bis)} vérifient les conditions {\rm (a bis) (b bis) (${\rm
a'}$ bis) (${\rm b'}$ bis) (c bis) (${\rm c'}$ bis)}, il existe une et
essentiellement une catégorie $F$ cofibrée sur $S$ et munie
d'isomorphismes $\fprod{F}{(S,i)}{S}=F_i$, $\fprod{F}{(S,p)}{S}=F_p$,
ces isomorphismes étant compatibles avec les identifications
$F_{iX}=F_{pX}$ $(X\in \Ob S)$ et tels que les isomorphismes $d(D)$ de
{\rm (iii bis)} deviennent les isomorphismes composés
${p'}^{\ast}j^{\ast}=(jp')^{\ast}=(pj')^{\ast}={j'}^{\ast}p^{\ast}$.
\end{enonce*}

On notera cependant que \hyperref[XVII.prop3.3.4]{3.3.4} n'est pas dual de \hyperref[XVII.prop3.3.2]{3.3.2},
car les hypothèses faites sur $(S,i)$ et $(S,p)$ ne sont pas autoduales.

\section{Résolutions. Application à la flèche de changement de
base}\etiquette[4]{XVII.sec4}\pageoriginale

\subsection{Résolutions plates}\etiquette[4.1]{XVII.subsec4.1}

Rappelons que si $f:(X,\Aa)\rightarrow (Y,\Bb)$ est un
morphisme de sites (ou de topos) annelés, l'image réciproque par $f$
d'un faisceau de $\Bb$-modules à gauche $\Mm$ est le
faisceau de $\Aa$-modules à gauche
$\Aa\otimes_{f^{\ast}\Bb} f^{\ast}\Mm$.

\begin{enonce*}{Proposition 4.1.1}\etiquette[4.1.1]{XVII.prop4.1.1}
Soient $f:(X,\Aa)\rightarrow (Y,\Bb)$ un morphisme de
sites annelés, $K$ un faisceau de $\Aa$-modules à droite, et
$0\rightarrow L\rightarrow M\rightarrow N\rightarrow 0$ une suite
exacte de faisceaux de $\Bb$-modules à gauche. Si $M$ et $N$
sont plats, alors
\begin{enumerate}
\renewcommand{\theenumi}{\roman{enumi}}
\renewcommand{\labelenumi}{(\theenumi)}
\item $L$ est plat,

\item la suite suivante est exacte:
\[
0\longrightarrow K\otimes_{\Aa}f^{\ast}L\longrightarrow
K\otimes_{\Aa}f^{\ast}M\longrightarrow
K\otimes_{\Aa}f^{\ast}N\longrightarrow 0{\quoi}.
\]
\end{enumerate}
\end{enonce*}

Factorisons $f$ en les flèches
$(X,\Aa)\xrightarrow{u}(X,f^{\sisi{\ast}{-1}}\Bb)\xrightarrow{v}(Y,\Bb)$.
Le foncteur $v^{\ast}$ est exact, et transforme Modules plats en
Modules plats (\sisi{IV 15}{\hyperrefext[V][V.sec1]{1}}) de sorte qu'il suffit de prouver \hyperref[XVII.prop4.1.1]{4.1.1} pour
$u$. L'isomorphisme
$$\etiquette[4.1.1.1]{XVII.eq4.1.1.1}
K\otimes_{\Aa}(\Aa\sisi{_{f^{\ast}\Bb}\otimes}{\otimes_{f^{-1}\Bb}} P)=K\otimes_{f^{\sisi{\ast}{-1}}\Bb}P\leqno(4.1.1.1)
$$
nous ramène alors au cas où $f$ est l'identité.

Soit une suite exacte $0\rightarrow R\rightarrow P\rightarrow
K\rightarrow 0$ avec $P$ plat\sisi{ (\hyperref[XVII.subsubsec4.1.2]{4.1.2})}{}. Les lignes et colonnes du
diagramme suivant sont exactes:
\[
\xymatrix{
 &      &                          &  0\ar[d]                        &  &\\
 &      & R\otimes L  \ar[r]\ar[d] & P\otimes L  \ar[r]\ar[d] & K\otimes L
 \ar[r]\ar[d]& 0\\
 & 0\ar[r]     & R\otimes M  \ar[r]\ar[d] & P\otimes M  \ar[r]\ar[d] & K\otimes M
 \ar[r]\ar[d]& 0\\
 & 0\ar[r]     & R\otimes N  \ar[r] \ar[d]& P\otimes N  \ar[r]\ar[d] & K\otimes N
 \ar[r] \ar[d]& 0\\
 &      & 0         &   0   &   0   &.
}
\]

Appliquons\pageoriginale le diagramme du serpent au deux premières
colonnes: on trouve
\[
0\xrightarrow{j}K\otimes L\rightarrow K\otimes M{\quoi},
\]
ce qui prouve (ii). Sans plus supposer $P$ plat, on sait que dans le
diagramme précédent les trois colonnes et deux lignes sont des suites
exactes courtes; la première ligne est donc exacte, \ie $L$ est plat.

\setcounter{subsubsection}{1}
\subsubsection{}\etiquette[4.1.2]{XVII.subsubsec4.1.2}
Soit $\Ss$ une sous-catégorie de la catégorie des sites
annelés.\nde{\label{XVII.ndesous}la théorie vaut plus généralement pour des catégories $\Ss$ formées de sites annelés.} Les faisceaux de modules sur les sites annelés de
$\Ss$ forment une catégorie $\Mm$, abélienne sur
$\Ss$ (\hyperref[XVII.def2.2.2]{2.2.2}), fibrée et cofibrée, dont les fibres sont les
opposées des catégories usuelles de faisceaux de modules sur un
site. On désignera par $K(\Ss)$ la catégorie $K(\Mm)$
(\hyperref[XVII.exem2.2.8]{2.2.8}) et par $K^{+}(\Ss)$ (\resp $K^{-}(\Ss)$) la
sous-catégorie pleine de $K(\Ss)$, dont les catégories fibres
sont les opposées des catégories $K^{+}(S)$ (\resp $K^{-}(S)$) pour
$S\in \Ob \Ss$. Les catégories $K(\Ss)$,
$K^{+}(\Ss)$ et $K^{-}(\Ss)$ sont triangulées sur
$\Ss$, fibrées et cofibrées.

Soit\pageoriginale $f:S\rightarrow S'$ une flèche de $\Ss$. Tout
injectif est acyclique (\hyperref[XVII.def1.2.3]{1.2.3}) pour le foncteur
$Rf_{\ast}$, et tout faisceau de modules plats est acyclique pour le
foncteur $Lf^{\ast}$ (\hyperref[XVII.prop4.1.1]{4.1.1} et \hyperref[XVII.prop1.2.6]{1.2.6}),
d'où des foncteurs dérivés (\hyperref[XVII.prop1.2.7]{1.2.7})
\begin{align*}
Rf_{\ast}:D^{+}(S) &\longrightarrow D^{+}(S')\tag{4.1.2.1}\\
Lf^{\ast}:D^{-}(S\sisi{}{'}) &\longrightarrow D^{-}(S\sisi{'}{}){\quoi}.\tag{4.1.2.2}
\end{align*}\etiquette[4.1.2.1]{XVII.eq4.1.2.1}
\etiquette[4.1.2.2]{XVII.eq4.1.2.2}%L'espace vertical suivant ce paragraphe est trop large.

Si $Rf_{\ast}$ (\resp $Lf^{\ast}$) est de dimension finie, on peut
dans \eqhyperref[XVII.eq4.1.2.1]{4.1.2.1} (\resp \eqhyperref[XVII.eq4.1.2.2]{4.1.2.2}) remplacer l'exposant $+$ (\resp $-$) par
$b$, par $-$ (\resp $+$) ou le supprimer (\hyperref[XVII.prop1.2.10]{1.2.10}).\nde{\label{XVII.ndeKS}en utilisant les résolutions de Spaltenstein, on peut supprimer les exposants inconditionnellement. Dans le cas commutatif, voir M.~Kashiwara et P.~Schapira, \emph{Categories and Sheaves}, Springer, 2010, cité [KS] dans la suite, 18.6.}

On désigne par $D(\Ss)$ (\resp $D^{+}(\Ss)$, \resp
$D^{-}(\Ss)$) la catégorie déduite de $K(\Ss)$ (\resp
$K^{+}(\Ss)$, \resp $K^{-}(\Ss)$) en inversant les
quasi-isomorphismes.

\begin{enonce*}{Scholie 4.1.3}\etiquette[4.1.3]{XVII.sch4.1.3}
(« Théorème trivial de dualité »).
\begin{enumerate}
\renewcommand{\theenumi}{\alph{enumi}}
\renewcommand{\labelenumi}{(\theenumi)}
\item On suppose que tous les foncteurs $Rf_{\ast}$ (\resp
$Lf^{\ast}$) sont de dimension finie. Alors:
\begin{itemize}
\item[(i)] La formation de $D(\Ss)$ commute à tout changement
de catégorie base $\Ss'\rightarrow \Ss$. La catégorie
$D(\Ss)$ est triangulée sur $\Ss$ \emph{(\hyperref[XVII.def2.2.3]{2.2.3})}, et ses
fibres s'identifient aux opposées des catégories $D(S)$ pour $S\in \Ob
\Ss$.

\item[(ii)] La catégorie $D(\Ss)$ est cofibrée (\resp fibrée)
sur $\Ss$. Soit $f:S\rightarrow S'$; si $Lf^{\ast}$ (\resp
$Rf_{\ast}$) est de dimension finie, alors tout objet de $D(S')$
(\resp $D(S)$) a une image réciproque (\resp directe) par $f$ au sens
de la catégorie $D(\Ss)$ sur $\Ss$. Les foncteurs
image réciproque et image directe s'identifient à $Lf^{\ast}$ et
$Rf_{\ast}$.

\item[(iii)] Tout\pageoriginale objet de $D^{-}(S')$ (\resp $D^{+}(S)$) a une
image réciproque (\resp directe) au sens de $D(\Ss)$ et le
foncteur image réciproque (\resp directe) s'identifie à $Lf^{\ast}$
(\resp $Rf_{\ast}$).\nde{les hypothèses de finitude de dimension dans (a) sont superflues. \Cf la N.D.E. précédente.}
\end{itemize}

\item \begin{itemize}
\item[(i)] La formation de $D^{-}(\Ss)$ (\resp
$D^{+}(\Ss)$) commute à tout changement de catégorie base
$\Ss'\rightarrow \Ss$. La catégorie
$D^{-}(\Ss)$ (\resp $D^{+}(\Ss)$ est triangulée sur
$\Ss$, et ses fibres s'identifient aux opposées des catégories
$D^{-}(S)$ (\resp $D^{+}(S)$).

\item[(ii)] La catégorie $D^{-}(\Ss)$ (\resp
$D^{+}(\Ss)$) est fibrée (\resp cofibrée) sur $\Ss$;
les foncteurs « image réciproque »  (\resp « image directe »)
s'identifient aux foncteurs dérivés $Lf^{\ast}$ (\resp $Rf_{\ast}$).

\item[(iii)] Si $f:S\rightarrow S'$ est tel que $Rf_{\ast}$ (\resp
$Lf^{\ast}$) soit de dimension finie, tout objet de $D^{-}(S)$ (\resp
$D^{+}(S')$) a une image directe au sens de $D^{-}(\Ss)$
(\resp une image réciproque au sens de $D^{+}(\Ss)$), et le
foncteur « image directe »  (\resp « image réciproque »)
s'identifie à $Rf_{\ast}$ (\resp $Lf^{\ast}$).
\end{itemize}
\end{enumerate}
\end{enonce*}

Il suffira de vérifier que dans chaque cas les hypothèses de
\hyperref[XVII.prop2.4.1]{2.4.1} ou
\hyperref[XVII.prop2.4.3]{2.4.3} sont remplies, les autres assertions se déduisant après un
changement de base $\{S\}\rightarrow \Ss$, ou $I\rightarrow
\Ss$, de \hyperref[XVII.prop2.3.5]{2.3.5} et \hyperref[XVII.prop1.2.10]{1.2.10}.
Dans le cas (b), les complexes de
faisceaux plats (\resp flasques) satisfont aux hypothèses de
\sisi{\hyperref[XVII.prop2.4.1]{2.4.1}}{\hyperref[XVII.prop2.4.3]{2.4.3} (\resp \hyperref[XVII.prop2.4.1]{2.4.1})}. Dans le cas (a), on utilisera \sisi{\hyperref[XVII.prop2.4.3]{2.4.3}}{\hyperref[XVII.prop2.4.1]{2.4.1} (\resp \hyperref[XVII.prop2.4.3]{2.4.3})}
de façon à n'avoir à
considérer que des $f$ tels que la dimension de $Rf_{\ast}$ (\resp
$Lf^{\ast}$) soit majorée par $N$ fixe. Prouvons qu'il existe « assez »  de complexes $K$ dont les composantes ont leurs images
directes (\resp réciproques) toutes acycliques pour les foncteurs
image directe (réciproque) considérés. Soit $K$ un
complexe\pageoriginale et $K^{\ast}$ (\resp $K_{\ast}$) une résolution
de Cartan-Eilenberg de $K$ par des complexes à composantes flasques
(\resp plates).\nde{on exige seulement que l'augmentation $K^p\to (K^*)^{p\bigdot}$ (\resp $(K_*)^{p\bigdot}\to K^p$) soit une résolution flasque (\resp plate) pour tout $p$.} Si $n>N$, on voit que
$f_{\ast}\tau''_{<n}(K^{\ast})=\tau''_{<\sisi{}{n}}f_{\ast}(K^{\ast})$ (\resp
$f^{\ast}\tau''_{>-n}(K_{\ast})=\tau''_{>-n}f^{\ast}(K_{\ast})$) où
$\tau\sisi{^{\ast}}{''}$ désigne un tronqué relativement à la seconde
différentielle. Pour $n>2N$, on en déduit que $\tau''_{<n}(K^{\ast})$
(\resp $\tau''_{>-n}(K_{\ast})$) vérifie l'hypothèse de \sisi{\hyperref[XVII.prop2.4.3]{2.4.3}}{\hyperref[XVII.prop2.4.1]{2.4.1} (\resp \hyperref[XVII.prop2.4.3]{2.4.3})}, et tout
complexe $K$ est la source (\resp le but) d'un quasi-isomorphisme de
but (\resp de source) un tel complexe.

\setcounter{subsubsection}{3}
\subsubsection{}\etiquette[4.1.4]{XVII.subsubsec4.1.4}
Considérons un carré commutatif de sites annelés \eqhyperref[XVII.eq4.1.4.1]{4.1.4.1}
(\cf~\hyperref[XVII.subsec0.11]{0.11})
$$\etiquette[4.1.4.1]{XVII.eq4.1.4.1}
\vcenter{
\xymatrix{
  X' \ar[r]^{g'}\ar[d]_{f'}& X\ar[d]^{f}\\
S' \ar[r]^{g} & S\rlap{.}
}}\leqno{(4.1.4.1)}
$$
On peut considérer la sous-catégorie de la catégorie des sites réduite
à ce diagramme;\nde{il faut plutôt considérer la catégorie formée de sites annelés définie par ce diagramme, qui n'est pas une sous-catégorie stricto sensu de la catégorie des sites annelés en général. \Cf \nderef{XVII.ndesous}.} on utilise alors \sisi{\hyperref[XVII.subsubsec2.1.2]{2.1.2}}{\hyperref[XVII.prop2.1.3]{2.1.3}} pour
définir un « morphisme de changement de base » de $Lg^{\ast}Rf_{\ast}$ dans $Rf'_{\ast}L{g'}^{\ast}$ lorsque \hyperref[XVII.sch4.1.3]{4.1.3} est
applicable:

\begin{enonce*}{Proposition 4.1.5}\etiquette[4.1.5]{XVII.prop4.1.5}
La flèche \sisi{\hyperref[XVII.subsubsec2.1.2]{2.1.2}}{\hyperref[XVII.prop2.1.3]{2.1.3}} de changement de base, de $Lg^{\ast}Rf_{\ast}$
dans $Rf'_{\ast}L{g'}^{\ast}$, est définie dans chacun des cas
suivants:
\begin{enumerate}
\renewcommand{\theenumi}{\roman{enumi}}
\renewcommand{\labelenumi}{(\theenumi)}
\item On travaille dans les catégories $D^{+}$, et $Lg^{\ast}$,
$L{g'}^{\ast}$ sont de dimension finie.

\item On travaille dans les catégories $D^{-}$, et $Rf_{\ast}$,
$Rf'_{\ast}$ sont de dimension finie.

\item On travaille dans les catégories dérivées entières, et
$Rf_{\ast}$, $Rf'_{\ast}$, $Lg^{\ast}$, $L{g'}^{\ast}$ sont de
dimension finie, ainsi que $Lf^{\ast}$ et $L{f'}^{\ast}$ (ou: ainsi
que $Rg_{\ast}$ et $Rg'_{\ast}$).\nde{les hypothèses de finitude de dimension dans (iii) sont superflues. \Cf \nderef{XVII.ndeKS}.}
\end{enumerate}
\end{enonce*}

On\pageoriginale notera que la dernière hypothèse de (iii) semble
canularesque. Le rédacteur ne sait pas s'en débarrasser sauf lorsque
les sites considérés ont assez de points (\hyperref[XVII.subsubsec4.2.12]{4.2.12}). Les propriétés
fonctorielles de cette flèche de changement de base s'obtiennent
aussitôt à partir de la description \sisi{\hyperref[XVII.subsubsec2.1.2]{2.1.2}}{\hyperref[XVII.prop2.1.3]{2.1.3}}.

\setcounter{subsubsection}{5}
\subsubsection{}\etiquette[4.1.6]{XVII.subsubsec4.1.6}
Soit $(S,\Aa)$ un site annelé, et considérons le foncteur « produit tensoriel »
\[
\otimes:\Mod\sisi{_d}{}(S,\Aa^{\circ})\times
\Mod\sisi{_g}{}(S,\Aa)\longrightarrow \Mod(S,C)
\]%Pour la première composante de la source, on prend ou bien les $\Aa$-modules à droite, ou bien les $\Aa^{\circ}$-modules à gauche.
où $C$ est le centre de $\Aa$.

La recette de \hyperref[XVII.subsubsec1.1.6]{1.1.6}, pour l'application de laquelle on définira le
complexe simple associé à un double complexe par une somme, permet
d'étendre ce foncteur en un foncteur encore noté $\otimes$:
\[
\otimes:K(S,\Aa^{\circ})\times K(S,\Aa)\longrightarrow
K(S,C){\quoi}.
\]

\begin{enonce*}{Proposition 4.1.7}\etiquette[4.1.7]{XVII.prop4.1.7}
Avec les notations précédentes:
\begin{enumerate}
\renewcommand{\theenumi}{\roman{enumi}}
\renewcommand{\labelenumi}{(\theenumi)}
\item Le foncteur produit tensoriel est dérivable à gauche (\hyperref[XVII.def1.2.1]{1.2.1})
en un foncteur « produit tensoriel total »
\[
{}\overset{L}{\otimes}{}:D^{-}(S,\Aa^{\circ})\times
D^{-}(S,\Aa)\longrightarrow D^{-}(S,C){\quoi}\text{;}
\]

\item un couple $(K,L)$ de complexes bornés supérieurement est déployé
pour $\overset{L}{\otimes}$ dès que $K$ ou $L$ est à
composantes plates.
\end{enumerate}
\end{enonce*}

La démonstration est standard.

\begin{enonce*}{Notation 4.1.8}\etiquette[4.1.8]{XVII.not4.1.8}
Si\pageoriginale $K\in D^{-}(S,\Aa^{\circ})$ et $L\in D^{-}(S,\Aa)$,
on désigne par $K\overset{L}{\otimes}L$ le produit
tensoriel total de $K$ et $L$. Le faisceau
$H_{k}(K\overset{L}{\otimes}L)=H^{-k}
(K\overset{L}{\otimes}L)$
se désigne par $\Tor_{k}(K,L)$ et s'appelle le $k$\textsuperscript{ième}
hypertor local de $K$ et $L$.
\end{enonce*}

Soient $K\in D^{-}(S,\Aa^{\circ})$ et $L\in
D^{-}(S,\Aa)$. Si $H^{i}(K)=0$ pour $i>k$ et $H^{i}(L)=0$ pour
$i>\ell$, on vérifie que
$H^{i}(K\overset{L}{\otimes}L)=0$ pour $i>k+\ell$ et
que $H^{k+\ell}(K\overset{L}{\otimes}L)=H^{k}(K)\otimes
H^{\ell}(L)$. Fixons $K$ et regardons
$K\overset{L}{\otimes}L$ comme un foncteur en $L$. Si
on applique la définition \hyperref[XVII.def1.2.9]{1.2.9} de l'amplitude cohomologique comme
intervalle de $\ZZ$, la borne supérieure de cet intervalle est
de nature triviale. C'est le plus grand entier $i$ tel que
$H^{i}(K)\neq 0$. La borne inférieure est plus intéressante:

\begin{enonce*}{Définition 4.1.9}\etiquette[4.1.9]{XVII.def4.1.9}
On dit qu'un complexe $K\in \Ob D^{b}(S,\Aa^{\circ})$ est \emph{de
\tTor-dimension $\leq n$} si les conditions équivalentes suivantes
sont vérifiées:
\begin{enumerate}
\renewcommand{\theenumi}{\roman{enumi}}
\renewcommand{\labelenumi}{(\theenumi)}
\item Pour tout $\Aa$-module $L$, $\Tor_{k}(K,L)=0$ pour
$k>n$.

\item Pour tout complexe de $\Aa$-modules $L$ tel que
$H^{i}(L)=0$ pour $i<\ell$, $H^{i}(K\overset{L}{\otimes}L)=0$
pour $i<\ell-n$.

\item Il existe un quasi-isomorphisme de but $K$ et de source un
complexe\sisi{}{ borné} $K'$, à composantes plates et nulles pour $i<-n$.
\end{enumerate}
\end{enonce*}

\setcounter{subsubsection}{9}
\subsubsection{}\etiquette[4.1.10]{XVII.subsubsec4.1.10}
Voici quelques variations possibles:
\begin{enumerate}
\renewcommand{\theenumi}{\roman{enumi}}
\renewcommand{\labelenumi}{(\theenumi)}
\item Si $K\in \Ob D^{b}(S,\Aa^{\circ})$ est de \ttor-dimension
finie et si $L\in \Ob D(S,\Aa)$, alors le foncteur
$\overset{L}{\otimes}$ est défini (\hyperref[XVII.def1.2.1]{1.2.1}) en $(K,L)$.

\item Si $K\in \Ob K^{b}(S,\Aa^{\circ})$ a des composantes de
\tTor-dimension finie, le foncteur $L\mapsto K\otimes L$ $(L\in
K(S,\Aa))$ est dérivable à gauche.

\item Si tout Module est de \tTor-dimension finie,\nde{cette condition est superflue. Dans le cas où $\Aa$ est commutatif, voir [KS], 18.6, cité dans \nderef{XVII.ndeKS}.} le foncteur $\otimes$
se dérive en un foncteur
\[
{}\overset{L}{\otimes}{}:D(S,\Aa^{\circ})\times
D(S,\Aa)\longrightarrow D(S,\Cc){\quoi}.
\]
\end{enumerate}

\subsection{Résolutions flasques de Godement}\etiquette[4.2]{XVII.subsec4.2}\pageoriginale

\subsubsection{}\etiquette[4.2.1]{XVII.subsubsec4.2.1}
Soient $S$ un topos et $X$ un ensemble. Il revient au même de se
donner une famille de point\sisi{}{s} de $S$ indexée par $X$ ou de se donner un
morphisme $k:\Top(X)\rightarrow S$ du topos défini par l'espace
topologique discret $X$ dans $S$. Pour que la dite famille soit
conservative (IV 6.4.1), il faut et il suffit que pour tout faisceau
$F$ de $S$, la flèche d'adjonction de $F$ dans $k_{\ast}k^{\ast}F$
soit un monomorphisme (\cf IV 6.4.0).

\begin{enonce*}{Définition 4.2.2}\etiquette[4.2.2]{XVII.def4.2.2}
Soit $k:\Top(X)\rightarrow S$ une famille conservative de points d'un
topos~$S$. On appelle \emph{résolution de Godement} ou
\emph{résolution flasque canonique} (relative à $X$) d'un faisceau
abélien $F$ de $S$, et on désigne par $\Cc^{\ast}_X(F)$ (ou
simplement $\Cc^{\ast}(F)$), la résolution à droite suivante
de $F$:
\begin{enumerate}
\renewcommand{\labelenumi}{\theenumi)}
\item $\Cc^{0}(F)=k_{\ast}k^{\ast}F$ et
$\epsilon:F\rightarrow \Cc^{0}(F)$ est la flèche
d'adjonction.

\item Posons $d^{-1}=\epsilon$; on définit par récurrence, pour $n\geq
0$

$\Cc^{n+1}(F)=\Cc^{0}(\coker d^{n-1})$, et
$d^{n}$ comme flèche composée
\[
d^{n}:\Cc^n(F)\rightarrow
\coker(d^{n-1})\rightarrow \Cc^{0}(\coker(d^{n-1})){\quoi}.
\]
\end{enumerate}
\end{enonce*}

\begin{enonce*}{Proposition 4.2.3}\etiquette[4.2.3]{XVII.prop4.2.3}
Sous les hypothèses de \emph{\hyperref[XVII.def4.2.2]{4.2.2}},
\begin{enumerate}
\renewcommand{\theenumi}{\roman{enumi}}
\renewcommand{\labelenumi}{(\theenumi)}
\item $\Cc^n(F)$ est un faisceau flasque \emph{\sisi{(V 3.7)}{\hyperrefext[V][V.def4.1]{4.1}}},

\item $\Cc^{n}(F)$ est un foncteur exact en $F$,

\item la fibre en un point $x\in X$ du complexe
$\Cc^{\ast}(F)$ est une résolution canoniquement scindée de $F_x$.
\end{enumerate}
\end{enonce*}

Les\pageoriginale faisceaux $\Cc^{n}(F)$ sont flasques en
tant qu'image directe de faisceaux (automatiquement flasques) sur
l'espace topologique discret $X$.

Les foncteurs $k^{\ast}$ et $k_{\ast}$ sont exacts. Le foncteur
$\Cc^{0}$ est donc exact. Prouvons par récurrence sur
$n\geq 0$ que le foncteur $\mathcal{Z}^{n}(F)=\coker(d^{n-1})$ est
exact, donc aussi
$\Cc^{n+1}=\Cc^{0}\mathcal{Z}^{n}$. Posons
$\Zz^{-1}(F)=F$. Une suite exacte de faisceaux
\[
0\longrightarrow F'\longrightarrow F\longrightarrow F''\longrightarrow
0
\]
définit pour $n\geq 0$ un diagramme
\[
\xymatrix{
        &   0   \ar[d]             &  0   \ar[d]            &  0 \ar[d]\\
0  \ar[r] & \Zz^{n-1}(F\sisi{^1}{'})  \ar[r] \ar[d]& \Zz^{n-1}(F)  \ar[r] \ar[d]& \Zz^{n-1}(F'')
\ar[r] \ar[d]& 0 \\
0  \ar[r] & \Cc^n(F')  \ar[r] \ar[d]& \Cc^n(F)  \ar[r] \ar[d]&
\Cc^n (F'')  \ar[r] \ar[d]& 0 \\
 0  \ar[r] & \Zz^{n}(F')  \ar[r]\ar[d] & \Zz^{n}(F)  \ar[r] \ar[d]& \Zz^{n}(F'')  \ar[r] \ar[d]&
0\\
        &  0               &  0             &  0      &,
}
\]
dont les colonnes sont exactes, ainsi, d'après l'hypothèse de
récurrence, que les 2 premières lignes. La 3\sisi{è}{\textsuperscript{ième}} ligne est donc exacte,
et (ii) est prouvé.

Les suites exactes courtes
\[
0\longrightarrow \mathcal{Z}^{n+1}(F)\longrightarrow
\Cc^{n}(F)\longrightarrow \mathcal{Z}^{n}(F)\longrightarrow 0
\]
sont des cas particuliers de la suite exacte courte
$$\etiquette[4.2.3.1]{XVII.eq4.2.3.1}
0\rightarrow F\xrightarrow{\epsilon}\Cc^{0}(F)\rightarrow
\mathcal{Z}^{0}(F)\rightarrow 0\leqno(4.2.3.1)
$$
(faire\pageoriginale la substitution $F\mapsto
\mathcal{Z}^{n-1}(F)$). Pour prouver (iii), il reste à prouver que la
suite image réciproque de \eqhyperref[XVII.eq4.2.3.1]{4.2.3.1} par $k$ est canoniquement
scindée. La flèche composée
\[
k^{\ast}F\xrightarrow{k^{\ast}(\epsilon)}k^{\ast}\Cc^{0}F=k^{\ast}(k_{\ast}k^{\ast})F=(k^{\ast}k_{\ast})k^{\ast}F\longrightarrow
k^{\ast}F
\]
est en effet l'identité.

\begin{enonce*}[remark]{Remarque 4.2.4}\etiquette[4.2.4]{XVII.rem4.2.4}
Si $(x_i)_{i\in X}$ est une famille conservative de points d'un site,
à cette famille correspond une famille conservative de points du topos
engendré, qui permet encore de définir les résolutions flasques canoniques.
\end{enonce*}

\begin{enonce*}[remark]{Remarque 4.2.5}\etiquette[4.2.5]{XVII.rem4.2.5}
Si $S$ est un topos annelé, les foncteurs $\Cc^{n}$
transforment Modules en Modules.
\end{enonce*}

\begin{enonce*}[remark]{Remarque 4.2.6}\etiquette[4.2.6]{XVII.rem4.2.6}
Soit $\Ss$ une catégorie formée de sites. Supposons donné,
pour chaque $S\in \Ob \Ss$, un ensemble $S^{d}$ et une famille
conservative $k_{S}:\Top(S^{d})\rightarrow S$ de points de $S$ indexés
par $S^{d}$. Supposons que $S^{d}$ soit fonctoriel en $S$ et que les
diagrammes
\[
\xymatrix{
\Top(S^d)   \ar[r]^{f^d} \ar[d]^{k_S}& \Top(T^d)\ar[d]^{k_T}\\
S        \ar[r]^{f}  & T
}
\]
soient commutatifs. Les résolutions flasques canoniques
correspondantes sont alors « fonctorielles en $S$ ».
\end{enonce*}

Si\pageoriginale $\Ss$ est la catégorie des espaces
topologiques, on peut prendre $S^{d}=$ ensemble sous-jacent.

Si $\Ss$ est la catégorie des sites étales de schémas $S$
localement de type fini sur un schéma $S_{0}$, et si pour chaque
$s\in S_{0}$, $\bar{k(s)}$ est une clôture algébrique de $k(s)$,
on peut prendre pour $S^{d}$ l'ensemble des points de $S$ à valeurs
dans l'un des $\bar{k(s)}$.

On peut remplacer la condition « de type fini »  par une condition
de cardinal si les clôtures algébriques $\bar{k(s)}$ sont remplacées
par des extensions algébriquement closes assez grandes.

\setcounter{subsubsection}{6}
\subsubsection{}\etiquette[4.2.7]{XVII.subsubsec4.2.7}
Un des intérêts des résolutions flasques canoniques est qu'elles
permettent de construire des complexes déployés à la fois pour des
foncteurs du type « produit tensoriel »  et pour des foncteurs du
type « image directe ». \sisi{Il}{Elles} permettent par là une construction
peut-être plus compréhensible des morphismes de changement de base
(\hyperref[XVII.prop4.1.5]{4.1.5}).

\begin{enonce*}{Lemme 4.2.8}\etiquette[4.2.8]{XVII.lem4.2.8}
Soient $f:(S,\Aa)\rightarrow (S',\Bb)$ un morphisme de
sites annelés, $K\in \Ob K(S,\Aa^{\circ})$ et $L\in \Ob
K(S',\Bb)$ des complexes et $k:\Top(X)\rightarrow S$ un
ensemble conservatif de points de $S$. On suppose remplie l'une des
conditions suivantes:
\begin{enumerate}
\renewcommand{\theenumi}{\alph{enumi}}
\renewcommand{\labelenumi}{(\theenumi)}
\item Les cohomologies de $K$ et $L$ sont bornés supérieurement;

\item $\Aa$ est de \tTor-dimension finie sur
$f^{\sisi{\ast}{-1}}\Bb$ et $K\in \Ob D^{b}(S,\Aa^{\circ})$ est
de \tTor-dimension finie;

\item $L\in \Ob D^{b}(S',\Bb)$ est de \tTor-dimension finie.
\end{enumerate}

Alors\pageoriginale le foncteur $(K,L)\mapsto K\otimes f^{\ast}L$ est
dérivable\sisi{}{ à gauche} \emph{(\hyperref[XVII.def1.2.1]{1.2.1})} en $(K,L)$; soit
${}\overset{\LL}{\otimes}\LL f^{\ast}$ son dérivé.
\begin{enumerate}
\renewcommand{\theenumi}{\roman{enumi}}
\renewcommand{\labelenumi}{(\theenumi)}
\item Si $x$ est un point de $S$, alors
$(K\overset{\LL}{\otimes}\LL f^{\ast}L)_x\simeq
K_x\overset{\LL}{\otimes}_{\Bb_{f(x)}}L_{f(x)}$.

\item Le couple $(K,L)$ est déployé à gauche pour
${}\overset{\LL}{\otimes}\LL f^{\ast}$ \eqhyperref[XVII.def1.2.1]{1.2.1} si et
seulement si, pour tout $x\in X$, le couple $(K_x,L_{f(x)})$ est
déployé à gauche pour le foncteur $\otimes_{\Bb_{f(x)}}$.

\item Sous les hypothèses $(a)$ ou $(b)$, le foncteur
$K\overset{\LL}{\otimes}\LL f^{\ast}L$ en $L$ est
d'amplitude cohomologique $\subset d$ \emph{(\hyperref[XVII.def1.2.9]{1.2.9})} si et seulement si,
$\forall x \in V$, le complexe $K_x$ de
$\Bb_{f(x)}$-modules est de \tTor-amplitude $\subset d$.\nde{les hypothèse (a) (b) et (c) sont toutes superflues. \Cf \nderef{XVII.ndeKS}.}
\end{enumerate}
\end{enonce*}

Ce lemme, dont la démonstration est standard, exprime la « nature
ponctuelle »  du foncteur considéré.

\setcounter{subsubsection}{8}
\subsubsection{}\etiquette[4.2.9]{XVII.subsubsec4.2.9}

Soit $L$ un complexe. On appelle \emph{résolution de Godement}
(relative à $X$) de $L$ le complexe simple (\emph{défini en terme de
sommes}) associé au complexe double $\sisi{C}{\Cc}^{\ast}L$
(\hyperref[XVII.subsubsec1.1.4]{1.1.4}). On appelle \emph{résolution de Godement
tronquée à l'ordre} $n$ $\tau''_{\leq n}\sisi{C}{\Cc}^{\ast}L$ le complexe
simple associé au complexe double déduit de $\sisi{C}{\Cc}^{\ast}L$ par troncature
(\hyperref[XVII.subsubsec1.1.15]{1.1.15}) dans le sens de la différentielle de
Godement. En vertu de \hyperref[XVII.prop4.2.3]{4.2.3} (iii), $L$ est point de $X$ par point de
$X$ homotope aux complexes $\tau''_{\leq n} \sisi{C}{\Cc}^{\ast}L$ et $\sisi{C}{\Cc}^{\ast}L$;
ces complexes sont donc « aussi bons que $L$ » vis-à-vis des
foncteurs du type « produit tensoriel » (\hyperref[XVII.lem4.2.8]{4.2.8}).

\begin{enonce*}{Proposition 4.2.10}\etiquette[4.2.10]{XVII.prop4.2.10}
Soient $(S,\Aa)$ un site annelé, $X$ un ensemble conservatif
de points de $S$ et $N$, $M\in \NN\cup \{\infty\}$ avec $N$ ou
$M<\infty$. Soit $A$ la sous-catégorie de $K(S)$ formée des complexes
$L$ tels que
\begin{enumerate}
\renewcommand{\theenumi}{\alph{enumi}}
\renewcommand{\labelenumi}{(\theenumi)}
\item Si $N=\infty$, alors $H^{i}(L)=0$ pour $i$ assez grand; si
$M=\infty$, alors $H^{i}(L)=0$ pour $i$ assez petit.

\item Quels\pageoriginale que soient $x\in X$ et le
$\Aa_x$-module $Q$ de \tTor-dimension $\leq N$, le complexe
$L_x$ est déployé à gauche pour le foncteur
$\sisi{M_x}{Q}\otimes_{\Aa_x}\ast$.

\item Quel que soit le morphisme de topos $f:\sisi{X^{\sim}}{S}\rightarrow E$,
de dimension cohomologique $\leq M$, le complexe $L$ est déployé à
droite pour le foncteur $f_{\ast}$.
\end{enumerate}

Alors, l'inclusion de $A$ dans $K(S)$ induit une équivalence entre la
catégorie déduite de $A$ en inversant les quasi-isomorphismes et la
catégorie $D^{\ast}(S)$, $\ast=+$ si $N<M=\infty$, $\ast=-$ si
$M<N=\infty$ et $\ast=$ blanc si $N$, $M<\infty$.
\end{enonce*}

Soit $K$ un complexe vérifiant (a). Si $N=\infty$, alors la flèche
canonique $a:K_1=\sisi{\sigma}{\tau}_{\leq i}K\rightarrow K$ est quasi-isomorphisme
pour $i$ assez grand. Si $N<\infty$, on pose $K_1=K$. Soit
$K^{\ast}_1$ une résolution plate de Cartan-Eilenberg et posons
$K_2=\tau''_{\geq \sisi{}{-}N}K_1$ (troncature dans le sens de la nouvelle
différentielle). La flèche composée de $K_2$ dans $K$ est un
quasi-isomorphisme, et $K_2$ vérifie (b).

Soit $K$ un complexe vérifiant (a) et (b). Si $M=\infty$, alors
$N<\infty$; pour $i$ assez grand, $b:K\rightarrow \sisi{\sigma}{\tau}_{\geq
i}K_1=K_3$ est un quasi-isomorphisme \sisi{de}{et} $K_3$ vérifie encore (a) et
(b). Si $M<\infty$, on pose $K=K_3$. Soit $K_4=\sisi{\sigma}{\tau}''_{\leq M}
\sisi{C}{\Cc}^{\ast}K_3$. Alors $K_4$ vérifie (a) \sisi{(c), et (b)}{(b) et (c)} par \hyperref[XVII.prop4.2.3]{4.2.3} (iii).

La proposition résulte donc de deux applications successives de \cite[I
2.4.2]{SGA4-17:CD}.

\begin{enonce*}[remark]{Variantes 4.2.11}\etiquette[4.2.11]{XVII.var4.2.11}
Si $N<\infty$ (\resp si $M<\infty$, \resp si $N$, $M<\infty$) alors on
aurait pu travailler de même avec des complexes bornés inférieurement
(\resp supérieurement, \resp bornés).
\end{enonce*}

\setcounter{subsubsection}{11}
\subsubsection{}\etiquette[4.2.12]{XVII.subsubsec4.2.12}
Voici\pageoriginale comment utiliser ces constructions pour définir le
\emph{morphisme de changement de base} \hyperref[XVII.subsubsec4.1.4]{4.1.4} sous
des hypothèses plus générales que celles de \loccit, lorsque les sites
considérés ont assez de points.

Soient un diagramme commutatif (\cf (\hyperref[XVII.subsec0.11]{0.11})) de sites
$$\etiquette[4.2.12.1]{XVII.eq4.2.12.1}
\vcenter{
\xymatrix{
X'          \ar[r]^{g'} \ar[d]^{f'}& X\ar[d]^{f}\\
S'          \ar[r]^{g}  & S\rlap{,}
}}\leqno{(4.2.12.1)}
$$
$\Aa$ un faisceau d'anneaux sur $S$, $\Aa'$ un
faisceau d'anneaux sur $S'$ et $L$ un complexe de $(\Aa',
g^{\ast}\Aa)$-bimodules sur $S'$. On suppose que:
\begin{enumerate}
\renewcommand{\theenumi}{\roman{enumi}}
\renewcommand{\labelenumi}{(\theenumi)}
\item les sites $X$ et $X'$ admettent des familles conservatives de
points, $k:Q\rightarrow X$ et $k':Q'\rightarrow X'$, s'insérant dans
un diagramme commutatif
$$\etiquette[4.2.12.2]{XVII.eq4.2.12.2}
\vcenter{
\sisi{\xymatrix{
X'         \ar[r]^{g'}  \ar[d]^{k'} & X     \ar[d]^{k}  & \\
Q'         \ar[r]^{g'_0} & Q &.
}}{\xymatrix{
Q'         \ar[r]^{g'_0}  \ar[d]^{k'} & Q     \ar[d]^{k}  & \\
X'         \ar[r]^{g'} & X \rlap{.}
}}}\leqno{(4.2.12.2)}
$$

\item L'une des hypothèses suivantes est remplie:
\begin{enumerate}
\renewcommand{\theenumi}{\alph{enumi}}
\renewcommand{\labelenumi}{(\theenumi)}
\item Les foncteurs $Rf_{\ast}$ et $Rf'_{\ast}$ sont de dimension
cohomologique finie, et $L$ est borné supérieurement. On pose $s=-$.

\item $L$ est borné inférieurement, à composantes de \tTor-dimension
uniformément bornée (comme $g^{\ast}\Aa$-modules). On pose
$s=+$.

\item $L$ est borné, et les hypothèses de (a) et (b) sont remplies. On
pose $s=$ blanc.
\end{enumerate}
\end{enumerate}

On\pageoriginale désigne par
$L\overset{\LL}{\otimes}\LL g^{\ast}:D^{s}(S,\Aa)\rightarrow
D^{s}(S', \Aa')$ le foncteur dérivé gauche (\hyperref[XVII.def1.2.1]{1.2.1}) du foncteur
$K\mapsto L\otimes_{g^{\ast}\Aa}g^{\ast}K$ de
$K^{s}(S,\Aa)$ dans $K^{s}(S',\Aa')$.

On se propose de définir une flèche de changement de base
$$\etiquette[4.2.12.3]{XVII.eq4.2.12.3}
\varphi:L\overset{\LL}{\otimes}\LL g^{\ast}Rf_{\ast}K\longrightarrow
Rf'_{\ast}({f'}^{\ast}L\overset{\LL}{\otimes}\LL{g'}^{\ast}K)\leqno(4.2.12.3)
$$
pour $K$ dans $D^{s}(X,\Aa)$.\nde{\label{XVII.ndecb}des adjonctions dans les catégories dérivées non bornées permettent de définir une telle flèche sans les hypothèses (i) et (ii). \Cf \nderef{XVII.ndeKS}.}

En vertu de \hyperref[XVII.prop4.2.10]{4.2.10}, il suffit de la définir, fonctoriellement en $K$,
pour les complexes appartenant à la catégorie $A$ de \hyperref[XVII.prop4.2.10]{4.2.10}, pour
\begin{align*}
N & \geq\sup_i \text{ (tor\sisi{}{-}dim de $L^{i}$ sur
$g^{\ast}$)}\\
M & \geq\text{dimension cohomologique de $Rf_{\ast}$}.
\end{align*}
Un tel complexe est déployé tant pour le foncteur $Rf_{\ast}$ que pour
le foncteur ${f'}^{\ast}L\overset{\LL}{\otimes}_{\Aa'}\sisi{}{\LL}{g'}^{\ast}$
(\hyperref[XVII.lem4.2.8]{4.2.8}). Les flèches canoniques
\begin{align*}
\alpha &: f_{\ast}K\longrightarrow Rf_{\ast}K\\
\beta &:
{f'}^{\ast}L\overset{\LL}{\otimes}\LL \sisi{g}{{g'}}^{\ast}K\longrightarrow
\sisi{}{{f'}^*}L\otimes \sisi{g}{{g'}}^{\ast}K
\end{align*}
sont donc des quasi-isomorphismes. Pour calculer le membre de gauche
$I$ de \eqhyperref[XVII.eq4.2.12.3]{4.2.12.3}, il suffit donc de prendre une résolution gauche
(\ie un quasi-isomorphisme)
\[
\gamma:K_{1}\longrightarrow f_{\ast}K{\quoi},
\]
avec $K_{1}$ déployé pour
$L\overset{\LL}{\otimes}\LL g^{\ast}$; on
a
\[
I\simeq L\otimes \sisi{\LL}{} g^{\ast}K_{1}{\quoi}.
\]

De\pageoriginale même, pour calculer le membre de droite $II$ de
\eqhyperref[XVII.eq4.2.12.3]{4.2.12.3}, il suffit de prendre une résolution droite
\[
\delta:{f'}^{\ast}L\otimes {g'}^{\ast}K\longrightarrow K_2{\quoi},
\]
avec $K_{2}$ déployé pour $Rf'_{\ast}$; on a
\[
II\simeq f'_{\ast}K_{2}{\quoi}.
\]
On définit alors la flèche de changement de base comme le composé
\[
I\simeq L\otimes g^{\ast}K_{1}\xrightarrow{\gamma}L\otimes
g^{\ast}f_{\ast}K\xrightarrow{c}f'_{\ast}({f'}^{\ast}L\otimes
{g'}^{\ast}K)\xrightarrow{\sisi{\alpha}{\delta}}f'_{\ast}K_{2}\simeq II{\quoi},
\]
où la flèche $c$ est la flèche de changement de base
usuelle, au niveau des vrais complexes.

Cette flèche de changement de base a la propriété caractéristique
suivante, qui rend en principe aisée la vérification de ses propriétés
de compatibilité.

\begin{enonce*}{Proposition 4.2.13}\etiquette[4.2.13]{XVII.prop4.2.13}
Sous les hypothèses précédentes, soient $K\in \Ob
C^{s}(X\sisi{'}{},\Aa)$, $K_{1}\in \Ob C^{s}(\sisi{X}{S},\Aa)$ et
$K_{2}\in \Ob C^{s}(\sisi{S}{X'},\Aa)$. Soient de plus $u$ un
$f$-morphisme de $K$ dans $K_{1}$, \ie $u\in
\Hom(K_{1},f_{\ast}K)\simeq \Hom(f^{\ast}K_{1},K)$, et $v$ un
morphisme de ${f'}^{\ast}L\otimes {g'}^{\ast}K$ dans $K_{2}$. On
déduit de $u$, $v$ et de la flèche de changement de base usuelle un
morphisme de complexes $c_{u,v}:L\otimes g^{\ast}K_{1}\rightarrow
f'_{\ast}\sisi{k}{K}_2$. Soient $u'$, $v'$ et $c'_{uv}$ les flèches composées
\begin{align*}
u' &: K_{1}\xrightarrow{u} f_{\ast}K\rightarrow Rf_{\ast}K\\
v' &: L\overset{\LL}{\otimes}\LL{g'}^{\ast}K\rightarrow
L\otimes {g'}^{\ast}K\xrightarrow{v} K_{2}\\
c'_{u,v} &:
L\overset{\LL}{\otimes}\LL{g'}^{\ast}K_{1}\rightarrow
L\otimes {g'}^{\ast}K_{1}
\xrightarrow{c_{uv}}f'_{\ast}K_{2}\rightarrow Rf'_{\ast}K_{2}{\quoi}.
\end{align*}
Le\pageoriginale diagramme suivant est commutatif
$$\etiquette[4.2.13.1]{XVII.eq4.2.13.1}
\vcenter{
\xymatrix{
L \overset{\LL}{\otimes} \sisi{L}{\LL}g^{\ast}Rf_{\ast}K
\ar[r]^{\varphi} & Rf'_{\ast}(f'^{\ast}L
\overset{\LL}{\otimes}\sisi{L}{\LL} g'^{\ast}K) \ar[d]^{Rf'_{\ast}(v')}&\\
L \overset{\LL}{\otimes} \sisi{L}{\LL}g^{\ast}K_1 \ar[u]_{L \overset{\LL}{\otimes} \sisi{L}{\LL}g^{\ast}(u')}
\ar[r]^{c'_{uv}} & Rf'_{\ast}K_2.
}}\leqno{(4.2.13.1)}
$$
\end{enonce*}

Supposons tout d'abord que $K_{1}=f_{\ast}K$ et que
$K_{2}={f'}^{\ast}L\otimes {g'}^{\ast}K$, et soit un diagramme
commutatif de complexes
\[
\xymatrix{
K^{a}  \ar[r] \ar[d]& K\ar[d]\\
K^{c}  \ar[r] & K^{b}
}
\]
tel que $K^{a}$ soit déployé pour
$\sisi{L}{}{f'}^{\ast}L\overset{\LL}{\otimes}\sisi{L}{\LL}{g'}^{\ast}$,
que $K^{b}$ soit déployé pour $Rf_{\ast}$ et que $K^{c}$ soit déployé
pour ces deux foncteurs: pour construire ce diagramme, on commence par
construire $K^{a}$, puis on applique une résolution flasque canonique
(éventuellement tronquée) à $K^{a}$ et $K$.

Appliquons le foncteur $f_{\ast}$ (\resp~$L\otimes {g'}^{\ast}$) à ce
diagramme, et résolvons les complexes du diagramme
$(K_{1},K^{a}_{1},K^{b}_{1}, K^{c}_{1})$ (\resp $(K_{2}, K^{a}_{2},
K^{b}_{2}, K^{c}_{2})$) de façon à obtenir un cube commutatif, où
$\approx$ désigne un quasi-isomorphisme:
\[
\xymatrix@C=2em{
          &         &K'_1 \ar[rdd] \ar[rrrrrr]^{\approx}&        && &         &
          & K_1\ar[rdd]\\  K_1^{a'}\ar[rru]\ar[rdd]\ar[rrrrrr]^(.75){\approx} &         &      &        && & K_1^{a}\ar[rru] \ar[rdd]&       &\\
  &         &      &K^{b'}_1 \ar[rrrrrr] ^(.25){\approx}&& &         &         &   & K^{b}_1\\
          &K_1^{c'} \ar[rru]^{\approx}\ar[rrrrrr] ^(.65){\approx}&      &        && &         &K_1^{c}\ar[rru]^{\approx} \\
}
\]
avec\pageoriginale $K'_{1}$, $K^{a'}_{1}$, $K^{b'}_{1}$, $K^{c'}_{1}$
déployés pour le foncteur
$L\overset{\LL}{\otimes}\sisi{L}{\LL}g^{\ast}$ (\resp
\[
\xymatrix@C=2em{
        &         &K_2 \ar[rdd] \ar[rrrrrr]^{\approx}&        && &         &
          & K'_2\ar[rdd]\\
K_2^{a}\ar[rru]\ar[rdd]^{\approx}\ar[rrrrrr]^(.75){\approx} &         &      &        && & K_2^{a'}\ar[rru] \ar[rdd]^(.65){\approx}&       &\\
  &         &      &K^{b}_2 \ar[rrrrrr] ^(.25){\approx}&& &         &         &   & K^{b'}_2\\
          &K_2^{c} \ar[rru]\ar[rrrrrr] ^(.65){\approx}&      &
        && &         &K_2^{c'}\ar[rru] & &, \\
}
\]
avec $K'_{2}$, $K^{a'}_{2}$, $K^{b'}_{2}$, $K^{c'}_{2}$ déployés pour
le foncteur $Rf'_{\ast}$). Utilisant les morphismes usuels de
changement de base, on obtient un prisme commutatif de complexes sur
$S$.
%%%%%%
\newcommand{\circlenum}[1]{\bigcirc \kern -.255cm \lower
  .05em\hbox{${#1}$}}

\[
\def\objectstyle{\scriptstyle}
\xymatrix@C=-.27cm{
        &         & \scriptstyle{L\otimes g^{\ast}K'_1}\ar[rrrrrr]_*+[o][F-]{\scriptscriptstyle{4}} \ar[rdd]^(.35)*+[o][F-]{\scriptscriptstyle{2}}&        && &         &
          & \scriptstyle{L\otimes g^{\ast}K_1} \ar[rdd] \ar[rrrrr]_*+[o][F-]{\scriptscriptstyle{4}}& & & & & \scriptstyle{f'_{\ast}K_2} \ar[rrrrr]_*+[o][F-]{\scriptscriptstyle{4}}\ar[rdd] & & & && \scriptscriptstyle{f'_{\ast}K'_2} \ar[rdd]\\
\scriptstyle{L\otimes g^{\ast}K_1^{a'}} \ar[rru]\ar[rdd]\ar[rrrrrr] &         &
        &        && & \scriptstyle{L\otimes g^{\ast}K_1\sisi{}{^a}} \ar[rru] \ar[rdd]\ar[rrrrr]&       &
        &&& \scriptstyle{f'_{\ast}K^{a}_2}\ar[rdd]\ar[rru] \ar[rrrrr]&
        & & & &
        \scriptstyle{f'_{\ast}K^{a'}_2}\ar[rru]_*+[o][F-]{\scriptscriptstyle{3}}\ar[rdd]^(.7)*+[o][F-]{\scriptscriptstyle{3}}^{\quad \raisebox{1cm}{}{\scriptscriptstyle{-1}}}_(.7){\text{\rotatebox{110}{$\approx$}}}\\
  &         &      & \scriptstyle{L\otimes g^{\ast}K^{\sisi{c}{b}'}_1}  \ar[rrrrrr]
        && &         &         &   & \scriptstyle{L\otimes g^{\ast}K^{b}_2}
        \ar[rrrrr]& & & & & \scriptstyle{f'_{\ast}K^{b}_2} \ar[rrrrr]& & & & & \scriptstyle{f'_{\ast}K^{b'}_2}\\
          &\scriptstyle{L\otimes g^{\ast}K^{c'}_1}
	\ar[rrrrrr]_*+[o][F-]{\scriptscriptstyle{1}}\ar[rru]_{\text{\rotatebox{50}{$\approx$}}}^*+[o][F-]{\scriptscriptstyle{2}}^{\qquad\qquad \raisebox{-.3cm}{}{\scriptscriptstyle{-1}}}
&      &        && &         &\scriptstyle{L\otimes g^{\ast}K^{c}_2} \ar[rru] \ar[rrrrr]_*+[o][F-]{\scriptscriptstyle{1}}& & & & &\scriptstyle{f'_{\ast}K^{c}_2}\ar[rru]\ar[rrrrr]_*+[o][F-]{\scriptscriptstyle{1}}&
        & & & & \scriptstyle{f'_{\ast}K^{c'}_{2}}\ar[rru]\\
}
\]
et\pageoriginale on conclut en notant que les flèches du diagramme
\eqhyperref[XVII.eq4.2.13.1]{4.2.13.1} s'identifient aux flèches composées suivantes:

pour $\varphi$: la composée des flèches \ding{192};

pour $L\overset{\LL}{\otimes}\sisi{L}{\LL}g^{\ast}(u')$: la
composée des flèches \ding{193}, dont une flèche inverse d'un
quasi-isomorphisme;

pour $Rf'_{\ast}(v')$: la composée des flèches \ding{194}, dont une
flèche inverse d'un quasi-isomorphisme;

pour $\sisi{C}{c}'_{uv}$: la composé des flèches \ding{195}.

Dans le cas général, posons $K'_{1}=f_{\ast}K$ et
$K'_{2}={f'}^{\ast}L\otimes g^{\ast}K$. On dispose alors d'un
diagramme
\[
\xymatrix{
&L \overset{\LL}{\otimes} \sisi{L}{\LL}g^{\ast}Rf_{\ast}K
\ar[r]^\varphi&
Rf'_{\ast}(f'^{\ast}L \overset{\LL}{\otimes}\sisi{L}{\LL}
g'^{\ast}K) \ar[d] \ar[rd]& \\
L \overset{\LL}{\otimes} \sisi{L}{\LL}g^{\ast}K_1 \ar[r]\ar[ru]& L
\overset{\LL}{\otimes} \sisi{L}{\LL}g^{\ast}K'_1 \ar[r] \ar[u]&
Rf'_{\ast}K'_2 \ar[r]&  Rf'_{\ast}K_2,
}
\]
dont les triangles sont commutatifs, ainsi que le rectangle intérieur
d'après ce qui précède. Le contour est donc commutatif, et ceci prouve
\hyperref[XVII.prop4.2.13]{4.2.13}.

\subsection{Le théorème de changement de base}\etiquette[4.3]{XVII.subsec4.3}\pageoriginale

On se propose de reformuler en termes de catégories dérivées le
théorème de changement de base pour les morphismes propres \hyperrefext[XII][XII.thm5.1]{5.1}
(iii).

\begin{enonce*}{Théorème 4.3.1}\etiquette[4.3.1]{XVII.thm4.3.1}
Soit un diagramme cartésien de schémas
\[
\xymatrix{
X' \ar[rr]^{g'} \ar[d]^{f'}&& X\ar[d]^{f}\\
S' \ar[rr]^{g}& & S
}
\]
avec $f$ propre. Soient $\Aa$ et $\Aa'$ des faisceaux
d'anneaux de torsion sur $S$ et $S'$ et $L$ un complexe borné
supérieurement de $(\Aa',g^{\ast}\Aa)$-bimodules. On
suppose que la dimension des fibres de $f$ est bornée. L'hypothèse
\hyperref[XVII.subsubsec4.2.12]{4.2.12} {\rm(a)} est alors remplie, et la flèche de changement de
base \eqhyperref[XVII.eq4.2.12.3]{4.2.12.3}, entre foncteurs de $D^{-}(X,f^{\ast}\Aa)$
dans $D^{-}(S',\Aa')$:
\[
\varphi_K:L\overset{\LL}{\otimes}\sisi{L}{\LL}g^{\ast}\sisi{}{Rf_*}K\longrightarrow
Rf'_{\ast}({f'}^{\ast}L\overset{\LL}{\otimes}\sisi{L}{\LL}{g'}^{\ast}K){\quoi},
\]
est un isomorphisme.
\end{enonce*}

Les foncteurs considérés sont triangulés et « way-out » \cite[I 7]{SGA4-17:RD},
de sorte que, pour vérifier que $\varphi_{K}$ est un isomorphisme, il
suffit de vérifier que les $\varphi_{ℋ^{i}(K)}$ sont des
isomorphismes. On peut donc supposer que $K$ est réduit à un faisceau
de $f^{\ast}\Aa$-modules (de torsion) placé en degré $0$.

Pour\pageoriginale vérifier que $\varphi_{K}$ est un isomorphisme,
il suffit de le vérifier après application du foncteur d'oubli de la catégorie des
$\Aa'$-Modules dans celle des
$\ZZ/n\ZZ$-Modules, où $n\Aa=n\Aa'=0$.

Ceci permet de se ramener au cas où
$\Aa'=\ZZ/n\ZZ$, avec $n\Aa=0$. On peut
alors remplacer $L$ par un complexe isomorphe dans la catégorie
dérivée, à composantes plates sur $g^{\ast}\Aa$, et borné
supérieurement. Les foncteurs considérés étant triangulés et « way
out »  en $L$, il suffit de vérifier que $\varphi_{K}$ est un
isomorphisme après avoir remplacé $L$ par une de ses composantes. On
peut donc supposer que $L$ est réduit à un faisceau de
$g^{\ast}\Aa$-modules plats placé en degré $0$. Il faut
prouver que les flèches
\[
\varphi^{q}:L\otimes g^{\ast}R^{q}f_{\ast}K\longrightarrow
R^{q}f'_{\ast}({f'}^{\ast}L\otimes g^{\ast}K)
\]
sont des isomorphismes. Soit $s$ un point géométrique de $S'$, $X_s$
la fibre de $X'$ en $s$ et $K_s$ les images réciproques de $K$ sur
$X_s$. En vertu de \hyperrefext[XII][XII.thm5.1]{5.1} (iii), la fibre en $s$ de $\varphi^{q}$ est
la flèche
\[
\varphi^{q}_s:L_{s}\otimes H^{q}(X_s,K)\longrightarrow H^{q}(X_s,
L_s\otimes K){\quoi}.
\]
En vertu du théorème de $D$. Lazard \cite{SGA4-17:La}, le $\Aa_s$-module
plat $L_s$ est limite inductive filtrante de $\Aa_s$-modules
libres $L_{s,i}$; le foncteur $H^{q}(X_s,\ast)$ commute aux limites
inductives filtrantes\sisi{}{ (\hyperrefext[VII][VII.prop3.3]{3.3})}. La flèche $\varphi^{q}$ est donc limite
inductive des flèches
\[
\varphi^{q}_i:L_{s,i}\otimes H^{q}(X_s,K)\longrightarrow H^{q}(X_s,
L_{s,i}\otimes K){\quoi}.
\]
Ces dernières sont évidemment des isomorphismes, donc aussi
$\varphi^{q}_s$, et $\varphi\sisi{}{^q}$ est un isomorphisme.

\begin{enonce*}[remark]{Variantes 4.3.2}\etiquette[4.3.2]{XVII.var4.3.2}
\begin{enumerate}
\renewcommand{\theenumi}{\roman{enumi}}
\renewcommand{\labelenumi}{(\theenumi)}
\item Lorsque\pageoriginale $L$ est un complexe borné de composantes
de \tTor-dimension finie sur $g^{\ast}\Aa$, les hypothèses (c) de
\hyperref[XVII.subsubsec4.2.12]{4.2.12} sont remplies, et on peut travailler dans
la catégorie dérivée entière.\nde{on peut travailler dans
la catégorie dérivée entière sans ces hypothèses. \Cf \nderef{XVII.ndeKS}.}

\item Si $L$ borné inférieurement, à composantes de \tTor-dimension \sisi{finie}{uniformément bornée}
sur $g^{\ast}\Aa$, les hypothèses (b) de \hyperref[XVII.subsubsec4.2.12]{4.2.12} sont remplies,
même si $\Aa$ et $\Aa'$ ne sont pas de torsion; la
flèche de changement de base $\varphi_{K}$ $(K\in
D^{+}(X,f^{\ast}\Aa)$) sera un isomorphisme si les faisceaux
$ℋ^{i}(K)$ sont de torsion.
\end{enumerate}
\end{enonce*}

\section{Les foncteurs image directe à support propre}\etiquette[5]{XVII.sec5}

\subsection{La construction fondamentale}\etiquette[5.1]{XVII.subsec5.1}

On se propose de suivre le programme indiqué en \hyperref[XVII.subsec3.1]{3.1}
et de définir le
foncteur $Rf_{!}$ pour $f$ un morphisme compactifiable
(\hyperref[XVII.subsec3.2]{3.2}). On
montrera en appendice comment traiter le cas d'un morphisme séparé de
type fini quelconque.\nde{\Cf \nderef{XVII.ndeConrad}.}

\subsubsection{}\etiquette[5.1.1]{XVII.subsubsec5.1.1}

Soient $X$ un topos, $U$ un ouvert de $X$ (\hyperrefext[IV][IV.def8.3]{8.3}) et $j$ le morphisme
d'inclusion de $U$ dans $X$. On a défini en \hyperrefext[IV][IV.subsec5.2]{5.2} un foncteur « \emph{prolongement par le vide} »  de la catégorie des faisceaux
d'ensembles sur $U$ dans celle des faisceaux d'ensembles sur $X$. Ce
foncteur $j_{!}$ est adjoint à gauche au foncteur de restriction
$j^{\ast}$.

On désigne encore par $j_{!}$ le foncteur de la catégorie des
faisceaux d'ensembles pointés sur $U$ (\resp des faisceaux en groupes,
\resp des faisceaux de $j^{\ast}\Aa$-modules pour
$\Aa$ faisceau d'anneaux sur $X$) dans\pageoriginale la
catégorie analogue sur $X$, adjoint à gauche au foncteur de
restriction $j^{\ast}$ (\sisi{IV\quad}{\hyperrefext[IV][IV.prop11.3.1]{11.3.1}}). Ces trois foncteurs $j_{!}$ sont
compatibles aux foncteurs d'oubli, de la catégorie des Modules dans
celle des Groupes, et de la catégorie des Groupes dans celle des
Ensembles pointés \footnote[1]{Mais pas dans celle des Ensembles
!}. Si $i:F\rightarrow X$ est l'inclusion du fermé complémentaire de
$U$, alors $j^{\ast}j_{!}$ est (canoniquement isomorphe à) l'identité
et $i^{\ast}j_{!}$ est le foncteur constant de valeur l'objet initial
(ou final, cela revient ici au même). Ceci caractérise d'ailleurs
$j_{!}$. On en déduit que les trois foncteurs $j_{!}$ envisagés sont
exacts et s'injectent dans les foncteurs $j_{\ast}$ correspondants. On
les appelle foncteurs de \emph{prolongement par zéro}. La formule
$(j_1j_2)^{\ast}=j^{\ast}_2 j^{\ast}_1$ se transpose en un
isomorphisme de transitivité
\[
(j_1 j_2)_{!}=j_{1!}j_{2!}{\quoi}.
\]
Les foncteurs $j_{!}$ sont compatibles aux changements de base:

\begin{enonce*}{Lemme 5.1.2}\etiquette[5.1.2]{XVII.lem5.1.2}
Soient $f:X\rightarrow Y$ un morphisme de topos, $V$ un ouvert de $Y$
et $U$ son image réciproque dans $X$:
\[
\xymatrix{
U \ar[rr]^{f'}  \ar[d]^{j'}& & V\ar[d]^{j}\\
X \ar[rr]^{f}   && Y
}
\]
Il existe alors un et un seul isomorphisme de changement de base
$f^{\ast}j_{!}\isomto j'_{!}{f'}^{\ast}$ rendant commutatif
le diagramme
$$\etiquette[5.1.2.1]{XVII.eq5.1.2.1}
\vcenter{
\xymatrix{
f^{\ast}j_{!}    \ar[r]^{\sim} \ar@{^{(}->}[d] & j'_{!}f'^{\ast} \ar@{^{(}->}[d]     \\
f^{\ast}j_{\ast}  \ar[r]^{\alpha}&j'_{\ast}f'^{\ast},
}}\leqno{(5.1.2.1)}
$$
où $\alpha$ est la flèche de changement de base.
\end{enonce*}

\begin{enonce*}[remark]{Preuve}
La\pageoriginale flèche de changement de base
$\varphi:j^{\ast}f_{\ast}\rightarrow f'_{\ast}{j'}^{\ast}$ est
trivialement un isomorphisme. Elle se transpose en un isomorphisme de
changement de base $\varphi':j'_{!}{f'}^{\ast}\isomto
f^{\ast}j_{!}$. Pour vérifier que le diagramme \eqhyperref[XVII.eq5.1.2.1]{5.1.2.1} est
commutatif, il suffit de voir que sa restriction à $U$ l'est, ce qui
est trivial.
\end{enonce*}

\setcounter{subsubsection}{2}
\subsubsection{}\etiquette[5.1.3]{XVII.subsubsec5.1.3}
Sous les hypothèses de \hyperref[XVII.subsubsec5.1.1]{5.1.1}, avec $X$ annelé par un faisceau
d'anneaux $\Aa$, le foncteur $j_{!}$ sur les Modules est
exact, donc passe trivialement aux catégories dérivées et définit un
foncteur, encore noté $j_{!}$, de $D(U,j^{\ast}\Aa)$ dans
$D(X,\Aa)$. On dispose encore d'isomorphismes de transitivité
$$\etiquette[5.1.3.1]{XVII.eq5.1.3.1}
(j_1 j_2)\sisi{}{!}=j_{1!}j_{2!}{\quoi}.\leqno(5.1.3.1)
$$

On a encore, au niveau des catégories dérivées, une formule
d'adjonction
$$\etiquette[5.1.3.2]{XVII.eq5.1.3.2}
\Hom_{D(X)}(j_{!}K,L)\isomto \Hom_{D(U)}
(K,j^{\ast}L){\quoi}.\leqno(5.1.3.2)
$$

\subsubsection{}\etiquette[5.1.4]{XVII.subsubsec5.1.4}
Soit $f:X\rightarrow Y$ un morphisme propre, $\Aa$ un faisceau
d'anneaux \emph{de torsion} sur $Y$, et annelons $X$ par le faisceau
d'anneaux image réciproque de $\Aa$. Supposons que la
dimension des fibres de $f$ soit majorée par un entier $n$ (ce qui est
le cas si $Y$ est quasi-compact). On sait alors que pour chaque
faisceau de $f^{\ast}\Aa$-modules $F$ sur $X$, les faisceaux
$R^{q}f_{\ast}F$ sont nuls pour $q>2n$ (\hyperrefext[XII][XII.corbis5.3]{5.3 bis}). Le foncteur
$f_{\ast}$ admet donc un foncteur dérivé (\hyperref[XVII.prop1.2.10]{1.2.10})
\[
Rf_{\ast}:D(X,f^{\ast}\Aa)\longrightarrow
D(Y,\Aa){\quoi}.
\]

Si\pageoriginale $f$ est le composé $gh$ de deux morphismes
propres, on dispose d'une formule de transitivité
$Rf_{\ast}=Rg_{\ast}Rh_{\ast}$, qu'on peut comme d'habitude exprimer
en terme de catégories cofibrées (Scholie \sisi{\hyperref[XVII.subsubsec5.1.3]{5.1.3}}{\hyperref[XVII.sch4.1.3]{4.1.3}}).

\subsubsection{}\etiquette[5.1.5]{XVII.subsubsec5.1.5}
Soit un diagramme commutatif de schémas
$$\etiquette[5.1.5.1]{XVII.eq5.1.5.1}
\vcenter{
\xymatrix{
U        \ar@{^{(}->}[r]_{j_1} \ar[d]^{g}      & X  \ar[d]^{f}\\
V        \ar@{^{(}->}[r]_{j_2}       & Y,
}}\leqno{(5.1.5.1)}
$$
dans lequel $j_1$ et $j_2$ sont des immersions ouvertes et $f$, $g$
des morphismes propres dont la dimension des fibres est majorée par un
entier $n$. Soit $\Aa$ un faisceau d'anneaux de torsion sur
$Y$, et annelons $U$, $V$ et $X$ par les images réciproques de
$\Aa$. En vertu de \eqhyperref[XVII.eq5.1.3.2]{5.1.3.2}, si $K\in D(U)$, définir une flèche
$$\etiquette[5.1.5.2]{XVII.eq5.1.5.2}
d:j_{2!}Rg_{\ast}K\longrightarrow Rf_{\ast}j_{1!}K\leqno{(5.1.5.2)}
$$
revient à définir une flèche
\[
d':Rg_{\ast}K\longrightarrow j^{\ast}_2 Rf_{\ast}j_{1!}K{\quoi}.
\]

Le produit fibré $X'=V\times {}_YX$ s'insère dans un diagramme
commutatif
$$\etiquette[5.1.5.3]{XVII.eq5.1.5.3}
\vcenter{
\xymatrix{
U \ar@{^{(}->}[r]^{k} \ar[d]^{g} & X'  \ar@{^{(}->}[r]_{j'_2} \ar[d]^{f'}& X \ar[d]^{f}          &\\
V         \ar@{=}[r] & V         \ar@{^{(}->}[r]_{j_2}  & Y   \rlap{,}
}}\leqno{(5.1.5.3)}
$$
et,\pageoriginale trivialement, on a $j^{\ast}_2Rf_{\ast}j_{1!}K\simeq
Rf'_{\ast}(j'_2{}^\ast j_{1!}K)\simeq Rf'_{\ast}k_{!}K$.

La flèche $k$ est une immersion ouverte propre; on a donc
$k_{!}=k_{\ast}$, d'où un isomorphisme $Rf'_{\ast}k_{!}K\simeq
Rf'_{\ast}k_{\ast}K\simeq Rg_{\ast}K$, qui nous fournit l'isomorphisme
$d'$ cherché et le morphisme $d$. On voit de plus que la restriction
de $d$ à $V$ est un isomorphisme.

\begin{enonce*}{Lemme 5.1.6}\etiquette[5.1.6]{XVII.lem5.1.6}
Le morphisme \eqhyperref[XVII.eq5.1.5.2]{5.1.5.2} est un isomorphisme.
\end{enonce*}

Il reste à vérifier que la restriction de $d$ à $Y-V=F$ est un
isomorphisme. Soit $X''=\fprod{X}{F}{Y}$ et appliquons le théorème de
changement de base pour un morphisme propre (\hyperref[XVII.thm4.3.1]{4.3.1}) au carré cartésien
\[
\xymatrix{
X'' \ar[r]_{i'} \ar[d]^{f''}& X\ar[d]^{f\sisi{'}{}}\\
F   \ar[r]_{i}& Y\rlap{,}
}
\]
pour $\Aa'=i^{\ast}\Aa$ et $L=\Aa'[0]$. On
trouve:
$i^{\ast}Rf_{\ast}j_{1!}K\simeq Rf''_{\ast}{i'}^{\ast}j_{1!}K\simeq
Rf''_{\ast}0\simeq 0$, de sorte que les deux membres de \eqhyperref[XVII.eq5.1.5.2]{5.1.5.2} sont
nuls en dehors de $U$ et que $d$ est un isomorphisme.

\setcounter{subsubsection}{6}
\subsubsection{}\etiquette[5.1.7]{XVII.subsubsec5.1.7}
Soit $S$ un schéma annelé par un faisceau d'anneaux de torsion
$\Aa$. Pour tout $S$-schéma $\varphi:X\rightarrow S$, on pose
$\Aa_X=\varphi^{\ast}\Aa$. Plaçons-nous dans la catégorie $(S)$ des
$S$-schémas quasi-compacts quasi-séparés et des morphismes
$S$-compactifiables (\hyperref[XVII.subsec3.2]{3.2}). En vertu de
\hyperref[XVII.prop3.2.3]{3.2.3}, cette catégorie, la sous-catégorie $(S,i)$ des
immersions ouvertes, et la sous-catégorie $(S,p)$ des morphismes
propres vérifient les hypothèses \hyperref[XVII.subsubsec3.2.4]{3.2.4}.

On\pageoriginale se propose d'appliquer la théorie \hyperref[XVII.prop3.3.2]{3.3.2} aux données
suivantes (notations de \hyperref[XVII.subsubsec3.3.1]{3.3.1})

\begin{itemize}
\item[$\alpha$)] $F(X)=D(X,\Aa_X)$ pour $\varphi:X\rightarrow
S$ dans $S$:
\begin{itemize}
\item[(i)] Pour $p:X\rightarrow Y$ propre, on prend pour foncteur
$p_{\ast}$ le foncteur $Rp_{\ast}$.

\item[(ii)] On utilise les isomorphismes de transitivité usuels
$R(pq)_{\ast}=Rp_{\ast}Rq_{\ast}$.

\item[${\rm(i')}$] Pour une immersion ouverte $i:X\rightarrow Y$, on
prend pour foncteur $i_{\ast}$ le foncteur $i_{!}$.

\item[${\rm(ii')}$] On utilise les isomorphismes de transitivité
\eqhyperref[XVII.eq5.1.3.1]{5.1.3.1}.

\item[${\rm(iii)}$] Pour tout diagramme commutatif \eqhyperref[XVII.eq5.1.5.1]{5.1.5.1} on prend
pour isomorphisme $d$ l'isomorphisme \eqhyperref[XVII.eq5.1.5.2]{5.1.5.2}.
\end{itemize}
\end{itemize}

On vérifie les conditions (a) (b) (c) ${\rm(a')}$ ${\rm(b')}$
${\rm(c')}$ de \hyperref[XVII.prop3.3.2]{3.3.2}, ce qui prouve le point (a) du:

\begin{enonce*}{Théorème 5.1.8}\etiquette[5.1.8]{XVII.thm5.1.8}
\begin{enumerate}
\renewcommand{\theenumi}{\alph{enumi}}
\renewcommand{\labelenumi}{(\theenumi)}
\item Dans la catégorie $(S)$ définie dans \hyperref[XVII.subsubsec5.1.7]{5.1.7}, on peut d'une, et
essentiellement d'une seule façon, définir pour toute flèche
$f:X\rightarrow Y$ un foncteur
\[
Rf_{!}:D(X,\Aa_X)\longrightarrow D(Y,\Aa_Y){\quoi},
\]
et pour tout morphisme composé
$f=gh:X\xrightarrow{h}Y\xrightarrow{g}Z$ des isomorphismes de
transitivité $c_{\sisi{f,g}{g,h}}$ entre $Rf_{!}$ et $Rg_{!}$ $Rh_{!}$ de sorte
que
\begin{enumerate}
\renewcommand{\theenumii}{\roman{enumii}}
\renewcommand{\labelenumii}{(\theenumii)}
\item Les $c_{f,g}$ vérifient la « condition de cocycles »
\[
(c_{f,g}\ast Rh_{!})\circ c_{fg,h}=(Rf_{!}\ast c_{g,h})\circ
c_{f,gh}{\quoi}.
\]

\item Lorsqu'on se restreint aux morphismes propres (\resp aux
immersions ouvertes) on retrouve la théorie de variance
\hyperref[XVII.subsubsec5.1.4]{5.1.4} (\resp~\hyperref[XVII.subsubsec5.1.3]{5.1.3}).

\item Pour\pageoriginale tout diagramme \eqhyperref[XVII.eq5.1.5.1]{5.1.5.1}, l'isomorphisme
$d$ \eqhyperref[XVII.eq5.1.5.2]{5.1.5.2} est l'isomorphisme composé $c_{f,j\sisi{'}{_1}} \circ c^{-1}_{j\sisi{,f'}{_2,g}}$.
\end{enumerate}

\item Les foncteurs $Rf_{!}$ sont « way out »  \emph{\cite[I 7]{SGA4-17:RD}} et
triangulés.
\end{enumerate}
\end{enonce*}

Tout foncteur $Rf_{!}$ est composé d'un foncteur $Rg_{\ast}$, pour $g$
propre de but quasi-compact, donc à fibres de dimension bornée, et
d'un foncteur $j_{!}$ pour $j$ immersion ouverte. Ces foncteurs sont
way out et triangulés, donc aussi $Rf_{!}$.

\setcounter{subsubsection}{8}
\paragraph{}\etiquette[5.1.8.1]{XVII.par5.1.8.1}

Les isomorphismes de transitivité \hyperref[XVII.thm5.1.8]{5.1.8} (a) sont souvent utilisés via
la suite spectrale « des foncteurs composés »  qui s'en déduit,
valable pour $K\in \Ob D(X,\Aa_X)$:
$$\etiquette[5.1.8.2]{XVII.eq5.1.8.2}
E^{pq}_2=R^{p}g_! R^{q}h_{!}(K)\Longrightarrow
R^{p+q}f_{!}(K){\quoi}.\leqno{(5.1.8.2)}
$$

Pour l'obtenir, on note que d'après \hyperref[XVII.thm5.1.8]{5.1.8} b) et \textsc{Verdier}
\cite{SGA4-17:V2}, pour tout
$L\in \Ob D(Y,\Aa_Y)$, on dispose d'une suite spectrale
$$\etiquette[5.1.8.3]{XVII.eq5.1.8.3}
E^{pq}_2=R^{p} g_!(ℋ^{q}(L))\Longrightarrow
R^{p+q}g_!(L){\quoi},\leqno{(5.1.8.3)}
$$
et \eqhyperref[XVII.eq5.1.8.2]{5.1.8.2} est le cas particulier de \eqhyperref[XVII.eq5.1.8.3]{5.1.8.3} pour $L=Rh_{!}K$.

\begin{enonce*}{Définition 5.1.9}\etiquette[5.1.9]{XVII.def5.1.9}
\begin{enumerate}
\renewcommand{\theenumi}{\roman{enumi}}
\renewcommand{\labelenumi}{(\theenumi)}
\item On appelle « \emph{foncteurs image directe à support propre
\sisi{totals}{totaux}} »  les foncteurs construits en \hyperref[XVII.thm5.1.8]{5.1.8}. On omettra
l'adjectif « total »  lorsqu'il n'y aura pas de risque de
confusion.

\item On appelle $q$\textsuperscript{ième} \emph{foncteur image directe à
support propre,} et on désigne par $R^{q}f_{!}$, le foncteur composé
$\Hh^{q}\circ Rf_{!}$.

\item Si $X$ est un schéma sur un corps $k$ séparablement clos,
immergeable dans un schéma propre sur $k$, on identifie faisceaux de
modules sur\pageoriginale $\Spec (k)$ et modules usuels, et on désigne
le foncteur $R^{q}f_{!}$ par l'une des notations $H^{q}_c(X,\quad)$,
$H^{q}_{!}(X,\quad)$, ou simplement $H^{q}_c(\quad)$ s'il n'y a pas danger
de confusion. On l'appelle alors « $q$\textsuperscript{ième} \emph{groupe
d'hypercohomologie à support propre} » , ou « $q$\textsuperscript{ième} \emph{groupe de cohomologie à support propre} »  si on
l'applique à un complexe réduit au degré $0$.
\end{enumerate}
\end{enonce*}

\setcounter{subsubsection}{9}
\setcounter{paragraph}{0}
\paragraph{}\etiquette[5.1.9.1]{XVII.par5.1.9.1}
La notation $Rf_{!}$ est abusive en ce que \emph{le foncteur $Rf_{!}$
n'est pas le dérivé droit du foncteur $f_{!}$} qui sera défini en
\hyperref[XVII.def6.1.2]{6.1.2}. Certains préfèrent la notation $R_{!}f$.

\subsubsection{}\etiquette[5.1.10]{XVII.subsubsec5.1.10}
Les ingrédients dans la démonstration de \hyperref[XVII.thm5.1.8]{5.1.8} sont le
théorème de
changement de base propre et l'identité entre $j_{\ast}$ et $j_{!}$
lorsque $j$ est une immersion ouverte propre. Cela permet d'énoncer
diverses variantes à \hyperref[XVII.thm5.1.8]{5.1.8}; on laisse au lecteur le soin de leur
donner un sens précis et de les vérifier.

\begin{enonce*}{Variante 5.1.11}\etiquette[5.1.11]{XVII.var5.1.11}
Dans la catégorie $(S)$, on peut d'une, et essentiellement d'une seule
façon, définir pour toute flèche $f:X\rightarrow Y$ un foncteur
$f_{!}$ de la catégorie des faisceaux d'ensembles pointés sur $X$ dans
la catégorie des faisceaux d'ensembles pointés sur $Y$ et pour tout
morphisme composé $f=gh$ des isomorphismes de transitivité vérifiant
des conditions {\rm (i)\sisi{,}{} (ii) (iii)} analogues à celles de \sisi{\hyperref[XVII.subsubsec5.1.4]{5.1.4}}{\hyperref[XVII.thm5.1.8]{5.1.8}}.
\end{enonce*}

Voir \hyperref[XVII.def6.1.2]{6.1.2} pour une description directe de ce foncteur.

\begin{enonce*}{Variante 5.1.12}\etiquette[5.1.12]{XVII.var5.1.12}
On\pageoriginale peut, pour tout morphisme compactifiable
$f:X\rightarrow Y$, définir un foncteur $R^{1}f_{!}$ de la catégorie
des faisceaux de groupes \iInd-finis sur $X$ dans la catégorie des
faisceaux d'ensembles pointés sur $Y$, avec pour toute
compactification $f=\bar{f}j$ de $f$, un isomorphisme
$R^{1}f_{!}\simeq R^{1}\bar{f}_{\ast}\circ f_{!}$.
\end{enonce*}

On ne postule ici aucune formule de transitivité; le seul problème est
de montrer que le foncteur $R^{1}f_{!}$ ne dépend pas de la
compactification choisie de $f$. Soit $p:(\bar{f}_1, i_{1})\rightarrow
(\bar{f}_{2}, i_{2})$ un morphisme de compactification, \ie un
diagramme
$$\etiquette[5.1.12.1]{XVII.eq5.1.12.1}
\vcenter{
\xymatrix{
X   \ar[rr]^{i_1} \ar[rrd]^{i_2} \ar[dd]_{f}& & \bar{X}_1 \ar[d]^{p}\\
    & &  \bar{X}_2 \ar[lld]^{\bar{f}_2}\\
Y   & & .
}}\leqno{(5.1.12.1)}
$$
où $\bar{f}_{1}=\bar{f}_2\circ p$, $p$ étant propre. Si $G$ est un
faisceau \sisi{au}{de} groupes sur $X$, on dispose d'une suite exacte d'ensembles
pointés (avec $\alpha$ injectif)
\[
0\rightarrow
R^{1}\sisi{f}{\bar{f}}_{2\ast}(p_{\ast}i_{1!}G)\xrightarrow{\alpha}R^{1}\sisi{f}{\bar{f}}_{1\ast}(i_{1!}G)\rightarrow
\bar{f}_{2\ast}R^{1}p_{\ast}(i_{1!}G){\quoi},
\]
et d'une flèche de $i_{2!}G$ dans $p_{\ast}i_{1!}G$, d'où par
composition une flèche $\beta$ de $R^{1}\sisi{f}{\bar{f}}_{2\ast}i_{2!}G$ dans
$R^{1}\sisi{f}{\bar{f}}_{1\ast}i_{1!}G$. Raisonnant comme en \hyperref[XVII.par3.3.2.1]{3.3.2.1}, on voit que le
problème est de prouver que cette flèche est toujours un isomorphisme
lorsque $G$ est \iInd-fini.

Si\pageoriginale on applique le théorème de changement de base propre
à $p$ et au faisceau d'ensembles $i_{1!}G$, on vérifie que
$p_{\ast}i_{1!}G=i_{2!}G$, d'où la suite exacte
\[
0\rightarrow R^{1}
\sisi{f}{\bar{f}}_{2\ast}(i_{2!}G)\xrightarrow{\beta}R^{1}\sisi{f}{\bar{f}}_{1\ast}(i_{1!}G)\rightarrow
\bar{f}_{2\ast}R^{1}p_{\ast}i_{1!}G{\quoi}.
\]
La flèche $\beta$ est injective (car $\alpha$ l'est). Si on applique
le théorème de changement de base propre à $p$ et au faisceau en
groupes $G$, supposé \iInd-fini, on trouve que
$R^{1}p_{\ast}i_{1!}G=0$. Donc $\beta$ est un épimorphisme, ce qui
achève de démontrer \hyperref[XVII.var5.1.12]{5.1.12}.

\begin{enonce*}{Définition 5.1.13}\etiquette[5.1.13]{XVII.def5.1.13}
Pour tout site $X$, on désigne par $D_{\tors}(X)$ la sous-catégorie
pleine de la catégorie dérivée de la catégorie des faisceaux abéliens
sur $X$, formée des complexes dont les faisceaux de cohomologie sont
de torsion.
\end{enonce*}

\begin{enonce*}{Variante 5.1.14}\etiquette[5.1.14]{XVII.var5.1.14}
Dans la catégorie $(S)$, on peut d'une, et essentiellement d'une seule
façon, définir pour toute flèche $f:X\rightarrow Y$ un foncteur
triangulé et « way out »  \emph{\cite[I 7]{SGA4-17:RD}}
\[
Rf_{!}:D^{+}_{\tors}(X)\longrightarrow D^{+}_{\tors}(Y)
\]
et pour tout morphisme composé $f=gh$ des isomorphismes de
transitivité vérifiant des conditions \mbox{\rm (i) (ii) (iii)}
analogues à celles de \hyperref[XVII.subsubsec5.1.4]{5.1.4}.
\end{enonce*}

\begin{enonce*}{Variante 5.1.\sisi{}{1}5}\etiquette[5.1.15]{XVII.var5.1.15}
Plaçons-nous dans la catégorie des $S$-schémas annelés par des
faisceaux de torsion, les flèches étant les morphismes de schémas
annelés induisant (\hyperref[XVII.subsec0.8]{0.8}) un morphisme $S$-compactifiable de
schémas. Pour toute flèche
\[
f:(X,\Aa)\longrightarrow (Y,\Bb),
\]
soit\pageoriginale $Rf_{!}$ le composé du foncteur d'oubli de la catégorie des
$\Aa$-modules dans celle des $f^{\sisi{\ast}{-1}}\Bb$-modules, et
de $Rf_{!}$. Ces foncteurs jouissent des propriétés de variance des
foncteurs $Rf_{!}$ usuels.
\end{enonce*}

On notera toutefois que la notation $Rf_{!}$ n'est raisonnable ici que
si $\Aa$ est fini sur $f^{\sisi{\ast}{-1}}\Bb$, ainsi qu'on s'en
rend compte en prenant $X=Y=\Spec(k)$, $k$ algébriquement clos.

\setcounter{subsubsection}{15}
\subsubsection{}\etiquette[5.1.16]{XVII.subsubsec5.1.16}
Soient $f:X\rightarrow S$ un morphisme compactifiable, $j:U\rightarrow
X$ un ouvert de $X$, $i:F\rightarrow X$ le fermé complémentaire,
$\Aa_S$ un faisceau d'anneaux de torsion sur $S$,
$\Aa_X$ son image réciproque sur $X$ et $K\in \Ob
D(X,\Aa_X)$.

La suite exacte de complexes
\[
0\longrightarrow j_{!}j^{\ast}K\longrightarrow K\longrightarrow
i_{!}i^{\ast}K\longrightarrow 0
\]
donne naissance à un triangle distingué dans $D(X,\Aa_X)$
\textsc{Verdier} \cite[\sisi{I 2.4}{II 1.1.5}]{SGA4-17:CD}. Son image par $Rf_{!}$ est un triangle distingué
(\hyperref[XVII.thm5.1.8]{5.1.8} (b))
$$\etiquette[5.1.16.1]{XVII.eq5.1.16.1}
R(fj)_{!}j^{\ast}K\longrightarrow Rf_{!}K\longrightarrow
R(fi)_{!}i^{\ast}K\longrightarrow R(fj)_{!}j^{\ast}K[1]\ldots\leqno{(5.1.16.1)}
$$
qui définit une suite exacte longue de cohomologie
$$\etiquette[5.1.16.2]{XVII.eq5.1.16.2}
\ldots \sisi{}{\to}R^{q}(fj)_{!}j^{\ast}K\rightarrow
R^{q}f_{!}K\rightarrow R^{q}(fi)_{!}i^{\ast}K \xrightarrow{j} R^{q+1}(fj)_{!}j^{\ast}K\rightarrow\ldots.\leqno{(5.1.16.2)}
$$

Dans le cas particulier où $S$ est spectre d'un corps séparablement
clos et où $K$ est réduit à un faisceau $𝒢$ placé en degré $0$, la
suite \eqhyperref[XVII.eq5.1.16.2]{5.1.16.2} s'écrit
$$\etiquette[5.1.16.3]{XVII.eq5.1.16.3}
\ldots H^{q}_c(U,𝒢)\rightarrow H^{q}_c (X,𝒢)\rightarrow
H^{q}_c(F,𝒢)\xrightarrow{j} H^{q+1}_c(U,𝒢)\ldots\leqno{(5.1.16.3)}
$$

\subsubsection{}\etiquette[5.1.17]{XVII.subsubsec5.1.17}

Soit\pageoriginale un diagramme de schémas $S$-compactifiables
\[
\xymatrix{
X \ar[rr]^{u} \ar[rd]_{f}& & Y\ar[ld]^{g}\\
  & S &,
}
\]
avec $u$ \emph{propre}. On annèle $X$ et $Y$ par l'image réciproque
d'un faisceau d'anneaux de torsion $\Aa_S$ sur $S$. Pour $K\in
D(Y,\Aa_Y)$, la flèche évidente (d'adjonction) $K\rightarrow
Ru_{\ast}u^{\ast}K$ définit par application du foncteur $Rg_{!}$ une
flèche
$$\etiquette[5.1.17.1]{XVII.eq5.1.17.1}
u^{\ast}:Rg_{!}K\longrightarrow
Rg_{!}Ru_{\ast}u^{\ast}K=Rg_{!}Ru_{!}u^{\ast}K\sisi{\eqsim}{\simeq} Rf_{!}u^{\ast}K\leqno(5.1.17.1)
$$
(\emph{contravariance de la cohomologie à support propre vis-à-vis des
morphismes} propres).

Pour $S$ spectre d'un corps algébriquement clos et $K$ réduit à un
faisceau $𝒢$ placé en degré $0$, cette flèche s'écrit
$$\etiquette[5.1.17.2]{XVII.eq5.1.17.2}
u^{\ast}:H^{q}_c(Y,𝒢)\longrightarrow H^{q}_c(X,𝒢)\leqno(5.1.17.2)
$$

\subsection{Le théorème de changement de base}\etiquette[5.2]{XVII.subsec5.2}

\subsubsection{}\etiquette[5.2.1]{XVII.subsubsec5.2.1}

Soient $f:X\rightarrow Y$ un morphisme de topos, $V$ un ouvert de $Y$,
$U=f^{-1}(V)$ son image réciproque dans $X$, $\Aa$ un faisceau
d'anneaux sur $Y$, $\Aa'$ un faisceau d'anneaux sur $X$ et $L$
un complexe borné supérieurement de $(\Aa',
f^{\ast}\Aa)$-bimodules:
$$\etiquette[5.2.1.1]{XVII.eq5.2.1.1}
\vcenter{
\xymatrix{
U \ar[rr]^{f'} \ar@{^{(}->}[d]^{j'}& & V \ar@{^{(}->}[d]^{j}\\
X \ar[rr]_{f} & & Y.
}}\leqno{(5.2.1.1)}
$$

On\pageoriginale dispose alors d'un isomorphisme de changement de base
\eqhyperref[XVII.eq5.1.2.1]{5.1.2.1} entre foncteurs de $D(V,\Aa)$ dans
$D(X,f^{\ast}\Aa):f^{\ast}j_{!}\isomto
j'_{!}{f'}^{\ast}$. De plus, on dispose d'un isomorphisme évident
entre foncteurs de $D^{-}(U,{f'}^{\ast}\Aa)$ dans
$D^{-}(X,\Aa'):L\overset{\LL}{\otimes}_{f^{\ast}\Aa}j'_{!}K\simeq
j'_{!}({j'}^{\ast}L\overset{\LL}{\otimes}_{f^{\ast}\Aa}K)$.
Pour le construire, il suffit de noter que si un complexe $K$ est
déployé pour le foncteur
${j'}^{\ast}L\overset{\LL}{\otimes}_{f^{\ast}\Aa}$,
alors le complexe $j'_{!}K$ est déployé pour le foncteur
$L\overset{\LL}{\otimes}_{f^{\ast}\Aa}$,
et qu'on dispose d'un isomorphisme du type précédent au niveau des
complexes. Composant ces deux isomorphismes, on obtient un
isomorphisme de changement de base entre foncteurs de
$D^{-}(V,\Aa)$ dans $D^{-}(X,\Aa')$:
$$\etiquette[5.2.1.2]{XVII.eq5.2.1.2}
L\overset{\LL}{\otimes}_{f^{\ast}\Aa}f^{\ast}j_{!}K\isomto j'_{!}({j'}^{\ast}L\overset{\LL}{\otimes}_{{f\sisi{'}{}}^{\ast}\Aa}{f'}^{\ast}K){\quoi}.\leqno(5.2.1.2)
$$

Si $L$ est borné à composante de \ttor-dimension finie sur
$f^{\ast}\Aa$,\nde{cette condition est superflue. \Cf \nderef{XVII.ndeKS}.} cet isomorphisme est encore défini pour $K\in
\Ob D(V,\Aa)$.

\subsubsection{}\etiquette[5.2.2]{XVII.subsubsec5.2.2}

Soient $f:X\rightarrow S$ et $g:S'\rightarrow S$ des morphismes de
schémas, et supposons que $f$ soit le composé d'une immersion ouverte
$j$ et d'un morphisme propre $p$, dont la dimension des fibres soit
bornée par un entier $n$ (automatique pour $S$ quasi-compact). Ces
morphismes s'insèrent dans un diagramme à carrés cartésiens
\[
\xymatrix{
X' \ar[rrrr]_{g'} \ar@{^{(}->}[rd]^{j'} \ar[dd]_{f'}& &         &   &
X\ar@{^{(}->}[rd]^{j} \ar[dd]_(.35){f} & \\
             &\bar{X}'\ar[rrrr]|(.75)\hole_{\bar{g}}\ar[ld]^{p'}  &&   &   & \bar{X}\ar[ld]^{p}\\
S' \ar[rrrr]_{g}  & & &&  S &.
}
\]
Soient\pageoriginale $\Aa$ un faisceau d'anneaux de torsion sur $S$,
$\Aa'$ un faisceau d'anneaux de torsion sur $S'$ et $L$ un complexe
borné supérieurement de $(\Aa',g^{\ast}\Aa)$-bimodules. On annèle $X$
et $\bar{X}$ (\resp $X'$ et $\bar{X}'$) par les images réciproques
$\Aa_X$ et $\Aa_{\bar{X}}$ (\resp $\Aa'_{X'}$ et $\Aa'_{\bar{X}'}$) de
$\Aa$ (\resp $\Aa'$) sur $X$ et $\bar{X}$ (\resp $X'$ et
$\bar{X}'$). On dispose alors d'isomorphismes de changement de base
\eqhyperref[XVII.eq5.2.1.2]{5.2.1.2} entre les foncteurs
${p'}^{\ast}L\overset{\LL}{\otimes}
\bar{g}^{\ast}j_{!}K$ et $j'_{!}({f'}^{\ast}L\sisi{\otimes}{\overset{\LL}{\otimes}} {g'}^{\ast}K)$ de
$D^{-}(X,\Aa_X)$ dans $D^{-}(\bar{X}', \Aa'_{\bar{X}'})$, et
d'isomorphismes de changement de base (\hyperref[XVII.thm4.3.1]{4.3.1}) entre les
foncteurs $L\overset{\LL}{\otimes}g^{\ast}Rp_{\ast}K$
et $Rp'_{\ast}({p'}^{\ast}L\overset{\LL}{\otimes}\bar{g}^{\ast}K)$
de $D^{-}(\bar{X}, \Aa_{\bar{X}})$ dans $D^{-}(S',\Aa')$. Par
composition, on en déduit un isomorphisme de changement de base entre
foncteurs de $D^{-}(X,\Aa_X)$ dans $D^{-}(S',\Aa')$
$$\etiquette[5.2.2.1]{XVII.eq5.2.2.1}
L\overset{\LL}{\otimes}g^{\ast}Rp_{\ast}j_{!}K\isomto
Rp'_{\ast}
j'_{!}({f'}^{\ast}L\overset{\LL}{\otimes}{g'}^{\ast}K){\quoi}.\leqno(5.2.2.1)
$$

Lorsque $L$ est borné à composantes de \ttor-dimension finie, les
arguments précédents et \hyperref[XVII.var4.3.2]{4.3.2} (i) permettent de définir l'isomorphisme
de changement de base pour $K\in \Ob D(X,\Aa_X)$.\nde{en combinant la N.D.E. précédente et \nderef{XVII.ndecb}, on peut définir cet isomorphisme sans condition sur $L$ ni l'hypothèse que la dimension des fibres de $p$ soit bornée.}

Pour pouvoir référer à \hyperref[XVII.thm5.1.8]{5.1.8}, supposons que $S$ soit quasi-compact
quasi-séparé. Les foncteurs $Rp_{\ast}j_{!}=Rf_{!}$ et
$Rp'_{\ast}j'_{!}=Rf'_{!}$ ne dépendent alors que de $f$ et $f'$, ce
qui donne un sens au

\begin{enonce*}{Lemme 5.2.3}\etiquette[5.2.3]{XVII.lem5.2.3}
L'isomorphisme de changement de base \eqhyperref[XVII.eq5.2.2.1]{5.2.2.1} ne dépend pas de la
compactification choisie $f=pj$ de $X$.
\end{enonce*}

Soit un morphisme de compactification $q:\bar{X}_{1}\rightarrow
\bar{X}_2$
$$\etiquette[5.2.3.1]{XVII.eq5.2.3.1}
\vcenter{
\xymatrix{
X \ar[rr]^{j_1} \ar[rrd]|(.65)\hole^{j_2} \ar[dd]_{f}& & \bar{X}_1 \ar[d]_{q} \ar[lldd]^{p_1}\\
  & &  \bar{X}_2 \ar[lld]^{p_2}\\
S & & .
}}\leqno{(5.2.3.1)}
$$
L'isomorphisme\pageoriginale de changement de base $c_{1}$ relatif à
$\bar{X}_{1}$ s'obtient en « composant »  ceux relatifs à $j_{1}$
et $p_{1}$, ou encore en composant ceux relatifs à $j_{1}$, $q$ et
$p_{2}$. L'isomorphisme $c_2$ relatif à $\bar{X}_2$ s'obtient en
composant ceux relatifs à $j_{2}$ et $p_{2}$. Pour vérifier que
$c_{1}=c_{2}$, il suffit de vérifier que l'isomorphisme $c'_{1}$
relatif à $j_{2}$ s'obtient comme le composé $c'_{2}$ de ceux relatifs
à $j_{1}$ et $q$. Désignons par un $'$ les flèches déduites par le
changement de base $g$ des flèches \eqhyperref[XVII.eq5.2.3.1]{5.2.3.1}; $c'_{1}$ et $c'_{2}$
sont des isomorphismes
\[
c'_{1}, c'_{2}:L\overset{\LL}{\otimes}\bar{g}_{2}^{\ast}
j_{2!}K\isomto  j'_{2!}({j'}_{2}^{\ast}
L\overset{\LL}{\otimes} {g'}^{\ast}K){\quoi}.
\]
D'après \eqhyperref[XVII.eq5.1.3.2]{5.1.3.2}, pour vérifier que $c'_{1}=c'_{2}$, il suffit de le
vérifier après s'être restreint à $\sisi{X}{\bar{X}}'_{2}$, ce qui est trivial.

Les isomorphismes de changement de base relatifs à deux
compactifications entre lesquelles il existe un morphisme sont donc
égaux, et on conclut par \hyperref[XVII.prop3.2.6]{3.2.6} (ii).

\begin{enonce*}{Lemme 5.2.4}\etiquette[5.2.4]{XVII.lem5.2.4}
Soit $f=uv:X\rightarrow Y \rightarrow S$ un composé de morphismes
$S$-compactifiables. Les isomorphismes de changement de base
\eqhyperref[XVII.eq5.2.2.1]{5.2.2.1} relatifs à $u$, $v$ et $f$ donnent lieu à un diagramme
commutatif (notations de \eqhyperref[XVII.eq5.2.4.1]{5.2.4.1} ci-dessous):
\[
\xymatrix{
L\overset{\LL}{\otimes} g^{\ast} R\sisi{uv}{(uv)}_{!} K
\ar@<-4pt>[rr]^{\sim} \ar@{=}[d]& & R\sisi{u'v'}{(u'v')}_{!}((u'v')^{\ast}L
\overset{\LL}{\otimes}g''^{\ast}K)\ar@{=}[d]\\
L\overset{\LL}{\otimes} g^{\ast}Ru_{!}Rv_{!}K \ar@<-4pt>[r]^(.45){\sim}& Ru'_{!}(u'^{\ast}L
\overset{\LL}{\otimes}g'^{\ast}Rv_{!}\sisi{}{K}) \ar@<-4pt>[r]^(.45){\sim}& Ru'_{!}Rv'_{!}(v'^{\ast}u'^{\ast}L\overset{\LL}{\otimes}g''^{\ast}K).
}
\]
\end{enonce*}

Choisissons\pageoriginale une compactification de $(u,v)$ (\hyperref[XVII.def3.2.5]{3.2.5}):
$$\etiquette[5.2.4.1]{XVII.eq5.2.4.1}
\vcenter{
\xymatrix{
X' \ar[rr]_{g''} \ar[d]_{\sisi{u}{v}'}& & X \ar@{^{(}->}[r]^{j_1} \ar[d]_{\sisi{u}{v}}& \bar{X}
\ar[ld]_{p_1}\ar@{^{(}->}[r]^{j_2} \ar[d]& \bar{\bar{X}} \ar[ld]^{p_3}& \\
Y' \ar[rr]_{g'} \ar[d]_{\sisi{v}{u}'}& & Y \ar@{^{(}->}[r]^{j_3}\ar[d]_{\sisi{v}{u}}& \bar{Y}\ar[ld]^{p_2} &  &\\
S' \ar[rr]_{g} &  & S &      &  .
}}\leqno{(5.2.4.1)}
$$

Le lemme \hyperref[XVII.lem5.2.4]{5.2.4} exprime que l'isomorphisme de changement de base
$c_{1}$ relatif à $f$ est le composé $c_{2}$ de ceux relatifs à $u$ et
$v$. L'isomorphisme $c_{1}$ s'obtient en composant les isomorphismes
de changement de base relatifs à $j_{1}$, $j_{2}$, $p_{3}$ et
$p_{2}$. Pour $c_{2}$, on compose ceux relatifs à $j_{1}$, $p_{1}$,
$j_{3}$, $p_{2}$. Il suffit donc de comparer les isomorphismes
$c'_{1}$, composé de ceux relatifs à $p_{3}$ et $j_{2}$, et $c'_{2}$,
composé de ceux relatifs à $j_{3}$ et $p_{1}$.

Soit le diagramme cartésien
\[
\xymatrix{
  \bar{X}' \ar[ddd]^{p'_1} \ar[rrrr]^{\bar{g}''} \ar@{^{(}->}[rrd]^{j'_2} &    &       &    & \bar{X}\ar@{^{(}->}[rrd]^{j_2} \ar[ddd]^{\sisi{}{p_1}}|(.38)\hole\\
                          & &\bar{\bar{X}}' \ar[rrrr]_(.3){\bar{\bar{g}}} \ar[ddd]^{p'_3}&   &&
& \bar{\bar{X}} \ar[ddd]^{\sisi{}{p_3}}\\
 & & &\\
  Y'  \ar@{^{(}->}[rrd]_{j'_3}\ar[rrrr]^(.7){g'}|(.52)\hole \ar[rdd]_{\sisi{v}{u}'} &    &&
& Y\ar[rdd]_(.65){\sisi{v}{u}} \ar@{^{(}->}[rrd]^{j_3}\\
  & &\bar{Y}' \ar[rrrr]_{\bar{g}'} \ar[ld]^{p'_2}&        &  & &\bar{Y}\ar[ld]^{p_2}\\
 &S'  \ar[rrrr]_{g} &&& & S &,
}
\]
$L'={p'_2}^{\ast}L$\sisi{}{, $K'=j_{1!}K$}\pageoriginale et
\[
c'_{1}, c'_{2}: L'\overset{\LL}{\otimes} {\bar{g}'}^{\ast}
R(j_{3}p_{1})_{!} K\sisi{}{'}\isomto  R(j'_{3} p'_{1})_{!} ((j'_{3}
p'_{1})^{\ast} L'\otimes {\bar{g}''}^{\ast}K\sisi{}{'}){\quoi}.
\]
En vertu de \eqhyperref[XVII.eq5.1.3.2]{5.1.3.2}, pour vérifier que $c'_{1}=c'_{2}$, il suffit de
vérifier que ${j'_{3}}^{\ast}(c'_{1})={j'_{3}}^{\ast}(c'_{2})$, ce qui
est trivial.

\begin{enonce*}{Lemme 5.2.5}\etiquette[5.2.5]{XVII.lem5.2.5}
Soient des morphismes de schémas annelés par des faisceaux d'anneaux
de torsion
\[
(S'', \Aa'')\xrightarrow{g_{2}} (S',
\Aa')\xrightarrow{g_{1}} (S,\Aa){\quoi},
\]
et $f:X\rightarrow S$ un morphisme compactifiable, d'où un diagramme:
\[
\xymatrix{
X'' \ar[r]_{g'_2} \ar[d]^{f''}& X' \ar[r]_{g'_1}\ar[d]^{f'}& X \ar[d]^{f}\\
S'' \ar[r]_{g_2}& S' \ar[r]_{g_1}& S \rlap{.}
}
\]
On\pageoriginale annèle $X$, $X'$, $X''$ par $f^{\ast}\Aa$,
${f'}^{\ast}\Aa'$ et ${f''}^{\ast}\Aa''$. Alors, le
diagramme
\[
\xymatrix{
L(g_1g_2)^{\ast}Rf_{!}  \ar[rr]^{\sim} \ar@{=}[d]&& Rf''_{!}L(g'_1g'_2)^{\ast}\ar@{=}[d]\\
Lg_{\sisi{1}{2}}^{\ast}Lg^{\ast}_{\sisi{2}{1}} R f_{!} \ar[r]^{\sim}& Lg^{\ast}_{\sisi{1}{2}} Rf'_{!}
L\sisi{g^\ast_2}{g'^\ast_1}  \ar[r]^{\sim} & Rf''_{!} Lg'^{\ast}_{\sisi{1}{2}}  L g'^{\ast}_{\sisi{2}{1}}
}
\]
est commutatif.
\end{enonce*}

Ce lemme se déduit aussitôt des résultats analogues pour les foncteurs
$j_{!}$ et $Rp_{\ast}$ pour $p$ propre. On laisse au lecteur le soin
de le moduler.

On résume les énoncés qui précédent en le

\begin{enonce*}{Théorème 5.2.6}\etiquette[5.2.6]{XVII.thm5.2.6}
Le foncteur $Rf_{!}$ commute aux changements de base.
\end{enonce*}

\begin{enonce*}[remark]{Variantes 5.2.7}\etiquette[5.2.7]{XVII.var5.2.7}
\begin{enumerate}
\renewcommand{\theenumi}{\roman{enumi}}
\renewcommand{\labelenumi}{(\theenumi)}
\item Les foncteurs $f_{!}$ (pour les faisceaux d'ensembles pointés:
\hyperref[XVII.var5.1.11]{5.1.11}) et les foncteurs $R^{1}f_{!}$ (pour les faisceaux en groupes:
        \hyperref[XVII.var5.1.12]{5.1.12}) commutent aux changement de base.

\item Les foncteurs $Rf_{!}$ (\hyperref[XVII.var5.1.14]{5.1.14}) commutent aux changements de
base.

\item Sous les hypothèses (\hyperref[XVII.subsubsec5.2.2]{5.2.2}), si $L$ est
borné et à composantes de \ttor-dimension finie,\nde{comme on a remarqué en \hyperref[XVII.subsubsec5.2.2]{5.2.2}, ces hypothèses sont superflues.} les foncteurs
$Rf_{!}:D(X)\rightarrow D(S)$, commutent aux changements de base.
\end{enumerate}
\end{enonce*}

Le\pageoriginale théorème \hyperref[XVII.thm5.2.6]{5.2.6} est essentiellement équivalent à la
conjonction des deux corollaires \hyperref[XVII.prop5.2.8]{5.2.8} et \hyperref[XVII.prop5.2.9]{5.2.9} suivants.

\begin{enonce*}{Proposition 5.2.8}\etiquette[5.2.8]{XVII.prop5.2.8}
Soient $f:X\rightarrow S$ un morphisme compactifiable, et
$\mathcal{F}$ un faisceau de torsion sur $X$. Pour tout point
géométrique $s$ de $S$, soient $X_s$ la fibre géométrique de $X$ en
$s$ et $\mathcal{F}_s$ le faisceau induit. On a canoniquement
\[
(R^{q}f_{!}(\mathcal{F}))_s=H^{q}_c(X_s,\mathcal{F}_s){\quoi}.
\]
\end{enonce*}

\begin{enonce*}{Corollaire 5.2.8.1}\etiquette[5.2.8.1]{XVII.cor5.2.8.1}
Si les fibres de $f$ sont de dimension $\leq d$, le foncteur $Rf_{!}$
est de dimension cohomologique $\leq 2d$ (\hyperref[XVII.def1.2.9]{1.2.9}): pour tout faisceau
de torsion $\sisi{F}{ℱ}$, on a $R^{i}f_{!}\sisi{F}{ℱ}=0$ pour $i>2d$. Le foncteur
$R^{2d}f_{!}$ est donc exact à droite sur la catégorie des faisceaux
de torsion.
\end{enonce*}

On se ramène par changement de base (\hyperref[XVII.prop5.2.8]{5.2.8}) au cas où $S$ est le
spectre d'un corps algébriquement clos. Si alors $\bar{X}$ est un
schéma propre sur $S$ contenant $X$ comme ouvert de Zariski dense,
$j:X\rightarrow \bar{X}$, on a $\dim(\bar{X})=\dim(X)\leq \sisi{D}{d}$. Pour
$ℱ$ faisceau de torsion sur $X$ et $i>2d$, on a donc, par \hyperrefext[X][X.cor4.3]{4.3},
\refnoexpose{X.cor5.2}
\[
H^{i}_c(X,ℱ)=H^{i}(\bar{X},j_{!}\sisi{F}{ℱ})=0{\quoi}.
\]

\begin{enonce*}{Proposition 5.2.9}\etiquette[5.2.9]{XVII.prop5.2.9}
Soient $f:X\rightarrow S$ un morphisme compactifiable, $\Aa$
un faisceau d'anneaux de torsion sur $S$, $K$ un complexe borné
supérieurement de $f^{\ast}\Aa$-modules à gauche et $L$ un
complexe borné supérieurement de $\Aa$-modules à droite. On a
canoniquement
\[
L\overset{\LL}{\underset{\Aa}{\otimes}}
Rf_{!}(K)=Rf_{!}(f^{\ast}L\overset{\LL}{\underset{f^{\ast}\Aa}{\otimes}}K){\quoi}.
\]
\end{enonce*}

Cette\pageoriginale proposition implique formellement le résultat
suivant:

\begin{enonce*}{Théorème 5.2.10}\etiquette[5.2.10]{XVII.thm5.2.10}
Soient $f:X\rightarrow S$ un morphisme compactifiable, $\Aa$
un faisceau d'anneaux de torsion sur $S$, $K\in
D^{-}(X,f^{\ast}\Aa)$ et $d\in \ZZ$. Si $K$ est de
\tTor-dimension finie $\leq d$ (\hyperref[XVII.def4.1.9]{4.1.9}), alors $Rf_{!}K$ est de
\tTor-dimension finie $\leq d$.
\end{enonce*}

Il suffit de vérifier que pour tout faisceau de $\Aa$-modules
à droite $L$ sur $S$, on a
\[
\sisi{H}{\Hh}^{-i} (L\overset{\LL}{\otimes}Rf_{!}K)=0\quad\mbox{\rm
pour}\quad i>d{\quoi}.
\]
D'après \hyperref[XVII.prop5.2.9]{5.2.9}, on sait que
\[
\sisi{H}{\Hh}^{-i}(L\overset{\sisi{L}{\LL}}{\otimes}Rf_{!}K)=R^{-i}f_{!}(f^{\ast}L\overset{\sisi{L}{\LL}}{\otimes}K){\quoi},
\]
de plus, par hypothèse,
\[
\sisi{H}{\Hh}^{-i}(f^{\ast}L\overset{\sisi{L}{\LL}}{\otimes}K)=0\quad\mbox{\rm pour}\quad
i>d{\quoi},
\]
et on conclut en remarquant que si un complexe $M$ de faisceaux de
modules sur $X$ satisfait à $\sisi{H}{\Hh}^{-i}(M)=0$ pour $i>d$, alors on a
encore $\sisi{H}{\Hh}^{-i}(Rf_{!}M)=0$ pour $i>d$.

Dans le cas particulier des faisceaux d'anneaux constants et
noethériens, on a le résultat plus général:

\begin{enonce*}{Théorème 5.2.11}\etiquette[5.2.11]{XVII.thm5.2.11}
Soient $f:X\rightarrow S$ un morphisme de sites, $A$ un anneau, $C$
son centre et $d\in\ZZ$. On suppose que
\begin{enumerate}
\renewcommand{\theenumi}{\roman{enumi}}
\renewcommand{\labelenumi}{(\theenumi)}
\item le foncteur $Rf_{\ast}:D^{+}(X,C)\rightarrow D^{+}(S,C)$ est de
dimension cohomologique finie.

\item $A$ est noethérien à droite.

\item $S$ a assez de points.
\end{enumerate}
Pour\pageoriginale tout complexe $K\in \Ob D^{-}(X,A)$ de faisceaux de
$A$-modules à gauche, si $K$ est de \tTor-dimension finie $\leq d$,
alors le complexe $Rf_{\ast}K\in \Ob D^{-}(S,A)$ est de \tTor-dimension
$\leq d$.
\end{enonce*}

\begin{enonce*}[remark]{Preuve}
En vertu de (i), le foncteur $f_{\ast}$ est dérivable en des foncteurs
\[
Rf_{\ast}:D(X,A)\longrightarrow D(S,A)
\]
et
\[
Rf_{\ast}:D(X,C)\longrightarrow D(S,C){\quoi}.
\]
\end{enonce*}

D'après \hyperref[XVII.lem4.2.8]{4.2.8}, il suffit de vérifier que pour tout point $s$ de $S$,
le complexe $(Rf_{\ast}K)_s$ de $A$-modules est de \tTor-dimension $\leq
d$; il suffit pour cela que pour tout $A$-module à droite $L$ de type
fini on ait
\[
H^{-i}(L\overset{L}{\otimes}Rf_{\ast}K)=0\quad\mbox{\rm pour}\quad
i>d{\quoi}.
\]

Raisonnant comme en \hyperref[XVII.thm5.2.10]{5.2.10}, on voit qu'il suffit pour cela que la « flèche de changement de base »
$$\etiquette[5.2.11.1]{XVII.eq5.2.11.1}
L\overset{L}{\otimes}Rf_{\ast}K\longrightarrow Rf_{\ast}(L\overset{L}{\otimes}K)\leqno(5.2.11.1)
$$
soit un isomorphisme.

Soit $L^{\ast}$ une résolution de $L$ par des $A$-modules libres de
type fini. Les foncteurs considérés étant way-out à droite, il suffit
de prouver que \eqhyperref[XVII.eq5.2.11.1]{5.2.11.1} est un isomorphisme après avoir remplacé $L$
par une quelconque composante de $L\sisi{_{\ast}}{^\ast}$, et on est ramené au cas
trivial où $L$ est libre de type fini.

\begin{enonce*}[remark]{Remarques 5.2.12}\etiquette[5.2.12]{XVII.rems5.2.12}
\begin{enumerate}
\renewcommand{\theenumi}{\roman{enumi}}
\renewcommand{\labelenumi}{(\theenumi)}
\item L'hypothèse\pageoriginale (iii) de \hyperref[XVII.thm5.2.11]{5.2.11} est sans doute
inutile.

\item Si on omet l'hypothèse noethérienne sur $A$, \hyperref[XVII.thm5.2.11]{5.2.11} reste
valable lorsque les foncteurs $R^{i}f_{\ast}$ commutent aux limites
inductives filtrantes.

\item L'hypothèse que $S$ soit annelé par un faisceau d'anneaux
\emph{constant} est essentielle. On obtiendrait un contre-exemple en
prenant pour $S$ le spectre d'un anneau de valuation discrète
strictement local, pour $f$ l'inclusion du point générique $\eta$,
pour faisceau d'anneaux $\Aa$ le faisceau valant
$\ZZ/n\ZZ$ au point générique et
$\ZZ/n^{2}\ZZ$ au point fermé $s$ ($n>1$ premier aux
caractéristiques résiduelles) et pour faisceau le module
$f^{\ast}\Aa$.
\end{enumerate}
\end{enonce*}

\subsection{Théorème de comparaison et théorème de finitude}\etiquette[5.3]{XVII.subsec5.3}

\subsubsection{}\etiquette[5.3.1]{XVII.subsubsec5.3.1}
À un schéma $X$ de type fini sur $\CC$ sont associés son site
étale $X_{\text{\rm et}}$, son site transcendant $X_{\cl}$ et un morphisme de
sites (\hyperrefext[XI][XI.sec4]{4})
\[
\epsilon:X_{\cl}\longrightarrow X_{\et}{\quoi}.
\]

Si $j:U\rightarrow X$ est une immersion ouverte, le diagramme
commutatif
$$
\xymatrix{
U_{\cl}  \ar[r]\ar[d]   & U_{\et} \ar[d]\\
X_{\cl}  \ar[r] & X_{\et}
}
$$
identifie $U^{\sim}_{\cl}$ à l'ouvert de $X^{\sim}_{\cl}$ image
réciproque de l'ouvert $U_{\widetilde{\text{\rm et}}}$ de
$X_{\widetilde{\text{\rm et}}}$.

Si $p:Y\rightarrow X$ est un morphisme propre, l'application continue
$p^{\an}:Y^{\an}\rightarrow X^{\an}$ est propre et séparée. Pour le
voir, on se ramène par le lemme de Chow au cas $p$ projectif.

Soit\pageoriginale $f:X\rightarrow S$ un morphisme compactifiable de
schémas de type fini sur $\CC$ et soit $f=pj$ une
compactification de $f$.
$$\etiquette[5.3.1.1]{XVII.eq5.3.1.1}
\vcenter{
\xymatrix{
X_{\cl} \ar[rrr]_{\epsilon} \ar[rd]^{j_{\cl}} \ar[dd]_{f_{\cl}}&  & & X \ar[rd]^{j}\ar[dd]_(.6){f}|\hole & \\
            &\bar{X}_{\cl}\ar[ld]^{p_{\cl}}  \ar[rrr]_{\epsilon} &  & &\bar{X}\ar[ld]^{p} \\
S_{\cl}  \ar[rrr]_{\epsilon} &                 &   &  S & .
}}\leqno{(5.3.1.1)}
$$
Le foncteur $p_{\cl{\ast}}j_{\cl !}$ n'est autre que le foncteur
image directe à support propre $f_{\cl !}$ (voir \textsc{Verdier} \cite{SGA4-17:V3}). Si
$S_{\cl}$, donc $X_{\cl}$ est séparé, le foncteur $j_{\cl !}$
transforme faisceaux mous sur $X_{\cl}$ en faisceau mous sur
$\bar{X}_{\cl}$ \cite[II 3.5.5]{SGA4-17:TF}. La question étant locale sur $S$, il
transforme en tout cas faisceaux flasques sur $X_{\cl}$ en faisceaux
sur $\sisi{X}{\bar{X}}_{\cl}$ acycliques pour le foncteur $p_{\cl \ast}$. On a donc
(\cf \textsc{Verdier} \cite{SGA4-17:V3})
$$
Rf_{\cl}=Rp_{\cl \ast} j_{\cl !}{\quoi}.
$$

\subsubsection{}\etiquette[5.3.2]{XVII.subsubsec5.3.2}

D'après \hyperref[XVII.lem5.1.2]{5.1.2}, on dispose d'un isomorphisme de changement de base,
entre foncteurs de $D(X)$ dans $D(\bar{X}_{\cl})$
$$
\epsilon^{\ast} j_{!}\isomto  j_{\cl !}
\epsilon^{\ast}{\quoi}.
$$

D'après\pageoriginale \hyperref[XVII.subsubsec4.2.12]{4.2.12} ou \hyperref[XVII.prop4.1.5]{4.1.5}, on dispose
d'un morphisme de changement de base entre foncteurs de
$D^{+}(\bar{X})$ dans $D^{+}(S_{\cl})$:
$$
\epsilon^{\ast}Rp_{\ast}\longrightarrow Rp_{\cl
\ast}\epsilon^{\ast}{\quoi}.
$$
Par composition, on en déduit un \emph{morphisme de changement de
base, entre foncteurs de} $D^{+}_{\tors}(X)$ dans
$D^{+}_{\tors}(S_{\cl})$ (\hyperref[XVII.def5.1.13]{5.1.13})
$$\etiquette[5.3.2.1]{XVII.eq5.3.2.1}
\epsilon^{\ast}Rf_{!}\longrightarrow Rf_{\cl
!}\epsilon^{\ast}{\quoi}.\leqno(5.3.2.1)
$$

On vérifie comme en \hyperref[XVII.lem5.2.3]{5.2.3} qu'il ne dépend pas de la compactification
choisie de $f$; il vérifie une formule de transitivité du type \hyperref[XVII.lem5.2.4]{5.2.4}.

\begin{enonce*}{Théorème 5.3.3}\etiquette[5.3.3]{XVII.thm5.3.3}
\emph{(Théorème de comparaison)}. Le morphisme \eqhyperref[XVII.eq5.3.2.1]{5.3.2.1} est un isomorphisme.
\end{enonce*}

Par construction, il suffit de traiter le cas où $f$ est propre. Les
foncteurs considérés sont triangulés et « way-out » \cite[I 7]{SGA4-17:RD}. Il
suffit donc de traiter le cas où $K$ est réduit à un faisceau de
torsion placé en degré $0$, et on conclut par \hyperrefext[XVI][XVI.thm4.1]{4.1} (i).

\begin{enonce*}[remark]{Variantes 5.3.4}\etiquette[5.3.4]{XVII.var5.3.4}
\begin{enumerate}
\renewcommand{\theenumi}{\roman{enumi}}
\renewcommand{\labelenumi}{(\theenumi)}
\item Si $\Aa$ est un faisceau d'anneaux de torsion sur
$S_{\text{\rm et}}$, on obtient de même un isomorphisme \eqhyperref[XVII.eq5.3.2.1]{5.3.2.1}
entre foncteurs de $D(X,f^{\ast}\Aa)$ dans $D(S,\epsilon^{\ast}\Aa)$.

\item En bas degré, on trouve les variantes habituelles pour les
faisceaux d'ensembles pointés et les faisceaux en groupes \iInd-finis.
\end{enumerate}
\end{enonce*}

\begin{enonce*}{Corollaire 5.3.5}\etiquette[5.3.5]{XVII.cor5.3.5}
Soient\pageoriginale $X$ un schéma séparé de type fini sur
$\CC$, qui puisse se compactifier en un schéma propre sur
$\CC$, et $F$ un faisceau de torsion sur $X$. On a alors
$$
H^{q}_{!}(X,F)\isomto
H^{q}_c(X_{\cl},\epsilon^{\ast}F){\quoi}.
$$
\end{enonce*}

\begin{enonce*}{Théorème 5.3.6}\etiquette[5.3.6]{XVII.thm5.3.6}
Soient $f:X\rightarrow S$ un morphisme compactifiable de présentation
finie et $A$ un anneau de torsion noethérien à gauche. Le foncteur
$Rf_{!}:D(X,A)\rightarrow D(S,A)$ transforme complexes à cohomologie
constructible en complexes à cohomologie constructible.
\end{enonce*}

L'assertion est locale sur $S$, qu'on peut supposer affine. Par
passage à la limite, on se ramène au cas $S$ noethérien. Les complexes
à cohomologie constructible forment une sous-catégorie triangulée de
$D(S,A)$. Le foncteur $Rf_{!}$ étant « way-out », on se ramène à
vérifier que si $\mathcal{F}$ est un faisceau constructible de
$A$-modules, les faisceaux $R^{q}f_{!}\mathcal{F}$ sont
constructibles. Soit $f=pj$, avec $p$ propre et $j$ une immersion
ouverte. On a
$R^{q}f_{!}\mathcal{F}=R^{q}p_{\ast}(j_{!}\mathcal{F})$. Le faisceau
$j_{!}\mathcal{F}$ est constructible et on conclut par \hyperrefext[XIV][XIV.thm1.1]{1.1}.

\setcounter{subsubsection}{6}
\subsubsection{}\etiquette[5.3.7]{XVII.subsubsec5.3.7}
En bas degrés, on dispose d'énoncés analogues pour les images directes
à support propre d'un faisceau constructible d'ensembles ou de groupes
\iInd-finis.

\begin{enonce*}{Corollaire 5.3.8}\etiquette[5.3.8]{XVII.cor5.3.8}
Soient $X$ un schéma séparé de type fini sur un corps séparablement
clos $k$ et $\Ff$ un faisceau constructible de torsion sur
$X$. On suppose que $X$ puisse se compactifier en un schéma propre sur
$k$.\nde{cette hypothèse est automatiquement vérifiée. Voir \nderef{XVII.ndeConrad}.} Alors, les groupes $H^{q}_c(X,\mathcal{F})$ sont finis.
\end{enonce*}

\subsection{Formule de Künneth}\etiquette[5.4]{XVII.subsec5.4}\pageoriginale

\subsubsection{}\etiquette[5.4.1]{XVII.subsubsec5.4.1}
Soit un diagramme commutatif (\hyperref[XVII.subsec0.11]{0.11}) de sites
$$\etiquette[5.4.1.1]{XVII.eq5.4.1.1}
\vcenter{
\xymatrix@R=.5cm@C=1cm{
             &  Z'\ar[ld]_{p'}\ar[rd]^{q'} \ar[dd]^{h}&    \\
X' \ar[dd]_{f}&                   &  Y' \ar[dd]^{g}\\
             &       Z\ar[ld]_{p} \ar[rd]^{q}   &    \\
X \ar[rd]           &                    & Y \ar[ld]\\
             &      S             & &.
}}\leqno{(5.4.1.1)}
$$
On suppose $S$ annelé par un faisceau d'anneaux $\Aa$, de
centre $\Cc$; on désignera encore par $\Aa$ et
$\Cc$ les images réciproques de ces faisceaux sur les sites de
\eqhyperref[XVII.eq5.4.1.1]{5.4.1.1}.

Si $K$ est un faisceau de $\Aa$-modules à droite sur $X$ et
$L$ un faisceau de $\Aa$-modules à gauche sur $Y$, on pose
\[
K\mathop{\boxtimes}_{\Aa}
L=p^{\ast}\sisi{}{K}\underset{\Aa}{\otimes} q^{\ast}L
\]
(produit tensoriel externe). De même pour $K$ (\resp $L$) sur $X'$
(\resp $Y'$). Soit alors $K$ un faisceau de $\Aa$-modules à
droite sur $X'$ et $L$ un faisceau de $\Aa$-modules à gauche
sur $Y'$. Les flèches de changement de base (\sisi{}{\hyperref[XVII.prop2.1.3]{2.1.3}, }\sisi{\hyperrefext[XII][XII.eq4.1]{4.1}}{\neweqhyperrefext[XII][XII.eq4.1]{4.1}}) définissent
$$\etiquette[5.4.1.2]{XVII.eq5.4.1.2}
\begin{cases}
p^{\ast}f_{\ast}K\longrightarrow h_{\ast}{p'}^{\ast}K\\
q^{\ast}g_{\ast}L\longrightarrow h_{\ast}{q'}^{\ast}L{\quoi},
\end{cases}\leqno(5.4.1.2)
$$
d'où,\pageoriginale par produit tensoriel, une flèche
$$\etiquette[5.4.1.3]{XVII.eq5.4.1.3}
f_{\ast}K\boxtimes_{\Aa} g_{\ast}L\longrightarrow
h_{\ast}{p'}^{\ast}K\otimes_{\Aa} h_{\ast}
{q'}^{\ast}L\longrightarrow h_{\ast} (K\boxtimes_{\Aa}
L){\quoi}.\leqno(5.4.1.3)
$$

Supposons que les foncteurs
\begin{align*}
Rf_{\ast} &: D^{+} (X',\Aa^{\circ})\longrightarrow
D^{+}(X,\Aa^{\circ})\\
Rg_{\ast} &: D^{+}(Y',\Aa)\longrightarrow
D^{+}(Y,\Aa)\\
Rh_{\ast} &: D^{+} (Z', \Cc)\longrightarrow D^{+}
(Z,\Cc)
\end{align*}
soient de dimension cohomologique finie, et se prolongent donc aux
catégories dérivées entières. Il s'impose alors, pour $K\in \Ob
D^{-}(X',\Aa^{\circ})$ et $L\in \Ob D^{-}(Y',\Aa)$, de
définir une \emph{flèche de K\"unneth}
$$\etiquette[5.4.1.4]{XVII.eq5.4.1.4}
Rf_{\ast}K\overset{\LL}{\boxtimes}_{\Aa}
Rg_{\ast}L\longrightarrow Rh_{\ast} (K
\overset{\LL}{\boxtimes}_{\Aa} L){\quoi}.\leqno(5.4.1.4)
$$

On ne la définira que si les sites considérés ont assez de points, de
façon à pouvoir utiliser \hyperref[XVII.prop4.2.10]{4.2.10} et \hyperref[XVII.lem4.2.8]{4.2.8}.
En vertu de \hyperref[XVII.prop4.2.10]{4.2.10}, il
suffit de la définir pour un couple $(K,L)$ déployé à gauche pour le
foncteur $(K,L)\mapsto K\boxtimes_{\Aa}L$, $K$ et $L$ étant
déployés à droite pour les foncteurs $f_{\ast}$ et $g_{\ast}$ et
bornés supérieurement.

Soient $K_{1}$ une résolution plate de $f_{\ast} K$, $L_{1}$ une
résolution plate de $g_{\ast}L$ et
$u:K\overset{\LL}{\boxtimes}\sisi{}{_\Aa} L\rightarrow M$ une
résolution de $K\overset{\LL}{\boxtimes}\sisi{}{_\Aa}L$,
déployée pour le foncteur $Rh_{\ast}$. On dispose de morphismes
\begin{gather*}
Rf_{\ast}K\overset{\LL}{\boxtimes}_{\Aa}Rg_{\ast}L\simeq
K_{1} \boxtimes_{\Aa} L_{1}\longrightarrow
f_{\ast}K\boxtimes_{\Aa} g_{\ast}L\longrightarrow\\
\longrightarrow h_{\ast}(K\boxtimes_{\Aa} L)\longrightarrow
h_{\ast}M\simeq Rh_{\ast}(K\overset{\LL}{\boxtimes}_{\Aa}L)
\end{gather*}
et on définit \eqhyperref[XVII.eq5.4.1.4]{5.4.1.4} comme leur composé.

Cette\pageoriginale flèche a la propriété caractéristique suivante, de démonstration
toute analogue à celle de \hyperref[XVII.prop4.2.13]{4.2.13}:

\setcounter{paragraph}{4}
\paragraph{}\etiquette[5.4.1.5]{XVII.par5.4.1.5}
Soient $K\in \Ob C^{-}(X', \Aa^{\circ})$ et $L\in \Ob
C^{-}(Y', \Aa)$. Posons $K_{1}=f_{\ast}K$, $L_{1}=g_{\ast}K$,
et $M_{2}=K\boxtimes\sisi{}{_\Aa} L$. Le diagramme suivant est alors commutatif
$$
\xymatrix@R=1.25cm{
Rf_{\ast}K
\overset{\LL}{\boxtimes}_{\Aa}Rg_{\ast}
L \ar[rrr]& & & Rh_{\ast}(K
\overset{\LL}{\boxtimes}_{\Aa}L) \ar[d]\\
K_1\overset{\LL}{\boxtimes}_{\Aa}L_{\sisi{k}{1}}
\ar[r] \ar[u]&
K_1 \boxtimes\sisi{}{_\Aa} L_1 \ar[r]&  h_{\ast}M_2 \ar[r]& Rh_{\ast}M_2.
}
$$

On a une commutativité analogue (\cf \hyperref[XVII.prop4.2.13]{4.2.13}) en partant
plutôt de
morphismes $K_{1}\rightarrow f_{\ast}K$, $L_{1}\rightarrow g_{\ast}K$
et $K\boxtimes\sisi{}{_\Aa} L\rightarrow M_{2}$.

La commutativité \hyperref[XVII.par5.4.1.5]{5.4.1.5} rend en principe aisée la vérification des
propriétés de compatibilité du morphisme de Künneth.

\subsubsection{}\etiquette[5.4.2]{XVII.subsubsec5.4.2}
Soient un diagramme commutatif \eqhyperref[XVII.eq5.4.1.1]{5.4.1.1} de schémas et $\Aa$
un faisceau de torsion sur $S$. Supposons que $f$ et $g$ sont
compactifiables (\hyperref[XVII.def3.2.1]{3.2.1}), que $Z=\fprod{X}{Y}{S}$ et que
$Z'=\fprod{X'}{Y'}{S}$. Factorisons $f$ et $g$ en un morphisme propre
et une immersion ouverte; on obtient un diagramme
$$
\xymatrix@C=1.75cm@R=1cm{
  X'      \sisi{\ar[r]^{p'}}{} \ar[d]^{j_1} \ar@/_1.75pc/[dd]_{f}  & Z'   \ar[d]^{j_3}\sisi{}{\ar[l]_{p'}\ar[r]^{q'}}
& \sisi{Y'\ar[l]_{q'}}{} \ar[d]^{j_2} \ar@/^1.75pc/[dd]^{g}\\
\bar{X}' \sisi{\ar[r]^{\bar{p}}}{} \ar[d]^{\bar{f}}  & \bar{Z}' \ar[d]^{\bar{h}}\sisi{}{\ar[l]_{\bar{p}}\ar[r]^{\bar{q}}} & \bar{Y}' \sisi{\ar[l]_{\bar{q}}}{}\ar[d]^{\bar{g}}\\
X            &    Z \ar[l]_{p} \ar[r]^{q}  & Y
}
$$
dans\pageoriginale lequel on a posé
$\bar{Z}'=\fprod{\bar{X}'}{\bar{Y}'}{S}$. On
dispose d'un isomorphisme de nature triviale, pour $K$ et $L$
complexes sur $X'$ et $Y'$
$$\etiquette[5.4.2.1]{XVII.eq5.4.2.1}
j_{3!}(K\boxtimes_{\Aa}L)\isomto
j_{1!}K\boxtimes_{\Aa} j_{2!}L{\quoi}.\leqno(5.4.2.1)
$$

En « composant »  l'inverse de cet isomorphisme et la flèche de
Künneth \eqhyperref[XVII.eq5.4.1.4]{5.4.1.4} relative à $(\bar{f},\bar{g},\bar{h}, j_{1!}K,
j_{2!}L)$, on obtient, \emph{pour $K\in \Ob D^{-}(X', \Aa)$ et
$L\in \Ob D^{-}(Y',\Aa)$, une flèche de Künneth}
$$\etiquette[5.4.2.2]{XVII.eq5.4.2.2}
Rf_{!}K\overset{\LL}{\boxtimes}\sisi{}{_\Aa} Rg_{!}L\longrightarrow
Rh_{!}(K\overset{\LL}{\boxtimes}\sisi{}{_\Aa} L){\quoi},\leqno(5.4.2.2)
$$
dont on vérifie qu'elle ne dépend pas des compactifications choisies.

\begin{enonce*}{Théorème 5.4.3}\etiquette[5.4.3]{XVII.thm5.4.3}
Sous les hypothèses de \hyperref[XVII.subsubsec5.4.2]{5.4.2}, la flèche de Künneth en cohomologie à
supports propres \eqhyperref[XVII.eq5.4.2.2]{5.4.2.2} est un isomorphisme.
\end{enonce*}

\begin{enonce*}[remark]{Preuve}
\begin{enumerate}
\renewcommand{\theenumi}{\alph{enumi}}
\renewcommand{\labelenumi}{(\theenumi)}
\item Par construction, il suffit de vérifier ce résultat lorsque $f$
et $g$ (donc aussi $h$) sont propres.

\item Soit un diagramme commutatif
$$\etiquette[5.4.3.1]{XVII.eq5.4.3.1}
\vcenter{
\xymatrix@C=1.75cm@R=1.25cm{
X_1 \ar[r]^{p_1} \ar[d]& S_1 \sisi{\ar[r]^{q_1}}{}\ar[d]& Y_1\ar[d]\sisi{}{\ar[l]_{q_1}}\\
X   \ar[r]^{p}& S  \sisi{\ar[r]^{q}}{} &  Y\sisi{}{\ar[l]_q}.
}}\leqno{(5.4.3.1)}
$$
et soient $Z_{1}=\fprod{X_{1}}{Y_{1}}{S_{1}}$,
$X'_{1}=\fprod{X'}{X_{1}}{X}$, $Y'_{1}=\fprod{Y'}{Y_{1}}{Y}$ et
$Z'_{1}=\fprod{X'_{1}}{Y'_{1}}{S_{1}}=\fprod{Z'}{Z\sisi{_{1}}{}}{Z_{1}}$, de
sorte que l'on dispose d'un morphisme de diagrammes \sisi{(\hyperref[XVII.subsubsec5.4.1]{5.4.1})}{\eqhyperref[XVII.eq5.4.1.1]{5.4.1.1}},
\[
(S_{1},X_{1},Y_{1},Z_{1},X'_{1}, Y'_{1},Z'_{1})\longrightarrow
(S,X,Y,Z,X',Y',Z')
\]
dont les flèches seront uniformément désignées par $C$.


On\pageoriginale admettra que:

\begin{enonce*}{Lemme 5.4.3.2}\etiquette[5.4.3.2]{XVII.lem5.4.3.2}
Le diagramme suivant est commutatif:
$$
\xymatrix{
C^{\ast}(Rf_{!} K \overset{L}{\boxtimes}_{\Aa}
Rg_{!} L) \ar@{=}[r] \ar[d]& C^{\ast}Rf_{!}K \overset{L}{\boxtimes}_{\Aa}
C^{\ast} Rg_{!} L \ar[r]^{\sim}&  Rf_{1!} C^{\ast}K
\overset{L}{\boxtimes}_{\Aa} Rg_{1!} C^{\ast}
L \ar[d]\\
C^{\ast}(Rh_{!}(K \overset{L}{\boxtimes}_{\Aa}
L)) \ar[r]^{\sim}& Rh_{\sisi{}{1}!}(C^{\ast}(K
 \overset{L}{\boxtimes}_{\Aa}  L))\ar@{=}[r] &Rh_{\sisi{}{1}!}(C^{\ast}K
 \overset{L}{\boxtimes}_{\Aa}  C^*L ).
}
$$
\end{enonce*}

Dans ce diagramme, les flèches horizontales sont des flèches de
changement de base et les flèches verticales les flèches de Künneth
pour les deux diagrammes \eqhyperref[XVII.eq5.4.1.1]{5.4.1.1} considérés.

Pour prouver \hyperref[XVII.thm5.4.3]{5.4.3}, il suffit de prouver que la flèche
de Künneth est
un isomorphisme en tout point géométrique de $Z$. D'après
\hyperref[XVII.lem5.4.3.2]{5.4.3.2}, il
suffit de le vérifier après un quelconque changement de base \eqhyperref[XVII.eq5.4.3.1]{5.4.3.1}
tel que $S_{1}$ soit le spectre d'un corps algébriquement clos et
$p_{1}$, $q_{1}$ des isomorphismes.

Après un tel changement de base, le diagramme \eqhyperref[XVII.eq5.4.1.1]{5.4.1.1}
se réduit à un
diagramme
$$\vcenter{\etiquette[5.4.3.3]{XVII.eq5.4.3.3}
\xymatrix@R=1.5cm@C=2cm{
    & \fprod{X}{Y}{S} \ar[ld]_{p} \ar[rd]^{q} \ar[dd]^{f\times g}&     \\
X \ar[rd]_{f}  &                   &    Y \ar[ld]^{g}\\
    &      S             &\sisi{.}{}
}}\leqno{(5.4.3.3)}
$$
et la flèche de Künneth à une flèche \eqhyperref[XVII.eq5.4.4.1]{5.4.4.1}
$$\etiquette[5.4.3.4]{XVII.eq5.4.3.4}
k:Rf_{!}K\overset{\LL}{\otimes}\sisi{}{_\Aa} Rg_{!} L\longrightarrow
R(f\times g)_{!} (K\overset{\LL}{\boxtimes}\sisi{}{_\Aa} L){\quoi}.\leqno(5.4.3.4)
$$

\item Soit\pageoriginale un diagramme \eqhyperref[XVII.eq5.4.3.3]{5.4.3.3} dans lequel $f$ et $g$ sont
propres. On admettra que:

\begin{enonce*}{Lemme 5.4.3.5}\etiquette[5.4.3.5]{XVII.lem5.4.3.5}
Le diagramme suivant est commutatif
$$
\xymatrix@C=1.75cm{
Rf_{!}K \overset{\sisi{L}{\LL}}{\underset{\Aa}{\otimes}} Rg_{!}L
\ar[r] \ar[d]& R(f\times g)_{!}(K
\overset{\sisi{L}{\LL}}{\underset{\Aa}{\boxtimes}} L) \ar@{=}[d]\\
Rf_{!}(K \overset{\sisi{L}{\LL}}{\sisi{\boxtimes}{\otimes}}_{\Aa}
f^{\ast}Rg_{!}L) \ar[r]& Rf_{!} Rp_{!}(p^{\ast}K \overset{\sisi{}{\LL}}{\otimes}\sisi{}{_\Aa} q^{\ast}L).
}
$$
\end{enonce*}

Dans ce diagramme, les flèches horizontales sont \sisi{les flèches}{la flèche} de
Künneth et la flèche de changement de base. Les flèches verticales
sont une flèche de changement de base par $f$ et un isomorphisme de
transitivité. Le théorème \sisi{\hyperref[XVII.subsubsec5.5.9]{5.5.9}}{\hyperref[XVII.thm5.4.3]{5.4.3}} est conséquence de ce lemme et de
\hyperref[XVII.thm4.3.1]{4.3.1}.

Le lecteur courageux pourra vérifier que ce lemme reste valable si $f$
et $g$ ne sont pas nécessairement propres.
\end{enumerate}
\end{enonce*}

\setcounter{subsubsection}{3}
\subsubsection{}\etiquette[5.4.4]{XVII.subsubsec5.4.4}
Soient $S$ un schéma, $f_{i}:X'_{i}\rightarrow X_{i}$ $(i\in I)$ une
famille finie de morphismes compactifiables de $S$-schémas, $X$ le
produit sur $S$ des $X_{i}$, $X'$ le produit sur $S$ des $X'_{i}$ et
$f:X'\rightarrow X$ le produit des $f_{i}$. Si $S$ est annelé par
$\Aa$, comme en \hyperref[XVII.subsubsec5.4.1]{5.4.1}, et si on annèle les $X_{i}$ et $X'_{i}$
par l'image réciproque de $\Aa$, on a plus généralement un
isomorphisme
$$\etiquette[5.4.4.1]{XVII.eq5.4.4.1}
\mathop{\boxtimes}^{\LL}_{i\in I} Rf_{i!} (K_{i})
\isomto  Rf_{!} (\mathop{\boxtimes}^{\LL}_{i\in I}
K_{i}){\quoi},\leqno(5.4.4.1)
$$
pour $K_{i}\in \Ob D^{-}(X'_{i}, \Aa)$. La définition de cet
isomorphisme, ainsi que celle des deux membres de \eqhyperref[XVII.eq5.4.4.1]{5.4.4.1} \emph{ne
nécessite pas} le choix d'un ordre total sur $I$.

\subsection{Cohomologie des puissances symétriques}\etiquette[5.5]{XVII.subsec5.5}
\pageoriginale

Le résultat principal de ce ${\rm n^{\circ}}$ est le théorème \hyperref[XVII.thm5.5.21]{5.5.21},
donnant en particulier la cohomologie à supports propres d'une
puissance symétrique d'un schéma $X$ sur un corps $k$ algébriquement
clos, pour des coefficients convenables (par exemple $\ZZ/n\ZZ$) en
terme de la cohomologie à supports propres de $X$. L'analogue de ce
théorème en topologie semi-simpliciale est bien connu, tout au moins
pour la cohomologie à coefficients dans un anneau commutatif $A$; il
se vérifie au niveau des cochaînes semi-simpliciales. En effet, si
$X_{\ast}$ est un ensemble semi-simplicial, il existe une bijection
continue canonique de la réalisation géométrique de la puissance
symétrique n\sisi{.ème}{-ième} $(X_{\ast})^{n}/\Sf_{n}$ de $X_{\ast}$ dans la
puissance symétrique $n$\textsuperscript{ième} de la réalisation
géométrique de $X_{\ast}$, et cette bijection induit un isomorphisme
sur la cohomologie.

\subsubsection{}\etiquette[5.5.1]{XVII.subsubsec5.5.1}
Rappelons quelques propriétés (dues à \textsc{D. Lazard}) des modules plats,
quelques propriétés (dues à \textsc{Roby}) des foncteurs $TS^{n}$ et $T^{n}$,
et quelques propriétés (dues à \textsc{Dold} et \textsc{Puppe}) des
foncteurs dérivés de foncteurs non additifs.

\paragraph{}\etiquette[5.5.1.1]{XVII.par5.5.1.1}
\emph{Tout module plat sur un anneau $A$ est limite inductive
filtrante de modules libres de type fini} (\textsc{Lazard} \cite[I 1.2]{SGA4-17:La}). Ce
résultat implique les suivants:

\paragraph{}\etiquette[5.5.1.2]{XVII.par5.5.1.2}
Le foncteur $\varinjlim$ est une équivalence de la catégorie des
\iInd-objets (\hyperrefext[I][I.sec8.1]{8.1}) de la catégorie des $A$-modules libres de type fini
(\resp projectifs de type fini) avec la catégorie des $A$-modules
plats.

\paragraph{}\etiquette[5.5.1.3]{XVII.par5.5.1.3}
Si\pageoriginale $A$ et $B$ sont deux anneaux, tout foncteur
$T:\Mod(A)\rightarrow \Mod(B)$ qui transforme module libre de type
fini en module plat et commute aux limites inductives filtrantes,
transforme module plat en module plat (\textsc{Lazard} \cite[I 1.5]{SGA4-17:La}).

\paragraph{}\etiquette[5.5.1.4]{XVII.par5.5.1.4}
Toute suite exacte de $A$-modules, avec $L''$ plat,
$$
0\longrightarrow L'\longrightarrow L\longrightarrow L''\longrightarrow
0
$$
est limite inductive de suites exactes scindées (\textsc{Lazard} \cite[I 2.3]{SGA4-17:La}).

\subsubsection{}\etiquette[5.5.2]{XVII.subsubsec5.5.2}
Soient $A$ un anneau commutatif et $M$ un $A$-module.

\paragraph{}\etiquette[5.5.2.1]{XVII.par5.5.2.1}
Pour tout entier $n$, on désigne par $TS^{n}(M)$ (ou $TS^{n}_A(M)$) le
module des tenseurs symétriques de degré $n$ de $M$. Désignant par
$\Sf_n$ le groupe symétrique de degré $n$, on a par définition
$$
TS^{n}(M)=(\bigotimes^{n}M)^{\Sf_n}{\quoi}.
$$

Si $B$ est une $A$-algèbre, alors $TS^{n}(B)$ est une sous-algèbre de
$\bigotimes\limits^{n}B$. Si de plus $B$ est commutatif,
alors $\Spec (TS^{n}(B))$ est par définition la puissance symétrique
$n$\textsuperscript{ième} de $\Spec(B)$ sur $\Spec(A)$:
$$
\Spec(TS^{n}_{A}(B))=\Sym^{n}_{\Spec(A)}(\Spec(B)\underset{\dfn}{=}(\Spec(B)/\Spec(A))^{n}/\Sf_n{\quoi}.
$$

\paragraph{}\etiquette[5.5.2.2]{XVII.par5.5.2.2}
Rappelons (\textsc{Roby} \cite[I \sisi{1 pg.226}{2 p.~219}]{SGA4-17:Ro}) qu'une \emph{loi polynôme} $f$ d'un
$A$-module $M$ dans un $A$-module $N$ est la donnée, pour toute
$A$-algèbre commutative $A'$, d'une application
$f_{A'}:M\otimes_{A}A'\rightarrow N\otimes_{A}A'$, ces applications
étant telles que pour tout homomorphisme de $A$-algèbres
$u:A'\rightarrow A''$, le diagramme
\[
\xymatrix{
M \otimes_A A'  \ar[r]^{f_{A'}} \ar[d]^{u}& N \otimes_A A' \ar[d]^{u}\\
M \otimes_A A'' \ar[r]^{f_{A''}}& N \otimes_A A''
}
\]
soit\pageoriginale commutatif.

Une loi polynôme $f$ est dite \emph{homogène de degré} $n$ (\textsc{Roby} \cite[I
\sisi{1}{8} p\sisi{g}{}.~226]{SGA4-17:Ro}) si, pour tout $a'\in A'$, on a identiquement
\[
f_{A'}(a'x)={a'}^{n}f_{A'}(x){\quoi}.
\]
Une telle loi polynôme s'appellera aussi \emph{application}
$n$\textsuperscript{ique} (quadratique, cubique, \ldots); on notera leur
ensemble
\[
\Hom n^{\text{\rm ique}} (M,N){\quoi}.
\]

Plus généralement, si $(M_{i})_{i\in I}$ est une famille finie de
$A$-modules et $\sheafn =(n_{i})_{i\in I}$ une famille d'entiers, une loi
polynôme multihomogène de multidegré $\sheafn $ des modules $M_{i}$ dans
$N$ est une loi polynôme de $\sisi{\sum}{\oplus} M_{i}$ dans $N$ vérifiant
identiquement, pour $a_{i}\in A'$,
\[
f_{A'}((a_{i}x_{i})_{i\in I})=\prod_{i\in I} a^{n_{i}}_i\cdot
f_{A'}((x_{i})_{i\in I}){\quoi}.
\]
On notera $\Hom \sheafn ^{\text{\rm ique}}((M_{i}),N)$ l'ensemble de ces
lois polynômes.

\paragraph{}\etiquette[5.5.2.3]{XVII.par5.5.2.3}
Pour tout module $M$, le foncteur covariant des $A$-modules dans les
ensembles
$$
N\longmapsto \Hom n^{\text{\rm ique}} (M,N)
$$
est représentable par un $A$-module $\Gamma^{n}(M)$ (ou
$\Gamma^{n}_{A}(M)$), muni d'une \emph{application\pageoriginale
$n$\textsuperscript{ique} universelle} $\gamma^{n}:M\rightarrow
\Gamma^{n}(M)$ (\textsc{Roby} \cite[IV 1 \sisi{pg.226}{p.~265}]{SGA4-17:Ro}). Etant donné $(M_{i})_{i\in I}$
et $\sheafn $ comme en \hyperref[XVII.par5.5.2.2]{5.5.2.2}, le foncteur covariant
$$
N\longmapsto \Hom \sheafn ^{\text{\rm ique}} ((M_{i}),N)
$$
est représenté par le module $\bigotimes\limits_{i\in
I}\Gamma^{n_i}(M_{i})$, muni de la \emph{loi polynôme multihomogène
(de multidegré $\sheafn $) universelle} $(m_{i})\mapsto
\bigotimes\limits_{i\in I}\gamma^{n_{i}}(m_{i})$.

\paragraph{}\etiquette[5.5.2.4]{XVII.par5.5.2.4}
Les foncteurs $TS^{n}$ et $\Gamma^{n}$ commutent aux limites
inductives filtrantes (pour le second, \textsc{Roby} \cite[IV \sisi{6 pg.275}{7 p.~277}]{SGA4-17:Ro}). Ils
transforment modules libres en modules libres, donc modules plats en
modules plats (\hyperref[XVII.par5.5.1.2]{5.5.1.2}).

\paragraph{}\etiquette[5.5.2.5]{XVII.par5.5.2.5}
L'application $n^{\text{\rm ique}}$ de $M$ dans $TS^{n}(M)$ donnée par
$m\mapsto \bigotimes\limits^{n}m$ définit une
application $\varphi$, dite canonique, de $\Gamma^{n}(M)$ dans
$TS^{n}(M)$, vérifiant
$$
\varphi(\gamma^{n}(m))=\bigotimes^{n}m{\quoi}.
$$
Cette application est un isomorphisme pour $M$ libre de type fini
(\textsc{Roby} \cite[IV \sisi{pg.272}{5 p.~272}]{SGA4-17:Ro}) donc aussi pour $M$ plat (\hyperref[XVII.par5.5.2.4]{5.5.2.4}
et \hyperref[XVII.par5.5.1.1]{5.5.1.1}).

Soit $M$ un $A$-module et $P$ une suite exacte
$$
P:P_{1}\xrightarrow{d} P_{0} \xrightarrow{\epsilon} M\longrightarrow
0{\quoi}.
$$
Soit $P'_{1}=P_{1}\times P_{0}$, et soient $j_{0}$ et $j_{1}$ les
morphismes de $P'_{1}$ dans $P_{0}$ de coordonnées $(d,\Id)$ et
$(0,\Id)$. La suite
$$
\xymatrix@C=1.75cm{
\Gamma^{n}(P'_1) \ar@<2pt>[r]^{\Gamma^{n}(j_0)} \ar@<-2pt>[r]_{\Gamma^{n}(j_1)}& \Gamma^{n}(P_0) \ar[r]& \Gamma^{n}(M) \ar[r]& 0
}
$$
est alors exacte (\textsc{Roby} \cite[\sisi{pg.284}{IV 10 p.~284}]{SGA4-17:Ro}). \emph{Pour $P$ une présentation
plate de} $M$,\pageoriginale on en déduit une description du foncteur
$\Gamma^{n}$ via le foncteur $TS^{n}$
$$\etiquette[5.5.2.5.1]{XVII.eq5.5.2.5.1}
TS^{n}(P'_1)\rightrightarrows TS^{n}(P_{0})\rightarrow
\Gamma^{n}(M)\rightarrow 0{\quoi}.\leqno(5.5.2.5.1)
$$

\paragraph{}\etiquette[5.5.2.6]{XVII.par5.5.2.6}
Si $M=M'\oplus M''$, l'application $n^{\text{\rm ique}}$
\[
m\longmapsto (\gamma^{n'}m\otimes \gamma^{n''}m)_{n'+n''=n}
\]
définit un isomorphisme, dit canonique (\textsc{Roby} \cite[III \sisi{4}{9 p.~262}]{SGA4-17:Ro})
$$\etiquette[5.5.2.6.1]{XVII.eq5.5.2.6.1}
\Gamma^{n}(M)\isomto \sisi{\sum}{\bigoplus}_{n'+n''=n} \Gamma^{n'}(M)\otimes \Gamma^{n''}(M){\quoi}.\leqno(5.5.2.6.1)
$$

Si $\Sigma$ est une suite exacte
\[
\Sigma:0\longrightarrow M'\longrightarrow M\longrightarrow
M''\longrightarrow 0{\quoi},
\]
le module $\Gamma^{n}(M)$ est muni d'une filtration canonique
(décroissante) et de morphismes
$$\etiquette[5.5.2.6.2]{XVII.eq5.5.2.6.2}
\Gamma^{i}(M'')\otimes \Gamma^{n-i}(M')\longrightarrow \Gr^{i}(\Gamma^{n}(M)){\quoi}.\leqno(5.5.2.6.2)
$$
Ces données sont caractérisées par les propriétés d'être fonctorielles
en $\Sigma$ et de vérifier une compatibilité évidente avec \eqhyperref[XVII.eq5.5.2.6.1]{5.5.2.6.1}
pour $\Sigma$ scindable. Les flèches \eqhyperref[XVII.eq5.5.2.6.2]{5.5.2.6.2} sont des
isomorphismes pour $\Sigma$ scindable, donc aussi pour $M''$ plat
(\hyperref[XVII.par5.5.1.4]{5.5.1.4}).

\paragraph{}\etiquette[5.5.2.7]{XVII.par5.5.2.7}
Le foncteur $\Gamma^{n}$ commute à l'extension des scalaires (\textsc{Roby} \cite[III \sisi{3 pg.262}{8 p.~262}]{SGA4-17:Ro}) (la propriété analogue pour $TS^{n}$ est en général
fausse).

\paragraph{}\etiquette[5.5.2.8]{XVII.par5.5.2.8}
Si $u:M_{1}\times M_{2}\rightarrow N$ est une application bilinéaire,
alors $\gamma^{n}\circ u$ est une loi polynôme bihomogène de bidegré
$(n,n)$ et définit donc (\hyperref[XVII.par5.5.2.3]{5.5.2.3}) $u^{n}:\Gamma^{n}(M_{1})\otimes
\Gamma^{n}(M_{2})\rightarrow \Gamma^{n}(N)$ vérifiant
$u^{n}(\gamma^{n}x\otimes
\gamma^{n}y)=\gamma^{n}u(x,y)$. En\pageoriginale particulier, si $B$
est une $A$-algèbre, le produit: $B\times B\rightarrow B$ induit un
produit sur $\Gamma^{n}(B)$. Ce produit fait de $\Gamma^{n}(B)$ une
$A$-algèbre, et l'application canonique de $\Gamma^{n}(B)$ dans
$TS^{n}(B)$ est un homomorphisme d'algèbres.

\subsubsection{}\etiquette[5.5.3]{XVII.subsubsec5.5.3}
Soit $\Aa$ une catégorie abélienne. On désignera par
$\text{\rm ssimpl}^{+}(\Aa)$ la catégorie des objets
co-semi-simpliciaux\nde{on les appelle objets \emph{cosimpliciaux} dans la terminologie moderne.} de $\Aa$, par $C^{\geq 0}(\Aa)$ la catégorie des
complexes différentiels $K$ avec $K^{i}=0$ pour $i<0$, et par $K^{\geq
  0}(\Aa)$ l'image de $C^{\geq 0}(\Aa)$ dans $K(\Aa)$. Si $K\in \Ob
\text{ssimpl}^{+}(\Aa)$, le complexe normalisé associé
$N(K)\in \Ob C^{\geq 0}(\Aa)$ est défini par
$$\etiquette[5.5.3.1]{XVII.eq5.5.3.1}
\begin{cases}
N(K)^{n} =\text{intersection des noyaux des opérateurs de
dégénérescence}\\
\qquad\qquad\text{de $K^{n}$ dans $K^{n-1}$}\\
d=\sum(-1)^{i}\partial_{i}{\quoi}.
\end{cases}\leqno(5.5.3.1)
$$

Si on désigne encore par $K$ le complexe de composantes $K^{n}$, et de
différentielle $d=\sum(-1)^{i}\partial_{i}$, alors $N(K)$ est un
sous-complexe facteur direct dans $K$, et $K/N(K)$ est homotope à
zéro.

D'après \textsc{Dold-Puppe} \cite{SGA4-17:DP}, le foncteur $N$ est une équivalence de
catégories
$$
N:\;\text{\rm ssimpl}^{+}(\Aa)\longrightarrow C^{\geq
0}(\Aa){\quoi}.
$$
Cette équivalence et son inverse transforment morphismes homotopes en
morphismes homotopes. De plus, pour tout $K\in \Ob\text{\rm
ssimpl}^{+}(\Aa)$, $K^{n}$ est somme de copies de $N(K)^{i}$
$(i\leq n)$; si $\PP$ est une propriété d'objets de
$\Aa$, stable par sommes finies et facteurs directs, on en
déduit que les $K^{i}$ ont la propriété $\PP$ pour $i\leq n$ si
et seulement si les $N(K)^{i}$ ont cette propriété pour $i\leq n$.

L'équivalence\pageoriginale $N$ est utilisée par \textsc{Dold} et
\textsc{Puppe} pour
étendre tout foncteur (non nécessairement additif)
$T:\Aa\rightarrow \Bb$ de $\Aa$ dans une
catégorie quelconque $\Bb$ en un foncteur $C^{\geq 0}
(T)$ (ou simplement $T$) de $C^{\geq 0}(\Aa)$ dans $C^{\geq
0}(\Bb)$
$$\etiquette[5.5.3.2]{XVII.eq5.5.3.2}
C^{\geq 0}(T)=NTN^{-1}{\quoi}.\leqno(5.5.3.2)
$$

Le foncteur $C^{\geq 0}(T)$ passe au quotient pour définir
$$\etiquette[5.5.3.3]{XVII.eq5.5.3.3}
T\quad\mbox{\rm ou}\quad K^{\geq 0}(T):K^{\geq
0}(\Aa)\longrightarrow K^{\geq 0}(\Bb){\quoi}.\leqno(5.5.3.3)
$$

\setcounter{paragraph}{3}
\paragraph{}\etiquette[5.5.3.4]{XVII.par5.5.3.4}
Si $T$ est \emph{de degré} $\leq n$, \ie si le $(n+1)$\textsuperscript{ième} « cross effect »  est nul (\textsc{Dold-Puppe} \cite[4.18 \sisi{pg.232}{p.~232}]{SGA4-17:DP})\sisi{}{,}
alors, pour tout $K\in \Ob C^{\geq 0}(\Aa)$ tel que $K^{i}=0$
pour $i>m$, on a $[T(K)]^{i}=0$ pour $i>nm$ (\textsc{Dold-Puppe} \cite[4.2\sisi{.}{}3
\sisi{pg.233}{p.~233}]{SGA4-17:DP}).

\paragraph{}\etiquette[5.5.3.5]{XVII.par5.5.3.5}
En particulier, si $T$ est un foncteur de degré fini de la catégorie
des $A$-modules dans celle des $B$-modules, qui transforme objet plat
en objet plat, alors un complexe plat borné a pour image un complexe
plat borné.

\begin{enonce*}{Lemme 5.5.4}\etiquette[5.5.4]{XVII.lem5.5.4}
Soit $\Aa$ une catégorie abélienne.
\begin{enumerate}
\renewcommand{\theenumi}{\roman{enumi}}
\renewcommand{\labelenumi}{(\theenumi)}
\item Les quasi-isomorphismes dans $K^{\geq 0}(\Aa)$
permettent un calcul des fractions à gauche et à droite.

\item Le foncteur évident est une équivalence entre la catégorie
déduite de $K^{\geq 0}(\Aa)$ en inversant les
quasi-isomorphismes et la catégorie $D^{\geq 0}(\Aa)$.
\end{enumerate}
\end{enonce*}

Le\pageoriginale foncteur $\tau_{\geq 0}:C(\Aa)\rightarrow
C^{\geq 0}(\Aa)$ (\hyperref[XVII.def1.1.13]{1.1.13}) passe au quotient pour définir un
foncteur $\tau_{\geq 0}:K(\Aa)\rightarrow K^{\geq
0}(\Aa)$, adjoint à gauche à l'inclusion de $K^{\geq
0}(\Aa)$ dans $\sisi{\Aa}{K(\Aa)}$. Le foncteur $\tau_{\geq 0}$
transforme quasi-isomorphismes en quasi-isomorphismes, et la flèche
d'adjonction $X\rightarrow \tau_{\geq 0}(X)$ est un quasi-isomorphisme
si et seulement si $H^{i}(X)=0$ pour $i<0$. Le lemme
\hyperref[XVII.lem5.5.4]{5.5.4} résulte
donc de \textsc{Verdier} \cite{SGA4-17:CD} et des points (i) et (ii) du lemme suivant, dont
la démonstration est laissée au lecteur:

\begin{enonce*}{Lemme 5.5.4.1}\etiquette[5.5.4.1]{XVII.lem5.5.4.1}
Soient $K$ une catégorie, $K'$ une sous-catégorie pleine, $S$ un
ensemble de morphismes de $K$ égal à l'image inverse de l'ensemble des
isomorphismes de $K(S^{-1})$ (\cf \cite[\sisi{chap 1}{\chap~1}]{SGA4-17:GZ}) et $S'$ sa
trace dans $K'$. On suppose que l'inclusion $i$ de $K'$ dans $K$ admet
un adjoint à gauche $\tau$, et que $\tau(S)\subset S'$.
\begin{enumerate}
\renewcommand{\theenumi}{\roman{enumi}}
\renewcommand{\labelenumi}{(\theenumi)}
\item Le foncteur évident de $K'({S'}^{-1})$ dans $K(S^{-1})$ est
pleinement fidèle, et admet $\tau$ pour adjoint à gauche.

\item Si $K(S^{-1})$ peut se calculer par calcul de fractions à gauche
(\resp à droite), alors $K'({S'}^{-1})$ a la même propriété.

Si l'inclusion $i$ n'a pas nécessairement d'adjoint à gauche, et si
$\tau$ désigne le foncteur non part\sisi{an}{ou}t défini adjoint à gauche à $i$,
on a:

\item Si $K(S^{-1})$ peut se calculer par calcul de fractions à gauche
(\resp à droite) et si pour tout $u\in S$, de but (\resp de source)
dans $K'$, $\tau(u)$ est défini et dans $S'$, alors $K'({S'}^{-1})$
admet un calcul de fractions à gauche (\resp à droite) et le foncteur
$K'({S'}^{-1})\rightarrow K(S^{-1})$ est pleinement fidèle.
\end{enumerate}
\end{enonce*}

\setcounter{subsubsection}{4}
\subsubsection{}\etiquette[5.5.5]{XVII.subsubsec5.5.5}
Soit\pageoriginale $T:\Aa\rightarrow \Bb$ un foncteur
entre catégories abéliennes. Comme en \hyperref[XVII.def1.2.1]{1.2.1}, on définit les dérivés de
$T$
\begin{align*}
LT &: D^{\geq 0}(\Aa)\longrightarrow \Pro(D^{\geq
0}(\Bb))\\
RT &: D^{\geq 0}(\Aa)\longrightarrow \Ind(D^{\geq 0}(\Bb))
\end{align*}
par les formules
\begin{align*}
LT(K) &= \quotprojlim_{s:K'\rightarrow K}T(K')\\
RT(K) &= \quotinjlim_{s:K\rightarrow K'}T(K'){\quoi},
\end{align*}
où $s$ parcourt la catégorie filtrante des quasi-isomorphismes de but
(\resp de source) $K$.

On dit que $T$ est \emph{dérivable à gauche (\resp à droite) en $K$ si
$LT(K)$} $(RT(K))$ est représentable. On dit que $K$ est \emph{déployé
à gauche (\resp à droite)} pour $T$ (ou, par abus de langage,
déployé pour $LT$ (\resp $RT$)) si la flèche canonique suivante est un
isomorphisme:
\[
LT(K)\isomto  TK\text{ (\resp $TK \isomto  RT(K)$)}{\quoi}.
\]

\begin{enonce*}{Proposition 5.5.6}\etiquette[5.5.6]{XVII.prop5.5.6}
Soit $F$ un foncteur de la catégorie des modules sur un anneau $A$
dans une catégorie abélienne $\Bb$. Supposons que les limites
inductives filtrantes dans $\Bb$ soient représentables et
soient exactes, et que $F$ commute aux limites inductives
filtrantes. Alors tout complexe borné $M^{\bigdot}\in \Ob K^{\geq
0}(A)$ de modules plats est déployé pour $LF$.
\end{enonce*}

La\pageoriginale proposition résulte immédiatement des définitions, du
lemme suivant, qui sera prouvé en \hyperref[XVII.par5.5.6.4]{5.5.6.4}, et du fait que les
quasi-isomorphismes $u:N^{\bigdot}\rightarrow M^{\bigdot}$, avec
$N^{\bigdot}\in \sisi{B}{C}^{b,\geq 0}(A)$ à composantes plates, forment un
ensemble cofinal dans la catégorie des quasi-isomorphismes de but
$M^{\bigdot}$.

\begin{enonce*}{Lemme 5.5.6.1}\etiquette[5.5.6.1]{XVII.lem5.5.6.1}
Soit $u:N^{\bigdot}\rightarrow M^{\bigdot}$ un quasi-isomorphisme
entre complexes plats, objets de $C^{b,\geq 0}(A)$. Alors, $F(u)$ est
un quasi-isomorphisme.
\end{enonce*}

\begin{enonce*}{Lemme 5.5.6.2}\etiquette[5.5.6.2]{XVII.lem5.5.6.2}
Soient $\Aa$ une catégorie additive et $d$ un intervalle borné
supérieurement de $\ZZ$. On désigne par $C^{d}$ (\resp
$C^{d,b}$, \resp $K^{d}$, \resp $K^{d,b}$) la catégorie des complexes
nuls en degrés $i\not\in d$ (\resp et bornés, \resp à homotopie
près). On a une équivalence
$$
\sisi{\sheafell }{\ell}:\Ind(C^{d,b}(\Aa)\sisi{}{)}\isomto
C^{d}(\Ind\Aa){\quoi};
$$
l'inverse de cette équivalence induit un foncteur
$$
\sisi{\sheafell }{\ell}':K^{d}(\Ind\Aa)\longrightarrow
\Ind(K^{d,b}(\Aa)){\quoi}.
$$
\end{enonce*}

\begin{enumerate}
\renewcommand{\theenumi}{\Alph{enumi}}
\renewcommand{\labelenumi}{\theenumi)}
\item Prouvons que $\sisi{\sheafell }{\ell}$ est pleinement fidèle. Soit donc $K$ un
complexe borné d'objets de $\Aa$ et $L=\quotinjlim
L_{\alpha}$ un \iInd-complexe. On a
\begin{align*}
\Hom(K,L)\overset{\dfn}{=} \varinjlim \Hom(K,L_{\alpha})  = & \varinjlim
Z^{0}(\Hom^{\bigdot} (K,L_{\alpha}))\simeq Z^{0}
\varinjlim_{\alpha} \Hom^{\bigdot} (K,L_{\alpha})\\
& \isomto
Z^{0} \Hom^{\bigdot} (K,\ell L)=\Hom(K,\ell L){\quoi},
\end{align*}
l'avant-dernier isomorphisme provenant du fait que $K$ est fini, de
sorte que la $\varinjlim$ commute au passage du double au simple
complexe. Ceci prouve l'assertion.

\item On\pageoriginale laisse au lecteur le soin (assez fastidieux) de
vérifier que $\sisi{\sheafell }{\ell}$ est essentiellement surjectif en construisant,
de gauche à droite, à partir d'un degré arbitrairement choisi, « suffisamment »  de complexes bornés s'envoyant dans un complexe
borné supérieurement de \iInd-objets $K$; et de développer de même le
point suivant:

\item Si $\ell f:\ell \varinjlim K_{\alpha}\rightarrow \ell \varinjlim
L_{\beta}$ est homotope à zéro: $f=dH+Hd$, alors, pour tout $\alpha$,
$H$ se relève en $H^{i}_{\alpha}:K^{i}_{\alpha}\rightarrow
L^{i\sisi{}{-1}}_{\beta(\alpha)}$ qui, pour $\beta$ assez grand, définit une
homotopie à zéro de $f_{\alpha}:K_{\alpha}\rightarrow L_{\beta(\alpha)}$.
\end{enumerate}

\setcounter{subsubsection}{6}
\setcounter{paragraph}{2}
\paragraph{}\etiquette[5.5.6.3]{XVII.par5.5.6.3}
Soient $A$ un anneau, $d$ un intervalle borné supérieurement de
$\ZZ$, $P(A)$ la catégorie des $A$-modules projectifs de type
fini, $\Pl(A)$ la catégorie des $A$-modules plats. Désignons par
$D^{\Tor\subset d}(A)$ l'image de $C^{d}(\Pl(A))$ dans $D(A)$ et par
$D^{\parf\subset d}(A)$ l'image de $C^{d,b}(P(A))$ dans
$D(A)$. D'après \hyperref[XVII.par5.5.1.2]{5.5.1.2}, on a $\Pl(A)\simeq \Ind
P(A)$. On a par
ailleurs $K^{d,b}(P(A))\simeq D^{\parf\subset d}(A)$. On déduit dès
lors de \hyperref[XVII.lem5.5.6.2]{5.5.6.2}, et du fait que $H^{i}$ commute aux
limites inductives
filtrantes, qu'il existe un diagramme commutatif de foncteurs
$$
\xymatrix@C=2.25cm{
C^{d}(\text{Pl}(A)) \ar[d]& \Ind C^{d,b}(P(A)) \ar[d]\ar[l]_{\substack{\varinjlim\\\sim}}\\
K^{d}(\text{Pl}(A)) \ar[d]\ar[r]^{\ell'}& \Ind K^{d,b}(P(A)) \ar[d]^{\sisi{}{\simeq}}\\
D^{\Tor\subset d}(A) \ar[r]^{\ell''}& \Ind D^{\parf \subset d}(A).
}
$$

\setcounter{subsubsection}{6}
\paragraph{}\etiquette[5.5.6.4]{XVII.par5.5.6.4}
Prouvons\pageoriginale \hyperref[XVII.lem5.5.6.1]{5.5.6.1}. Avec les notations de
\hyperref[XVII.par5.5.6.3]{5.5.6.3}, si $d$ est un intervalle $[0,k]$, on dispose
(car $F$ commute aux $\varinjlim$) de diagrammes commutatifs
$$
\xymatrix@C=2.5cm{
K^{d}(\text{Pl}(A))          \ar[r]^{\ell'} \ar[d]_{K^{\geq 0}(F)} & \Ind K^d(P(A))  \ar[d]_{\Ind(K^{\geq 0}(F))}             \\
K^{+}(\Bb) \ar[d]_{H^i}                        & \Ind K^{+}(\Bb)     \ar[d]_{\Ind(H^i)}       \\
\Bb                  & \Ind \Bb  \ar[l]_{\varinjlim}.
}
$$
Si $u$ est comme en \hyperref[XVII.lem5.5.6.1]{5.5.6.1}, alors, d'après
\hyperref[XVII.par5.5.6.3]{5.5.6.3}, $\ell'(u)$ est un isomorphisme, et
$H^{i}K^{\geq 0}(F)(u)$ est donc un isomorphisme, ce qu'il fallait
prouver.

\subsubsection{}\etiquette[5.5.7]{XVII.subsubsec5.5.7}
Soient $X$ un schéma sur un schéma $S$ et $k$ un entier. On suppose
que la puissance symétrique $\Sym^{k}_{S}(X)$ est définie. Le lecteur
qui voudrait ne pas faire cette hypothèse pourra se placer dans la
catégorie des espaces algébriques (\textsc{Artin}\sisi{ [1]}{\nde{dans le texte original,
on cite ici un ouvrage non précisé, qui est probablement l'un des deux ouvrages d'Artin cités dans les notes en pages \pageref{XVII.noteArtin1} et \pageref{XVII.noteArtin2}.}}). On annèle $S$ par un
faisceau d'anneaux \emph{commutatifs} $\Aa$. À partir de
\hyperref[XVII.subsubsec5.5.19]{5.5.19}, $\Aa$ sera supposé \emph{de torsion}. On
\emph{désignera encore par} $\Aa$, par abus de notations,
l'image réciproque de $\Aa$ sur un quelconque $S$-schéma.

Soit $ℱ$ un faisceau de modules sur $X$; le faisceau
$\mathop{\boxtimes}\limits^{k} ℱ$, produit tensoriel externe
de $n$ copies de $ℱ$, est un faisceau sur $X^{k}/S$, équivariant par
l'action du groupe symétrique $\Sf_{k}$ sur $X^{k}/S$. Par
passage\pageoriginale aux invariants, il définit un faisceau
$TS^{k}_{\ext}(ℱ)$ sur $\Sym^{k}_{S}(X)$, appelé \emph{puissance
tensorielle symétrique externe de $ℱ$}\sisi{}{.}

Pour raisonner « point par point », il est utile de travailler
avec les points de $\Sym^{k}_{S}(X)$ décrits comme suit.

\paragraph{}\etiquette[5.5.7.1]{XVII.par5.5.7.1}
Soit $s$ un point géométrique de $S$. Un point $x$ de $\Sym_{S}(X)$ à
valeurs dans $k(s)$ définit un point géométrique de $\Sym_{S}(X)$. On
désignera par $(x_{1}\ldots x_{k})$ un point de $(X/S)^{k}$ à valeurs
dans $k(s)$ d'image $x$, par $y_{1}\ldots y_{\ell}$ les éléments
distincts de la suite $x_{1}\ldots x_{k}$ et par $n_{i}$ $(1\leq i\leq
\ell)$ la multiplicité de $y_{i}$ dans cette suite $(\sum n_{i}=k)$.

Avec les notations de \hyperref[XVII.par5.5.7.1]{5.5.7.1}, et désignant par $\Sf$ un
groupe symétrique, on a
$$\etiquette[5.5.7.2]{XVII.eq5.5.7.2}
TS^{k}_{\ext}(F)_x\simeq (\bigotimes^{\ell}_{1}
(\bigotimes^{n_{i}} F_{y_{i}}))^{\prod^{\ell}_{1}\Sf_{
n_{i}}}{\quoi},\leqno(5.5.7.2)
$$
les produits tensoriels étant pris sur $\Aa_{S}$. En
particulier, \emph{pour $F$ plat}, on a, grâce à \hyperref[XVII.par5.5.2.5]{5.5.2.5}:
$$\etiquette[5.5.7.3]{XVII.eq5.5.7.3}
TS^{k}_{\ext}(F)_x\isomfrom  \bigotimes^{\ell}_{1}
\Gamma^{n_{i}} F_{y_{i}}{\quoi},\leqno(5.5.7.3)
$$
en utilisant un argument par récurrence sur $\ell$, utilisant la
platitude des $\Gamma^{n_{i}}F_{y_{i}}$ (\hyperref[XVII.par5.5.2.4]{5.5.2.4}) et le fait suivant:
si $G$ (\resp $H$) est un groupe fini opérant sur un $A$-module
$M$\sisi{}{ }(\resp $N$), et si $M$ et $N^{H}$ sont plats, alors
$M^{G}\otimes_{A}N^{H}\isomto  (M\otimes_{A}N)^{G\times H}$.

\subsubsection{}\etiquette[5.5.8]{XVII.subsubsec5.5.8}
Soit $F$ un Module sur $X$, et $P$ une présentation plate de $F$
$$
P:P_{1}\longrightarrow P_{0}\longrightarrow F\longrightarrow 0{\quoi}.
$$ On\pageoriginale définit $P'_{1}$, $j_{0}$ et $j_{1}$ comme en
\hyperref[XVII.par5.5.2.5]{5.5.2.5}, et on définit $\Gamma^{k}_{\ext}(F)$
(\emph{foncteur $\Gamma$ externe}) comme le conoyau de la double
flèche $(TS^{k}_{\ext}(j_{0}),TS^{k}_{\ext}(j_{1}))$. Avec les
notations de \hyperref[XVII.par5.5.7.1]{5.5.7.1}, on déduit de \eqhyperref[XVII.eq5.5.7.3]{5.5.7.3} et
\hyperref[XVII.par5.5.2.5]{5.5.2.5} des isomorphismes
$$\etiquette[5.5.8.1]{XVII.eq5.5.8.1}
\Gamma^{k}_{\ext}(F)_{x}\simeq \bigotimes^{\ell}_{1} \Gamma^{n_{i}}(F_{y_{i}}){\quoi}.\leqno(5.5.8.1)
$$

De cette formule, on déduit que $\Gamma^{k}_{\ext}(F)$, qui a priori
dépend fonctoriellement de $P$, ne dépend, à isomorphisme unique près,
que de $F$.

\subsubsection{}\etiquette[5.5.9]{XVII.subsubsec5.5.9}
On dispose d'un morphisme fonctoriel
$$
\alpha:\Gamma^{k}_{\ext}(F)\longrightarrow TS^{k}_{\ext}(F){\quoi}.
$$
Le couple $(\Gamma,\alpha)$ est caractérisé par les deux propriétés:
\begin{enumerate}
\renewcommand{\theenumi}{\roman{enumi}}
\renewcommand{\labelenumi}{(\theenumi)}
\item Pour $F$ plat, $\alpha$ est un isomorphisme

\item Avec les notations de \hyperref[XVII.par5.5.2.5]{5.5.2.5}, pour toute présentation
$$
P_{1}\xrightarrow{d} P_{0} \xrightarrow{\epsilon} F\longrightarrow 0
$$
de $F$, la suite
$$
\Gamma^{k}_{\ext}(P'_{1})\rightrightarrows
\Gamma^{k}_{\ext}(P_{0})\rightarrow \Gamma^{k}_{\ext}(F)\rightarrow 0
$$
est exacte.

On a de plus:

\item Les foncteurs $TS^{k}$ et $\Gamma^{k}$ commutent aux limites
inductives filtrantes (d'après \sisi{\hyperref[XVII.eq5.5.7.2]{5.5.7.2}}{\eqhyperref[XVII.eq5.5.7.2]{5.5.7.2}}, \sisi{\hyperref[XVII.eq5.5.8.1]{5.5.8.1}}{\eqhyperref[XVII.eq5.5.8.1]{5.5.8.1}},
\hyperref[XVII.par5.5.2.4]{5.5.2.4}).

\item Le foncteur $\Gamma^{k}$ commute à toute extension
$\Aa\rightarrow \Aa'$ de l'anneau de base.

\item Le\pageoriginale foncteur $\Gamma^{k}$ commute à tout changement
de base $S'\rightarrow S$: pour tout carré cartésien
$$
\xymatrix@C=2.5cm@R=1.5cm{
X' \ar[r]^{g'}\ar[d]^{f'}& X \ar[d]^{f}\\
S' \ar[r]^{g}& S,
}
$$
on a
$$
\Sym^{k}_{g}(g')^{\ast}\Gamma^{k}_{\ext}\simeq
\Gamma^{k}_{\ext}{g'}^{\ast}{\quoi}.
$$
(\sisi{\hyperref[XVII.eq5.5.8.1]{5.5.8.1}}{\eqhyperref[XVII.eq5.5.8.1]{5.5.8.1}} et \hyperref[XVII.par5.5.2.7]{5.5.2.7}).

\item Le foncteur $\Gamma^{k}$ commute aux images inverses par tout
morphisme \emph{radiciel} $i:X'\rightarrow X$ (par exemple une
immersion fermée): \sisi{O}{o}n a
$$
\Sym^{k}_{S}(i)^{\ast} \Gamma^{k}_{\ext}\simeq
\Gamma^{k}_{\ext}i^{\ast}
$$
\eqhyperref[XVII.eq5.5.8.1]{5.5.8.1}.

\item Pour $j:X'\rightarrow X$ une \emph{immersion}, on a
$$
\Sym^{k}_{S}(j)_{!} \Gamma^{k}_{\ext}\simeq \Gamma^{k}_{\ext} j_{!}
$$
\eqhyperref[XVII.eq5.5.8.1]{5.5.8.1}.
\end{enumerate}

\subsubsection{}\etiquette[5.5.10]{XVII.subsubsec5.5.10}

Soient $k_{1}$, $k_{2}$ et $k=k_{1}+k_{2}$ trois entiers $\geq 0$. On
dispose alors d'un morphisme fini (\quad):\nde{je n'ai pas trouvé de références pour la finitude de $\vee_{k_1k_2}$ en général.}
$$
\vee_{k_{1}k_{2}}: \Sym_{S}^{k_{1}}(X)\times_{S} \Sym_{S}^{k_{2}}(X)\longrightarrow \Sym_{S}^{k}(X){\quoi}.
$$
Si $F_{i}$ $(i=1,2)$ est un Module sur $\Sym_{S} {}^{k_{i}}(X)$, on
désigne par $F_{1}\vee F_{2}$ le faisceau sur $\Sym_{S}^{k}(X)$
$$
F_{1}\vee F_{2}=\vee_{k_{1}k_{2}\ast}(F_{1}\boxtimes F_{2}){\quoi}.
$$

\subsubsection{}\etiquette[5.5.11]{XVII.subsubsec5.5.11}
Soit\pageoriginale une suite exacte de Modules sur $X$
\[
\Sigma:0\longrightarrow F'\longrightarrow F\longrightarrow
F''\longrightarrow 0{\quoi}.
\]
Avec les notations \hyperref[XVII.par5.5.7.1]{5.5.7.1}, on dispose, grâce à
\eqhyperref[XVII.eq5.5.8.1]{5.5.8.1} et \hyperref[XVII.par5.5.2.6]{5.5.2.6}, d'une filtration sur
$\Gamma^{k}_{\ext}(F)_{x}$, et de morphismes:
\[
\sum_{\sum^{\ell}_{1} i_{\alpha}=i}
\bigotimes^{\ell}_{1}
\Gamma^{i_{\alpha}}(F'_{y_{\alpha}})\otimes
\bigotimes^{\ell}_{1}
\Gamma^{n_{i}-i_{\alpha}}(F''_{y_{\alpha}})\longrightarrow
\Gr^{i}(\Gamma^{k}_{\ext}(F)_{x}){\quoi}.
\]
Le premier membre n'est autre que $(\Gamma^{i}_{\ext}(F')\vee
\Gamma^{k-i}_{\ext}(F''))_{x}$. On vérifie facilement que ces données,
construites point par point, proviennent d'une (unique) filtration en
$k+1$ crans de $\Gamma^{k}_{\ext}(F)$, et de morphismes (uniques)
$$\etiquette[5.5.11.1]{XVII.eq5.5.11.1}
\Gamma^{i}_{\ext}(F')\vee \Gamma^{k-i}_{\ext}(F'')\longrightarrow \Gr^{i}(\Gamma^{k}_{\ext}(F)){\quoi}.\leqno(5.5.11.1)
$$

Si la suite exacte $\Sigma$ est scindée, on a plus précisément
(vérification point par point à partir de \sisi{\hyperref[XVII.eq5.5.2.6.2]{5.5.2.6.2}}{\eqhyperref[XVII.eq5.5.2.6.1]{5.5.2.6.1}}).
$$\etiquette[5.5.11.2]{XVII.eq5.5.11.2}
\Gamma^{k}_{\ext}(F'\oplus F'')\simeq
\sisi{\sum}{\bigoplus}_{k_{1}+k_{2}=k}\Gamma^{k_{1}}_{\ext}(F')\vee
\Gamma^{k}_{\ext}(F''){\quoi}.\leqno(5.5.11.2)
$$

Si la suite exacte $\Sigma$ est scindable point par point, ou si $F''$
est plat, les morphismes \eqhyperref[XVII.eq5.5.11.1]{5.5.11.1} sont des isomorphismes (\cf
\hyperref[XVII.par5.5.1.4]{5.5.1.4}).

\subsubsection{}\etiquette[5.5.12]{XVII.subsubsec5.5.12}
Soit $u:X\rightarrow Y$ un morphisme de $S$-schémas. On suppose que
les puissances symétriques $\Sym_{S}^{k}(X)$ et $\Sym_{S}^{k}(Y)$
existent. Soient $F$ un Module sur $X$, $G$ un Module sur $Y$ et $v\in
\Hom(u^{\ast}G,F)\underset{\sisi{n}{}}{\simeq}\Hom(G,u_{\ast}F)$
un $u$-homomorphisme de $F$ dans $G$ (0.2). Si
$\prod\limits^{n} u$ est la puissance cartésienne
$n$\textsuperscript{ième} de $u$, l'homomorphisme $v$ induit un
$\prod\limits^{n} u$-homomorphisme
$\mathop{\boxtimes}\limits^n v$ de
$\mathop{\boxtimes}\limits^n F$ dans
$\mathop{\boxtimes}\limits^n G$. Par\pageoriginale passage
aux invariants, celui-ci définit
\[
TS^{n}(v):TS^{n}_{\ext}(F)\xrightarrow[\Sym_{S}^{n}(u)]{}
TS^{n}_{\ext}(G){\quoi}.
\]

On déduit de la définition \hyperref[XVII.subsubsec5.5.8]{5.5.8} qu'il existe un et un seul
$\Sym^{n}(u)$-homomorphisme $\Gamma^{n}_{\ext}(v)$, fonctoriel en
$(F,G,v)$,
$$\etiquette[5.5.12.0]{XVII.eq5.5.12.0}
\Gamma^{n}_{\ext}(v):\Gamma^{n}_{\ext}(F)\xrightarrow[\Sym^{n}_{S}(u)]{} \Gamma^{n}_{\ext}(G)\leqno(5.5.12.0)
$$
rendant commutatif le diagramme suivant (\cf (0.2) pour le sens
(inhabituel) des flèches):
$$
\xymatrix@C=3.5cm@R=2cm{
\Gamma^{n}_{\ext}(F) \ar[r]^{\Gamma^{n}_{\ext}(v)}_{\text{sur $\Sym^{n}_S(u)$}}& \Gamma^{n}_{\ext}(G)\\
TS^{n}_{\ext}(F)   \ar[r]^{TS^{n}_{\ext}(v)}_{\text{sur $\Sym^{n}_S(u)$}}\ar[u]^{\alpha}_{\text{sur $\Id$}}  & TS^{n}_{\ext}(G)\ar[u]^{\alpha}_{\text{sur $\Id$}}.
}
$$

On utilisera cette flèche surtout comme « flèche de Künneth
symétrique »: pour $G=u_{\ast}F$, et $v$ la flèche canonique, la
flèche $\Gamma^{n}_{\ext}(v)$ s'identifie à un homomorphisme\sisi{s}{} de
Modules
$$\etiquette[5.5.12.1]{XVII.eq5.5.12.1}
k^{n}:\Gamma^{n}_{\ext}(u_{\ast}F)\longrightarrow \Sym^{n}(u)_{\ast}\Gamma^{n}_{\ext}(F){\quoi}.\leqno(5.5.12.1)
$$

\subsubsection{}\etiquette[5.5.13]{XVII.subsubsec5.5.13}
Soit la condition:

\paragraph{}\etiquette[5.5.13.1]{XVII.par5.5.13.1}
$K\in \Ob K^{b,\geq 0}(X,\Aa)$ et, pour tout point géométrique
$x$ de $X$, le complexe $K_{x}$ est homotope à un complexe de modules
plats $L_{x}\in \Ob C^{b,\geq 0}(\Aa_{x})$.

\begin{enonce*}{Lemme 5.5.13.2}\etiquette[5.5.13.2]{XVII.lem5.5.13.2}
La\pageoriginale sous-catégorie $D^{b,\tor\leq 0}(X,\Aa)$ de
$D^{b}(X,\Aa)$ formée des complexes de \tTor-dimension $\leq 0$ \emph{\hyperref[XVII.def4.1.9]{4.1.9}}
est équivalente à la catégorie qui se déduit de la catégorie des
complexes $K\in \Ob K^{b,\geq 0}(X,\Aa)$ vérifiant
\emph{(\hyperref[XVII.par5.5.13.1]{5.5.13.1})} en inversant les quasi-isomorphismes. Cette
construction peut se faire indifféremment par calcul de fractions à gauche ou à droite.
\end{enonce*}

Soit $P$ la sous-catégorie de $K^{-}(X,\Aa)$ formée des
complexes point par point homotopes à un complexe de modules plats. On
sait (\textsc{Verdier} \sisi{\cite{SGA4-17:CD}}{\cite[I 2.4.2]{SGA4-17:CD}}) que $D^{-}(X,\Aa)$ se déduit de $P$ en
inversant les quasi-isomorphismes, par calcul de fractions à droite ou
à gauche. On applique alors \hyperref[XVII.lem5.5.4.1]{5.5.4.1} à l'inclusion de $K^{b,\geq
0}(X,\Aa)\cap P$ dans $P$. En effet, si $K$ est de
\tTor-dimension $\leq 0$, et point par point homotope à des complexes de
modules plats $L_{x}$, alors $\tau_{\geq 0}(K)$ est point par point
homotope aux $\tau_{\geq 0}(L_{x})$, et ceux-ci sont encore des
complexes de modules plats, de sorte que $\tau_{\geq 0}(K)$ vérifie
\hyperref[XVII.par5.5.13.1]{5.5.13.1}.

\subsubsection{}\etiquette[5.5.14]{XVII.subsubsec5.5.14}
D'après \hyperref[XVII.lem5.5.6.1]{5.5.6.1}, \sisi{\hyperref[XVII.eq5.5.7.2]{5.5.7.2}}{\eqhyperref[XVII.eq5.5.7.2]{5.5.7.2}} ou \sisi{\hyperref[XVII.eq5.5.8.1]{5.5.8.1}}{\eqhyperref[XVII.eq5.5.8.1]{5.5.8.1}} et
\hyperref[XVII.par5.5.2.4]{5.5.2.4}, tout
quasi-isomorphisme $v:K'\rightarrow K$ entre complexes vérifiant
(\hyperref[XVII.par5.5.13.1]{5.5.13.1}) induit des quasi-isomorphismes
\begin{align*}
TS^{k}(v) &: TS^{k}_{\ext}(K')\xrightarrow{\approx} TS^{k}_{\ext}(K)\\
\Gamma^{k}(v) &: \Gamma^{k}_{\ext}(K')\xrightarrow{\approx} \Gamma^{k}_{\ext}(K){\quoi}.
\end{align*}
Il en résulte que les foncteurs $\Gamma^{k}$ et $TS^{k}$ sont
dérivables (\hyperref[XVII.subsubsec5.5.5]{5.5.5}) sur la sous-catégorie $D^{b,\tor\leq
0}(X,\Aa)$ de $D^{b}(X,\Aa)$, et que les complexes
vérifiant \hyperref[XVII.par5.5.13.1]{5.5.13.1} sont déployés à gauche pour ces
foncteurs. D'après \hyperref[XVII.par5.5.2.4]{5.5.2.4},\pageoriginale $\Gamma^{k}$
et $TS^{k}$ induisent des foncteurs
\begin{align*}
LTS^{k} &: D^{b,\tor \leq 0}(X,\Aa)\longrightarrow D^{b,\tor
\leq 0}(\Sym^{k}_{S}(X),\Aa)\tag{5.5.14.1}\\
L\Gamma^{k} &: D^{b,\tor\leq 0} (X,\Aa)\longrightarrow
D^{b,\tor \leq 0}(\Sym^{k}_{S}(X),\Aa){\quoi}.\tag{5.5.14.2}
\end{align*}\etiquette[5.5.14.1]{XVII.eq5.5.14.1}
\etiquette[5.5.14.2]{XVII.eq5.5.14.2}%L'enlèvement de ce commentaire introduira un espace blanc non désiré.
De plus, le morphisme canonique $\alpha$ (\hyperref[XVII.subsubsec5.5.9]{5.5.9})
induit d'après \hyperref[XVII.subsubsec5.5.9]{5.5.9} (i) un isomorphisme entre les
foncteurs \eqhyperref[XVII.eq5.5.14.1]{5.5.14.1} et \eqhyperref[XVII.eq5.5.14.2]{5.5.14.2}
$$\etiquette[5.5.14.3]{XVII.eq5.5.14.3}
\alpha:L\Gamma^{k}\isomto  LTS^{k}{\quoi}.\leqno(5.5.14.3)
$$

\setcounter{paragraph}{3}
\paragraph{}\etiquette[5.5.14.4]{XVII.par5.5.14.4}
Les isomorphismes de \hyperref[XVII.subsubsec5.5.9]{5.5.9} (iv) (v) (vi) (vii)
passent à la catégorie dérivée et définissent des isomorphismes:
\begin{itemize}
\item[${\rm(iv')}$] Pour $\Aa'$ une $\Aa$-algèbre, on a
$$
L\Gamma^{k}_{\ext_{\Aa}}(K)\overset{\LL}{\otimes}
\Aa'\simeq
L\Gamma^{k}_{\ext_{\Aa'}}(K\overset{\LL}{\underset{\Aa}{\otimes}}\Aa'){\quoi}.
$$

\item[${\rm(v')}$] Pour tout changement de base $g:S'\rightarrow S$,
on a
$$
\Sym^{k}_{g}(g')^{\ast}L\Gamma^{k}_{\ext}K\simeq
L\Gamma^{k}_{\ext}({g'}^{\ast}K){\quoi}.
$$

\item[${\rm (vi')}$] Pour tout morphisme \emph{radiciel}
$i:X'\rightarrow X$, on a
$$
\Sym^{k}_{S}(i)^{\ast}L\Gamma^{k}_{\ext}K\simeq
L\Gamma^{k}_{\ext}i^{\ast}K{\quoi}.
$$

\item[${\rm (vii')}$] Pour $j:X'\rightarrow X$ une \emph{immersion},
on a:
$$
\Sym^{k}_{S}(j)_{!} L\Gamma^{k}_{\ext}\simeq
L\Gamma^{k}_{\ext}j_{!}{\quoi}.
$$
\end{itemize}

\subsubsection{}\etiquette[5.5.15]{XVII.subsubsec5.5.15}
Soit une suite exacte de complexes plats de $C^{\geq
0}(X,\Aa)$,
$$
0\longrightarrow K'\longrightarrow K\longrightarrow K''\longrightarrow
0{\quoi},
$$
et appliquons le foncteur $\Gamma^{k}_{\ext}$
(\hyperref[XVII.subsubsec5.5.3]{5.5.3}). On trouve
(comme en \hyperref[XVII.subsubsec5.5.11]{5.5.11}) que le Module semi-simplicial
$\Gamma^{k}_{\ext}(N^{-1}K)$ sur $\Sym^{k}_{S}(X)$ admet une
filtration dont\pageoriginale les quotients successifs sont les
Modules semi-simpliciaux $\Gamma^{k'}_{\ext}(N^{-1}K')\vee
\Gamma^{k''}_{\ext}(N^{-1}K'')$ $(k'+k''=k)$.

Passant à la catégorie dérivée on trouve que $L\Gamma^{k}_{\ext}(K)$
est représenté par un complexe filtré dont les quotients successifs
représentent les objets
$L\Gamma^{i}_{\ext}(K'\sisi{}{)}\overset{\LL}{\vee}L\Gamma^{k-i}_{\ext}
(K'')$ de la catégorie dérivée (appliquer \hyperref[XVII.par5.5.2.4]{5.5.2.4} et le
théorème de
\textsc{Eilenberg-Zilber-Cartier} (isomorphisme dans $\mathcal{K}$ d'un objet
semi-simplicial double et de sa diagonale: \textsc{Dold\sisi{ }{-}Puppe} \cite[2.9 \sisi{pg.213}{p.~213}]{SGA4-17:DP}).

Dans la catégorie dérivée, cette construction peut déjà se faire à
partir d'une suite exacte de complexes vérifiant \hyperref[XVII.par5.5.13.1]{5.5.13.1}.

${}^{\ast}$ En termes plus intrinsèques, soit $\Delta$ un « \emph{vrai triangle} »  d'objets de $D^{b,\tor\leq
0}(X,\Aa)$, \ie un objet de la « catégorie dérivée filtrée » avec $\Gr^{i}(\Delta)=0$ pour $i\neq 0$, $1$ et $\Gr^{i}(\Delta)\in
\Ob D^{b,\tor\leq 0}(X,\Aa)$ (\textsc{Deligne}, in Appendice à \textsc{Verdier}
\cite{SGA4-17:V2}, ou \refvb[VI B nᵒ]{Vbis.subsubsec2.5.1}, \refnoexpose{Vbis.subsubsec2.5.2}).

La construction précédente fournit alors un objet
$L\Gamma^{k}_{\ext}(\Delta)$ de la catégorie dérivée filtrée de
$\Sym^{k}_{S}(X)$, des isomorphismes
\[
\Gr^{i}(L\Gamma^{k}_{\ext}(\Delta))\simeq
L\Gamma^{i}_{\ext}(\Gr^{1}(\Delta))\overset{\LL}{\vee}L\Gamma^{k-i}_{\ext}
(\Gr^{0}(\Delta)\sisi{}{)}
\]
et, désignant par $\lvert\quad \rvert$ le complexe sous-jacent à un
complexe filtré, un isomorphisme
\[
\lvert L\Gamma^{k}_{\ext}(\Delta))\rvert \simeq L\Gamma^{k}_{\ext}(\lvert\Delta
\rvert){\quoi}.{}_{\ast}
\]

\begin{enonce*}{Lemme 5.5.16}\etiquette[5.5.16]{XVII.lem5.5.16}
Soit $N$ un entier. La catégorie $D^{b,\tor\leq 0}(X,\Aa)$ se
déduit de la catégorie des complexes $K\in \Ob K^{b,\geq
0}(X,\Aa)$ vérifiant les conditions {\rm (i)} et {\rm (ii)}
ci-dessous en inversant les quasi-isomorphismes:
\begin{enumerate}
\renewcommand{\theenumi}{\roman{enumi}}
\renewcommand{\labelenumi}{(\theenumi)}
\item condition\pageoriginale (\hyperref[XVII.par5.5.13.1]{5.5.13.1}).

\item Le complexe $K$ est déployé pour tout foncteur image directe de
dimension cohomologique $\leq N$, relatif à un isomorphisme de topos
$f:\widetilde{X}_{\text{\rm et}}\rightarrow T$.
\end{enumerate}
\end{enonce*}

Qu'il y ait un calcul des fractions à droite résulte de \hyperref[XVII.subsubsec5.5.13]{5.5.13} et du
fait que, via les résolutions flasques canoniques tronquées (\hyperref[XVII.subsubsec4.2.9]{4.2.9}),
tout complexe $K$ vérifiant (i) s'envoie dans un complexe vérifiant
(i) et (ii).

\setcounter{subsubsection}{16}

\subsubsection{}\etiquette[5.5.17]{XVII.subsubsec5.5.17}
Soit un morphisme de $S$-schémas $u:X\rightarrow Y$. On suppose que
$\Sym^{n}_{S}(X)$ et $\Sym^{n}_{S}(Y)$ existent, qu'il existe un
entier $N$ tel que le foncteur $u_{\ast}$ soit de dimension $\leq N$
sur la catégorie des $\Aa$-Modules, et que le foncteur
$Ru_{\ast}$ transforme complexes bornés de \tTor-dimension $\leq 0$ en
complexes de \tTor-dimension $\leq 0$. Ces deux dernières conditions
sont automatiques si $u$ est propre, $\Aa$ de torsion et $Y$
cohérent (\ie quasi-compact, quasi-séparé) (\hyperref[XVII.cor5.2.8.1]{5.2.8.1}, \hyperref[XVII.thm5.2.10]{5.2.10}).\nde{la cohérence de $Y$ peut être remplacée par la condition plus faible que la dimension des fibres de $u$ est bornée.}

Soit $K\in \Ob K^{b,\geq 0}(X,\Aa)$ un complexe vérifiant les
conditions (i) et (ii) de \hyperref[XVII.lem5.5.16]{5.5.16}, avec $N$ comme ci-dessus. D'après
\hyperref[XVII.subsubsec5.5.14]{5.5.14}, la flèche
$$
\beta:L\Gamma^{n}_{\ext}(K)\isomto  \Gamma^{n}_{\ext}(K)
$$
est un quasi-isomorphisme. De même, la flèche
$$
\gamma:u_{\ast}K\isomto  Ru_{\ast}K
$$
est un quasi-isomorphisme.

On\pageoriginale définit la \emph{flèche de K\"unneth symétrique} de
$L\Gamma^{n}_{\ext} Ru_{\ast}(K)$ dans
$R\Sym^{n}(u)_{\ast}L\Gamma^{n}_{\ext}(K)$ comme le composé
\begin{multline*}
k^{n}: L\Gamma^{n}_{\ext}Ru_{\ast}K\xleftarrow[\gamma]{\sim}
L\Gamma^{n}u_{\ast}K\rightarrow
\Gamma^{n}_{\ext}u_{\ast}K\xrightarrow[5.5.12.1]{}\Sym^{n}_{S}(u)_{\ast}\Gamma^{n}_{\ext}(K)\\
\longrightarrow R
\Sym^{n}_{S}(u)_{\ast}\Gamma^{n}_{\ext}(K)\xleftarrow[\beta]{\sim} R \Sym^{n}_{S}(u)_{\ast}L\Gamma^{n}_{\ext}(K){\quoi}.\tag{5.5.17.1}
\end{multline*}\etiquette[5.5.17.1]{XVII.eq5.5.17.1}%L'espace vertical suivant ce paragraphe est trop large.

Il résulte de \hyperref[XVII.lem5.5.16]{5.5.16} que les morphismes \eqhyperref[XVII.eq5.5.17.1]{5.5.17.1} définissent un
morphisme entre foncteurs de $D^{b,\tor\leq 0}(X,\Aa)$ dans
$D^{+}(\Sym^{n}_{S}(Y),\Aa)$
$$\etiquette[5.5.17.2]{XVII.eq5.5.17.2}
k^{n}:L\Gamma^{n}_{\ext}Ru_{\ast}\longrightarrow R \Sym^{n}_{S}u_{\ast}L\Gamma^{n}_{\ext}{\quoi},\leqno(5.5.17.2)
$$
appelé encore \emph{morphisme de K\"unneth symétrique}.

Ce morphisme est caractérisé par la propriété \hyperref[XVII.prop5.5.18]{5.5.18} ci-dessous.

\setcounter{paragraph}{2}
\paragraph{}\etiquette[5.5.17.3]{XVII.par5.5.17.3}
Soient $S$, $X$, $Y$ et $u$ comme plus haut, $K\in \Ob C^{b,\geq
0}(X,\Aa)$ un complexe de \tTor-dimension $\leq 0$, $K_{1}\in
\Ob C^{b,\geq 0}(Y,\Aa)$ et $K_{2}\in \Ob
C^{b}(\Sym^{n}_{S}(X),\Aa)$. Soit $v$ un $u$-homomorphisme
(0.2) de $K$ dans $K_{1}$, et soit $t$ un homomorphisme (au sens
ordinaire, comme tous ceux qui suivent) de $\Gamma^{n}_{\ext}(K)$ dans
$K_{2}$. On dispose de morphismes
\begin{enumerate}
\renewcommand{\theenumi}{\alph{enumi}}
\renewcommand{\labelenumi}{\theenumi)}
\itemsep=8pt
\item $\Gamma^{n}_{\ext}(v):\Gamma^{n}_{\ext}(K_{1})\rightarrow
\Sym^{n}_{S}(u)_{\ast} \Gamma^{n}_{\ext}(K)$\qquad \eqhyperref[XVII.eq5.5.12.0]{5.5.12.0}

\item
$\Sym^{n}_{S}(u)_{\ast}(t):\Sym^{n}_{S}(u)_{\ast}\Gamma^{n}_{\ext}(K)\rightarrow
\Sym^{n}_{S}(u)_{\ast}K_{2}$,
\end{enumerate}
de composé
$$
\varphi':\Gamma^{n}_{\ext}(K_{1})\longrightarrow
\Sym^{n}_{S}(u)_{\ast}K_{2}{\quoi},
$$
définissant, par composition
$$
\varphi:L\Gamma^{n}_{\ext}(K_{1})\longrightarrow
R\Sym^{n}(u)_{\ast}K_{2}{\quoi}.
$$
De\pageoriginale plus, $v$ et $t$ induisent:
$$
\begin{cases}
v': K_{1}\longrightarrow Ru_{\ast}K\\
t': L\Gamma^{n}_{\ext}K\longrightarrow K_{2}{\quoi}.
\end{cases}
$$

On a alors:

\begin{enonce*}{Proposition 5.5.18}\etiquette[5.5.18]{XVII.prop5.5.18}
Sous les hypothèses \hyperref[XVII.par5.5.17.3]{5.5.17.3}, le diagramme suivant est commutatif
dans $D^{+}(\Sym_{S}(Y),\Aa)$:
\[
\xymatrix@C=4cm@R=1.75cm{
  L\Gamma^{n}_{\ext} Ru_{\ast}K \ar[r]^{k^n}_{\text{\textsuperscript{(K\"unneth symétrique)}}}& R \Sym^{n}_{S}(u)_{\ast} L
\Gamma^{n}_{\ext}K \ar[d]_{R
  \Sym^{n}_{S}(u)_{\ast}(t')}\\
L\Gamma^{n}_{\ext} K_1  \ar[r]^{\varphi}\ar[u]^{L\Gamma^{n}_{\ext}(v')} & R \Sym^{n}_{S}(u)_{\ast}K_2.
}
\]
\end{enonce*}

La démonstration est toute pareille à celle de \hyperref[XVII.par5.4.1.5]{5.4.1.5} et laissée au
lecteur.

\setcounter{subsubsection}{18}
\subsubsection{On suppose dorénavant $\Aa$ de
  torsion}\etiquette[5.5.19]{XVII.subsubsec5.5.19}

Soit $u:X\rightarrow Y$ un morphisme compactifiable de $S$-schémas:
\[
\xymatrix@C=3cm@R=1.75cm{
   & \bar{X} \ar[rd]^{\bar{u}}&\\
X \ar@{^{(}->}[ru]^{j} \ar[rr]^{u} \ar[rd]&         & Y \ar[ld]\\
   &   S     & .
}
\]
On\pageoriginale suppose que $\Sym^{n}_{S}(X)$,
$\Sym^{n}_{S}(\bar{X})$ et $\Sym^{n}_{S}(Y)$ existent. Pour tout
Module $F$ sur $X$, on a alors
$$\etiquette[5.5.19.1]{XVII.eq5.5.19.1}
\Gamma^{n}_{\ext}(j_{!}F)\simeq \Sym^{n}_{S}(j)_{!}\Gamma^{n}_{\ext}(F){\quoi},\leqno(5.5.19.1)
$$
et cet isomorphisme passe trivialement aux catégories dérivées\sisi{}{ (\hyperref[XVII.par5.5.14.4]{5.5.14.4} (vii'))}. On
définit le \emph{morphisme de K\"unneth symétrique} entre foncteurs de
$D^{b,\tor\leq 0}(X,\Aa)$ dans $D^{b,\tor\leq
0}(\Sym^{n}_{S}(Y),\Aa)$\nde{Comment montrer que $\Sym_S(u)$ est de type fini en général ? Cette condition est nécessaire pour que $R\Sym_S(u)_!$ soit défini.}
$$\etiquette[5.5.19.2]{XVII.eq5.5.19.2}
k^{n}_{u}\quad\mbox{\rm ou}\quad
k^{n}:L\Gamma^{n}_{\ext}Ru_{!}\longrightarrow
R\Sym_{S}(u)_{!}L\Gamma^{n}_{\ext}\leqno(5.5.19.2)
$$
comme étant le composé
\begin{align*}
L\Gamma^{n}_{\ext}Ru_{!}\stackrel{\dfn}{\simeq}L\Gamma^{n}_{\ext}R\bar{u}_{\ast}j_{!}
& \xrightarrow{\eqhyperref[XVII.eq5.5.17.2]{5.5.17.2}}
R\Sym^{n}_{S}(\bar{u})_{\ast}L\Gamma^{n}_{\ext j!}\\
& \xrightarrow[\sisi{}{\text{(\hyperref[XVII.par5.5.14.4]{5.5.14.4} (vii'))}}]{\sisi{\substack{\eqhyperref[XVII.eq5.5.19.1]{5.5.19.1}\\\sim}}{\sim}}
R\Sym^n_s(\bar{u})_*\Sym^{n}_{S}(j)_{!}L\Gamma^{n}_{\ext}\\
 & \xrightarrow{\sisi{\text{(5.5.13.5 (vii'))}}{}}
 R\Sym^{n}_{S}(u)_{!}L\Gamma^{n}_{\ext}
\end{align*}

\setcounter{subsubsection}{19}
\subsubsection{}\etiquette[5.5.20]{XVII.subsubsec5.5.20}
On laisse au lecteur le soin de faire les vérifications suivantes,
faciles en principe en utilisant \hyperref[XVII.prop5.5.18]{5.5.18}.

\paragraph{}\etiquette[5.5.20.1]{XVII.par5.5.20.1}
(Caractère intrinsèque). Le morphisme \eqhyperref[XVII.eq5.5.19.2]{5.5.19.2} est indépendant du
choix de la compactification $\bar{u}$ de $u$.

\paragraph{}\etiquette[5.5.20.2]{XVII.par5.5.20.2}
(Transitivité). Soit un diagramme de $S$-schémas
$$
X\xrightarrow{u}Y\xrightarrow{v}Z{\quoi},
$$
avec\pageoriginale $u$ et $v$ quasi-projectifs; on suppose que $\Sym^{n}_{S}(Z)$
existe (donc aussi $\Sym^{n}_{S}(Y)$ et $\Sym^{n}_{S}(X)$ (\SGA 1 V
1.7, \EGA II 4.5.4)). Le diagramme de foncteurs suivant est alors
commutatif
$$
\xymatrix@C=1.5cm@R=1.75cm{
L\Gamma^{n}_{\ext} R(vu)_{!} \ar@{=}[d] \ar[rr]^{k^n_{vu}}& & R \Sym^{n}_S(vu)_{!}L\Gamma^{n}_{\ext}\ar@{=}[d]\\
L\Gamma^{n}_{\ext} Rv_{!} Ru_{!} \ar[r]^{k^n_v}& R \Sym^{n}_S(v)_!
L\Gamma^{n}_{\ext} Ru_{!} \ar[r]^{k^n_u}& R \Sym^{n}_S(v)_{!} R \Sym^{n}_S(u)_{!}L\Gamma^{n}_{\ext}.
}
$$

\paragraph{}\etiquette[5.5.20.3]{XVII.par5.5.20.3}
Le morphisme de K\"unneth symétrique est compatible aux changements de
base ou d'espace du type 5.5.9 (iv) (v) (vi). De façon plus précise,
on a les compatibilités suivantes avec les morphismes \sisi{5.5.13.5}{\hyperref[XVII.par5.5.14.4]{5.5.14.4}}
${\rm(iv')}$ ${\rm (v')}$ ${\rm (vi')}$.
\begin{itemize}
\item[${\rm(iv'')}$] Pour toute $\Aa$-algèbre $\Aa'$,
et tout $S$-morphisme quasi-projectif $u:X\rightarrow Y$ le diagramme
\[\def\objectstyle{\scriptstyle}
\xymatrix@C=.8cm@R=1.75cm{
L\Gamma^{n}_{\ext}Ru_{!}(K)\overset{\LL}{\otimes}_{\Aa}\Aa'
\ar@{-}[r]^{\sim} \ar[d]^{k^n_u}&
L\Gamma^{n}_{\ext}(Ru_{!}(K)\overset{\LL}{\otimes}_{\Aa}\Aa')
\ar@{->}[r]^{\sim}&
L\Gamma^{n}_{\ext}Ru_{!}(K\overset{\LL}{\otimes}_{\Aa}\Aa') \ar[d]^{k^n_u}\\
R\Sym^{n}_{S}(u)_{!}L\Gamma^{n}_{\ext}(K)\overset{\LL}{\otimes}_{\Aa}\Aa'
\ar@{->}[r]^{\sim}&
R\Sym^{n}_{S}(u)_{!}(L\Gamma^{n}_{\ext}(K)\overset{\LL}{\otimes}_{\Aa}\Aa')
\ar@{-}[r]^{\sim}&
R\Sym^{n}_{S}(u)_{!}L\Gamma^{n}_{\ext}(K\overset{\LL}{\otimes}_{\Aa}\Aa')
}
\]
est commutatif.

\item[${\rm(v'')}$] Pour\pageoriginale tout changement de base
$g:S'\rightarrow S$, donnant lieu à un diagramme \sisi{cartésien}{à carrés cartésiens}
$$
\xymatrix{
X' \ar[rrr]^{g'}\ar[rd]^{u'}\ar[dd]&   & & X \ar[rd]^{u}\ar[dd]|\hole &   \\
   & Y' \ar[ld]\ar[rrr]^{g''}&  & & Y \ar[ld]\\
S' \ar[rrr]_{g}&  &  & S &   ,
}
$$
le diagramme de foncteurs, tout analogue à ${\rm(iv'')}$, est
commutatif:
$$\def\objectstyle{\scriptstyle}
\xymatrix@R=2.25cm{
\Sym^{n}_g(g'')^{\ast}L\Gamma^{n}_{\ext}Ru_{!}  \ar@{-}[r]^{\sim} \ar[d]^{k^n_u}&
L\Gamma^{n}_{\ext}g''^{\ast}Ru_{!} \ar[r]^{\sim}& L\Gamma^{n}_{\ext}Ru'_{!}
g'^{\ast}\ar[d]^{k^n_{u'}} \\
\Sym^{n}_g(g'')^{\ast}R \Sym^{n}_S(u)_{!}L\Gamma^{n}_{\ext} \ar[r]^{\sim}& R
\Sym^{n}_S(u')_{!} \Sym^{n}_{g}(g')^{\ast}L\Gamma^{n}_{\ext} \ar@{-}[r]^\sim& R \Sym^{n}_{S}(u')_{!}L\Gamma^{n}_{\ext}g'^{\ast}.
}
$$

\item[${\rm(vi'')}$] Pour tout changement de base \emph{radiciel}
$g:Y'\rightarrow Y$, donnant lieu à un diagramme cartésien de
$S$-schémas
$$
\xymatrix@C=2.5cm@R=1.5cm{
X' \ar[r]^{g'}\ar[d]^{u'}& X \ar[d]^{u}\\
Y' \ar[r]^{g} & Y,
}
$$
le diagramme de foncteurs suivant est commutatif:
\[
\xymatrix{
\Sym^{n}_S(g)^{\ast}L \Gamma^{n}_{\ext} Ru_{!} \ar@{-}[r]^{\sim} \ar[d]^{k^n_u}&
L\Gamma^{n}_{\ext}g^{\ast} Ru_{!} \ar[r]^{\sim}&  L\Gamma^{n}_{\ext}
Ru'_{!} g'^{\ast} \ar[d]^{k^n_{u'}}\\
\Sym^{n}_{S}(g)^{\ast}R \Sym^{n}_{S}(u\sisi{_!)}{)_!L\Gamma^n_{\ext}} \ar[r]^{\sim}& R
\Sym^{n}_{S}(u')_{!}{g'}^{\ast}L\Gamma^{n}_{\ext} \ar@{-}[r]^{\sim}& R \Sym^{n}_{S}(u')_{!}L\Gamma^{n}_{\ext}g'^{\ast}
}
\]
\end{itemize}

\begin{enonce*}{Théorème 5.5.21}\etiquette[5.5.21]{XVII.thm5.5.21}
{\rm(Formule\pageoriginale de K\"unneth symétrique)}. Soient $S$ un schéma cohérent (=
quasi-compact quasi-séparé) annelé par un faisceau d'anneaux
commutatifs de torsion $\Aa$, à fibres des anneaux
noethériens. On annèle chaque $S$-schéma par l'image réciproque de
$\Aa$. Soit $u:X\rightarrow Y$ un morphisme de $S$-schémas
quasi-projectifs. Quel\sisi{}{s} que soient l'entier $n\geq 0$ et $K\in \Ob
D^{b,\tor\leq 0}(X,\Aa)$, la flèche de K\"unneth symétrique
\eqhyperref[XVII.eq5.5.17.2]{5.5.17.2}
$$
k^{n}:L\Gamma^{n}_{\ext}Ru_{!}K\longrightarrow
R\Sym^{n}_{S}(u)_{!}L\Gamma^{n}_{\ext}K
$$
est un isomorphisme dans $D^{b,\tor \leq
  0}(\Sym^{n}_{S}(Y),\Aa)$.
\end{enonce*}

La démonstration de ce théorème occupera toute la suite de 5.5.

\begin{enonce*}[remark]{Remarque 5.5.21.1}\etiquette[5.5.21.1]{XVII.rem5.5.21.1}
L'hypothèse de quasi-projectivité sert seulement à assurer que $u$
admet une compactification $\bar{u}$
$$
X\xrightarrow{j}\bar{X}\xrightarrow{\bar{u}}Y
$$
telle que les schémas $\Sym^{n}_{S}(X)$, $\Sym^{n}_{S}(\bar{X})$ et
$\Sym^{n}_{S}(Y)$ existent pour tout $n$. Le langage des espaces
algébriques devrait permettre de remplacer cette hypothèse par la
seule hypothèse que $u$ est compactifiable.
\end{enonce*}

La\pageoriginale première partie de la démonstration consiste à
ramener par dévissage \hyperref[XVII.thm5.5.21]{5.5.21} à un cas « irréductible »:

\begin{enonce*}{Lemme 5.5.22}\etiquette[5.5.22]{XVII.lem5.5.22}
Le théorème \hyperref[XVII.thm5.5.21]{5.5.21} est impliqué par le cas particulier suivant:
\begin{itemize}
\item[${(\ast)}$] Pour tout corps algébriquement clos $k$, toute courbe
projective et lisse $X$ sur $k$ et tout nombre premier $\ell$, les
flèches de K\"unneth symétriques (obtenues pour $Y=S=\Spec(k)$,
$\Aa=\ZZ/\ell$, $K=\ZZ/\ell$; voir explications
ci-dessous)
$$\etiquette[5.5.22.1]{XVII.eq5.5.22.1}
L\Gamma^{n}(R\Gamma(X,\ZZ/\ell))\longrightarrow
R\Gamma(\Sym^{n}_{k}(X),\ZZ/\ell)\leqno(5.5.22.1)
$$
sont des isomorphismes $(n\geq 0)$.
\end{itemize}
\end{enonce*}

Dans \sisi{\hyperref[XVII.eq5.5.22.1]{5.5.22.1}}{\eqhyperref[XVII.eq5.5.22.1]{5.5.22.1}}, le symbole $\Gamma$ a un malencontreux
double usage:
dans l'expression $R\Gamma$, il désigne le foncteur « sections
globales »; dans l'expression $L\Gamma^{n}$, il désigne le foncteur
\hyperref[XVII.par5.5.2.3]{5.5.2.3}. On utilise dans cette formule l'isomorphisme
$\Gamma^{n}_{\ext}(\ZZ/\ell)\simeq \ZZ/\ell$.

Soient donnés $S$, $X$, $Y$, $u$ et $\Aa$ comme en \hyperref[XVII.thm5.5.21]{5.5.21}.

\begin{enonce*}{Lemme 5.5.22.2}\etiquette[5.5.22.2]{XVII.lem5.5.22.2}
Soit une suite exacte de complexes dans $C^{b,\geq 0}(X,\Aa)$,
de \ttor-dimension $\leq 0$:
$$
0\longrightarrow K'\longrightarrow K\longrightarrow K''\longrightarrow
0{\quoi}.
$$
Si les flèches de K\"unneth $k^{n}\;(n\geq 0)$ relatives à deux des
complexes $K'$, $K$ et $K''$ sont toutes des isomorphismes, alors les
flèches de K\"unneth relatives au troisième sont aussi des isomorphismes.
\end{enonce*}

On\pageoriginale se ramène au cas où $u$ est projectif, et où les
complexes $K'$, $K$, $K''$ vérifient les conditions (i) et (ii) de
\hyperref[XVII.lem5.5.16]{5.5.16}, pour $N$ au moins égal au double de la dimension des flèches
de $u$. Les flèches de Künneth symétriques sont alors données par
\eqhyperref[XVII.eq5.5.17.1]{5.5.17.1}.

Montrons par récurrence sur $n$ que si pour $i\leq n$ et deux des
complexes $K'$, $K$, $K''$, les flèches $k^{i}$ sont des
isomorphismes, alors les flèches $k^{i}\;(i\leq n)$ sont des
isomorphismes pour le troisième.

Pour le voir, on note que grâce à \sisi{\hyperref[XVII.subsubsec5.5.14]{5.5.14}}{\hyperref[XVII.subsubsec5.5.15]{5.5.15}}, le morphisme
\begin{equation*}
k^{n}:L\Gamma^{n}_{\ext}Ru_{\ast}K\longrightarrow R\Sym^{n}(u)_{\ast}L\Gamma^{n}_{\ext}K\tag{a}
\end{equation*}
peut se représenter par un morphisme de complexes filtrés
$$
{k'}^{n}:K_{0}\longrightarrow K_{1}
$$
de filtration finie, ce morphisme induisant sur les quotients
successifs des flèches isomorphes aux flèches composées
\begin{equation*}
L\Gamma^{i}_{\ext}Ru_{\ast}K'\overset{\LL}{\vee}
L\Gamma^{n-i}_{\ext} Ru_{\ast}K''\longrightarrow
R\Sym^{n}(u)_{\ast}(L\Gamma^{i}_{\ext}K'\overset{\sisi{}{\LL}}{\vee}
L\Gamma^{n-i}_{\ext}K'')\tag{$\text{b}_i$}
\end{equation*}
suivantes:
\begin{gather*}
L\Gamma^{i}_{\ext}Ru_{\ast}K'\overset{\LL}{\vee}
L\Gamma^{n-i}_{\ext}Ru_{\ast}K''\xrightarrow{k^{i}\overset{\LL}{\vee}k^{n-i}}
R\Sym^{i}(u)_{\ast}L\Gamma^{i}_{\ext}K'\overset{\LL}{\vee}
R\Sym^{n-i}(u)_{\ast}L\Gamma^{n-i}_{\ext}(K'')\\
\underset{\defn}{=}
\vee_{i,n-i {\ast}}(R\Sym^{i}(u)_{\ast}L\Gamma^{i}_{\ext}K'\overset{\LL}{\boxtimes}
R\Sym^{n-i}(u)_{\ast}L\Gamma^{n-i}_{\ext}(K''))\\
\xrightarrow{\substack{\hyperref[XVII.thm5.4.3]{5.4.3}\\\sim}}\vee_{i,n-i {\ast}}\sisi{(R\Sym^{i}(u)_{\ast}\times R \Sym^{n-i}(u)_{\ast})}{(R\Sym^i(u)\times R\Sym^{n-i}(u))_\ast}
(L\Gamma^{i}_{\ext}K'\overset{\LL}{\boxtimes}
L\Gamma^{n-i}_{\ext}K'')\\
=R\Sym^{n}(u)_{\ast}\vee_{i,n-i{\ast}}(L\Gamma^{i}_{\ext}K'\overset{\LL}{\boxtimes}L\Gamma^{n-i}_{\ext}K'')\overset{\dfn}{=} R\Sym^{n}(u)_{\ast}(L\Gamma^{i}_{\ext}K'\overset{\LL}{\vee}L\Gamma^{n-i}_{\ext}K''){\quoi}.
\end{gather*}

D'après\pageoriginale l'hypothèse de récurrence (et la formule de
Künneth), ces flèches sont des quasi-isomorphismes pour $i\neq 0$,
$n$. Deux quelconques des assertions suivantes impliquent donc la
troisième: (a) est un quasi-isomorphisme; $(\text{b}_0)$ est un
quasi-isomorphisme; $(\text{b}_n)$ est un quasi-isomorphisme. Il s'agit
là des trois flèches qui étaient en question, \cqfd

\setcounter{subsubsection}{22}
\setcounter{paragraph}{2}
\paragraph{}\etiquette[5.5.22.3]{XVII.par5.5.22.3}
Supposons vraie l'assertion $(\ast)$ envisagée dans \hyperref[XVII.lem5.5.22]{5.5.22}, et
prouvons \hyperref[XVII.thm5.5.21]{5.5.21}. Soient donc $S$, $X$, $Y$, $\Aa$, $K$ et $u$
comme en \hyperref[XVII.thm5.5.21]{5.5.21}. Le problème est local sur $S$, qu'on peut supposer
affine. Tout complexe $K$ est extension successive de complexes de la
forme $i_{!}i^{\ast}K$ pour $i$ inclusion dans $X$ d'une partie affine
localement fermée. Par dévissage (\hyperref[XVII.lem5.5.22.2]{5.5.22.2}), on peut donc supposer $K$
de la forme $i_{!}K_{1}$, pour une immersion $i:X_{1}\rightarrow X$ de
source affine. La compatibilité \hyperref[XVII.par5.5.20.2]{5.5.20.2} permet de remplacer $X$ par
$X_{1}$, $K$ par $K_{1}$ et de supposer donc $u$ affine. Factorisons
$u$ comme un composé $u=u_{1}\ldots u_{k}$ de morphismes à fibres de
dimension $\leq 1$. La compatibilité \hyperref[XVII.par5.5.20.2]{5.5.20.2} montre qu'il suffit de
prouver que les flèches $k^{n}$ relatives aux $u_{i}$ et à un
quelconque complexe $K$ sont des isomorphismes. On est ainsi ramené au
cas $S$ \emph{affine, $u$ de dimension relative $\leq 1$}.

\paragraph{}\etiquette[5.5.22.4]{XVII.par5.5.22.4}
Pour que $k^{n}$ soit un isomorphisme, il suffit que ses restrictions
aux flèches géométriques de $\Sym^{n}_{S}(Y)$ le soient; la
compatibilité \hyperref[XVII.par5.5.20.3]{5.5.20.3} ${\rm(v'')}$ nous ramène donc à supposer $S$
\emph{spectre d'un corps algébriquement clos $k$.}

\paragraph{}\etiquette[5.5.22.5]{XVII.par5.5.22.5}
Pour\pageoriginale tout point rationnel $y$ de $\Sym^{n}_{k}(Y)$, il
existe une partie réduite finie $Y'$ de $Y$ telle que $\sisi{x}{y}$ soit dans
l'image de $\Sym^{n}_{k}(Y')$. La compatibilité \hyperref[XVII.par5.5.20.3]{5.5.20.3} ${\rm(vi'')}$
nous permet donc de supposer que $Y$ \emph{est somme d'un nombre fini
de copies de} $\Spec(k)$.

\paragraph{}\etiquette[5.5.22.6]{XVII.par5.5.22.6}
Si $Y=\amalg Y_{i}$, alors tout complexe $K$ sur $X$ est somme de
complexes concentrés sur l'un des schémas $X_{i}=u^{-1}(Y_{i})$. Par
dévissage (\hyperref[XVII.lem5.5.22.2]{5.5.22.2}) et d'après la compatibilité
\hyperref[XVII.par5.5.20.3]{5.5.20.3}
${\rm (vi'')}$, il suffit de prouver \hyperref[XVII.thm5.5.21]{5.5.21} pour chacun des morphismes
$u_{i}:X_{i}\rightarrow Y_{i}$. Par \hyperref[XVII.par5.5.22.5]{5.5.22.5}, ceci nous ramène au cas
où
$$
\text{\rm $Y=S=\Spec(k)$, $k$ algébriquement clos, $\dim X\leq 1$.}
$$

Le faisceau $\Aa$ se réduit maintenant à un anneau noethérien
$A$. La proposition suivante permettra de se ramener au cas
où $A$ est un corps.

\begin{enonce*}{Proposition 5.5.23}\etiquette[5.5.23]{XVII.prop5.5.23}
Soient $A$ un anneau noethérien (commutatif), $K\in \Ob D^{-}(A)$ un
complexe borné supérieurement de $A$-modules et $n$ un entier. Les
conditions suivantes sont équivalentes
\begin{enumerate}
\renewcommand{\theenumi}{\roman{enumi}}
\renewcommand{\labelenumi}{(\theenumi)}
\item $K$ est de \tTor-dimension $\leq n$

\item Pour tout $x\in \Spec (A)$ et tout $i<-n$,
$H^{i}(K\overset{\LL}{\underset{A}{\otimes}} k(x))=0$.
\end{enumerate}
\end{enonce*}

On a (i) $\Rightarrow$ (ii) trivialement\sisi{,}{.} Supposons (ii). Pour prouver
(i), nous prouverons par récurrence noethérienne sur le fermé $F$ de
$\Spec(A)$ que si un module $N$ vérifie $N_{x}=0$ pour $x\not\in F$,
alors $H^{i}(K\overset{\LL}{\otimes}_{A}N)=0$
pour $i<-n$.

Soit\pageoriginale $I$ l'unique idéal qui définisse $F$ tel que $A/I$
soit réduit, et soit $N_{k}$ le sous-module de $N$ annulé par
$I^{k}$. On a $N=\varinjlim N_{k}$. Puisque les foncteurs Tor
commutent aux limites inductives filtrantes, on se ramène au cas où il
existe $k$ tel que $N=N_{k}$, de sorte que $N$ est extension
successive de modules tués par $I$, et la suite exacte longue de
cohomologie nous ramène au cas où $N$ est déjà un $A/I$-module.

Soit $\eta$ un point maximal de $F$, et définissons des $A/I$-modules
$M_{1}$ et $M_{2}$ par la suite exacte
$$\etiquette[5.5.23.1]{XVII.eq5.5.23.1}
0\longrightarrow M_{1}\longrightarrow N \xrightarrow{i} N\otimes
k(\eta)\longrightarrow M_{2} \longrightarrow 0{\quoi}.\leqno(5.5.23.1)
$$

On a $(M_{j})_{\eta}=0$ $(j=1,2)$. Si pour chaque fermé $G\subset F$
avec $\eta\not\in G$ on désigne par $M_{j,G}$ le sous-module de
$M_{j}$ formé des sections de $M_{j}$ à support dans $G$, on a
$$
M_{j}=\varinjlim_{\eta\not\in G} M_{j,G}{\quoi}.
$$
D'après l'hypothèse de récurrence et un passage à la limite, on a donc
$$
H^{i}(K\overset{\LL}{\otimes}_{A} M_j)=0\quad\mbox{\rm
pour}\quad i< -n{\quoi}.
$$

Le module $N\otimes k(\eta)$ est un vectoriel sur $k(\eta)$, donc une
somme de copies de $k(\eta)$. Par hypothèse, on a
$$
H^{i}(K\overset{\LL}{\otimes} (N\otimes
k(\eta))\sisi{}{)}=0\quad\mbox{\rm pour}\quad i<-n{\quoi}.
$$

Si $I$ est l'image de la flèche $i$ de \eqhyperref[XVII.eq5.5.23.1]{5.5.23.1}, on dispose de
suites exactes
\begin{align*}
& H^{i-1}(K\overset{\LL}{\otimes}M_{2})\longrightarrow
H^{i}(K\overset{\LL}{\otimes}I)\longrightarrow
H^{i}(K\overset{\LL}{\otimes}(N\otimes k(\eta)))\\
& H^{i}(K\overset{\LL}{\otimes}M_{1})\longrightarrow
H^{i}(K\overset{\LL}{\otimes}N)\longrightarrow
H^{i}(K\overset{\LL}{\otimes}I){\quoi},
\end{align*}
dont\pageoriginale on tire que
$H^{i}(K\overset{\LL}{\otimes}N)=0$ pour $i<-n$,
\cqfd

\begin{enonce*}{Corollaire 5.5.24}\etiquette[5.5.24]{XVII.cor5.5.24}
Pour qu'une flèche $u:K\rightarrow L$ dans $D^{-}(A)$ soit un
isomorphisme, il suffit que pour tout $x\in \Spec(A)$, la flèche
$$
u_{x}:K\overset{\LL}{\otimes}_{A}k(x)\longrightarrow
L\overset{\LL}{\otimes}_{A}k(x)
$$
soit un isomorphisme dans $D^{-}(k(x))$.
\end{enonce*}

\sisi{\noindent}{}Soit $M$ le cône (mapping cylinder) de $u$. Par hypothèse, on a
$M\overset{\LL}{\otimes}_{A}k(x)=0$ pour tout
$x\in\Spec(A)$. D'après \hyperref[XVII.prop5.5.23]{5.5.23}, on a donc $M=0$, et $u$ est un
isomorphisme.

\setcounter{subsubsection}{25}
\setcounter{paragraph}{0}
\paragraph{}\etiquette[5.5.25.1]{XVII.par5.5.25.1}
Poursuivons la démonstration de \hyperref[XVII.lem5.5.22]{5.5.22}, interrompue
en \hyperref[XVII.prop5.5.23]{5.5.23}. D'après
la compatibilité \hyperref[XVII.par5.5.20.3]{5.5.20.3} ${\rm(iv'')}$ et
\hyperref[XVII.cor5.5.24]{5.5.24}, le morphisme $k^{n}$
relatif à $K\in \Ob D^{b,\tor\geq 0}(X,A)$ est un isomorphisme si et
seulement si les morphismes $k^{n}$ relatifs aux
$K\overset{\LL}{\otimes}_{A}k(x)\in \Ob
D^{b,\tor\geq 0}(X,k(x))$ sont des isomorphismes, pour $x\in
\Spec(A)$. Ceci nous ramène au cas où $A$ \emph{est un corps de
caractéristique} $\ell > 0$.

\paragraph{}\etiquette[5.5.25.2]{XVII.par5.5.25.2}
La suite exacte (où $C$ désigne un cône)
$$
0\longrightarrow K[-1]\longrightarrow CK[-1]\longrightarrow
K\longrightarrow 0
$$
montre que les flèches $k^{n}$ sont des isomorphismes pour $K$ si et
seulement si elles le sont pour $K[-1]$ (\hyperref[XVII.lem5.5.22.2]{5.5.22.2}).

\paragraph{}\etiquette[5.5.25.3]{XVII.par5.5.25.3}
Les\pageoriginale suites exactes
$$
0\longrightarrow \sigma_{\geq n} K\longrightarrow K\longrightarrow
\sigma_{\leq n} K\longrightarrow 0
$$
montrent que pour les flèches de Künneth soient un isomorphisme, pour
$K$ complexe plat, il suffit qu'elles le soient pour les composantes
de $K$, placées en degré $0$ (\hyperref[XVII.lem5.5.22.2]{5.5.22.2}\sisi{)}{} et \sisi{(}{}\hyperref[XVII.par5.5.25.2]{5.5.25.2}).

\paragraph{}\etiquette[5.5.25.4]{XVII.par5.5.25.4}
Soit donc $\sisi{F}{K}$ un faisceau de $A$-vectoriels. Ce faisceau est limite
inductive de faisceaux constructibles (\hyperrefext[IX][IX.cor2.7.2]{2.7.2}), automatiquement
plats puisque $A$ est un corps, et par passage à la limite, on se
ramène au cas $K$ \emph{constructible}, réduit au degré zéro.

\paragraph{}\etiquette[5.5.25.5]{XVII.par5.5.25.5}
Si un complexe $K$ est facteur direct d'un complexe $L$, le morphisme
$k^{n}$ relatif à $K$ est facteur direct du morphisme $k^{n}$ relatif
à $L$, donc un isomorphisme si ce dernier en est un. La méthode de la
trace (\hyperrefext[IX][IX.rem5.8]{5.8}) s'applique donc et il suffit de traiter le cas où $F$
est de la forme $j_{!}A_{X'}$, pour $j:X'\rightarrow X$ quasi-fini. La
compatibilité \hyperref[XVII.par5.5.20.2]{5.5.20.2} permet de remplacer $X$ par $X'$, d'où la
réduction
$$
\mbox{$X$ \emph{est de dimension $\leq 1$ et} $K=A_{X}${\quoi}.}
$$

\paragraph{}\etiquette[5.5.25.6]{XVII.par5.5.25.6}
On peut supposer $X$ réduit. Si $j:U\rightarrow X$ est l'ouvert de
lissité de $X$, de complément $i:F\rightarrow X$, la suite exacte
$$
0\longrightarrow j_{!}j^{\ast}A_{X}\longrightarrow
A_{X}\longrightarrow i_{!}i^{\ast}A_{X}\longrightarrow 0
$$
nous ramène par dévissage aux cas où $X$ \emph{est soit une courbe
lisse, soit un ensemble fini.} De plus, la formule
$$
A_{X}=A\overset{\LL}{\otimes}_{\ZZ/\ell}(\ZZ/\ell)_{X}
$$
ramène par \hyperref[XVII.par5.5.20.3]{5.5.20.3} ${\rm(iv'')}$ au cas où $A_{X}=\ZZ/\ell$.

Si\pageoriginale $X$ est une courbe, on plonge $X$ dans la courbe
projective et lisse correspondante, et un nouveau dévissage ramène à
l'un des deux cas suivants
\begin{enumerate}
\renewcommand{\theenumi}{\roman{enumi}}
\renewcommand{\labelenumi}{(\theenumi)}
\item $X$ est une courbe projective et lisse, $Y=S=\Spec(k)$ avec $k$
algébriquement clos, $\Aa=\ZZ/\ell$, $K=\ZZ/\ell$.

\item $X$ est une somme d'un nombre fini de copies de $\Spec(k)$,
$Y=S=\Spec(k)$ ($k$ algébriquement clos),
$\Aa=\ZZ/\ell$, $K=\ZZ/\ell$.
\end{enumerate}

Un ultime dévissage de $\sisi{}{(}\ZZ/\ell\sisi{}{)}_{X}$ en les faisceaux
$\sisi{}{(}\ZZ/\ell\sisi{}{)}_{X\sisi{}{'}}$, pour $X'$ composante connexe de $X$, ramène dans
(i) et (ii) à supposer $X$ connexe. Le cas (ii) est alors
trivial. Ceci prouve \hyperref[XVII.lem5.5.22]{5.5.22}.

\begin{enonce*}{Lemme 5.5.26}\etiquette[5.5.26]{XVII.lem5.5.26}
Le théorème \emph{\hyperref[XVII.thm5.5.21]{5.5.21}} est conséquence des deux cas particuliers
suivants de l'assertion $(\ast)$ de \emph{\hyperref[XVII.lem5.5.22]{5.5.22}}:
\begin{itemize}
\item[\emph{Cas} 1.] On suppose que $k=\CC$.

\item[\emph{Cas} 2.] On suppose que $k$ est de caractéristique $p>0$
et que $\ell=p$.
\end{itemize}
\end{enonce*}

Avec les notations de \hyperref[XVII.lem5.5.22]{5.5.22}, désignons par $p$ la caractéristique de
$k$. Si $p=0$, le principe de Lefschetz (s'appuyant ici sur la
compatibilité \hyperref[XVII.par5.5.20.3]{5.5.20.3} ${\rm(v'')}$, et sur XII 5.4) permet de se
ramener au cas où $k=\CC$.

Reste à traiter le cas où $p>0$ et où $\ell\neq p$. La courbe $X$ se
remonte en une courbe $X'$ sur l'anneau $W(k)$ des vecteurs de Witt
sur $k$ (\SGA 1 III 7.4). Posons $S=\Spec(W(k))$, soit le diagramme
$$
\xymatrix@C=2.25cm@R=1.5cm{
X'_{\eta} \ar[d]^{f_\eta}          \ar[r]              & X' \ar[d]^{f'}& X\ar[l]\ar[d]^{f}\\
\eta               \ar[r]         & S   & \Spec(k), \ar[l],
}
$$
et\pageoriginale soient $f_{n}$, $f'_{n}$ et $f_{\eta n}$ les
projections respectives de $\Sym^{n}_{k}(X)$, $\Sym^{n}_{S}(X')$ et
$\Sym^{n}_{\eta}(X'_{\eta})$ sur $\Spec(k)$, $S$ et $\eta$. On sait
que $f'_{n}$ est un morphisme projectif et lisse (la projectivité est
évidente; pour la lissité, le plus simple est de se ramener au cas où
$X'=\PP^{1}_{S}$, auquel cas $\Sym^{n}_{S}(X')\simeq
\PP^{n}$ comme il résulte du théorème des fonctions symétriques
élémentaires). D'après le théorème de spécialisation \hyperrefext[XVI][XVI.cor2.2]{2.2}, les
faisceaux $R^{i}f'_{n{\ast}}$ $(\ZZ/\ell)$ sont localement
constants sur $S$. Il en est de même des faisceaux
$L^{i}\Gamma^{n}(Rf'_{\ast}\ZZ/\ell)$. Pour vérifier que la
flèche de \eqhyperref[XVII.eq5.5.22.1]{5.5.22.1} relative à $X$ sur $k$ est un isomorphisme, il
suffit donc d'après \hyperref[XVII.par5.5.20.3]{5.5.20.3} ${\rm(v'')}$ de prouver l'énoncé analogue
pour la fibre générique géométrique $X_{\bar{\eta}}$ de $X$ sur
$S$. Le corps $k(\bar{\eta})$ est de caractéristique $0$, et, vu la
première réduction, ceci prouve \hyperref[XVII.lem5.5.26]{5.5.26}.

Pour prouver \hyperref[XVII.thm5.5.21]{5.5.21} dans les deux cas particuliers
\hyperref[XVII.lem5.5.26]{5.5.26}, on
comparera le morphisme de Künneth symétrique d'une part à son analogue
transcendant (\hyperref[XVII.subsubsec5.5.27]{5.5.27} et
\hyperref[XVII.subsubsec5.5.28]{5.5.28}) d'autre part à son analogue cohérent
(\hyperref[XVII.subsubsec5.5.29]{5.5.29} à \hyperref[XVII.lem5.5.37]{5.5.37}).

\setcounter{subsubsection}{26}
\subsubsection{}\etiquette[5.5.27]{XVII.subsubsec5.5.27}
Si $p:X\rightarrow S$ est une application continue entre espaces
topologiques, avec $S$ annelé par un faisceau d'anneaux commutatifs
$\Aa$, les constructions \hyperref[XVII.subsubsec5.5.7]{5.5.7},
\hyperref[XVII.subsubsec5.5.8]{5.5.8} s'appliquent telles
quelles pour définir un foncteur $\Gamma^{n}_{\ext}$ des modules sur
$X$ dans les modules sur $\Sym^{n}_{S}(X)=X^{n}/\Sf_{n}$. Les
sorites \hyperref[XVII.subsubsec5.5.7]{5.5.7} à \hyperref[XVII.prop5.5.18]{5.5.18} se
transposent. De même, si $S$, $X$ et $Y$
sont des espaces topologiques localement compacts séparés
$$
\xymatrix@C=1.5cm@R=1.25cm{
X \ar[rr]^{u}\ar[rd]&  &  Y\ar[ld]\\
            & S &,
}
$$
et\pageoriginale si les fibres de $u$ sont de dimension cohomologique
bornée, le sorite \hyperref[XVII.subsubsec5.5.19]{5.5.19},
\hyperref[XVII.subsubsec5.5.20]{5.5.20} du morphisme de Künneth symétrique en
cohomologie à support propre se transpose.

Si $p:X\rightarrow S$ est un morphisme de schémas de type fini sur
$\CC$, les foncteurs $\Gamma^{n}$ et $L\Gamma^{n}$ commutent aux
foncteurs « passage au transcendant ». Enfin, si $X$, $Y$, $S$,
$\Aa$, $K$ sont comme en \hyperref[XVII.thm5.5.21]{5.5.21}, avec $S$ séparé de type fini
sur $\CC$, et si on désigne génériquement par $\epsilon$ les
morphismes de sites de type $\epsilon:X_{\cl}\rightarrow
X_{\et}$, le diagramme suivant est commutatif:
$$
\vcenter{\etiquette[5.5.27.1]{XVII.eq5.5.27.1}
\xymatrix@C=1.75cm@R=1.5cm{
\epsilon^{\ast} L\Gamma^{n}_{\ext} Ru_{!}K  \ar[r]^{k^{\sisi{k}{n}}} \ar[d]&
\epsilon^{\ast} R \Sym^{n}_{S}(u)_{!} L \Gamma^{n}_{\ext} K \ar[d]\\
L\Gamma^{n}_{\ext} R u^{\an}_{!} \epsilon^{\ast}K  \ar[r]^{k^{n \an}}
& R \Sym^{n}_{S}(u^{\an})_{!} L\Gamma^{n}_{\ext} \epsilon^{\ast}K.
}}\leqno{(5.5.27.1)}
$$

\subsubsection{}\etiquette[5.5.28]{XVII.subsubsec5.5.28}
D'après la compatibilité \eqhyperref[XVII.eq5.5.27.1]{5.5.27.1}, et un cas
élémentaire du théorème
de comparaison (\hyperrefext[XI][XI.thm4.4]{4.4} ou \hyperrefext[XVI][XVI.thm4.1]{4.1}),
pour traiter le cas 1 de \hyperref[XVII.lem5.5.26]{5.5.26},
il suffit de vérifier que si $X$ est une courbe projective non
singulière complexe, alors l'analogue topologique de \eqhyperref[XVII.eq5.5.22.1]{5.5.22.1} est un
isomorphisme.

L'espace topologique $X^{\an}$ est triangulable (ce fait est ici
trivial, car $X$ est un revêtement ramifié de la sphère de
Riemann). L'analogue topologique du lemme de dévissage \hyperref[XVII.lem5.5.22.2]{5.5.22.2} nous
ramène alors à prouver
\begin{itemize}
\item[$(\ast)$] Si\pageoriginale $D$ est un simplexe
fermé (de dimension 0, 1 ou 2), la flèche de Künneth symétrique
topologique
$$\etiquette[5.5.28.1]{XVII.eq5.5.28.1}
k^{n}:L\Gamma^{n} R\Gamma(D,\ZZ/\ell)\longrightarrow R\Gamma(\Sym^{n}(D\sisi{}{)},\ZZ/\ell)\leqno(5.5.28.1)
$$
est un isomorphisme.
\end{itemize}

On a
$$
R\Gamma(D,\ZZ/\ell)=\ZZ/\ell\;\mbox{\rm placé en degré }
0{\quoi},
$$
et donc
$$
L\Gamma^{n}R\Gamma(D,\ZZ/\ell)=\ZZ/\ell\; \mbox{\rm
placé en degré $0$}{\quoi}.
$$

D'autre part, $\Sym^{n}(D)$ est contractible, et on a encore
$$
R\Gamma(\Sym^{n}(D), \ZZ/\ell)=\ZZ/\ell\;\mbox{\rm placé
en degré $0$}{\quoi}.
$$

Les $H^{i}$ des deux membres de \eqhyperref[XVII.eq5.5.28.1]{5.5.28.1} sont donc nuls, et la flèche
induite sur le $H^0$ s'identifie à l'identité de
$\ZZ/\ell$.

\subsubsection{}\etiquette[5.5.29]{XVII.subsubsec5.5.29}
Pour l'application que nous avons en vue, il y a intérêt à regarder un
faisceau quasi-cohérent sur un schéma $X$ comme étant un faisceau
quasi-cohérent de $\Oo$-modules sur le site étale de $X$.

Soit $p:X\rightarrow \sisi{X}{S}$ un morphisme quasi-projectif de schémas (cf
\hyperref[XVII.rem5.5.21.1]{5.5.21.1} \footnote[1]{Tout ce qui suit, y inclus
  \hyperref[XVII.prop5.5.34]{5.5.34}, est valable
en supposant seulement que $X$ et $Y$ admettent des faisceaux
inversibles amples relativement à $S$ (sans supposer $X$ et $Y$ de
type fini sur $S$).}). Si $ℱ$ est un faisceau de $\Oo_{X}$-modules
sur $X_{\text{\rm et}}$, alors
\begin{enumerate}
\renewcommand{\theenumi}{\alph{enumi}}
\renewcommand{\labelenumi}{\theenumi)}
\item On désigne par $\bigotimes\limits^{n}ℱ$ le
faisceau de modules sur $(X/S)^{n}_{\text{\rm et}}$ produit tensoriel
des $\Oo_{(X/S)^n}$-modules images réciproques de $ℱ$ sur
$(X/S)^{n}_{\et}$;

\item le faisceau $\bigotimes\limits^{n}ℱ$ est un
faisceau $\Sf_{n}$-équivariant. Par passage aux invariants, il
définit un faisceau $TS^{n}_{\text{\rm ext}}(ℱ)$ sur $\Sym^{n}_{S}(X)$:
si $\sigma$ est la projection de $(X/S)^{n}$ sur $\Sym^{n}_{S}(X)$, on a
$$\etiquette[5.5.29.1]{XVII.eq5.5.29.1}
TS^{n}_{\ext}(ℱ)=(\sigma_{\ast}\bigotimes^{n}ℱ)^{\Sf^{n}}{\quoi}.\leqno(5.5.29.1)
$$
\end{enumerate}

\emph{Supposons\pageoriginale que $p$ soit plat}; comme en
\hyperref[XVII.subsubsec5.5.8]{5.5.8}, on définit alors le foncteur
$\Gamma^{n}_{\text{\rm ext}}$ comme étant le $0$\sisi{.ème}{\textsuperscript{ième}} dérivé gauche du
foncteur $TS^{n}_{\ext}$. Avec les notations de \hyperref[XVII.par5.5.2.5]{5.5.2.5} et
\hyperref[XVII.subsubsec5.5.8]{5.5.8}, si
$$
P_{1}\longrightarrow P_{0}\longrightarrow ℱ\longrightarrow 0
$$
est une présentation plate de $ℱ$, alors la suite
$$
\xymatrix@C=1.75cm@M=.25cm{
TS^{n}_{\ext} (P'_1) \ar@<3pt>[r]^{j_0} \ar@<-3pt>[r]_{j_1}& TS^{n}_{\ext}(P_0)\ar[r] & \Gamma^{n}_{\ext}(\sisi{F}{ℱ})
}
$$
est exacte. Si $x$ est un point géométrique de $\Sym^{n}_{S}(X)$ comme
en \hyperref[XVII.par5.5.7.1]{5.5.7.1}, on a (\cf \eqhyperref[XVII.eq5.5.8.1]{5.5.8.1})
\begin{align*}
& \Gamma^{n}_{\text{\rm
ext}}(ℱ)_{x}=\Oo_{x}\otimes_{C}(\bigotimes^{\ell}_{1}
\Gamma_{\Oo_{S}}^{n_{i}}(\sisi{F}{ℱ}_{y_{i}})){\quoi},\tag{5.5.29.2}\\
& C=\bigotimes^{\ell}_{1} \Sym^{n_{i}}_{\Oo_{S}}(\Oo_{y_{i}}){\quoi}.\tag*{avec}
\end{align*}\etiquette[5.5.29.2]{XVII.eq5.5.29.2}%L'enlèvement de ce commentaire introduira un espace blanc non désiré.
Il suffit en effet de vérifier cette formule dans le cas trivial où
les fibres de $\sisi{F}{ℱ}$ sont libres. On notera que $\Oo_{x}$ est plat sur
$C$, étant l'hensélisé strict de $C$ au point fermé image de $x$.

\subsubsection{}\etiquette[5.5.30]{XVII.subsubsec5.5.30}
Supposons $X$ et $S$ affines, $p$ plat et $ℱ$ quasi-cohérent. Soient
$S=\Spec(A)$, $X=\Spec(B)$ et soit $F=H^0(X,ℱ)$, de sorte que
$ℱ$ est le faisceau $F^{\sim}$ défini par le $B$-module $F$. On lit
alors sur \eqhyperref[XVII.eq5.5.29.1]{5.5.29.1} que
$$
TS^{n}_{\ext}(ℱ)=TS^{n}_{A}(F)^{\sim}{\quoi},
$$
et, par dérivation, on a donc
$$\etiquette[5.5.30.1]{XVII.eq5.5.30.1}
\Gamma^{n}_{\text{\rm
ext}}(ℱ)=\Gamma^{n}_{A}(F)^{\sim}{\quoi}.\leqno(5.5.30.1)
$$
Le\pageoriginale foncteur $\Gamma^{n}_{\text{\rm ext}}$ apparaît donc
comme un foncteur $\Gamma^{n}$ « relatif à $A$ », ce que
\eqhyperref[XVII.eq5.5.29.2]{5.5.29.2} exprimait déjà sous forme locale.

\subsubsection{}\etiquette[5.5.31]{XVII.subsubsec5.5.31}
On désignera par $D^{b,\tor_{S}\leq 0}(X,\Oo_{X})$ la sous-catégorie
de $D^{b}(X,\Oo_{X})$ formé des complexes de $\Oo_{X}$-modules
qui, en tant que complexes de $p^{-1}(\Oo_{S})$-modules sont de
\tTor-dimension $\leq 0$.

Soit la condition

\paragraph{}\etiquette[5.5.31.1]{XVII.par5.5.31.1}
$K\in \Ob C^{b,\geq 0}(X,\sisi{O}{\Oo}_{X})$ et, en chaque point géométrique $x$
de $X$, $K_{x}$ est homotope à un complexe de $\Oo_{X,x}$-modules, à
degrés positifs, plat sur $\Oo_{S,p(x)}$.

Le morphisme $p$ étant plat, de sorte qu'il y a « assez »  de
$\Oo_{X}$-modules plats sur $S$, on vérifie comme en \hyperref[XVII.lem5.5.13.2]{5.5.13.2} que
$D^{b,\tor_{S}\leq 0}(X,\Oo_{X})$ se déduit de la catégorie des
complexes $K\in \Ob K^{b,\geq 0}(X,\Oo_{X})$ vérifiant (\hyperref[XVII.par5.5.31.1]{5.5.31.1}) en
inversant les quasi-isomorphismes.

On lit sur \eqhyperref[XVII.eq5.5.29.2]{5.5.29.2} et \hyperref[XVII.lem5.5.6.1]{5.5.6.1} que si $u:K\rightarrow L$ est un
quasi-isomorphisme entre complexes vérifiant (\hyperref[XVII.par5.5.31.1]{5.5.31.1}), alors
$\Gamma^{n}_{\text{\rm ext}}(u)$ (au sens de \textsc{Dold-Puppe}) est encore un
quasi-isomorphisme. Le foncteur $\Gamma^{n}_{\text{\rm ext}}$ se
dérive dès lors en
$$\etiquette[5.5.31.2]{XVII.eq5.5.31.2}
L\Gamma^{n}_{\text{\rm ext}}:D^{b,\tor_{S}\leq
0}(X,\Oo_{X})\longrightarrow D^{b,\tor_{S}\leq 0}(\Sym^{n}_{S}(X), \Oo_{\Sym^{n}_{S}(X)})\leqno(5.5.31.2)
$$

\subsubsection{}\etiquette[5.5.32]{XVII.subsubsec5.5.32}
On désignera par $D^{b,\tor_{S}\leq 0}_{\text{\rm qcoh}}(X,\Oo_{X})$
la sous-catégorie de $D^{b,\tor_{S}\leq 0}(X,\Oo_{X})$ formée des
complexes à cohomologie quasi-cohérente.

Des\pageoriginale arguments standards permettent, comme en
\hyperref[XVII.subsubsec5.5.17]{5.5.17}, de
définir une flèche de Künneth symétrique du type suivant:

Soit $u:X\rightarrow Y$ un morphisme de $S$-schémas quasi-projectifs
et plats sur $S$; pour $K\in \Ob D^{b,\tor_{S}\leq 0}_{\text{\rm
qcoh}}(X,\Oo_{X})$, la flèche de Künneth symétrique $k^{n}$ dans
$D^{b,\tor_{S}\leq 0}_{\text{\rm qcoh}}(\Sym^{n}_{S}(Y),\Oo)$ est
une flèche
$$\etiquette[5.5.32.1]{XVII.eq5.5.32.1}
k^{n}:L\Gamma^{n}_{\text{\rm ext}}Ru_{\ast}K\longrightarrow
R\Sym^{n}_{S}u_{\ast}L\Gamma^{n}_{\text{\rm ext}} K{\quoi}.\leqno(5.5.32.1)
$$

On laisse au lecteur le soin de vérifier les propriétés formelles de
\eqhyperref[XVII.eq5.5.32.1]{5.5.32.1} qui seront utilisées, notamment la suivante:

\begin{enonce*}{Lemme 5.5.33}\etiquette[5.5.33]{XVII.lem5.5.33}
Si $S$ est de caractéristique $p>0$, et si $u:X\rightarrow S$ est
propre et plat, le diagramme de flèches de Künneth symétriques
\eqhyperref[XVII.eq5.5.17.2]{5.5.17.2} et \eqhyperref[XVII.eq5.5.32.1]{5.5.32.1} est commutatif ($Y=S, K=\sisi{\sheafZZ}{\ZZ}/p$ ou
$\Oo_{X}$):
\[
\xymatrix{
L\Gamma^{n}_{\sisi{\sheafZZ}{\ZZ}/p}Ru_{\ast}(\sisi{\sheafZZ}{\ZZ}/p)
\ar[r]^{k^n} \ar[d]&R
\Sym^{n}_S(u)_{\ast} (\sisi{\sheafZZ}{\ZZ}/p) \ar[d]&\\
L\Gamma^{n}_{\Oo} Ru_{\ast}(\Oo) \ar[r]^{k^n}& R
\Sym^{n}_{S}(u)_{\ast}\Oo &.
}
\]
\end{enonce*}


\begin{enonce*}{Proposition 5.5.34}\etiquette[5.5.34]{XVII.prop5.5.34}
Sous les hypothèses de \hyperref[XVII.subsubsec5.5.32]{5.5.32}, la flèche de Künneth symétrique
\eqhyperref[XVII.eq5.5.32.1]{5.5.32.1} est un isomorphisme.
\end{enonce*}

\begin{enumerate}
\renewcommand{\theenumi}{\alph{enumi}}
\renewcommand{\labelenumi}{\theenumi)}
\item La question est locale sur $S$, qu'on peut supposer affine.

\item Pour tout point $x$ de $\Sym^{n}_{S}(Y)$, il existe un ouvert
affine $Y_{1}$ de $Y$ tel que $x$ soit dans $\Sym^{n}_{S}(Y_{1})$. Par
localisation, on peut donc supposer $Y$ affine.

\item Tout\pageoriginale complexe $K$ du type considéré peut se
représenter par un complexe $K_{0}$ de modules quasi-cohérents sur
$S$, nul en degrés $i<0$. Si $\mathcal{U}$ est un recouvrement affine
de $X$, alors $K$ admet une résolution finie par des sommes de
complexes de la forme $j_{\ast}j^{\ast}K$ pour $j$ inclusion d'un
ouvert affine (complexe de $\check{C}$ech alterné). Utilisant
l'analogue quasi-cohérent du lemme de dévissage \hyperref[XVII.lem5.5.22.2]{5.5.22.2}, ceci nous
ramène au cas où il existe un ouvert affine $j:X_{1}\hookrightarrow X$
de $X$ et $K_{1}\in D^{b,\tor_{S}\leq 0}_{\text{\rm
qcoh}}(X_{1},\Oo)$ tel que $K=Rj_{\ast}K_{1}$.

\item Vu l'analogue quasi-cohérent de la transitivité \hyperref[XVII.par5.5.20.2]{5.5.20.2}, il
suffit de prouver \hyperref[XVII.prop5.5.34]{5.5.34} pour $uj$ et $j$. La réduction b), appliquée
à $j$, nous ramène enfin au cas où $X$, $Y$ et $S$ sont affines. On
peut alors représenter $K$ par un complexe $K_{2}\in C^{b,\geq
0}(X,\Oo_{X})$ à composantes \sisi{}{$S$-}plates et quasi-cohérentes (\sisi{\quad}{\hyperref[XVII.def4.1.9]{4.1.9}}). Les
foncteurs $Ru_{\ast}$ et $L\Gamma^{n}$ se calculent alors sans plus
devoir résoudre $K_{2}$, et il reste à noter que, d'après \eqhyperref[XVII.eq5.5.30.1]{5.5.30.1},
pour $ℱ$ quasi-cohérent sur $X$, on a
\[
\Gamma^{n}_{\text{\rm ext}} u_{\ast}ℱ\isomto
\Sym^{n}_{S}(u)_{\ast} \Gamma^{n}_{\text{\rm ext}}ℱ{\quoi},
\]
les deux membres s'identifiant à
\[
\Gamma^{n}_{A}(F)^{\sim} \quad\mbox{\rm pour}\quad
ℱ=F^{\sim}{\quoi},
\]
\cqfd
\end{enumerate}

\setcounter{subsubsection}{34}
\subsubsection{}\etiquette[5.5.35]{XVII.subsubsec5.5.35}
Traitons le cas $2$ de \hyperref[XVII.lem5.5.26]{5.5.26}.
Avec les notations de \hyperref[XVII.lem5.5.22]{5.5.22}, le
diagramme suivant est commutatif (\hyperref[XVII.lem5.5.33]{5.5.33}):
$$\vcenter{\etiquette[5.5.35.1]{XVII.eq5.5.35.1}
\xymatrix@C=2cm{
  L\Gamma^{n} (R\Gamma(X,\ZZ/p)) \ar[r]^{k^n}_*+[o][F-]{\scriptstyle{1}}
  \ar[d]& R\Gamma(\Sym^{n}_{k}(X),\ZZ/p)\ar[d] \\
L\Gamma^{n}(R\Gamma(X,\Oo)) \ar[r]^{k^n}_*+[o][F-]{\scriptstyle{2}} & R\Gamma(\Sym^{n}_{k}(X),\Oo),
}}\leqno{(5.5.35.1)}
$$
et\pageoriginale d'après \hyperref[XVII.prop5.5.34]{5.5.34}, la flèche \ding{193} est un
isomorphisme. L'endomorphisme de Frobenius $F:x\rightarrow x^{p}$ de
$\Oo$ agit $p$-linéairement sur $R\Gamma(X,\Oo)$ et sur
$R\Gamma(\Sym^{n}_{k}(X),\Oo)$. Il agit donc aussi $p$-linéairement
sur $L\Gamma^{n}(R\Gamma(X,\Oo))$. On vérifie que $k^{n}$ et $F$
commutent.

La théorie d'Artin-Schreier nous fournit des suites exactes
$$
0\longrightarrow H^{i}(X,\ZZ/p)\longrightarrow
H^{i}(X,\Oo)\xrightarrow{F-I} H^{i}(X,\Oo)\longrightarrow
\sisi{\Oo}{0}{\quoi},
$$
et de même pour $\Sym^{n}_{S}(X)$. Pour déduire que \ding{192} est
un isomorphisme du fait que \ding{193} en soit un, il nous reste à
vérifier le

\begin{enonce*}{Lemme 5.5.36}\etiquette[5.5.36]{XVII.lem5.5.36}
Soient $k$ un corps parfait de caractéristique $p>0$, $K\in C^{\geq
0}(\ZZ/p)$ un complexe de $\ZZ/p$-vectoriels de
dimension finie, $L\in C^{\geq 0}(k)$ un complexe de $k$-vectoriels de
dimension finie, $F$ un endomorphisme $p$-linéaire de $L$ et
$u:K\rightarrow L$ un morphisme. On suppose que
$H^{\ast}(Fu)=H^{\ast}(u)$, et que les suites
$$
0\longrightarrow H^{i}(K)\longrightarrow
H^{i}(L)\xrightarrow{F-I}H^{i}(L)\longrightarrow 0
$$
sont exactes. Alors,\pageoriginale les suites
$$
0\longrightarrow H^{i}(\Gamma^{n}K)\longrightarrow
H^{i}(\Gamma^{n}L)\xrightarrow{F-I}H^{i}(\Gamma^{n}L)\longrightarrow 0
$$
sont exactes.
\end{enonce*}

Puisque $H^{\ast}(Fu)=H^{\ast}(u)$ et que $\ZZ/p$ est un corps,
il existe une homotopie
$$
Fu-u=dH+Hd{\quoi}.
$$
Le corps $k$ étant parfait, on peut écrire
$$
FH_{0}-H_{0}=H{\quoi}.
$$
Remplaçant $u$ par $u-(dH_{0}+H_{0}d)$, on peut supposer que
$Fu=u$.

Le foncteur $S$ qui à chaque $k$-vectoriel $V$ muni d'un endomorphisme
$p$-linéaire $F$ associe le $\ZZ/p$-vectoriel $V^{S}$ des $v\in
V$ tels que $Fv=v$ est un foncteur exact. L'hypothèse signifie donc
que le morphisme
$$
K\longrightarrow L^{S}
$$
est un quasi-isomorphisme. Reste à prouver que le foncteur $S$ commute
au foncteur $\Gamma^{n}$. Ce point résulte aisément du lemme bien
connu suivant, dont la démonstration est laissée au lecteur:

\begin{enonce*}{Lemme 5.5.37}\etiquette[5.5.37]{XVII.lem5.5.37}
Un vectoriel de dimension finie $V$ sur un corps parfait $k$, muni
d'un endomorphisme $p$-linéaire $F$, admet une et une seule
décomposition
$$
V=V_{0}+V_{n}
$$
avec $V^{S}\otimes k\isomto  V_{0}$, $F(V_{n})\subset V_{n}$
et $F|V_{n}$ nilpotent.
\end{enonce*}

Ceci achève la démonstration de \hyperref[XVII.thm5.5.21]{5.5.21}.

\section[Le foncteur $f_{!}$]{Le foncteur $f_{!}$}\etiquette[6]{XVII.sec6}\pageoriginale

On se propose de donner une construction directe du foncteur $f_{!}$
de \hyperref[XVII.var5.1.11]{5.1.11}. On utilisera ensuite cette construction pour définir le « morphisme trace »  (\hyperref[XVII.thm6.2.3]{6.2.3}) pour les morphismes plats quasi-finis de
présentation finie.

\subsection[Nouvelle construction du foncteur $f_{!}$]{Nouvelle construction du foncteur $f_{!}$}\etiquette[6.1]{XVII.subsec6.1}

Soit $f:X\rightarrow S$ un morphisme localement de type fini. Une
partie $P$ de $X$ est dite \emph{propre} sur $S$ si elle est l'image
d'un sous-schéma de $X$ propre sur $S$.

\begin{enonce*}{Proposition 6.1.1}\etiquette[6.1.1]{XVII.prop6.1.1}
Soit $f:X\rightarrow S$ un morphisme localement de type fini et
quasi-séparé.
\begin{enumerate}
\renewcommand{\theenumi}{\roman{enumi}}
\renewcommand{\labelenumi}{(\theenumi)}
\item Une partie de $X$ propre sur $S$ est localement fermée, et est
fermée si $f$ est séparé.

\item Si $f$ est séparé, la réunion de deux parties de $X$ propres sur
$S$ est propre sur $S$.

\item Pour tout $S$-schéma $U$, désignons par $\Phi_{X}(U)$ l'ensemble
des parties de $\fprod{U}{X}{S}$ propres sur $U$. Le foncteur
$\Phi_{X}$ est un faisceau pour la topologie \emph{fpqc}. Si $f$ est
séparé, $\Phi\sisi{}{_X}$ est un faisceau d'ensembles ordonnés filtrants pour la
relation d'inclusion.
\end{enumerate}
\end{enonce*}

Les\pageoriginale assertions (i) et (ii) sont triviales, et la seconde
assertion de (iii) résulte de (ii). Le foncteur $U\mapsto
\mathcal{P}(\fprod{U}{X}{S})$ est un faisceau pour toute topologie
pour laquelle les recouvrements soient surjectifs. Il résulte de \EGA
IV 2.3.14 et \EGA IV 2.7.1 (vii) que $\Phi_{X}$ en est un
sous-faisceau.

Si $\Ff $ est un faisceau étale d'ensembles pointés sur un
schéma $X$, rappelons (IV 8.5.2) que le support d'une section $s$ de
$\Ff $ est le complémentaire du plus grand ouvert où $s=0$,
$0$ désignant la section marquée.

\begin{enonce*}{Définition 6.1.2}\etiquette[6.1.2]{XVII.def6.1.2}
Soient $f:X\rightarrow S$ un morphisme localement de type fini et
séparé, et $\Ff $ un faisceau d'ensembles pointés sur $X$. On
désigne par $f_{!}\Ff $ le sous-faisceau de
$f_{\ast}\Ff $ dont les sections sur $U$ (étale sur $S$) sont
les sections de $\Ff $ sur $\fprod{U}{X}{S}$ à support propre
sur $U$.
\end{enonce*}

Il résulte de \hyperref[XVII.prop6.1.1]{6.1.1} que $f_{!}\Ff $ est effectivement un
faisceau.

On vérifie aisément que:

\begin{enonce*}{Lemme 6.1.3}\etiquette[6.1.3]{XVII.lem6.1.3}
Si $f$ est le composé $gh$ de deux morphismes séparés localement de
type fini, il existe un et un seul isomorphisme de foncteur
$f_{!}\simeq g_{!}h_{!}$ rendant commutatif le diagramme
$$
\xymatrix@C=1cm@R=1.25cm@M=.2cm{
f_{!}  \ar@{-}[r]^{\sim} \ar@{^{(}->}[d]& g_{!}h_{!} \ar@{^{(}->}[d]\\
f_{\ast} \ar@{=}[r]      & g_{\ast} h_{\ast}.
}
$$
\end{enonce*}

Si\pageoriginale $f$ est une immersion ouverte (\resp un morphisme
propre) le foncteur $f_{!}$ est le foncteur de prolongement par zéro
(\resp le foncteur $f_{\ast}$). D'après \hyperref[XVII.lem6.1.3]{6.1.3}, les
définitions \hyperref[XVII.def6.1.2]{6.1.2} et \hyperref[XVII.var5.1.11]{5.1.11} coïncident donc dans leur
domaine commun de validité.

Les propriétés du foncteur $f_{!}$ vérifiées dans le \S 5 restent
valable dans le cas plus général \hyperref[XVII.def6.1.2]{6.1.2}.

\begin{enonce*}{Proposition 6.1.4}\etiquette[6.1.4]{XVII.prop6.1.4}
Soit $f:X\rightarrow S$ un morphisme séparé localement de type fini.
\begin{enumerate}
\renewcommand{\theenumi}{\roman{enumi}}
\renewcommand{\labelenumi}{(\theenumi)}
\item Le foncteur $f_{!}$ est exact à gauche (\ie commute aux limites
projectives finies) et commute aux limites inductives filtrantes.

\item Si $f$ est quasi-fini, le foncteur $f_{!}$ est un foncteur exact
de la catégorie des faisceaux abéliens sur $X$ dans celle des
faisceaux abéliens sur $S$.

\item Pour tout diagramme cartésien
$$
\xymatrix@C=1.5cm@R=1.25cm{
X'          \ar[r]^{g'} \ar[d]^{f'}& X \ar[d]^{f}\\
S'          \ar[r]^{g}  & S
}
$$
il existe un isomorphisme de changement de base rendant commutatif le
diagramme
$$
\xymatrix@C=1.5cm@R=1.25cm{
g^{\ast}f_{!}  \sisi{\ar@{=}[r]}{\ar@{-}[r]^\sim} \ar@{^{(}->}[d]& f'_{!}g'^{\ast}\ar@{^{(}->}[d]\\
g^{\ast}f_{\ast}  \ar[r]    & f'_{\ast}g'^{\ast},
}
$$
dans lequel la seconde flèche horizontale est la flèche de changement
de base \emph{\neweqhyperrefext[XII][XII.eq4.1]{4.1}}.
\end{enumerate}
\end{enonce*}

\begin{enonce*}[remark]{Preuve}
Quel\pageoriginale que soit $U\in \Ob S_{\et}$, on a
\[
f_{!}\Ff (U)=\varinjlim_{P\in
\Phi(U)}\Gamma_{\sisi{p}{P}}(\Ff ){\quoi},
\]
et pour $U$ quasi-compact quasi-séparé, le foncteur
$\Ff \mapsto f_{!}\Ff (U)$ commute aux limites
projectives finies et aux limites inductives filtrantes, comme limite
inductive filtrante (\hyperref[XVII.prop6.1.1]{6.1.1}\sisi{)}{} (iii)\sisi{}{)} de foncteurs ayant ces propriétés.

Ceci prouve (i). L'assertion (ii) résulte aussitôt de (iii) (prendre
comme changements de bases les points géométriques de $S$).

L'assertion (iii) est de nature locale sur $S$, qu'on peut supposer
affine. Soit $I$ l'ensemble ordonné filtrant des ouverts de $X$ de
type fini sur $S$. Si
$U\overset{j}{\hookrightarrow}X$
est un ouvert de $X$,
$(fj)_{!}(\Ff |U)=f_{!}(j_{!}(\Ff |U))$ s'envoie dans
$f_{!}(\Ff\sisi{|U}{})$. De plus
$$
f_{!}(\Ff )=\varinjlim_{U\in I}
(fj\sisi{}{)}_{!}(\Ff |U){\quoi}.
$$
Ceci permet de supposer $X$ de type fini sur $S$ et d'appliquer le
lemme de Chow (\EGA II 2\sisi{è}{\textsuperscript{ième}} éd.)\nde{voir [C, 2.6], cité dans \nderef{XVII.ndeConrad}.} pour trouver un morphisme projectif et
surjectif: $p^{0}:X^{0}\rightarrow X$ tel que $fp\sisi{_{0}}{^0}$
soit compactifiable. Soit $X^{1}=\fprod{X^{0}}{X^{0}}{X}$ et $p^{1}$
la flèche naturelle de $X^{1}$ dans $X$. La suite
$$
0\rightarrow \Ff \rightarrow
p^{0}_{\ast}p^{0{\ast}}\Ff \rightrightarrows
p^{1}_{\ast}p^{1{\ast}}\Ff
$$
est exacte (\sisi{VII 9.2}{\hyperrefext[VIII][VIII.cor9.2]{9.2}}).

Compte tenu de la fonctorialité de la flèche de changement de base, et
de (i), il suffit de vérifier (iii) pour les faisceaux
$p^{0}_{\ast}p^{0^{\ast}}\Ff $ et
$p^{1}_{\ast}p^{1^{\ast}}\Ff $. Par le changement de notations
$(p^{1}\mapsto p^{0})$ on\pageoriginale se ramène donc à supposer
$\Ff $ de la forme $p^{0}_{\ast}\mathcal{G}$. L'assertion
(iii) pour le couple $(f,p^{0}_{\ast}\mathcal{G})$ équivaut à
l'assertion (iii) pour les couples $(p^{0}_{\ast}\sisi{}{,}\mathcal{G})$ et
$(fp^{0}_{\ast},\mathcal{G})$. On conclut par \hyperref[XVII.var5.2.7]{5.2.7} (i), applicable
car $p^{0}$ et $fp^{0}$ sont compactifiables.
\end{enonce*}

\begin{enonce*}[remark]{Contre-exemple 6.1.6}\etiquette[6.1.6]{XVII.contrexem6.1.6}
Le foncteur $Rf_{!}$ ne coïncide en général pas avec le foncteur
dérivé du foncteur $f_{!}$. Si $X$ est affine sur un corps
algébriquement clos $k$, on a en effet
$$
f_{!}\Ff =\sisi{\sum}{\bigoplus}_{x\in X(k)}H^{0}_{x}(\Ff ){\quoi},
$$
et le $i$\textsuperscript{ème} foncteur dérivé de $f_{!}$ vaut donc en
$\Ff $
$$
\sisi{\sum}{\bigoplus}_{x\in X(k)}H^{i}_{x}(\Ff ){\quoi}.
$$
Par contre, si $X$ est lisse et irréductible, et de dimension 1,
$$
H^{2}_{c}(X,\mu_{\ell^{n}})\simeq \ZZ/\ell^{n}
$$
pour $\ell$ premier à la caractéristique de $k$.
\end{enonce*}

\subsection{Le morphisme trace en dimension relative 0}\etiquette[6.2]{XVII.subsec6.2}

Le résultat principal de ce ${\rm n^{\circ}}$ est le théorème
\hyperref[XVII.thm6.2.3]{6.2.3}. Il est recommandé au lecteur de ne pas lire les préliminaires
techniques \hyperref[XVII.subsubsec6.2.1]{6.2.1},
\hyperref[XVII.prop6.2.2]{6.2.2}. Ceux-ci permettent, dans la démonstration de
\hyperref[XVII.thm6.2.3]{6.2.3}, de ne considérer que des morphismes $f$ finis, réduction qui,
via \EGA IV 18.12\sisi{.}{,} est de toute façon assez évidente.

\subsubsection{}\etiquette[6.2.1]{XVII.subsubsec6.2.1}
Soit\pageoriginale $f:X\rightarrow S$ un morphisme \emph{quasi-fini
séparé.} Quel que soit $U\in \Ob S_{\et}$, désignons par
$\Psi_{f}(U)$ la catégorie des couples formés d'un ensemble pointé
fini $I$, de point marqué noté $0$, et d'une décomposition de
$f^{\ast}U$ en parties ouvertes et fermées $f^{\ast}U=\coprod_{i\in
I}V_{i}$, avec $V_{i}$ fini sur $U$ pour $i\neq 0$. Un morphisme de
$\varphi=(V_{i})_{i\in I}$ dans $\varphi'=(V'_{i})_{i\in I'}$ est un
morphisme $\sigma$ de $I'$ dans $I$ tel que pour $i\in I$, on ait
$V_{i}=\bigcup_{j\in \sigma^{-1}(i)}V'_{j}$. On dit que $\varphi'$
\emph{raffine} $\varphi$ s'il existe un morphisme de $\varphi$ dans
$\varphi'$.

Soient $\Ff $ un faisceau de groupes abéliens sur $S$, $U\in
\Ob S_{\text{\rm et}}$, $\varphi=(V_{i})_{i\in I}\in \Ob\Psi_{f}(U)$,
$I^{\ast}=I\backslash \{0\}$, et $f_{i}$ la projection de $V_{i}$ sur
$U$. On a
$$
f_{!}f^{\ast}\Ff |U=f_{0!}f_{0}^{\ast}(\Ff |U)\oplus
\bigoplus_{i\neq 0} f_{i\ast}f_{i}^{\ast}(\Ff |U){\quoi},
$$
d'où une flèche composée
$$\etiquette[6.2.1.1]{XVII.eq6.2.1.1}
\tau_{\varphi}:\Ff ^{I^{\ast}}|U\longrightarrow
\bigoplus_{i\neq 0}
f_{i\ast}f_{i}^{\ast}(\Ff |U)\longrightarrow f_{!}f^{\ast}\Ff |U{\quoi}.\leqno(6.2.1.1)
$$

Si $\sigma:I\rightarrow J$ est un morphisme d'ensembles pointés, on
désigne par $\leftexp{t}{\sigma}:\Ff ^{J^{\ast}}\rightarrow
\Ff ^{I^{\ast}}$ la flèche de coordonnées $\leftexp{t}{\sigma}^i_{\phantom{i}j}=\Id$
si $\sigma(i)=j$, et $\leftexp{t}{\sigma}^i_{\phantom{i}j}=0$ si $\sigma(i)\neq
j$. Si $\sigma:I'\rightarrow I$ est un morphisme de $\varphi$ dans
$\varphi'$, le diagramme
$$
\vcenter{\etiquette[6.2.1.2]{XVII.eq6.2.1.2}
\xymatrix@R=1.25cm{
\Ff^{I^{\ast}}|U \ar[rr]^{\leftexp{t}{\sigma}}   \ar[rd] \ar[rd]_{\tau_{\varphi}}&&
\Ff^{I'^{\ast}}|U     \ar[ld]^{\tau_{\varphi'}}\\
      &          f_{!}f^{\ast}\Ff|U  &
}}\leqno{(6.2.1.2)}
$$
est\pageoriginale commutatif. De plus, ${}^{t}(\sigma\rho)={}^{t}\rho
{}^{t}\sigma$. Les $\Psi_{f}(U)$ forment une catégorie fibrée
$\Psi_{f}$ sur $S_{\text{\rm et}}$; les constructions précédentes
commutent à la localisation et définissent donc un $S_{\text{\rm
et}}$-foncteur $\varphi=(V_{i})_{i\in I} \in \Psi_{f}(U)\mapsto
\Ff ^{I^{\ast}}|U$ de $\Psi$ dans la $S_{\text{\rm
et}}$-catégorie des faisceaux sur des ouverts étales de $S$, et un
morphisme de foncteurs
$\tau_{\varphi}:\Ff ^{I^{\ast}}|U\rightarrow
f_{!}f^{\ast}\Ff |U$.

\begin{enonce*}{Proposition 6.2.2}\etiquette[6.2.2]{XVII.prop6.2.2}
\begin{enumerate}
\renewcommand{\theenumi}{\roman{enumi}}
\renewcommand{\labelenumi}{(\theenumi)}
\item La catégorie $\Psi_{f}$ est localement filtrante à droite
\emph{(\sisi{IV 14}{\hyperrefext[V][V.def8.1.1]{8.1.1}})} sur $S_{\text{\rm et}}$.

\item Les flèches \eqhyperref[XVII.eq6.2.1.1]{6.2.1.1} définissent un isomorphisme\sisi{s}{}
\[
\tau:\varinjlim_{\varphi=(V_{i})_{i\in I}\in
\Psi_{f}(U)}\Ff ^{I^{\ast}}|U\isomto  f_{!}f^{\ast}\Ff\sisi{}{|U} {\quoi}.
\]
\end{enumerate}
\end{enonce*}

Si $S$ est strictement local de point fermé $s$, le $S$-schéma $X$
admet une décomposition $\varphi\in\Psi_{f}(S)$ en parties ouvertes et
fermées $X=\coprod_{i\in I}X_{i}$ indexée par un ensemble pointé $I$
telle que la fibre de $X_{0}$ en $s$ soit vide, que les $X_{i}$ $(i\neq
0)$ soient finis sur $S$ et que la fibre de $X_{i}$ $(i\neq 0)$ en $s$
soit réduite à un point (\EGA IV 18.12.1). Une telle décomposition est
un élément final de $\Psi\sisi{}{_f}(S)$. De plus, en vertu de \hyperref[XVII.prop6.1.4]{6.1.4} (iii) (ou
parce que c'est évident) la fibre en $s$ de la flèche $\tau_{\varphi}$
\eqhyperref[XVII.eq6.2.1.1]{6.2.1.1} est un isomorphisme.

Dans le cas général, soit $x$ un point géométrique de $S$ et
$\Psi_{x}$ la fibre de $\Psi_{f}$ en $x$,
\[
\Psi_{x}=\varinjlim_{U(x)} \Psi_{f}(U){\quoi},
\]
pour $U$ voisinage étale de $x$. Soit $S_{x}$ le localisé strict de
$S$ en $x$, $X_{x}=\fprod{X}{S_{x}}{S}$ et $f_{x}:X_{x}\rightarrow
S_{x}$. Les sorites de passage à la limite montrent\pageoriginale que
la catégorie $\Psi_{x}$ est équivalente à la catégorie
$\Psi_{f_{x}}(S_{x})$, donc filtrante à droite. Il en résulte que
$\Psi\sisi{}{_f}$ est localement filtrante à droite. De plus, la fibre en $x$ de
$\tau$ s'identifie à la flèche
$$
\tau_{x}:\varinjlim_{\Psi_{x}}\Ff _{x}^{I_{x}}\longrightarrow
(f_{!}f^{\ast}\Ff )_{x}
$$
et, d'après ce qui précède, cette flèche est un isomorphisme.

\begin{enonce*}{Théorème 6.2.3}\etiquette[6.2.3]{XVII.thm6.2.3}
On peut d'une et d'une seule façon définir, pour tout morphisme
$f:X\rightarrow S$, séparé, plat de présentation finie et quasi-fini,
et tout faisceau abélien $\Ff $ sur $S$, un morphisme trace
$$
\Tr_{f}:f_{!}f^{\ast}\Ff \longrightarrow \Ff
$$
de sorte que les conditions suivantes soient vérifiées:
\begin{itemize}
\item[\textsuperscript{(Var 1)}] $\Tr_{f}$ est fonctoriel en $\Ff $.

\item[\textsuperscript{(Var 2)}] $\Tr_{f}$ est compatible à tout changement de base, \ie
pour tout diagramme cartésien
$$
\xymatrix{
X' \ar[rr]^{g'} \ar[d]^{f'}& & X\ar[d]^{f}\\
S' \ar[rr]^{g} & &  S
}
$$
le diagramme de foncteurs suivant est commutatif:
$$
\xymatrix{
f'_{!} f'^{\ast} g^{\ast} \ar@{=}[r] \ar[d]^{\Tr_{f'}\ast g^{\ast}}& f'_{!}g'^{\ast} f^{\ast}\ar@{=}[r] &
g^{\ast}f_{!}f^{\ast} \ar[d]^{g^{\ast}\ast
  \Tr_f}\\
g^{\ast} \ar@{=}[rr]&  & g^{\ast} &.
}
$$

\item[\textsuperscript{(Var 3)}] $\Tr_{f}$\pageoriginale est compatible avec la
composition, \ie si $f$ est le composé $f=gh$ de deux morphismes
séparés, plats, de présentation finie et quasi-finis, le diagramme
suivant est commutatif:
$$
\xymatrix{
f_{!}f^{\ast}\ar[rrdd]_{\Tr_f} \ar@{=}[rr]& & g_{!}h_{!} h^{\ast}g_{\ast}\ar[d]^{g_{!}\ast\mathord{\Tr_{h}}\ast g^\ast}\\
              & &  g_{!}g^{\ast} \ar[d]^{\Tr_g}\\
              & & \Id .
}
$$

\item[\textsuperscript{(Var 4)}] (Normalisation). Si $f$ est fini de rang constant $n$,
l'homomorphisme composé
$$
\Ff \rightarrow
f_{\ast}f^{\ast}\Ff =f_{!}f^{\ast}\Ff \xrightarrow{\Tr_{f}}\Ff
$$
est la multiplication par $n$.
\end{itemize}
\end{enonce*}

\sisi{\emph{Notons tout d'abord que les conditions \emph{(Var 2)} à \emph{(Var 4)}
impliquent que}}{Notons tout d'abord que les conditions (Var 2) à (Var 4)
impliquent que}

\begin{enonce*}{Lemme 6.2.3.1}\etiquette[6.2.3.1]{XVII.lem6.2.3.1}
Si $X=X'\amalg X''$\sisi{}{,} alors, posant $f'=f|X'$ et $f''=f|X''$, le
morphisme $\Tr_{f}$:
$$
f_{!}f^{\ast}\Ff =f'_{!}{f'}^{\ast}\Ff \oplus
f''_{!}{f''}^{\ast}\Ff \longrightarrow \Ff
$$
est somme de $\Tr_{f'}$ et $\Tr_{f''}$.
$$
\xymatrix@C=2.25cm@R=.75cm{
  X' \ar@{^{(}->}[rd]^{j'}&           &\\
   & X   \ar[r]^{f}  \ar[r]^{f}    & S\\
X'' \ar@{^{(}->}[ru]_{j''}& &.
}
$$
\end{enonce*}

Si\pageoriginale on applique (Var 2) à $j'$, pour le changement de
base par $j'$, et (Var 4) pour l'application identique, on voit que
$$
\Tr_{j'}:j'_{!}{j'}^{\ast}f^{\ast}\Ff \longrightarrow
f^{\ast}\Ff
$$
est l'application canonique définie par la décomposition de
$f^{\ast}\Ff $ en somme directe
$$
f^{\ast}\Ff =j'_{!}{j'}^{\ast}f^{\ast}\Ff \oplus
j''_{!}{j''}^{\ast}f^{\ast}\Ff {\quoi}.
$$

Si on applique (Var 3) à $f'=f\circ j'$, on voit que le diagramme
suivant est commutatif
\[
\xymatrix@R=1.25cm{
f'_{!} f'^{\ast}\Ff  \ar[rr] \ar[rd]_{\Tr_{f'}}& &f_{!}f^{\ast}\Ff \ar[ld]^{\Tr_f}\\
                                   & \Ff    &    ,
}
\]
ce qui, joint à l'énoncé analogue pour $f''$, n'est autre que \hyperref[XVII.lem6.2.3.1]{6.2.3.1}.

Soit $f:X\rightarrow S$ un morphisme quasi-fini séparé plat de
présentation finie. Quel que soit $U\in \Ob S_{\text{\rm et}}$, soit
$\Psi'_{f}(U)$ la sous-catégorie pleine de $\Psi_{f}(U)$ (\sisi{\hyperref[XVII.prop6.1.1]{6.1.1}}{\hyperref[XVII.subsubsec6.2.1]{6.2.1}})
formée des décompositions $(V_{i})_{i\in I}$ telle que, pour $i\neq
0$, $V_{i}$ soit localement libre de rang constant sur $U$. Les
catégories\pageoriginale $\Psi'_{f}(U)$ forment une sous-catégorie
fibrée de $\Psi_{f}$, cofinale dans $\Psi_{f}$.

Soient $\Ff $ et $\mathcal{G}$ des faisceaux abéliens sur
$S$. D'après \hyperref[XVII.prop6.2.2]{6.2.2} (ii), il revient au même de se donner un morphisme
$g:f_{!}f^{\ast}\Ff \rightarrow \mathcal{G}$ ou de se donner,
pour chaque décomposition $\varphi=(V_{i})_{i\in I}\in\Psi'_{f}(U)$,
un morphisme $g_{\varphi}:\Ff ^{I^{\ast}}|U\rightarrow
\mathcal{G}|U$, compatible aux morphismes ${}^{t}\sigma$ de
\eqhyperref[XVII.eq6.2.1.2]{6.2.1.2}. Si $f_{i}$ est la projection de $V_{i}$ sur $U$,
$g_{\varphi}$ est la flèche de coordonnées les flèches
composées:
\[
g^{i}_{\varphi}:\Ff |U\rightarrow
f_{i\ast}f^{\ast}_{i}\Ff |U\rightarrow
f_{!}f^{\ast}\Ff |U\xrightarrow{g} \mathcal{G}|U{\quoi}.
\]

On en déduit qu'il existe une et une seule flèche $\Tr_{f}$, compatible
à la localisation étale et vérifiant (Var 4) et \hyperref[XVII.lem6.2.3.1]{6.2.3.1}. Il est clair
que cette flèche vérifie (Var 1) et (Var~2).

Pour vérifier (Var 3), on se ramène par changement de base au cas
facile où $S$ est le spectre d'un corps algébriquement clos.

\setcounter{subsubsection}{3}
\subsubsection{}\etiquette[6.2.4]{XVII.subsubsec6.2.4}
Une \emph{pondération} d'un morphisme quasi-fini séparé
$f:X\rightarrow S$ est la donnée, pour chaque point géométrique $x$ de
$X$ d'un entier $n(x)$, ne dépendant que du point de $X$ image de $x$,
cette donnée vérifiant:
\begin{itemize}
\item[$(\ast)$] \emph{Quels que soient $U\in \Ob S_{\text{\rm et}}$ et
la partie ouverte et fermée $X_{1}$ de $f^{-1}(U)$, finie sur $U$, la
fonction qui à chaque point géométri\sisi{e}{que} $s$ de $U$ associe le nombre
$\sum_{x\in X_{1}(s)}n(x)$ est localement constante.}
\end{itemize}

Si $f=gh$ est le composé de deux morphismes pondérés, on pondère $f$
en posant $n_{f}(x)=n_{h}(x)\sisi{.}{\cdot}n_{g}(h(x))$.

Si\pageoriginale le morphisme $f'$ se déduit par un changement de base
$g$ d'un morphisme pondéré $f$, on pondère $f'$ en posant
$n_{f'}(x)=n_{f}(g'(x))$.

Les arguments qui précèdent prouvent en fait:

\begin{enonce*}{Proposition 6.2.5}\etiquette[6.2.5]{XVII.prop6.2.5}
On peut d'une et d'une seule façon définir, pour tout morphisme
quasi-fini séparé pondéré $f:X\rightarrow S$ et tout faisceau abélien
$\Ff $ sur $S$, un morphisme trace
$\Tr_{f}:f_{!}f^{\ast}\Ff \rightarrow \Ff $, qui soit
fonctoriel en $\Ff $, compatible aux changements de base et à
la composition, et vérifie:
\begin{itemize}
\item[\textsuperscript{(Var $\sisi{\leftexp{\prime}{4}}{4'}$)}] Si $f$ est fini et s'il existe un entier $n$ tel
que, pour tout point géométrique $s$ de $S$, on ait $\sum_{x\in
X(s)}n(x)=n$, alors la flèche composée
$$
\Ff \longrightarrow
f_{\ast}f^{\ast}\Ff =f_{!}f^{\ast}\Ff \xrightarrow{\Tr_{f}}\Ff
$$
est la multiplication par $n$.
\end{itemize}
\end{enonce*}

\begin{enonce*}[remark]{Exemple 6.2.6}\etiquette[6.2.6]{XVII.exem6.2.6}
Soient $S$ un schéma \sisi{}{réduit et }noethérien géométriquement unibranche (par
exemple normal), et $f:X\rightarrow S$ un morphisme quasi-fini
séparé. Pour tout point géométrique $x$ de $X$ d'image $s$ dans $S$,
soit $\Oo_{x}$ (\resp $\Oo_{s}$) l'anneau strictement local de $X$
en $x$ (\resp de $S$ en $s$). Le schéma $\Spec(\Oo_{s})$ est réduit
et a un unique point maximal $\eta$. De plus, $\Oo_{x}$ est fini sur
$\Oo_{s}$. On pose $n(x)=\rg_{\eta}(\Oo_{x})$. On vérifie que
$n(x)$ est une pondération de $f$, d'où un morphisme trace $\Tr_{f}$.
\end{enonce*}

\setcounter{subsubsection}{6}
\subsubsection{}\etiquette[6.2.7]{XVII.subsubsec6.2.7}
Soit\pageoriginale $u:X\rightarrow Y$ un morphisme séparé quasi-fini
plat de présentation finie. Le morphisme trace (\hyperref[XVII.thm6.2.3]{6.2.3}) étant
fonctoriel, il définit, pour tout complexe de faisceaux abéliens $K$
sur $Y$, un morphisme de complexes
$$\etiquette[6.2.7.1]{XVII.eq6.2.7.1}
\Tr_{u}:u_{!} u^{\ast}K\longrightarrow K{\quoi}.\leqno(6.2.7.1)
$$
On appellera encore, au choix, \emph{morphisme trace, morphisme de
Gysin} ou \emph{morphisme} $u_{\ast}$ tout morphisme déduit de
\eqhyperref[XVII.eq6.2.7.1]{6.2.7.1} ou \hyperref[XVII.thm6.2.3]{6.2.3} par application d'un foncteur. Par exemple les
suivants.

Soit un diagramme commutatif
$$
\xymatrix@C=1.5cm@R=1.75cm{
X  \ar[rr]^{u} \ar[rd]_{f}&   & Y\ar[ld]^{g}\\
                          & S &
}
$$
avec $g$ séparé localement de type fini. Pour tout faisceau abélien
$\Ff $ sur $Y$, le morphisme trace définit un morphisme
$$\etiquette[6.2.7.2]{XVII.eq6.2.7.2}
\Tr_{u}\quad\mbox{\rm ou}\quad u_{\ast}:f_{!}u^{\ast}\Ff \simeq
g_{!}u_{!}u^{\ast}\Ff \xrightarrow{g_{!}(\Tr_{u})} g_{!}\Ff {\quoi}.\leqno(6.2.7.2)
$$

Si $g$ est compactifiable, et si $K$ est un complexe borné
inférieurement à cohomologie de torsion, ou un complexe de
$g^{\ast}\Aa$-modules pour $\Aa$ faisceau d'anneaux de
torsion sur $S$, le morphisme trace \eqhyperref[XVII.eq6.2.7.1]{6.2.7.1} induit de même
\begin{align*}
\Tr_{u}\quad\mbox{\rm ou}\quad u_{\ast} &: Rf_{!}u^{\ast}K\simeq
Rg_{!}u_{!}u^{\ast}K\xrightarrow{Rg_{!}(\Tr_{u})}Rg_{!}K{\quoi},\tag{6.2.7.3}\\
\Tr_{u}\quad\mbox{\rm ou}\quad u_{\ast} &: R^{q}f_{!}u^{\ast}K\simeq
R^{q}g_{!}u_{!}u^{\ast}K\xrightarrow{R^{q}g_{!}(\Tr_{u})}
R^{q}g_{!}K{\quoi}.\tag{6.2.7.4}
\end{align*}\etiquette[6.2.7.3]{XVII.eq6.2.7.3}
\etiquette[6.2.7.4]{XVII.eq6.2.7.4}%L'espace vertical suivant ce paragraphe est trop large.

Pour\pageoriginale $S$ spectre d'un corps algébriquement clos,
\eqhyperref[XVII.eq6.2.7.4]{6.2.7.4} est un morphisme
$$\etiquette[6.2.7.5]{XVII.eq6.2.7.5}
u_{\ast}:\HH^{q}_{c}(X,u^{\ast}K)\longrightarrow \HH^{q}_{c}(Y,K){\quoi}.\leqno(6.2.7.5)
$$

On voit donc que la \emph{cohomologie à supports propres présente un
caractère covariant} vis-à-vis des morphismes quasi-finis.

\subsubsection{}\etiquette[6.2.8]{XVII.subsubsec6.2.8}
Soient $u:X\rightarrow Y$ comme en \hyperref[XVII.subsubsec6.2.7]{6.2.7}, $\Delta_{n}=[0,n]$
l'ensemble type à $\sisi{}{(}n+1\sisi{}{)}$-éléments et $X_{n}$ le produit fibré itéré
$(n+1)$\textsuperscript{uple} de $X$ au-dessus de $Y$. Les $X_{n}$ forment
un schéma simplicial augmenté vers $Y$ (\cf \sisi{VI${}_{\text{B}}$}{VBIS}). On
désignera par $u_{n}:X_{n}\rightarrow Y$ le morphisme
d'augmentation. Si $\Ff $ est un faisceau abélien sur $Y$, les
faisceaux $\Ff _{n}=u_{n!}u_{n}^{\ast}\Ff $ forment,
en vertu \eqhyperref[XVII.eq6.2.7.2]{6.2.7.2} (appliqué après le changement de langage
$S\sisi{\rightarrow}{\mapsto} Y$, $X\sisi{\rightarrow}{\mapsto} X_{n}$, $Y\sisi{\rightarrow}{\mapsto} Y$ ou $X_{m}$),
un faisceau \emph{simplicial} augmenté vers $\Ff $. Le
faisceau différentiel des « \emph{chaînes non dégénérées} »  de ce
faisceau simplicial a par définition pour $n$\textsuperscript{ième}
composante le quotient de $\Ff _{n}$, conoyau de la somme
directe des opérateurs de dégénérescence de but $\Ff _{n}$ et
des morphismes $\sigma-\epsilon(\sigma)$ pour $\sigma$ opérateur de
symétrie.

\begin{enonce*}{Proposition 6.2.9}\etiquette[6.2.9]{XVII.prop6.2.9}
Sous les hypothèses précédentes, pour $u$ étale séparé de type fini
surjectif et pour tout faisceau abélien $\Ff $ sur $Y$, on a
\begin{enumerate}
\renewcommand{\theenumi}{\roman{enumi}}
\renewcommand{\labelenumi}{(\theenumi)}
\item Les faisceaux $\Ff _{n}$ forment une résolution (gauche)
de $\Ff $.

\item Le complexe différentiel des chaînes non dégénérées du faisceau
simplicial $(\Ff _{n})$ est une résolution (gauche) de
$\Ff $. Si les fibres géométriques de $u$ ont toutes au plus
$d$ points, cette résolution est de longueur $\leq d$.
\end{enumerate}
\end{enonce*}

Par\pageoriginale changement de base, il suffit de vérifier cet énoncé
pour $Y$ le spectre d'un corps algébriquement clos, auquel cas
\hyperref[XVII.subsubsec6.2.8]{6.2.8} se ramène au calcul de l'homologie simplicial
d'un simplexe.

Soient $(U_{i})_{i\in I}$ un recouvrement ouvert fini du schéma $Y$,
$\Ff $ un faisceau abélien sur $Y$ et, pour chaque partie $P$
de $I$, soit $j_{P}$ l'inclusion de $U_{P}=\cap_{i\in P}U_{i}$ dans
$\sisi{X}{Y}$. Supposons choisi un ordre total sur $I$. Si $X=\coprod_{i\in
I}U_{i}$ et si $u$ est la projection de $X$ sur $Y$, la résolution
\hyperref[XVII.prop6.2.9]{6.2.9} (ii) de $\Ff $ est donnée par
$$\etiquette[6.2.9.1]{XVII.eq6.2.9.1}
(\Ff _{n})_{n\,\text{dég}}=\bigoplus_{|P|=n+1}j_{P!}j^{\ast}_{P}\Ff {\quoi}.\leqno(6.2.9.1)
$$

Cette résolution garde un sens pour $I$ infini.

\setcounter{subsubsection}{9}
\subsubsection{}\etiquette[6.2.10]{XVII.subsubsec6.2.10}
Sous les hypothèses de \hyperref[XVII.prop6.2.9]{6.2.9}, soient $\Aa$ un faisceau
d'anneaux de torsion sur $Y$ et $K\in \Ob^{-}(Y,\Aa)$. On
annèle $X$ par $u^{\ast}\Aa$. La résolution \hyperref[XVII.prop6.2.9]{6.2.9} (i) (\resp
\hyperref[XVII.prop6.2.9]{6.2.9} (ii)), étant fonctorielles en $\Ff $, définit une
résolution de $K$ par un double complexe $K_{1}$. Cette résolution
s'appelle \emph{la résolution de Cech} extraordinaire (\resp
\emph{extraordinaire alternée}) de $K$.

Soit un diagramme commutatif
$$
\xymatrix@C=1.5cm@R=1.75cm{
X  \ar[rr]^{u} \ar[rd]_{f}&   & Y\ar[ld]^{g}\\
                          & S & ,
}
$$
avec $g$ compactifiable. Avec les notations de
\hyperref[XVII.subsubsec6.2.8]{6.2.8}, si $f_{n}$ est la projection de $X_{n}$ sur
$S$, on a $Rf_{n!}=Rg_{!}\sisi{U}{u}_{n!}$. Le complexe $K_{1}$, filtré
par\pageoriginale le degré semi-simplicial, définit donc une suite
spectrale
$$\etiquette[6.2.10.1]{XVII.eq6.2.10.1}
E^{-p,q}_{1}=R^{q}f_{p!}(u^{\ast}_{p}K)\Longrightarrow R^{-p+q} g_{!}
K{\quoi}.\leqno(6.2.10.1)
$$
Les différentielles $d_{1}$ proviennent des flèches simpliciales, de
sorte que
$$\etiquette[6.2.10.2]{XVII.eq6.2.10.2}
E^{-p,q}_{2}=\check{\Hh}^{p}(R^{q}f_{p!}(u^{\ast}_{p}K)\sisi{}{)}\Longrightarrow
R^{-p+q}g_{!}K{\quoi}.\leqno(6.2.10.2)
$$

Sous les hypothèses de \eqhyperref[XVII.eq6.2.9.1]{6.2.9.1}, on trouve de même une suite
spectrale
$$\etiquette[6.2.10.3]{XVII.eq6.2.10.3}
E^{-p,q}_{1}=\bigoplus_{|P|=p+1}R^{q}f_{p!}(j^{\ast}_{P}K)\Longrightarrow R^{-p+q}g_{!}K{\quoi}.\leqno(6.2.10.3)
$$

Les suites spectrales \eqhyperref[XVII.eq6.2.10.1]{6.2.10.1} et \eqhyperref[XVII.eq6.2.10.3]{6.2.10.3} s'appellent les
\emph{suites spectrales de localisation en haut.}

\begin{enonce*}{Proposition 6.2.11}\etiquette[6.2.11]{XVII.prop6.2.11}
Si $f$ est un morphisme étale, séparé de type fini, le foncteur
$f_{!}$ est un adjoint à gauche du foncteur $f^{\ast}$, et admet le
morphisme trace pour flèche d'adjonction.
\end{enonce*}

On sait que si $X$ est un objet d'un site $S$, le foncteur $f^{\ast}$
de restriction des faisceaux abéliens à $X$ admet un adjoint à gauche
$f_{!}$. Si $F$ est un faisceau abélien sur $X$, $f_{!}F$ est le
faisceau engendré par le préfaisceau
$$
V\longrightarrow \bigoplus_{\xi \in \Hom(V,X)}F(\xi){\quoi}.
$$
Dans le cas étale, notons provisoirement ce foncteur $f\sisi{''}{^0}_{!}$. De sa
propriété d'adjonction résulte aussitôt que sa formation commute à
tout changement de base. Le morphisme trace définit un morphisme $t$
de $f_{!}\Ff $ dans $f^{0}_{!}\Ff $,
compatible\pageoriginale à tout changement de base. Il suffit de voir
que $t$ est un isomorphisme fibre par fibre, \ie pour $S$ spectre d'un
corps séparablement clos, où cela est trivial.

\setcounter{subsubsection}{11}
\paragraph{}\etiquette[6.2.11.1]{XVII.par6.2.11.1}
On prendra garde que cette interprétation de $f^{0}_{!}$ est
spécifique au cas du site étale et ne s'étend pas au site fpqf d'un
préschéma, par exemple.

\subsection{Variante de la trace pour des coefficients continus}\etiquette[6.3]{XVII.subsec6.3}

On démontre dans ce nᵒ les assertions de \textsc{Grothendieck}
\cite[nᵒ 6]{SGA4-17:Gr}, et on applique ces résultats à la définition
d'un « morphisme trace » pour des faisceaux fppf convenables.

\subsubsection{}\etiquette[6.3.1]{XVII.subsubsec6.3.1}
Soit $A$ un anneau commutatif. Si $L$ et $M$ sont deux $A$-modules
projectifs de type fini, l'application $u\mapsto
\bigwedge\limits^{m}u$, de $\Hom(L,M)$ dans
$\Hom(\bigwedge\limits^{m}L,\bigwedge\limits^{m}M)$, est
compatible à tout changement de base et $m$\sisi{-ique}{\textsuperscript{ique}} (\sisi{\hyperref[XVII.subsubsec5.5.2]{5.5.2}}{\hyperref[XVII.par5.5.2.2]{5.5.2.2}}). D'après
\sisi{\hyperref[XVII.subsubsec5.5.2]{5.5.2}}{\hyperref[XVII.par5.5.2.3]{5.5.2.3}}, elle définit un morphisme
$$\etiquette[6.3.1.1]{XVII.eq6.3.1.1}
\bigwedge^{m}:\Gamma^{m}\Hom(L,M)\longrightarrow
\Hom(\bigwedge^{m} L, \bigwedge^{m}M)\leqno(6.3.1.1)
$$
tel que, après tout changement de base, pour $u\in \Hom(L,M)$, on ait
$\bigwedge\limits^{m}(\gamma^{m}(u))=\bigwedge\limits^{m}u$.
Puisque $\Hom(L,M)$ est plat, ce morphisme s'identifie (\sisi{\hyperref[XVII.subsubsec5.5.2]{5.5.2}}{\hyperref[XVII.par5.5.2.5]{5.5.2.5}}) à un
morphisme
$$\etiquette[6.3.1.2]{XVII.eq6.3.1.2}
\bigwedge^{m}:TS^{m}\Hom(L,M)\longrightarrow
\Hom(\bigwedge^{m}L, \bigwedge^{m}M){\quoi}.\leqno(6.3.1.2)
$$
N.B. On peut aussi définir un tel homomorphisme dès que 2 est
inversible dans $A$: un élément symétrique de
$\bigotimes\limits^{m}\Hom(L,M)$ définit un
homomorphisme symétrique de $\bigotimes\limits^{m}L$
dans $\bigotimes\limits^{m}M$, et celui-ci passe au
quotient pour définir un homomorphisme de
$\bigwedge\limits^{m}L$ dans
$\bigwedge\limits^{m}M$.

Etant\pageoriginale donnés trois $A$-modules projectifs de type fini,
le diagramme
$$
\vcenter{\etiquette[6.3.1.3]{XVII.eq6.3.1.3}
\xymatrix{
\Gamma^{m} \Hom(L,M) \otimes \Gamma^{m} \Hom(M,N)  \ar[r] \ar[d]^{\bigwedge\limits^{m}\otimes \bigwedge\limits^{m}}& \Gamma^{m}
\Hom(L,N) \ar[d]^{\bigwedge\limits^{m}}\\
\Hom(\bigwedge\limits^{m} L,
\bigwedge\limits^{m} M) \otimes
\Hom(\bigwedge\limits^{m} M,
\bigwedge\limits^{m} N)         \ar[r]     &
\Hom(\bigwedge\limits^{m} L,
\bigwedge\limits^{m} N),
}}\leqno{(6.3.1.3)}
$$
dans lequel les flèches horizontales sont déduites des flèches de
composition (et, pour la première, de (\hyperref[XVII.par5.5.2.8]{5.5.2.8})), est commutatif. Ceci
traduit la formule
$\bigwedge\limits^{m}(vu)=\bigwedge\limits^{m}v\circ
\bigwedge\limits^{m}u$.

Si $L$ est un $A$-module libre de rang $n$, \eqhyperref[XVII.eq6.3.1.2]{6.3.1.2} définit un
morphisme canonique
$$\etiquette[6.3.1.4]{XVII.eq6.3.1.4}
\det:\Gamma^{n}\End_{A}(L)=TS^{n}\End_{A}(L)\longrightarrow
A{\quoi},\leqno(6.3.1.4)
$$
 qui est un homomorphisme d'algèbres d'après
 \eqhyperref[XVII.eq6.3.1.3]{6.3.1.3}. Par définition,
 on a
$$\etiquette[6.3.1.5]{XVII.eq6.3.1.5}
\det (u^{\otimes n})=\det (u\text{; }L){\quoi}.\leqno(6.3.1.5)
$$

Si $L$ est de plus un module sur une $A$-algèbre $B$, \eqhyperref[XVII.eq6.3.1.4]{6.3.1.4}
définit par composition un homomorphisme d'algèbres
$$\etiquette[6.3.1.6]{XVII.eq6.3.1.6}
{\det}_{L}:TS^{n}_{A}(B)\longrightarrow A{\quoi}.\leqno(6.3.1.6)
$$

\begin{enonce*}{Lemme 6.3.2}\etiquette[6.3.2]{XVII.lem6.3.2}
Supposons que $L$ s'insère dans une suite exacte
\[
\sisi{F}{E}:0\longrightarrow L'\longrightarrow L\longrightarrow
L''\longrightarrow 0
\]
avec $L'$ (\resp $L''$) localement libre de rang $n_{1}$ (\resp
$n_{2}$). Soient $n=n_{1}+n_{2}$, $\alpha$ le morphisme « gradué
associé »  du groupe $\End(E)$ des endomorphismes\pageoriginale de
l'extension $E$ dans $\End(L')\oplus \End(L'')$ et, pour $X$ et $Y$
deux $A$-modules plats, soit $\beta$ le morphisme de $TS^{n}(X\oplus
Y)$ dans $TS^{n_{1}}(X)\otimes TS^{n_{2}}(Y)$ déduit de l'isomorphisme
(\hyperref[XVII.par5.5.2.6]{5.5.2.6})
\[
TS^{n}(X\oplus Y)\simeq \sisi{\sum}{\bigoplus}_{n'+n''=n} TS^{n'\sisi{_{1}}{}}(X)\otimes
TS^{n''}(Y){\quoi}.
\]
Le diagramme suivant est commutatif:{\footnotesize
\[
\xymatrix@C=1cm@R=1.5cm{
  TS^{n}(\End\sisi{)}{} (E)) \ar[d] \ar[r]^(.4){TS^{n}(\alpha)} & TS^{n}(\End(L')\oplus
\Ens(L''))  \ar[r]^(.45){\beta} & TS^{n_1} (\End(L')) \otimes TS^{n_2}
(\End(L'')) \ar[d]^{\det \otimes \det}\\
TS^{n}(\sisi{L}{\End(L)})         \ar[rr]^{\det} &                 & A.
}
\]
}
\end{enonce*}

D'après la propriété universelle \hyperref[XVII.subsubsec5.5.2]{5.5.2} de $TS^{n}$, il suffit de
vérifier que pour $u\in \End(E)$ les deux images de $u^{\otimes n}$
coïncident. Si $\alpha(u)=(u',u'')$, ces images sont respectivement
$\det(u)$ et $\det(u').\det(u'')$, d'où la conclusion.

\begin{enonce*}{Lemme 6.3.3}\etiquette[6.3.3]{XVII.lem6.3.3}
Soient $A$ un anneau\sisi{}{ commutatif}, $B$ une $A$-algèbre\sisi{}{ commutative}, $C$ une $B$-algèbre non
nécessairement commutative, $E$ un $B$-module libre de rang $n$ sur
$A$, et $F$ un $C$-module libre de rang $m$ sur $B$. On suppose $C$
plat sur $A$. D'après la propriété universelle \hyperref[XVII.par5.5.2.2]{5.5.2.2} de
$\Gamma^{nm}$ et \hyperref[XVII.par5.5.2.3]{5.5.2.3}, il existe une unique flèche
$\sigma:TS^{nm}_{A}(C)\rightarrow TS^{n}_{A}TS^{\sisi{n}{m}}_{B}(C)$ telle
qu'après tout changement d'anneaux $A\rightarrow A'$, on ait
$\sigma(c^{\otimes nm})=(c^{\otimes m})^{\otimes n}$. Le diagramme
suivant est alors commutatif:
$$
\xymatrix@C=2cm@R=1.5cm{
TS^n_A (B) \ar[d]^{\det_E}  &  & TS^{n}_A TS^{m}_B (C) \ar[ll]_{TS^n_A
  (\det_F)}\\
A    & & TS^{nm}_A (C) \ar[ll]_{\det_{E\otimes_B F}} \ar[u]_{\sigma}.
}
$$
\end{enonce*}

\begin{enonce*}[remark]{Preuve}
D'après\pageoriginale \hyperref[XVII.par5.5.2.2]{5.5.2.2} et \hyperref[XVII.par5.5.2.3]{5.5.2.3},
il suffit de montrer que les
deux images d'un élément $u^{\otimes nm}$ de $TS^{nm}_{A}(C)$
coïncident. Ces images sont $\det_{E}(\det_{F}(u))$ et
$\det_{E\otimes_{B}F}(u)$; elles s'expriment en termes de $A$, $B$,
$E$, $F$ et $u\in \End_{B}(F)$ et, pour vérifier qu'elles coïncident,
on peut au préalable remplacer $B$ par son image (commutative) $B_{1}$
dans $\End(E)$, et $F$ par $F\otimes_{B}B_{1}$. La formule à prouver
est alors la formule de calcul par bloc des déterminants (\textsc{Bourbaki},
\sisi{Alg.}{\emph{Alg.}} \sisi{Ch.}{\chap}~8 \S~12 lemme 2 \sisi{p.140}{p.~140}).
\end{enonce*}

Lorsque $C$ n'est plus nécessairement plat sur $A$, la même
démonstration fournit plus généralement un diagramme commutatif
d'algèbres
$$\vcenter{\etiquette[6.3.3.1]{XVII.eq6.3.3.1}
\xymatrix@C=3cm@R=1.5cm{
\Gamma^{n}_A (B)  \ar[d]^{\det_E}& \Gamma^{n}_A
\Gamma^n_{B}(C) \ar[l]_{\Gamma_A^{n}(\det_F)}\\
A            & \ar[l]_{\det_{E\otimes_B F}}  \Gamma^{nm}_A (C) \ar[u]_{\sigma}.
}}\leqno{(6.3.3.1)}
$$

\setcounter{subsubsection}{3}
\subsubsection{}\etiquette[6.3.4]{XVII.subsubsec6.3.4}
Si $B$ est une $A$-algèbre commutative, $\Spec(TS^{n}(B))$ n'est autre
que la puissance symétrique $n$\textsuperscript{ième} de $\Spec(B)$ sur
$\Spec(A)$, \ie le quotient de $(\Spec(B)/\Spec(A))^{n}$ par le groupe
symétrique (\SGA 1 \sisi{IV 3}{V 1}).

Si\pageoriginale $X$ est un $S$-schéma séparé\sisi{}{ sur $S$}, on dira qu'un
$\Oo_{X}$-module quasi-cohérent $\Ff $ sur $X$ est
\emph{fini localement libre} (\resp et de rang $n$) \emph{sur} $S$ si
son support est fini sur $S$ et si son image directe est localement
libre de rang fini (\resp de rang $n$) sur $S$. Par \emph{support} de
$\Ff $, on entendra encore le \emph{sous-schéma} de $X$ défini
par l'annulation de $\Ff $ (qui est un Idéal quasi-cohérent,
$\Ff $ étant quasi-cohérent de type fini).

Si $\Ff $ est fini localement libre sur $S$, de rang $n$ et de
support $E$, le schéma $E$ est affine sur $S$; par globalisation sur
$S$, $\Ff $ définit, via \eqhyperref[XVII.eq6.3.1.6]{6.3.1.6}, une section
$\sigma_{E}(\Ff )$ de $\Sym^{n}_{S}(E)$, d'où, par
fonctorialité, une section $\sigma(\Ff )$ de l'espace
algébrique \footnote[1]{\label{XVII.noteArtin1}M. \textsc{Artin}, Algebraization of formal moduli
I\sisi{}{, dans \emph{Global Analysis: Papers in Honor of K.~Kodaira}, Princeton University Press, 1969, 21--71}. On n'aura à utiliser la construction précédente que lorsque
$\Sym^{n}_{S}(X)$ existe en tant que schéma, ce qui est le cas pour
$X$ quasi-projectif sur $S$.} $\Sym^{n}_{S}(X)$,
$$\etiquette[6.3.4.1]{XVII.eq6.3.4.1}
\sigma(\Ff )\in \Gamma(\Sym^{n}_{S}(X)/S){\quoi}.\leqno(6.3.4.1)
$$

En particulier, si $f:S'\rightarrow S$ est fini localement libre de
rang $n$, on obtient, pour $\Ff =\Oo_{S'}$, une section
canonique
$$\etiquette[6.3.4.2]{XVII.eq6.3.4.2}
t\in\Gamma(\Sym^{n}_{S}(S')/S){\quoi}.\leqno(6.3.4.2)
$$

\begin{enonce*}{Proposition 6.3.5}\etiquette[6.3.5]{XVII.prop6.3.5}
Soient $X$ un schéma séparé sur $S$ et $\Ff _{1}$,
$\Ff _{2}$ deux $\Oo_{X}$-modules finis localement libres de
rang $n_{1}$ et $n_{2}$ sur $S$. On pose $n=n_{1}+n_{2}$. Soit $d$ la
flèche canonique $d:\Sym^{n_{1}}_{S}(X)\times
\Sym^{n_{2}}_{S}(X)\rightarrow \Sym^{n}_{S}(X)$. Si $\Ff $ est
une extension de $\Ff _{1}$ par $\Ff _{2}$, on a
$$
\sigma(\Ff )=d(\sigma(\Ff _{1}),\sigma(\Ff _{2})){\quoi}.
$$
\end{enonce*}

On\pageoriginale peut remplacer $X$ par le support (fini sur $S$) de
$\Ff _{1}\oplus \Ff _{2}$ et supposer $S$
affine. Traduite en terme d'anneaux, la proposition résulte alors de
\hyperref[XVII.lem6.3.2]{6.3.2}.

\begin{enonce*}{Lemme 6.3.6}\etiquette[6.3.6]{XVII.lem6.3.6}
Si $X$ est un schéma séparé sur $S$, $\Ff $ un
$\Oo_{X}$-module fini localement libre de rang $n$ sur $S$ et si
$\Ll$ est un faisceau inversible sur $X$, on a
$$
\sigma(\Ff )=\sigma(\Ff \otimes\Ll){\quoi}.
$$
\end{enonce*}

Remplaçons $X$ par le support (\hyperref[XVII.subsubsec6.3.4]{6.3.4})
de $\Ff $\sisi{ l}{. L}ocalement
sur $S$, $\Ll$ est alors isomorphe à $\Oo_{X}$, et
l'assertion devient triviale.

Le lecteur vérifiera:

\begin{enonce*}{Proposition 6.3.7}\etiquette[6.3.7]{XVII.prop6.3.7}
Pour tout diagramme de $S$-schémas séparés sur $S$
\[
\xymatrix{
X        \ar[rr]^{u} \ar[rd]_{f} &        &  Y  \ar[ld]^{g}\\
                                & S      &,
}
\]
et tout $\Oo_{X}$-module $\Ff $ fini localement libre de
rang $n$ sur $S$, le $\Oo_{Y}$-module $u_{\ast}\Ff $ est
fini localement libre de rang $n$ sur $S$, et
\[
\sisi{u}{(\Sym^n_S u)}(\sigma(\Ff ))=\sigma(u_{\ast}\Ff ){\quoi}.
\]
\end{enonce*}

\setcounter{subsubsection}{7}
\paragraph{}\etiquette[6.3.7.1]{XVII.par6.3.7.1}
Notons enfin que la construction des morphismes $\sigma$ \eqhyperref[XVII.eq6.3.4.1]{6.3.4.1} et
$t$ \eqhyperref[XVII.eq6.3.4.2]{6.3.4.2} est \emph{compatible à tout changement de base}, ce qui
a un sens d'après \hyperref[XVII.par5.5.2.7]{5.5.2.7}.

\subsection*{Application 1. Le cas des courbes lisses}\pageoriginale

\subsubsection{}\etiquette[6.3.8]{XVII.subsubsec6.3.8}
Soit $f:X\rightarrow S$ une courbe lisse et quasi-projective sur $S$,
\ie un morphisme lisse quasi-projectif purement de dimension relative
1. La théorie des schémas de Hilbert \sisi{[}{}\textsc{Grothendieck} \cite{SGA4-17:Gr}) montre que le
   faisceau fpqc $\sheaf{\Div}^{n}_{X/S}$ qui, à chaque $S$-schéma $S'$,
   associe l'ensemble des diviseurs relatifs sur $X'=\fprod{X}{S'}{S}$
   qui sont \emph{finis} sur $S$ et de degré $n$, est un faisceau
   représentable.

Pour chaque diviseur relatif $E$ sur $X$, fini de degré $n$ sur $S$,
on dispose d'un élément canonique \sisi{(6.3.4.3)}{\eqhyperref[XVII.eq6.3.4.2]{6.3.4.2}} $t(E)\in \Sym^{n}_{S}(E)$,
d'image $\sigma(E)$ dans $\Sym^{n}_{S}(X)$. Cette construction, étant
compatible à tout changement de base, définit
$$\etiquette[6.3.8.1]{XVII.eq6.3.8.1}
\sigma:\sheaf{\Div}^{n}_{X/S}\longrightarrow
\Sym^{n}_{S}(X):E\longrightarrow \sigma(\Oo_{E}){\quoi}.\leqno(6.3.8.1)
$$

Chaque section $s$ de $f$ définit un diviseur relatif fini sur $S$ de
degré 1, le diviseur $s(S)$. Le morphisme
\[
(X/S)^{n}\longrightarrow \sheaf{\Div}^{n}_{X/S}:(s_{i})_{0<i\leq
n}\longmapsto \sum s_{i}(S)
\]
est symétrique, et induit donc un morphisme
$$\etiquette[6.3.8.2]{XVII.eq6.3.8.2}
\alpha:\Sym^{n}_{S}(X)\longrightarrow \sheaf{\Div}^{n}_{X/S}{\quoi}.\leqno(6.3.8.2)
$$

\begin{enonce*}{Proposition 6.3.9}\etiquette[6.3.9]{XVII.prop6.3.9}
Les morphismes $\sigma$ et $\alpha$ de \hyperref[XVII.subsubsec6.3.8]{6.3.8} sont des isomorphismes
inverses l'un de l'autre.
\end{enonce*}

\begin{enonce*}[remark]{Preuve}
\begin{enumerate}
\renewcommand{\theenumi}{\alph{enumi}}
\renewcommand{\labelenumi}{\theenumi)}
\item $\sigma\alpha=\Id$. Il suffit de prouver que la flèche composée
$$
(X/S)^{n}\xrightarrow{\pi}\Sym^{n}_{S}(X)\xrightarrow{\alpha}\sheaf{\Div}^{n}_{X/S}\xrightarrow{\sigma}\Sym^{n}_{S}(X)
$$
est\pageoriginale $\pi$. Pour $n=1$, l'assertion est claire; le cas
général s'en déduit via le

\begin{enonce*}{Lemme 6.3.9.1}\etiquette[6.3.9.1]{XVII.lem6.3.9.1}
Avec les notations de \hyperref[XVII.prop6.3.5]{6.3.5} et \eqhyperref[XVII.eq6.3.8.1]{6.3.8.1}, on a
\[
\sigma(E+F)=d(\sigma(E),\sigma(F)){\quoi}.
\]
\end{enonce*}

Ce lemme résulte de \hyperref[XVII.prop6.3.5]{6.3.5}, \hyperref[XVII.lem6.3.6]{6.3.6} et de la suite exacte
\[
0\longrightarrow \Oo(-F)\otimes \Oo_{E}\longrightarrow
\Oo_{E+F}\longrightarrow \Oo_{F}\longrightarrow 0{\quoi}.
\]

\item On sait que $\sheaf{\Div}^{n}_{X/S}$ est lisse sur $S$ (il n'y a pas
d'obstruction à relever infinitésimalement des diviseurs sur une
courbe (\cf \SGA 3 XI 1)). Le morphisme
$\alpha\pi:(X/S)^{n}\rightarrow \sheaf{\Div}^{n}_{X/S}$ est un morphisme
quasi-fini surjectif entre schéma\sisi{}{s} lisses sur $S$ de source à fibres
équidimensionnelles. Il est donc plat (\EGA ${\rm 0_{IV}}$ 17.3.5),
donc $\alpha\pi$ est un épimorphisme de faisceaux fppf, donc $\alpha$
l'est aussi. Compte tenu de a), ceci achève la démonstration de \hyperref[XVII.prop6.3.9]{6.3.9}.
\end{enumerate}
\end{enonce*}

\begin{enonce*}{Lemme 6.3.10}\etiquette[6.3.10]{XVII.lem6.3.10}
Soient $k$ un anneau commutatif, $M$ un $k$-module libre et $u$ un
endomorphisme de $M$. Désignons par $N_{\det(u-T)=0}$ la norme de
$k[T]/\det(u-T)$ à $k$. Pour $P\in k[T]$, on a alors
\[
N_{\det(u-T)=0} (P(T))=\det(P(u),M){\quoi}.
\]
\end{enonce*}

Il s'agit de vérifier une identité algébrique entre les coefficients
de la matrice de $u$ et les coefficients de $P$. Pour prouver \hyperref[XVII.lem6.3.10]{6.3.10},
il suffit donc (grâce à des arguments familiers) de traiter le cas où
$k$ est un corps algébriquement clos et où $u$ est semi-simple, de
valeurs propres $a_{i}$.

On\pageoriginale a alors
\begin{align*}
& \det(u-T)=\Pi(a_{i}-T)\\
& N_{\det(u-T)=0}(P(T))=\Pi P(a_{i})\quad\mbox{\rm et}\\
& \det(P(u),M)=\Pi P(a_{i}){\quoi}.
\end{align*}

\setcounter{subsubsection}{11}
\setcounter{paragraph}{-1}
\paragraph{}\etiquette[6.3.11.0]{XVII.par6.3.11.0}
Soient $f:X\rightarrow S$ un morphisme quasi-projectif, de
présentation finie, plat à fibres purement de dimension 1 et
Cohen-Macaulay; ces conditions sont donc vérifiées si $X$ est une
courbe lisse quasi-projective sur $S$. Soit $\Ff $ un
$\Oo_{X}$-module, fini localement libre \emph{sur} $S$. Le module
$\Ff $ étant nul sur un ouvert dense fibre par fibre de $X$,
on dispose d'une flèche canonique $\psi:\Oo_{X}\rightarrow
\det^{\ast}(\Ff )$ (\SGA 6 XI ou \textsc{Mumford} \sisi{[1 Ch.~5 \S~3]}{\cite[\chap~5 \S~3]{SGA4-17:Mu}}). On
désignera par $\delta(\Ff )$ le diviseur défini par l'idéal
annulateur de $\coker(\psi)$ (Ibidem); ce diviseur est fini sur $S$.

\begin{enonce*}{Proposition 6.3.11.1}\etiquette[6.3.11.1]{XVII.prop6.3.11.1}
Sous les hypothèses \hyperref[XVII.par6.3.11.0]{6.3.11.0} et avec les notations de
\hyperref[XVII.subsubsec6.3.8]{6.3.8} et \eqhyperref[XVII.eq6.3.4.1]{6.3.4.1}, si $\Ff $ est de
rang $n$, alors $\delta(\Ff )$ est de degré $n$ et
\[
\sigma(\Ff )=\sigma(\delta(\Ff )){\quoi}.
\]
\end{enonce*}

\begin{enonce*}[remark]{Démonstration}
\begin{enumerate}
\renewcommand{\theenumi}{\alph{enumi}}
\renewcommand{\labelenumi}{\theenumi)}
\item \emph{Cas $f$ lisse.} On se ramène de façon standard au cas
$S$ local artinien à corps résiduel algébriquement clos. L'additivité
(\hyperref[XVII.prop6.3.5]{6.3.5}) de $\sigma$ permet de supposer le support de $\Ff $
concentré en un point rationnel $x$ de $X$. On vérifie aussitôt que si
$T$ est une uniformisante en $x$, définissant $p:X\rightarrow
\sisi{A}{\AAA}^{1}_{S}$, alors $p_{\ast}(\Oo_{\sisi{S}{\delta}(\Ff )})\simeq
\Oo_{\delta(p_{\ast}\Ff )}$ (car les constructions
effectuées commutent au passage au complété). On se ramène par là et
\hyperref[XVII.prop6.3.7]{6.3.7} au cas où $X=\sisi{A'}{\AAA^1}_{S}$.

Si\pageoriginale $X=\AAA\sisi{'}{^1}_{S}$ et si $S$ est affine,
$S=\Spec(k)$, alors $\Ff $ s'identifie à un $k$-module
localement libre $M$, muni d'un endomorphisme $u$ (l'action de
$T$). Le module $M$, en tant que $k[T]$-module, admet la résolution
\[
0\rightarrow
k[T]\otimes_{k}M\xrightarrow{\varphi}k[T]\otimes_{k}M\xrightarrow{\epsilon}
M\rightarrow 0{\quoi},
\]
où $\epsilon(P\otimes m)=P(u)(m)$ et où $\varphi(1\otimes m)=1\otimes
um-T\otimes m$. De la formule $\varphi=u-T$, on déduit que le diviseur
$\delta(\Ff )$ admet pour équation $\det(u-T,M)=0$.

Pour vérifier \sisi{6.3.11}{\hyperref[XVII.prop6.3.11.1]{6.3.11.1}}, il suffit de montrer que, après tout changement
de base et pour toute $\varphi$ section de $\Oo_{X}$, on a
\[
{\det}_{k}(\varphi,M)={\det}_{k}(\varphi,\Oo_{\delta(\Ff )})(=N_{\delta(\Ff )/S}(\varphi)){\quoi}.
\]
Pour $\sisi{f}{\varphi}=P(T)$, c'est ce qu'affirme \hyperref[XVII.lem6.3.10]{6.3.10}.

\item Cas général. Explicitons d'abord l'énoncé à prouver sous la
forme plus concrète suivante:

\begin{enonce*}{Corollaire 6.3.11.2}\etiquette[6.3.11.2]{XVII.cor6.3.11.2}
Sous les hypothèses \hyperref[XVII.par6.3.11.0]{6.3.11.0} sur $f:X\rightarrow S$, soient $M$ un
module localement libre de type fini sur $X$, et $u$ un endomorphisme
de $M$ tel que $\det u$ soit une section de $\Oo_{X}$ régulière sur
chaque fibre $X_{s}$ (d'où résulte que $\Coker u=M/u(M)$ et $\Coker
\det u=\Oo_{X}/\det u$ sont tous deux plats de présentation finie
sur $S$ et à \sisi{fibres}{supports} quasi-finies). Supposons que le support commun\nde{le support de $M/uM$ au sens de \hyperref[XVII.subsubsec6.3.4]{6.3.4} est un sous-schéma du support de $\Oo/\det u$ défini par un Nilidéal.} de
ces faisceaux soit fini sur $S$, de sorte qu'ils définissent des
faisceaux localement libres de type fini sur $S$. Alors pour toute
section $\varphi$ de $\Oo_{X}$, les déterminants par rapports à
$\Oo_{S}$ de $\varphi$ opérant sur $\Coker u$ et sur $\Coker \det u$
sont égaux:
$$\etiquette[6.3.11.3]{XVII.eq6.3.11.3}
{\det}_{S}(\varphi,M/u(M))={\det}_{S}(\varphi,\Oo_{X}/\det
u){\quoi}.\leqno(6.3.11.3)
$$
\end{enonce*}

On\pageoriginale notera que cet énoncé implique l'égalité des rangs
des deux Modules localement libres envisagés sur $S$, par raison
d'homogénéité (lorsqu'on y remplace $\varphi$ par $T\varphi$, $T$ une
section « indéterminée »  de $\Oo_{S}$). D'ailleurs, rempla\c
cant $\varphi$ par $1+T\varphi$, où $\sisi{j}{T}$ est une telle indéterminée (de
façon précise, on fait un changement de base $S'=S[T]\rightarrow
S$), on déduit de \eqhyperref[XVII.eq6.3.11.3]{6.3.11.3} la formule en apparence plus générale
$$\etiquette[6.3.11.4]{XVII.eq6.3.11.4}
\Pol_{S}(\varphi,M/uM)(T)=\Pol_{S}(\varphi,\Oo_{X}/\det u)(T)\leqno(6.3.11.4)
$$
sur les polynômes caractéristiques.

Pour prouver \eqhyperref[XVII.eq6.3.11.3]{6.3.11.3}, nous allons nous ramener au cas lisse déjà
prouvé. Par localisation étale sur $S$, on se ramène au cas où $S$ est
strictement local de point fermé $s$, puis par additivité on se ramène
au cas où le support de $M/u(M)$ a un seul point sur $X_{s}$, soit
$x$. Quitte à faire une extension plate de la base, on peut supposer
l'extension résiduelle $k(x)/k(s)$ triviale, et quitte à remplacer $X$
par un voisinage assez petit de $x$, et à ajouter à $\varphi$ une
fonction qui induise l'opération nulle sur $M/u(M)$ et sur
$\Oo_{X}/ \det u$ (ce qui ne modifie pas la formule à prouver), on
voit aisément qu'on peut supposer que $\varphi$ induit une section de
$\Oo_{X_{s}}$ régulière en $x$, donc (quitte à rapetisser $X$)
définit un morphisme quasi-fini et plat
$$
p:X\longrightarrow X'=S[T]{\quoi}.
$$
Par localisation étale sur $X$ et sur $X'$, on peut donc supposer
qu'on a un morphisme \emph{fini} et \emph{plat} $p$ comme dessus, avec
$X'$ lisse sur $S$, et $\varphi=p^{\ast}(\varphi')$, où $\varphi'$ est
une section de $\Oo_{X'}$. Posons alors
\begin{align*}
& M'=p_{\ast}(M)\quav u'=p_{\ast}(u){\quoi},\\
& P=\Coker u\quav Q=\Coker \det u{\quoi},\\
& P'=\Coker u'\quav Q'=\Coker \det u'{\quoi}.
\end{align*}
On\pageoriginale sait qu'on a
\[
\det u'=N_{\sisi{X'/X}{X/X'}}\det u,\quad\mbox{\rm donc}\quad Q'=\Coker
N_{\sisi{X'/X}{X/X'}}(\det u){\quoi}.
\]
On a trivialement
\[
{\det}_{S}(\varphi,P)={\det}_{S}(\varphi',P\sisi{}{'}){\quoi},
\]
et d'après la formule \eqhyperref[XVII.eq6.3.11.3]{6.3.11.3} prouvée dans le cas lisse,
\[
{\det}_{S}(\varphi',P')={\det}_{S}(\varphi',Q'){\quoi},
\]
de sorte qu'il reste à prouver la formule
\[
{\det}_{S}(\varphi',Q')={\det}_{S}(\varphi,Q){\quoi},
\]
qui \emph{équivaut} à la formule \eqhyperref[XVII.eq6.3.11.3]{6.3.11.3} à établir. On voit donc
que la validité de celle-ci ne dépend que de la section $\det u$ de
$\Oo_{X}$, ou encore qu'elle équivaut à la même formule dans le cas
où $M$ est remplacé par $\Oo_{X}$, et $\varphi$ par $\det u$. Mais
dans le cas $M=\Oo_{X}$ la formule \eqhyperref[XVII.eq6.3.11.3]{6.3.11.3} est tautologique !
Cela achève la démonstration.
\end{enumerate}
\end{enonce*}

\begin{enonce*}[remark]{Remarque 6.3.11.5}\etiquette[6.3.11.5]{XVII.rem6.3.11.5}
Dans\pageoriginale le cas où $S$ est artinien, l'égalité des rangs sur
$S$ de $\Coker u$ et $\Coker \det u$ (notations de \hyperref[XVII.cor6.3.11.2]{6.3.11.2}) est un
cas particulier de \EGA IV 21.10.17.3, dont \eqhyperref[XVII.eq6.3.11.3]{6.3.11.3} peut être
considéré comme une version « relative ».
\end{enonce*}

\begin{enonce*}{Proposition 6.3.12}\etiquette[6.3.12]{XVII.prop6.3.12}
Soient $X$ et $Y$ deux courbes quasi-projectives de présentation
finie, plates à fibres Cohen-Macaulay sur $S$, et $u:X\rightarrow Y$
un morphisme quasi-fini et plat.

Le diagramme suivant
\[
\xymatrix@C=4em{
\sheaf{\Div}^{n}_{X/S}  \ar[r]^{u_{\ast}} \ar[d]_{\sigma_X} &
\sheaf{\Div}^{n}_{Y/S} \ar[d]_{\sigma_Y}\\
\Sym^{n}_S(X)     \ar[r]^{\Sym^{n}(u)} & \Sym^{n}_S(Y)
}
\]
est alors commutatif \footnote[1]{\label{XVII.note192}Le morphisme $u_{\ast}$ est défini
parce qu'on se limite aux diviseurs \emph{finis} sur $S$. On prend par
exemple pour définition $u_{\ast}(D)=\delta(u_{\ast}\Oo_{D})$,
formule qui coïncide avec \EGA IV 21.5.5 pour $u$ fini.}.
\end{enonce*}

\sisi{\noindent}{}%
Puisque $u_{\ast}(D)=\delta(u_{\ast}\Oo_{D})$, \hyperref[XVII.prop6.3.12]{6.3.12} résulte de
\hyperref[XVII.prop6.3.7]{6.3.7} et \hyperref[XVII.prop6.3.11.1]{6.3.11.1}.

\subsection*{Application 2. Le morphisme trace}

\setcounter{subsubsection}{12}
\subsubsection{}\etiquette[6.3.13]{XVII.subsubsec6.3.13}
Soit $G$ un faisceau abélien sur le grand site fppf d'un schéma $S$
(\hyperref[XVII.subsec0.10]{0.10}). Supposons que $G$ vérifie la condition suivante.

\paragraph{}\etiquette[6.3.13.1]{XVII.par6.3.13.1}
\emph{Pour tout morphisme affine plat de présentation finie
$f:X\rightarrow Y$ de $S$-schémas, et pour tout entier $n\geq 1$, le
morphisme}
\[
\Hom_{S}(\Sym^{n}_{Y}(X),G)\longrightarrow \Hom_{S}((X/Y)^{n},G)
\]
\emph{est\pageoriginale injectif et a pour image l'ensemble des
morphismes symétriques de $(X/Y)^{n}$ dans $G$.}

Cette condition est stable par changement de base.

Soit $f:T\rightarrow S$ un morphisme fini localement libre de rang
$n$. Soit $\pr_{i}$ la $i$\textsuperscript{ème} projection de $(T/S)^{n}$
dans $T$, et $t$ la section \sisi{(6.3.4.3)}{\eqhyperref[XVII.eq6.3.4.2]{6.3.4.2}} de $\Sym^{n}_{S}(T)$. Si $u\in
G(T)$, alors $\sum \pr^{\ast}_{i}(u)$ est un élément symétrique de
$G\sisi{}{(}(T/S)^{n}$) et définit donc
\[
u_{n}\in \sisi{}{G(}\Sym^{n}_{S}(T)\sisi{}{)}{\quoi}.
\]
On pose
\[
\Tr_{f}(u)=t^{\ast}(u_{n}){\quoi}.
\]

Lorsque $f$ est fini localement libre, on définit $\Tr_{f}(u)$,
localement sur $S$, par la formule précédente. Cette construction est
compatible à tout changement de base et définit donc un
\emph{morphisme trace}
$$\etiquette[6.3.13.2]{XVII.eq6.3.13.2}
\Tr_{f}:f_{\ast}f^{\ast}G\longrightarrow G{\quoi}.\leqno(6.3.13.2)
$$

\subsubsection{}\etiquette[6.3.14]{XVII.subsubsec6.3.14}
La condition (\hyperref[XVII.par6.3.13.1]{6.3.13.1}) est vérifiée dans les cas suivants
\begin{enumerate}
\renewcommand{\theenumi}{\alph{enumi}}
\renewcommand{\labelenumi}{\theenumi)}
\item Si $G$ est représentable, par définition des puissances
symétriques.

\item Pour $G$ l'image réciproque sur $S_{\text{\rm fppf}}$ d'un
faisceau sur le petit site étale de $S$ (\hyperref[XVII.subsec0.10]{0.10}).

\item Pour $G$ défini par un $\Oo_{S}$-module quasi-cohérent $G_{0}$.
\end{enumerate}

L'énoncé c) se ramène à (\hyperref[XVII.par5.5.2.7]{5.5.2.7}).

\begin{enonce*}{Proposition 6.3.15}\etiquette[6.3.15]{XVII.prop6.3.15}
Pour\pageoriginale $G$ vérifiant (\hyperref[XVII.par6.3.13.1]{6.3.13.1}), on a:
\begin{enumerate}
\renewcommand{\theenumi}{\roman{enumi}}
\renewcommand{\labelenumi}{(\theenumi)}
\item Le morphisme $\Tr_{f}$ est fonctoriel en $G$, donc additif.

\item Sa formation est compatible à tout changement de base
$S'\rightarrow S$.

\item Si $T=\coprod_{i}T_{i}$, de sorte que $f_{\ast}G\simeq \sum
f_{i\ast}G$, on a $\Tr_{f}=\sum \Tr_{f_{i}}$.

\item Si $T$ est de rang constant $n$ sur $S$, le morphisme composé
\[
G\longrightarrow f_{\ast}f^{\ast}G\xrightarrow{\Tr_{f}} G
\]
est la multiplication par $n$.

\item Le morphisme trace est compatible à la composition \emph{(\cf
\hyperref[XVII.thm6.2.3]{6.2.3} Var 3).}
\end{enumerate}
\end{enonce*}

Les assertions (i) et (ii) sont claires, et (iii) est conséquence
facile de \hyperref[XVII.prop6.3.5]{6.3.5}. Avec les notations
de \hyperref[XVII.subsubsec6.3.13]{6.3.13}, si $u\in G(S)$, alors
$u_{n}$ est $n$ fois l'image réciproque de $u$, de sorte que
$\Tr_{f}(f^{\ast}u)=t^{\ast}(u_{n})=nu$.

Pour prouver (v), on se ramène, par localisation étale sur $S$ et
utilisant (iii), à prouver que pour
\[
f=gh:S''\xrightarrow{h} S'\xrightarrow{g}S{\quoi},
\]
avec $\sisi{f}{g}$ et $\sisi{g}{h}$ finis localement libres de \emph{rangs constants} $n$
et $m$, le diagramme
$$\vcenter{\etiquette[6.3.15.1]{XVII.eq6.3.15.1}
\xymatrix{
G(S'')    \ar[rr]^{\Tr_h}\ar[rd]_{\Tr_f} &        &G(S') \ar[ld]^{\Tr_g}\\
                &  G(S)         &
}}\leqno{(6.3.15.1)}
$$
est commutatif. Pour $u\in G(S'')$, considérons le diagramme
$$\vcenter{\etiquette[6.3.15.2]{XVII.eq6.3.15.2}
\xymatrix@C=2.5cm@R=1.75cm{
\Sym^{n}_S(S')  \ar[rr]^{\Sym^{n}(t)} \ar[rd]^{\Tr_{\sisi{g}{h}}(u)_n} & & \Sym^{n}_S
\Sym^{m}_{S'} (S'')  \ar[dd] \ar[ld]^{(u_m)_n}\\
                                          & G &\\
          S   \ar@<3pt>[ru]^{\Tr_{\sisi{f}{g}}\Tr_{\sisi{g}{h}}(u)} \ar@<-3pt>[ru]_{\Tr_{\sisi{fg}{gh}}(u)}  \ar[rr]^{t}  \ar[uu]_{t}   &        & \Sym^{nm}_S(S'') \ar[lu]_{u_{nm}}.
}}\leqno{(6.3.15.2)}
$$
La\pageoriginale flèche verticale de droite est définie car $S\sisi{'}{''}$ est
plat sur $S$; chaque triangle est commutatif et (v) sera donc
conséquence de la commutativité du bord extérieur.

Traduisant en terme d'anneaux, c'est une conséquence du cas
particulier $B=E$ et $C=F$ de \hyperref[XVII.lem6.3.3]{6.3.3}.

\setcounter{subsubsection}{15}
\subsubsection{}\etiquette[6.3.16]{XVII.subsubsec6.3.16}
Soit maintenant $f:T\rightarrow S$ un morphisme \sisi{séparé }{}quasi-fini plat de
présentation finie. Si $ℱ$ est un faisceau fppf sur $T$, on définit
son image directe à support propre $f_{!}\sisi{T}{ℱ}$ sur $S$ comme le
préfaisceau qui à $S'/S$ associe l'ensemble des sections de $\sisi{F}{ℱ}$ sur
$T'$ dont le support est propre sur $S'$. Ce préfaisceau est un
faisceau d'après \hyperref[XVII.prop6.1.1]{6.1.1}. Avec la notation de
\hyperref[XVII.prop6.1.1]{6.1.1}, on peut aussi
l'exprimer comme limite inductive locale (\sisi{IV 14.\quad}{\hyperrefext[V][V.subsec8.2]{8.2}}).
\[
f_{!}\sisi{F}{ℱ(U)}=\varinjlim_{V\in\Phi (U)}f_{\ast}(\sisi{F}{ℱ}|V){\quoi},
\]
formule\pageoriginale qui rend compte du fait assez évident que les
propriétés de $f_{!}$ se ramène à celles de $f'_{\ast}$ pour $f'$
fini. On trouve ainsi:

\begin{enonce*}{Proposition 6.3.17}\etiquette[6.3.17]{XVII.prop6.3.17}
On peut d'une et d'une seule façon définir, pour tout morphisme
$f:T\rightarrow S$ \sisi{}{séparé }quasi-fini plat de présentation finie et pour tout
faisceau fppf $G$ sur $S$ vérifiant (\hyperref[XVII.par6.3.13.1]{6.3.13.1}), un morphisme trace
$$
\Tr_{f}:f_{!}f^{\ast}G\longrightarrow G
$$
vérifiant:
\begin{enumerate}
\renewcommand{\theenumi}{\roman{enumi}}
\renewcommand{\labelenumi}{(\theenumi)}
\item Pour $f$ fini on retrouve le morphisme \eqhyperref[XVII.eq6.3.13.2]{6.3.13.2}.

\item Les conditions \emph{(Var 1)} à \emph{(Var 4)} de \hyperref[XVII.thm6.2.3]{6.2.3} sont
remplies, mutadis mutandis.
\end{enumerate}
\end{enonce*}

Le morphisme quasi-fini $f$ induit un diagramme commutatif (\cf
(\hyperref[XVII.subsec0.11]{0.11})) de sites (dans lequel $\alpha_{\ast}$ est le foncteur « restriction »)
$$
\xymatrix{
T_{\text{\rm fppf}} \ar[rr]^{\alpha} \ar[d]^{f}& & T_{\text{et}} \ar[d]^{f_{\et}}\\
S_{\text{\rm fppf}} \ar[rr]^{\alpha}& & S_{\et}.
}
$$
On a, pour $G$ sur $T_{\text{\rm fppf}}$,
$$
f_{!}\alpha_{\ast}\simeq \alpha_{\ast}f_{!}{\quoi}.
$$
On\pageoriginale dispose aussi d'un morphisme « de changement de
base »
$$
c:f^{\ast}_{\text{\rm et}}\alpha_{\ast}\longrightarrow
\alpha_{\ast}f^{\ast}{\quoi}.
$$
On laisse au lecteur le soin de vérifier la compatibilité suivante
entre les morphisme traces \hyperref[XVII.thm6.2.3]{6.2.3} et
\hyperref[XVII.prop6.3.17]{6.3.17}. Pour $G$ sur $S_{\text{\rm
fppf}}$ vérifiant (\hyperref[XVII.par6.3.13.1]{6.3.13.1}), le diagramme suivant est commutatif:
$$\label{XVII.eq6.3.17.1}
\vcenter{\xymatrix{
f_{\et!}f^{\ast}_{\et}\alpha_{\ast} G \ar[r]^{f_{!}c} \ar[d]_{\Tr_f}^{(6.2.3)}&
f_{\et!}\alpha_{\ast}f^{\ast}G  \ar@{<->}[r]^{\sim}&
\alpha_{\ast}f_!f^{\ast}G \ar[d]_{\alpha_{\ast}\Tr_f}^{(6.3.17)}\\
\alpha_{\ast}G                   \ar@{=}[rr] &        & \alpha_{\ast}G.
}}\leqno(6.3.17.1)
$$

En termes plus concrets, ceci signifie que si $G$ est un faisceau fppf
sur $S$, qui par restriction au site étale définit un faisceau étale
$\alpha_{\ast}G$, et si $u$ est une section à support propre  de
l'image réciproque $f^{\ast}_{\text{\rm et}}(\alpha_{\ast}G)$, d'image
$u_{1}$ dans $f^{\ast}G$, les éléments $\Tr_{f}(u)$ (au sens \hyperref[XVII.thm6.2.3]{6.2.3}) et
$\Tr_{f}(u_{1})$ (au sens 6.2.17) coïncident dans $G(S)$.

\begin{enonce*}[remark]{Exemple 6.3.18}\etiquette[6.3.18]{XVII.exem6.3.18}
Soit $f:S'\rightarrow S$ un morphisme fini localement libre de rang
$n$ et $u\in \Gamma(S',\Oo_{S'})=\Hom_{S}(S',\EE^{1}_{S})$.\nde{ici $\EE^1_S$ désigne la droite affine sur $S$.} Le morphisme $u$
définit un morphisme
\[
\Sym(u):\Sym^{n}_{S}(S')\longrightarrow
\Sym^{n}_{S}(\EE^{1}_{S})=\Spec(S[\sigma_{1}\ldots
\sigma_{n}]\sisi{)}{\quoi},
\]
où les $\sigma_{i}$ sont des fonctions symétriques élémentaires (\cf
\textsc{Bourbaki}, \emph{Alg.} 5, App. I nᵒ 2\nde{référence à vérifier. Probablement IV \S~6 nᵒ 1 dans la nouvelle version.}). On a alors (par un
calcul facile laissé au lecteur):
$$\etiquette[6.3.18.1]{XVII.eq6.3.18.1}
[\Sym(u)\circ t]^{\ast}(\sigma_{i})=\sigma_{i}(u){\quoi},\leqno(6.3.18.1)
$$
où\pageoriginale $t\in \Gamma$ $(\Sym^{n}_{S}(S')/S)$ est l'élément
canonique \sisi{(6.3.4.3)}{\eqhyperref[XVII.eq6.3.4.2]{6.3.4.2}}, et où on désigne par $\sigma_{i}(u)$ le
$i$\textsuperscript{ième} coefficient du polynôme caractéristique de la
multiplication par $u$ dans $f_{\ast}\Oo_{S'}$. On en déduit que les
applications « trace »
\[
\begin{cases}%raghu
\GG_{a}(S')\longrightarrow \GG_{a}(S)\\
\GG_{m}(S')\longrightarrow \GG_{m}(S)
\end{cases}
\]
sont respectivement la \emph{trace} et la \emph{norme} usuel\sisi{}{le}s; en
effet, le morphisme $\Sym^{n}_{S}(\EE^{1}_{S})\rightarrow
\EE^{1}_{S}$ déduit de la loi additive (\resp multiplicative)
sur $\EE^{1}_{S}$ est donné par la fonction $\sigma_{1}$ (\resp
$\sigma_{n}$) sur $\Sym^{n}_{S}(\EE^{1}_{S})\simeq \Spec
(S[\sigma_{1},\ldots,\sigma_{n}])$.

\sisi{\noindent}{}%
On déduit alors de (6.3.17.1) que pour $n$ inversible sur $S$ le
diagramme de faisceaux étales
$$\vcenter{\etiquette[6.3.18.2]{XVII.eq6.3.18.2}
\xymatrix{
0 \ar[r]& f_{\ast}\mu_n \ar[r] \ar[d]_{\Tr_f}^{(6.2.3)}& f_{\ast}\sisi{G}{\GG}_m
\ar[r] \ar[d]^{N_f}
& f_{\ast}\sisi{G}{\GG}_m \ar[d]^{N_f} \ar[r]& 0\\
0 \ar[r]& \mu_n \ar[r]        & \relax\sisi{G}{\GG}_m    \ar[r]     & \relax\sisi{G}{\GG}_m  \ar[r]       & 0
}}\leqno{(6.3.18.2)}
$$
est commutatif.
\end{enonce*}

\begin{enonce*}{Proposition 6.3.19}\etiquette[6.3.19]{XVII.prop6.3.19}
Soit un diagramme commutatif
\[
\xymatrix{
X \ar[rr]^{u} \ar[rd]_{p}&   & Y\ar[ld]^{q}\\
  & S &
}
\]
dans\pageoriginale lequel $X$ et $Y$ sont des courbes lisses
quasi-projectives sur $S$ et $u$ un morphisme quasi-fini.
\begin{enumerate}
\renewcommand{\theenumi}{\roman{enumi}}
\renewcommand{\labelenumi}{(\theenumi)}
\item Si $D$ est un diviseur sur $X$ fini sur $S$ (voir note
\emph{p.~\pageref{XVII.note192}}) et $g$ une section de $G_{Y}$ sur $Y$, on a
\[
\Tr_{D/S}(u^{\ast}g)=\Tr_{(u_{\ast}D)/S}(g){\quoi}.
\]

\item Si $E$ est un diviseur sur $Y$ fini sur $S$, $g$ une section de
$G_{X}$ sur $X$ et si $u$ est fini, on a
\[
\Tr_{E/S}(\Tr_{X/Y}(g))=\Tr_{(u\sisi{_{\ast}}{^\ast}E)/S}(\sisi{G}{g}){\quoi}.
\]
\end{enumerate}
\end{enonce*}

\begin{enonce*}[remark]{Preuve}
\begin{enumerate}
\renewcommand{\theenumi}{\roman{enumi}}
\renewcommand{\labelenumi}{(\theenumi)}
\item Soit $g_{n}$ la section de $G$ sur $\Sym^{n}_{S}(Y)$ déduite de
$g$ (\hyperref[XVII.subsubsec6.3.13]{6.3.13}). D'après \hyperref[XVII.prop6.3.12]{6.3.12}, on a
\begin{align*}
\Tr_{D/S}(u^{\ast}g)\overset{\dfn}{=}
\sigma(D)^{\ast}(u^{\ast}g)_{n} & =\sigma(D)^{\ast}(\sisi{u}{\Sym^n_S(u)}^{\ast}g_{n})=\sisi{u}{(\Sym^n_S(u)}\circ
\sigma(D))^\ast
(g_{n})\\
& =\sigma(u_{\ast}D)^\ast(g_{n})=\Tr_{u_{\ast}D}(g_{n}){\quoi}.
\end{align*}

\item La formule (ii) résulte de \hyperref[XVII.prop6.3.15]{6.3.15} (ii) et (v):
\[
\Tr_{E/S}(\Tr_{X/Y}(g))=\Tr_{E/S}(\Tr_{u^{\ast}E/E}(g))=\Tr_{u^{\ast}E/S}(g){\quoi}.
\]
\end{enumerate}
\end{enonce*}

\subsection*{Application 3. Trace d'un torseur}

\setcounter{subsubsection}{19}
\subsubsection{}\etiquette[6.3.20]{XVII.subsubsec6.3.20}
Soit $G$ un faisceau sur le grand site fppf de $S$ qui vérifie
(\hyperref[XVII.par6.3.13.1]{6.3.13.1}). Soit $u:X\rightarrow Y$ un morphisme fini localement libre
de $S$-schémas et $K$ un $G_{X}$-torseur sur $X$. Supposons que,
localement fppf sur $Y$, le torseur $K$ sur $X$ soit trivial. Tel est
le cas si $K$ est trivial localement pour la topologie étale.

\begin{enonce*}{Proposition 6.3.21}\etiquette[6.3.21]{XVII.prop6.3.21}
Sous\pageoriginale les hypothèses \hyperref[XVII.subsubsec6.3.20]{6.3.20}, il existe à isomorphisme
unique près un et un seul torseur $\Tr_{u}(K)$ sous $G_{Y}$ sur $Y$,
muni d'un morphisme $\Tr_{u}$ ou $\Tr^{K}_{u}:u_{\ast}K\rightarrow
\Tr_{u}(K)$, qui, pour $g$ section locale (au sens fppf) de $G$ sur
$Y$ et $k$ section locale de $u_{\ast}K$, vérifie
$$
\Tr^{K}_{u}(gk)=\Tr_{u}(g)\cdot\Tr_{u}\sisi{}{^K}(k){\quoi}.
$$
\end{enonce*}

Le problème est local sur $Y$; on peut donc supposer $K$ trivial:
$K=G_{X}$. Une solution du problème est de prendre $K=G_{Y}$ et pour
morphisme trace le morphisme \eqhyperref[XVII.eq6.3.13.2]{6.3.13.2}. On vérifie aussitôt que c'est
la seule.

Le foncteur trace \sisi{\hyperref[XVII.subsubsec6.3.20]{6.3.20}}{\hyperref[XVII.prop6.3.21]{6.3.21}} est défini en particulier
lorsque $G$ vérifie la condition suivante

\setcounter{subsubsection}{21}
\paragraph{Pour tout $S$-schéma $X$, tout torseur sous $G_{X}$ sur $X$
est localement trivial pour la topologie étale}\etiquette[6.3.21.1]{XVII.par6.3.21.1}

Cette condition est vérifiée pour $G$ un groupe lisse, pour $G$ image
réciproque d'un faisceau sur le petit site étale de $S$, ou pour $G$
défini par un faisceau quasi-cohérent sur $S$.

On a par construction
$$\etiquette[6.3.21.2]{XVII.eq6.3.21.2}
\Tr_{u}(K)=u_{\ast}(K) \times_{u_{\ast}G_{X}}G_{Y}{\quoi}.\leqno(6.3.21.2)
$$

\subsubsection{}\etiquette[6.3.22]{XVII.subsubsec6.3.22}
Pour définir un foncteur trace sous des hypothèses plus générales, le
point clef est, pour $f:X\rightarrow T$ un morphisme quasi-projectif
plat de présentation finie de $S$-schémas, de définir un foncteur
raisonnable $K\mapsto K_{n}$ des $G_{X}$-torseurs sur $X$ dans les
$G_{Z}$-torseurs sur $Z=\Sym^{n}_{T}(X)$. On\pageoriginale a tout
d'abord, $f_{n}$ désignant la projection de $\Sym^{n}_{T}(X)$ sur $T$,
le lemme suivant, qui se démontre comme \hyperref[XVII.prop6.3.21]{6.3.21}:

\begin{enonce*}{Lemme 6.3.22.1}\etiquette[6.3.22.1]{XVII.lem6.3.22.1}
Avec les notations précédentes et si, localement sur $T$, le torseur
$K$ est trivial, alors il existe, à isomorphisme unique près, un et un
seul torseur $K_{n}$ sur $\Sym^{n}_{T}(X)$, muni d'un morphisme
$\alpha:f_{\ast}K\rightarrow f_{n{\ast}}(K_{n})$, qui pour $g$
section locale \emph{(fppf)} de $f_{\ast}G_{X}$ et $k$ section locale
de $\sisi{u}{f}_{\ast}K$ vérifie $\alpha(gk)=g_{n}\sisi{\circ}{\cdot} \alpha(k)$ (notation de
\hyperref[XVII.subsubsec6.3.13]{6.3.13} pour $g_{n}$; la définition de $g_{n}$ ici utilise
l'hypothèse de platitude \etc).
\end{enonce*}

Le cas général se traitera par une descente un peu canulée. Tout
d'abord:

\begin{enonce*}{Proposition 6.3.23}\etiquette[6.3.23]{XVII.prop6.3.23}
Soit un diagramme commutatif
$$\vcenter{\etiquette[6.3.23.1]{XVII.eq6.3.23.1}
\xymatrix{
X \ar[rr]^{u} \ar[rd]& & Y\ar[ld]\\
  & T &
}}\leqno{(6.3.23.1)}
$$
de $T$-schémas plats de présentation finie et quasi-projectifs, avec
$u$ étale surjectif. Alors, le système semi-simplicial des schémas
$\Sym^{n}_{T}((X/Y)^{k+1})$ est un hyperrecouvrement pour la topologie
étale du schéma $\Sym^{n}_{T}(Y)$.
\end{enonce*}

\begin{enonce*}[remark]{Preuve}
\begin{enumerate}
\renewcommand{\theenumi}{\alph{enumi}}
\renewcommand{\labelenumi}{\theenumi)}
\item On aura à utiliser le résultat auxiliaire.

\setcounter{subsubsection}{23}\setcounter{paragraph}{1}
\paragraph{}\etiquette[6.3.23.2]{XVII.par6.3.23.2}
\emph{Soient $(X,Y,T)$ comme plus haut, $t\in T$, $Z$ un
$T$-schéma somme d'un nombre fini de $T$-schémas $\Spec(k_{i})$, pour
$k_{i}$ extension finie de $k_{t}$,\nde{ici $k_t$ désigne le corps résiduel de $T$ en $t$.} et\pageoriginale un diagramme
commutatif de $T$-schémas}
$$\vcenter{\etiquette[6.3.23.3]{XVII.eq6.3.23.3}
\xymatrix{
              & Z \ar@{^{(}->}[ld]_{x} \ar@{^{(}->}[rd]^{y}     \\
X_t \ar[rr]^{u}&    & Y_t,
}}\leqno{(6.3.23.3)}
$$
\emph{tel que $x$ et $y$ soient des immersions fermées. Soit $z\in
\Sym^{n}_{T}(Z)$, et soient $x(\sisi{Z}{z})$ et $y(\sisi{Z}{z})$ ses images dans
$\Sym^{n}_{T}(X)$ et $\Sym^{n}_{T}(Y)$. Alors, $\Sym^{n}_{T}(u)$ est
étale en $x(\sisi{Z}{z})$.}

La formation du schéma $\Sym^{n}_{T}(X)$ est compatible à tout
changement de base $\sisi{S'}{T'}\rightarrow \sisi{S}{T}$ (\hyperref[XVII.par5.5.2.7]{5.5.2.7}). Par réduction au cas
noethérien, on en déduit que $\Sym^{n}_{T}(X)$ est de présentation
finie sur $\sisi{S}{T}$; il est par ailleurs plat sur $\sisi{S}{T}$ (\hyperref[XVII.par5.5.2.4]{5.5.2.4}). Grâce au
critère de platitude fibre par fibre, on se ramène pour prouver
(\hyperref[XVII.par6.3.23.2]{6.3.23.2}) au cas $\sisi{S}{T}$ spectre d'un corps. Le complété de
$\Sym^{n}_{T}(X)$ en $\sisi{Z}{z}$ a alors pour algèbre affine la limite
projective des algèbres affines localisées en $x(z)$ des schémas
$\Sym^{n}_{T}(Z_{k})$, pour $Z_{k}$ voisinage infinitésimal de $Z$
dans $X$. Par hypothèse, $u$ induit des isomorphismes entre les
voisinages infinitésimaux de $Z$ dans $X$ et dans $Y$. Le complété de
$\Sym^{n}_{T}(u)$ en $x(z)$ et $y(z)$ est donc un isomorphisme, et
ceci prouve (\hyperref[XVII.par6.3.23.2]{6.3.23.2}).

On déduit de \hyperref[XVII.par6.3.23.2]{6.3.23.2} que le schéma somme, pour $z\in
\Sym^{n}_{T}(Z)$, des hensélisés en $y(z)$ de $\Sym^{n}_{T}(Y)$ ne
dépend que de $Z$ et du $T$-schéma $X^{h}_{Z}$ somme des hensélisés de
$X$ en les divers points de $Z$. On désignera ce schéma par la
notation $\Sym^{h\,n}_{T}(X^{h}_{Z})$; il est fonctoriel en
$X^{h}_{Z}$.

\item Par\pageoriginale passage à la limite, pour prouver \hyperref[XVII.prop6.3.23]{6.3.23}, il
suffit de montrer que pour $y\in \Sym^{n}_{T}(Y)$, fermé dans sa
fibre, l'image réciproque sur l'hensélisé de $\Sym^{n}_{T}(Y)$ en y du
schéma simplicial $\Sym^{n}_{T}((X/Y)^{k+1})$ est un hyperrecouvrement
pour la topologie étale. Ce problème est local sur $T$ pour la
topologie étale; on se ramène par là au cas où il existe un diagramme
\eqhyperref[XVII.eq6.3.23.3]{6.3.23.3} avec $y=x(z)$ (en tuant les extensions résiduelles par une
extension séparable de $k_{t}$). Le système semi-simplicial considéré
se déduit alors du système
$$
\Sym^{h\,n}_{T}((X^{h}_{u^{-1}(Z)}/Y^{h}_{Z})^{k+1}),\quad\mbox{\rm
augmenté vers}\quad \Sym^{h\,n}_{T}(Y^{h}_{Z})
$$
par changement de base, de sorte qu'il suffit de considérer ce
dernier. Le système semi-simplicial des
$(X^{h}_{u^{-1}(Z)}/Y^{h}_{Z})^{k+1}$ est semi-simplicialement
homotope au système constant $Y^{h}_{Z}$, car
$X^{h}_{u^{-1}(Z)}/Y^{h}_{Z}$ admet une section. Le système
$\Sym^{h\,n}_{T}((X^{h}_{u^{-1}(Z)}/Y^{h}_{Z})^{h+1})$ est donc de même
homotope à $\Sym^{h\, n}_{T} (Y^{h}_{Z})$, et en particulier en est
un hyperrecouvrement. Ceci prouve \sisi{6.3.2.3}{\hyperref[XVII.prop6.3.23]{6.3.23}}.
\end{enumerate}
\end{enonce*}

\subsubsection{}\etiquette[6.3.24]{XVII.subsubsec6.3.24}
Soient $f:Y\rightarrow T$ un morphisme quasi-projectif plat de
présentation finie de $S$-schémas, et $K$ un $G_{Y}$-torseur sur
$Y$. Considérons le problème suivant:

\setcounter{paragraph}{-1}
\paragraph{}\etiquette[6.3.24.0]{XVII.par6.3.24.0}
Définir, pour tout morphisme $u:X\rightarrow Y$, avec $X$
quasi-projectif plat de présentation finie sur $T$, un torseur
$(u^{\ast}K)_{n}$ sur $\Sym^{n}_{T}(X)$, et pour tout $v:X'\rightarrow
X$ un isomorphisme $c_{v}:(v^{\ast}u^{\ast}K)_{n}\simeq
\Sym^{n}_{T}(v)^{\ast}\sisi{K}{(u^K)}_{n}$, de sorte que
\begin{enumerate}
\renewcommand{\theenumi}{\roman{enumi}}
\renewcommand{\labelenumi}{(\theenumi)}
\item les\pageoriginale isomorphismes $c_{v}$ sont compatibles à la
  composition;

\item pour $u^{\ast}K$ localement trivial sur $T$, on retrouve le
torseur de \hyperref[XVII.lem6.3.22.1]{6.3.22.1}, et sa variance.
\end{enumerate}

D'après \hyperref[XVII.prop6.3.23]{6.3.23}, la question d'existence et d'unicité d'une solution à
\hyperref[XVII.par6.3.24.0]{6.3.24.0}, pour tout torseur déduit de $K$ par une image réciproque du
type envisagé en \hyperref[XVII.par6.3.24.0]{6.3.24.0}, est une question locale pour la topologie
étale sur $Y$. En particulier, si $G$ vérifie (\hyperref[XVII.par6.3.21.1]{6.3.21.1}), le problème
\hyperref[XVII.par6.3.24.0]{6.3.24.0} admet une et une seule solution.

Supposons maintenant que $G$ vérifie:

\paragraph*{(6.3.24.1)}\label{XVII.par6.3.24.1}
\emph{Pour tout entier $n$ et pour tout diagramme de $S$-schémas
affines
\[
\xymatrix{
X \ar[rr]^{u} \ar[rd]_{f}& & Y\ar[ld]^{g}\\
  &T&,
}
\]
avec $f$ et $g$ plats de présentation finie et $u$ fini
localement libre surjectif, le schéma semi-simplicial tronqué
$(\Sym^{n}_{T}((X/Y)^{k+1}))_{0\leq k\leq 2}$, augmenté vers
$\Sym^{n}_{T}(Y)$, est de descente \emph{(\textsc{Giraud} \cite{SGA4-17:Gi})} pour la
catégorie des $G$-torseurs.}

\paragraph*{6.3.24.1 bis}\etiquette[6.3.24.1 bis]{XVII.par6.3.24.1bis}
On déduit de façon standard de \hyperref[XVII.prop6.3.23]{6.3.23} que la
condition  (\ref{XVII.par6.3.24.1})
implique la condition analogue, où on suppose seulement $f$ et $g$
plats de présentation finie quasi-projectifs et $u$ plat
surjectif. Localement pour la topologie étale, $u$ admet en effet des
quasi-sections finies localement libres et on conclut par \hyperref[XVII.prop6.3.23]{6.3.23} et
\textsc{Giraud} \cite{SGA4-17:Gi} appliqué à un diagramme formant carré de données de descente.

Etant\pageoriginale donné un $G_{Y}$-torseur $K$ sur un $T$-schéma $Y$
quasi-projectif plat de présentation finie, il existe par définition
un morphisme fidèlement plat de présentation finie $u:X\rightarrow Y$
tel que $u^{\ast}K$ soit trivial.\nde{on peut supposer $T$ affine. Alors il existe $u$ comme énoncé qu'on peut même supposer affine pour pouvoir ensuite appliquer l'alinéa précédent.} Si $u_{i}$ désigne la projection de
$(X/Y)^{i+1}$ sur $Y$ $(i\geq 0)$, les torseurs $(u^{\ast}_{i}K)_{n}$
(\hyperref[XVII.lem6.3.22.1]{6.3.22.1}) sur les schémas $\Sym^{n}_{T}((X/Y)^{i+1})$ forment une
donnée de descente $(i=0,1,2)$ sur ce schéma semi-simplicial augmenté
vers $\Sym^{n}_{T}(Y)$. Supposons que $G$ vérifie la condition:

\paragraph*{(6.3.24.2)}\label{XVII.par6.3.24.2}
\emph{Les données de descentes construites plus haut sont toutes
effectives} (\cf \ref{XVII.par6.3.24.1}).

Alors, on désignera par $K_{n}$ le torseur sur $\Sym^{n}_{T}(Y)$
(déterminé à isomorphisme unique près) défini par cette donnée de
descente.

\subsubsection{}\etiquette[6.3.25]{XVII.subsubsec6.3.25}
Supposons que $G$ vérifie (\hyperref[XVII.par6.3.21.1]{6.3.21.1}) ou les conditions
(\ref{XVII.par6.3.24.1}) (\ref{XVII.par6.3.24.2}). On laisse au
lecteur le soin de donner un sens au formulaire suivant et de le vérifier:

\paragraph{}\etiquette[6.3.25.1]{XVII.par6.3.25.1}
\emph{Changement de base}. Soit un carré commutatif
\[
\xymatrix{
X' \ar[rr]^{g'} \ar[d]_{f'}& & X \ar[d]^{f}\\
T' \ar[rr]_{g}& & T.
}
\]
Désignons par $g'_{n}$ la flèche de $\Sym^{n}_{T'}(X')$ dans
$\Sym^{n}_{T}(X)$. Pour $K$ torseur sous $G$ sur $X$, on a
\[
{g'_n}^{\ast} K_{n}\simeq ({g'}^{\ast}K)_{n}{\quoi}.
\]

Soit\pageoriginale maintenant $f:X\rightarrow T$ et $K$ un
$G_{X}$-torseur sur $X$.

\paragraph{}\etiquette[6.3.25.2]{XVII.par6.3.25.2}
\emph{Changement de groupe.} Soit $r:G_{1}\rightarrow G_{2}$ un
morphisme entre groupes vérifiant tous deux (\hyperref[XVII.par6.3.21.1]{6.3.21.1})
ou (\ref{XVII.par6.3.24.1}), (\ref{XVII.par6.3.24.2}). On a
\[
r(K_{n})\simeq r(K)_{n}{\quoi}.
\]

\paragraph{}\etiquette[6.3.25.3]{XVII.par6.3.25.3}
\emph{Additivité en $K$}\sisi{:}{.} On a
\[
(K'+K'')_{n}\simeq K'_{n}+K''_{n}{\quoi}.
\]

\paragraph{}\etiquette[6.3.25.4]{XVII.par6.3.25.4}
\emph{Additivité en $n$}. Soient $n_{1}$ et $n_{2}$ deux entiers,
$n=n_{1}+n_{2}$ et
\[
\alpha:\fprod{\Sym_{T}{}^{n_{1}}(X)}{\Sym_{T}{}^{n_{2}}(X)}{T}\longrightarrow
\Sym^{n}_{T}(X)
\]
la flèche canonique; on a alors
\[
\alpha^{\ast}K_{n}\simeq \pr^{\ast}_{1}
K_{n_{1}}+\pr^{\ast}_{2}K_{n_{2}}{\quoi}.
\]

\paragraph{}\etiquette[6.3.25.5]{XVII.par6.3.25.5}
\emph{Multiplicativité en $n$.} Soient $n$ et $m$ deux entiers et
$\alpha$ la flèche canonique de $\Sym^{n}_{T}(\Sym^{m}_{T}(X))$ dans
$\Sym^{nm}_{T}(X)$ (\cf \hyperref[XVII.lem6.3.3]{6.3.3}). On a
\[
\alpha^{\ast}(K_{nm})\simeq (K_{m})_{n}{\quoi}.
\]

\subsubsection{}\etiquette[6.3.26]{XVII.subsubsec6.3.26}
Supposons que $G$ vérifie (\hyperref[XVII.par6.3.21.1]{6.3.21.1}) ou
(\ref{XVII.par6.3.24.1}),
(\ref{XVII.par6.3.24.2}). Si
$f:X\rightarrow T$ est un morphisme fini localement libre de rang $n$
entre $S$-schémas, si $t$ est la section (6.3.4.3) de
$\Sym^{n}_{T}(X)$, et si $K$ est un $G_{X}$-torseur sur $X$, on
définit un $G_{T}$-torseur $\Tr_{f}(K)$ sur $T$ par la formule
$$\etiquette[6.3.26.1]{XVII.eq6.3.26.1}
\Tr_{f}(K)=t^{\ast}(K_{n}){\quoi}.\leqno(6.3.26.1)
$$

Si\pageoriginale $f$ est fini localement libre, on définit $\Tr_{f}(K)$,
localement sur $T$, par la même formule. Cette définition, via un
isomorphisme canonique et fonctoriel, est équivalente à la définition
\hyperref[XVII.prop6.3.21]{6.3.21} dans leur domaine commun de validité
(la définition \hyperref[XVII.prop6.3.21]{6.3.21}
fournissant une solution au problème de descente posé). On déduit
aussitôt de (\hyperref[XVII.par6.3.25.1]{6.3.25.1}), (\hyperref[XVII.par6.3.25.2]{6.3.25.2})
et (\hyperref[XVII.par6.3.25.3]{6.3.25.3}) que le foncteur
$\Tr_{f}$ est \emph{compatible aux changements de base $T'\rightarrow
T$, compatible aux changements de groupe $r:G_{1}\rightarrow G_{2}$,
compatible à l'addition des torseurs.}

La commutativité du bord extérieur du diagramme \eqhyperref[XVII.eq6.3.15.2]{6.3.15.2}, et
(\hyperref[XVII.par6.3.25.1]{6.3.25.1}) (\hyperref[XVII.par6.3.25.4]{6.3.25.4}) fournissent
une \emph{compatibilité à la composition des morphismes}:
$$
\Tr_{fg}(K)\simeq \Tr_{f}(\Tr_{g}(K)\sisi{}{)}{\quoi}.
$$

\subsubsection{}\etiquette[6.3.27]{XVII.subsubsec6.3.27}
Dans le cas particulier des courbes relatives, on a de plus:

\begin{enonce*}{Construction 6.3.27.1}\etiquette[6.3.27.1]{XVII.const6.3.27.1}
\emph{Si} $f:X\rightarrow S$ est une courbe lisse et quasi-projective
sur $S$, $D_{1}$ et $D_{2}$ deux diviseurs relatifs finis sur $S$ et
$K$ un $G$-torseur sur $X$, on a, posant
$\Tr_{D/S}(K)=\Tr_{D/S}(K|D)$, un isomorphisme canonique
$$
\Tr_{D_{1}+D_{2}}(K)\simeq \Tr_{D_{1}}(K)+\Tr_{D_{2}}(K){\quoi}.
$$
\end{enonce*}

\sisi{\noindent}{}%
Cet isomorphisme se déduit aussitôt de l'isomorphisme (\hyperref[XVII.par6.3.25.4]{6.3.25.4}) et de
\hyperref[XVII.lem6.3.9.1]{6.3.9.1}. Enfin les arguments de \hyperref[XVII.prop6.3.19]{6.3.19} fournissent encore:

\begin{enonce*}{Construction 6.3.27.2}\etiquette[6.3.27.2]{XVII.const6.3.27.2}
Soit un diagramme commutatif
$$
\xymatrix{
X   \ar[rr]^{u} \ar[rd]_{p}&     &  Y \ar[ld]^{q}\\
     &   S &,
}
$$
dans\pageoriginale lequel $X$ et $Y$ sont des courbes lisses
quasi-projectives sur $S$ et $u$ un morphisme quasi-fini.
\begin{enumerate}
\renewcommand{\theenumi}{\roman{enumi}}
\renewcommand{\labelenumi}{(\theenumi)}
\item Si $D$ est un diviseur sur $X$ fini sur $S$ et $K$ un torseur
sous $G$ sur $Y$, on a \emph{(\cf note p.~\pageref{XVII.note192})}:
$$
\Tr_{D/S}(u^{\ast}K)\simeq \Tr_{(u_{\ast}D)/S}(K){\quoi}.
$$

\item Si $E$ est un diviseur sur $Y$ fini sur $S$, $K$ un torseur sous
$G$ sur $X$ et si $u$ est fini, on a
$$
\Tr_{E/S}(\Tr_{X/Y}(K))\simeq \Tr_{u^{\ast}E/S}(K){\quoi}.
$$
\end{enumerate}
\end{enonce*}

\begin{enonce*}{Proposition 6.3.28}\etiquette[6.3.28]{XVII.prop6.3.28}
Soient $A$ et $B$ deux $R$-algèbres (commutatives) et $u:A\rightarrow
B$ un homomorphisme qui fasse de $B$ un $A$-module localement libre de
type fini et fidèle. Alors le complexe \sisi{}{co-}semi-simplicial d'algèbres
$(TS^{n}_{R}(\bigotimes\limits^{k+1}_{A}B))_{k\geq 0}$
est, en tant que $TS^{n}_{R}(A)$-module différentiel \sisi{}{co-}semi-simplicial,
une résolution scindée de $TS^{n}_{R}(A)$.
\end{enonce*}

\begin{enonce*}[remark]{Preuve}
Le complexe de modules associé au module \sisi{}{co-}semi-simplicial
$TS^{n}_{R}(\bigotimes\limits^{k+1}_{A}B)$ ne dépend que
du $A$-module $B$, muni d'un homomorphisme de modules $u:A\rightarrow
B$. En tant que $A$-module, $B$ peut s'écrire $B=A\oplus L$, et ceci
permet d'expliciter un opérateur d'homotopie.
\end{enonce*}

Rappelons que:

\begin{enonce*}{Proposition 6.3.29}\etiquette[6.3.29]{XVII.prop6.3.29}
Soient $A$ un anneau, $B^{\ast}$ une $A$-algèbre co-semi-simpliciale,
et $c:A\rightarrow B^{\ast}$ le morphisme structural. Pour que le
foncteur image réciproque de modules ($=$ extension des scalaires)
$\epsilon^{\ast}$ soit fidèle (\resp pleinement fidèle) il suffit que
la suite de $A$-modules
$$
0\longrightarrow A\longrightarrow B^{0}
$$
(\resp\pageoriginale $0\rightarrow A\rightarrow
B^{0}\xrightarrow{j_{0}-j_{1}}B^{1}$)
soit exacte, et le reste après tensorisation par un quelconque $A$-module.
\end{enonce*}

La démonstration est triviale et laissée au lecteur.

On déduit de façon standard de (\hyperref[XVII.prop6.3.28]{6.3.28}) (\hyperref[XVII.prop6.3.29]{6.3.29}) le

\begin{enonce*}{Corollaire 6.3.30}\etiquette[6.3.30]{XVII.cor6.3.30}
La condition (\ref{XVII.par6.3.24.1}) est vérifiée par tout groupe affine $G$.
\end{enonce*}

(Utiliser qu'un torseur sous $G$ est représentable par un schéma
affine sur $T$.)

Il ne devrait pas être difficile d'étendre \hyperref[XVII.cor6.3.30]{6.3.30} à tout groupe $G$
représentable par un espace algébrique \footnote[1]{\label{XVII.noteArtin2}Algebraic Spaces
par M.~\textsc{Artin}. Yale University, Whittemore Lectures, May 1969.}. Je
conjecture que la condition (\ref{XVII.par6.3.24.2}) est vérifiée par tout groupe
représentable par un espace algébrique, et plat. On a l'énoncé
canularesque:

\begin{enonce*}{Proposition 6.3.31}\etiquette[6.3.31]{XVII.prop6.3.31}
Soit une suite exacte sur $S$
$$
0\rightarrow G\xrightarrow{r}G'\xrightarrow{u}H\rightarrow 0{\quoi}.
$$ Supposons que $G'$ et $H$, donc $G$ vérifient
(\hyperref[XVII.par6.3.13.1]{6.3.13.1}), que $G'$ vérifie (\hyperref[XVII.par6.3.21.1]{6.3.21.1}) et
que $G$ vérifie (\ref{XVII.par6.3.24.1}). Alors, $G$ vérifie
(\ref{XVII.par6.3.24.2}).
\end{enonce*}

Ceci s'applique par exemple à un groupe affine noyau d'un épimorphisme
de groupes représentables et lisses.

\begin{enonce*}{Preuve de 6.3.31}
\emph{Soit\pageoriginale $K$ un $G$-torseur sur
$X$. Par \hyperref[XVII.prop6.3.23]{6.3.23}, la question est locale sur $X$ pour la topologie
étale; il suffit donc de traiter le cas où $r(K)$ est trivial. Les
$G$-torseurs $K$ tels que $r(K)$ soit trivial s'identifient aux
$G'$-torseurs triviaux $K'$ munis d'un isomorphisme
$k:H\isomto u(K')$. Posons $Z=\Sym^{n}_{T}(X)$. Le morphisme
de torseurs triviaux $k$ induit $k_{n}:H_{Z}\isomto
u(K'_{n})$, et le $G$-torseur correspondant sur $\Sym^{n}\sisi{}{_T}(X)$
représente la donnée de descente considérée.}
\end{enonce*}

\section{Appendice}\etiquette[7]{XVII.sec7}\pageoriginale

\setcounter{subsection}{-1}
\subsection{Préliminaires}\etiquette[7.0]{XVII.subsec7.0}

Nous utiliserons dans cet appendice les résultats de \sisi{${\rm VI_{B}}$}{VBIS},
dont nous conservons les notations.

\subsubsection{}\etiquette[7.0.1]{XVII.subsubsec7.0.1}
Soit $\Sch_{!}$ la catégorie dont les objets sont les schémas
quasi-compacts et quasi-séparés, les morphismes étant les morphismes
de schémas séparés de type fini: on désigne par $\Ee$ la
catégorie bifibrée au-dessus de $\Sch_{!}$, à catégories fibres des
catégories opposées à des topos, définie par les faisceaux étales sur
les schémas \sisi{${\rm VI_{B}}$ (4.3.0)}{(\hyperrefext[Vbis][Vbis.subsubsec4.3.0]{4.3.0})}. On se fixe une fois pour toutes un
objet $R$ de $\Sch_{!}$ et on désigne par $\Ee_{R}$ la
catégorie bifibrée au-dessus de $\Sch_{!R}$ déduite de
$\Ee$. Soit enfin $\mathcal{A}$ un anneau du topos
$\sheafGamma(\Ee^{\circ}_{R})$ qui vérifie les conditions
suivantes:
\begin{enumerate}
\renewcommand{\theenumi}{\roman{enumi}}
\renewcommand{\labelenumi}{(\theenumi)}
\item $\mathcal{A}$ est une section \emph{cocartésienne} de
$\Ee^{\circ}_{R}$ au-dessus de $\Sch^{\circ}_{!R}$.

\item Pour tout schéma $X$ de $\Sch_{!R}$, $\mathcal{A}_{X}$ est un
faisceau de \emph{torsion}.
\end{enumerate}

Ceci posé, nous travaillerons avec la catégorie
$\sheaf{\Mod}(\sisi{}{\sheafGamma(}\Ee^{\circ}_{R}\sisi{}{)},\mathcal{A})$ (\refvb[VI B ]{Vbis.subsubsec1.3.5}).

On se propose de définir, pour tout morphisme $f:X\rightarrow Y$ de
$\Sch_{!R}$, un foncteur $Rf_{!}:D(X,\mathcal{A}_X)\rightarrow
D(Y,\mathcal{A}_{Y})$, coïncidant avec le foncteur précédemment défini
lorsque $f$ est compactifiable, et faisant des catégories
$D(X,\mathcal{A}_{X})$ pour $X$ variable une catégorie cofibrée
au-dessus de $\Sch_{!R}$.\nde{\label{XVII.ndeapp}Tout morphisme de $\Sch_{!}$ est compactifiable d'après le théorème de compactification de Nagata (voir \nderef{XVII.ndeConrad}). Cet appendice montre comment éviter l'usage du théorème de compactification dans la définition de $Rf_!$ en s'appuyant sur le lemme de Chow (\hyperref[XVII.lem7.0.6]{7.0.6}).}

\subsubsection{}\etiquette[7.0.2]{XVII.subsubsec7.0.2}
Bien que tous les résultats de cet appendice soient énoncés dans le
cadre explicité en \hyperref[XVII.subsubsec7.0.1]{7.0.1}, on pourrait les transcrire avec des
modifications évidentes dans le cadre suivant: on prend pour section
de $\sheafGamma(\Ee^{\circ}_{R})$ la section constante de
valeur le faisceau constant $\ZZ$ et on se limite pour tout
schéma $X$ à la catégorie dérivée $D^{\ast}_{\tors}(X)$ définie par
les complexes de faisceaux de groupes abéliens dont les faisceaux de
cohomologie sont de torsion (\refvb[VI B ]{Vbis.subsubsec4.3.1}).

\begin{enonce*}{Lemme 7.0.3}\etiquette[7.0.3]{XVII.lem7.0.3}
Soit\pageoriginale un diagramme commutatif de schémas
$$\vcenter{\etiquette[7.0.3.1]{XVII.eq7.0.3.1}
\xymatrix{
X' \ar[rr]^{j'} \ar[d]_{f'}& & X \ar[d]^{f}\\
Y' \ar[rr]_{j}& & Y.
}}\leqno{(7.0.3.1)}
$$
où $f$, $f'$ sont propres et $j$, $j'$ sont des immersions. Alors pour
tout complexe $K^{\bigdot}$ de $D^{+}_{\tors}(X')$, la flèche
canonique\nde{la définition de cette flèche est semblable à celle de \eqhyperref[XVII.eq5.1.5.2]{5.1.5.2} modulo un dévissage facile.}
$$
j_{!}\circ R^{+}f'\sisi{}{_*}(K^{\bigdot})\longrightarrow R^{+}f\sisi{}{_*}\circ
j_{!}(K^{\bigdot})
$$
est un isomorphisme.
\end{enonce*}

En vertu du « lemma on way-out functors »  (\cite{SGA4-17:RD} I \sisi{(}{}7.1), il
suffit de montrer que pour tout entier $i$ et pour tout faisceau $F$
de torsion, la flèche
$$\etiquette[7.0.3.2]{XVII.eq7.0.3.2}
j_{!}\circ R^{i}f'\sisi{}{_*}(F)\longrightarrow R^{i}f\sisi{}{_*}\circ j_{!}(F)\leqno(7.0.3.2)
$$
est un isomorphisme.

En prenant les images fermées de $j$ et $j'$ on est ramené au cas où
$j$ et $j'$ sont des immersions ouvertes dont les images sont partout
denses: la propreté de $f$ et de $f'$ entraîne alors que le diagramme
\eqhyperref[XVII.eq7.0.3.1]{7.0.3.1} est cartésien. Dans ces conditions, le théorème du
changement de base pour un morphisme propre, sous la forme \hyperrefext[XII][XII.cor5.2]{5.2},
montre que, pour tout point géométrique $\xi$ de $Y$, la flèche
$(j_{!}\circ R^{i}f'\sisi{}{_*}(F))_{\xi}\rightarrow (R^{i}f\sisi{}{_*}\circ
j_{!}(F))_{\xi}$ est un isomorphisme.

\begin{enonce*}{Construction 7.0.4}\etiquette[7.0.4]{XVII.const7.0.4}
\emph{Soient $D$ une catégorie, $X$ et $Y$ deux $D$-objets de
$\Sch_{!R}$ tels que pour toute flèche $\alpha\rightarrow \beta$ de
$D$, les morphismes $X_{\beta}\rightarrow X_{\alpha}$ et
$Y_{\beta}\rightarrow Y_{\alpha}$ soient propres. Soit $j:X\rightarrow
Y$ un morphisme tel que pour tout objet $\alpha$ de $D$,
$j_{\alpha}:X_{\alpha}\rightarrow Y_{\alpha}$ soit une immersion:
l'isomorphisme \eqhyperref[XVII.eq7.0.3.2]{7.0.3.2} pour $i=0$ vérifie les conditions de
compatibilités usuelles et permet de construire un foncteur
\emph{exact}:}
\[
j_{!}:\Mod(\sheafGamma(\bar{X}),\mathcal{A})\longrightarrow
\Mod(\sheafGamma(\bar{Y}),\mathcal{A})
\]
\emph{(\cf \sisi{(}{}\refvb{Vbis.subsubsec1.2.5} pour les notations)\sisi{.}{}
qui\pageoriginale induit un foncteur triangulé, encore noté $j_{!}$:}
\[
j_{!}:D(\sheafGamma(\bar{X}),\mathcal{A})\longrightarrow
D(\sheafGamma(\bar{Y},\mathcal{A}))
\]

\emph{Alors, pour tout objet $\alpha$ de $D$, on a un diagramme commutatif,
où $j_{\alpha !}$ désigne le foncteur « prolongement par $0$ »
usuel:}
\[
\xymatrix{
D(\sheafGamma(\bar{X}), \sisi{}{\Aa}) \ar[rr]^{j_!} \ar[d]_{R e^{\sisi{+}{*}}_{\alpha}}&
&D(\sheafGamma(\bar{Y}),\Aa) \ar[d]^{R e^{\sisi{+}{*}}_{\alpha}}\\
D(X_{\alpha}, \Aa_{X_{\alpha}})  \ar[rr]_{j_{\alpha !}} & &
  D(Y_{\alpha},\Aa_{Y_\alpha}).
}
\]
\end{enonce*}

\begin{enonce*}{Construction 7.0.5}\etiquette[7.0.5]{XVII.constru7.0.5}
Soit un diagramme commutatif dans $\Sch_{!R}$
$$\vcenter{\etiquette[7.0.5.1]{XVII.eq7.0.5.1}
\xymatrix{
X'   \ar[rr]^{i'} \ar[d]_{f'}&    & X \ar[d]^{f}\\
Y'   \ar[rr]^{i} &    & Y
}}\leqno{(7.0.5.1)}
$$
où $i$, $i'$ sont des immersions, $f$ un morphisme propre et $f'$
un morphisme propre \emph{surjectif}. Il définit un diagramme commutatif
d'objets simpliciaux de $\Sch_{!R}$
$$\vcenter{\etiquette[7.0.5.2]{XVII.eq7.0.5.2}
\xymatrix{
[X'|Y']      \ar[d]_{u_{f'}}  \ar[r]^{j} & [X|Y] \ar[d]^{u_f}\\
C^{\Delta}_{Y'}  \ar[r]_{i} & C^{\Delta}_{Y}
}}\leqno{(7.0.5.2)}
$$
où $j_{n}:[X'|Y']_{n}\rightarrow [X|Y]_{n}$ est une immersion
pour tout $n$. Avec ces notations, on a un isomorphisme canonique
\[
i_{!}\isomto  R^{+}\bar{u}_{f\sisi{+}{*}}\circ j_{!}\circ
L^{+}\bar{u}_{f'}\sisi{}{^*}{\quoi}.
\]
\end{enonce*}

En prenant les images fermées de $i$ et $i'$ on peut supposer que $i$
et $i'$ sont des immersions ouvertes dont les images sont partout
denses. Le diagramme \eqhyperref[XVII.eq7.0.5.1]{7.0.5.1} est alors cartésien et $f$ est
\emph{surjectif}. De plus, pour tout entier $n$, $j_{n}$ est
une\pageoriginale immersion ouverte et le diagramme
\[
\xymatrix{
[X'|Y']_n        \ar[r]^{j_n} \ar[d]& [X|Y]_n \ar[d]\\
Y'  \ar[r]_{i} & Y
}
\]
est cartésien.

Dès lors $L^{+}\bar{u}_{f}^{\sisi{+}{*}}\circ i_{!}$ est isomorphe à $j_{!}\circ
L^{+}\bar{u}_{f'}^{\ast}$: utilisant le fait que $u_{f}$ est une
augmentation de descente cohomologique (\sisi{${\rm VI_{B}}$ (\refnoexpose{Vbis.prop4.3.2})}{\refvb{Vbis.prop4.3.2}}), on
trouve une suite d'isomorphismes
\[
i_{!}\isomto R^{+}\bar{u}_{f\sisi{+}{*}}\circ
L^{+}(\bar{u}^{\sisi{+}{*}}_{f})\circ i_{!}\isomto
R^{+}\bar{u}_{f\ast}\circ j_{!}\circ L^{+}\bar{u}_{f'}^{\ast}
\]
d'où l'isomorphisme voulu.\nde{plus directement, sans supposer $f'$ surjectif, on a encore une suite de flèches canoniques
\[i_!\to i_! R^+\bar{u}_{f'*}L^+\bar{u}^*_{f'}\isomto  R^+\bar{u}_{f*} j_! L^+u^*_{f'},\]
dont la deuxième est un isomorphisme (\hyperref[XVII.lem7.0.3]{7.0.3}). Lorsque $f'$ est surjectif la première en est aussi d'après la descente cohomologique et le composé coïncide avec l'isomorphisme construit dans le texte.
}

Enfin, nous utiliserons la variante suivante du lemme de Chow, dont la
démonstration se trouvera dans \EGA II 2\sisi{ème}{\textsuperscript{ième}} édition,\nde{Voir [C, 2.6], cité dans \nderef{XVII.ndeConrad}.} et qui est une
amélioration de (\sisi{\hyperref[XVII.subsec7.1]{7.1}}{\hyperrefext[XII][XII.lem7.1]{7.1}}):

\begin{enonce*}{Lemme 7.0.6}\etiquette[7.0.6]{XVII.lem7.0.6}
Soient $S$ un schéma quasi-compact et quasi-séparé et $f:X\rightarrow
S$ un morphisme séparé de type fini. Alors il existe un diagramme
commutatif
\[
\xymatrix{
X' \ar[rr]^{p} \ar[d]_{f'}& & X \ar[lld]^{f}\\
S                         & &
}
\]
vérifiant les conditions suivantes:
\begin{enumerate}
\renewcommand{\theenumi}{\roman{enumi}}
\renewcommand{\labelenumi}{(\theenumi)}
\item $p$ est propre et \emph{surjectif}

\item il existe un ouvert dense $U$ de $X$ tel que $p$ induise un
isomorphisme de $p^{-1}(U)$ sur $U$

\item $f'$ est quasi-projectif.
\end{enumerate}
\end{enonce*}

\subsection{Définition de $R^{+}f_{!}$}\etiquette[7.1]{XVII.subsec7.1}\pageoriginale

\begin{enonce*}{Définition 7.1.1}\etiquette[7.1.1]{XVII.def7.1.1}
Soit $f:X\rightarrow S$ un morphisme de schémas. Une
\emph{pseudo-compactification} de $f$ consiste en la donnée d'un
diagramme commutatif
$$\vcenter{
\xymatrix{
X_1 \ar[r]^{\sisi{}{j_1}} \ar[d]_{p_1}& T_1\ar[ldd]^{f_1}\\
X   \ar[d]_{f}& \\
S   &
}}\leqno{(7.1.1.0)}
$$\etiquette[7.1.1.0]{XVII.eq7.1.1.0}
où $p_{1}$ est propre et surjectif, $j_{1}$ est une immersion et
$f_{1}$ est propre.
\end{enonce*}

\setcounter{subsubsection}{1}
\subsubsection{}\etiquette[7.1.2]{XVII.subsubsec7.1.2}
Soit $f:X\rightarrow S$ un morphisme de $\Sch_{!R}$: d'après le
lemme de Chow (\hyperref[XVII.lem7.0.6]{7.0.6}), $f$ possède une pseudo-compactification
$(p_{1},j_{1},f_{1})$. Le diagramme \eqhyperref[XVII.eq7.1.1.0]{7.1.1.0} définit un morphisme
\[
[X_{1}|X]\xrightarrow{j_{1}}[T_{1}|S]
\]
encore noté $j_{1}$, tel que $j_{1n}$ soit une immersion pour tout
$n$. On pose alors $R^{+(1)}f_{!}=\sisi{L^{+}\bar{u}_{p_{1}}\cdot
j_{1!}\circ R^{+}\bar{u}_{f_{1}}}{R^+\bar{u}_{f_1*}\circ j_{1!} \circ L^+\bar{u}_{p_1}^*}$ de sorte que l'on obtient un
foncteur \emph{triangulé} de $D^{+}(X,\mathcal{A}_{X})$ dans
$D^{+}(S,\mathcal{A}_{S})$.

\begin{enonce*}{Construction 7.1.3}\etiquette[7.1.3]{XVII.constru7.1.3}
Pour tout couple $((p_{1},j_{1},f_{1}),(p_{2},j_{2},f_{2}))$ de
pseudo-compactifications de $f$, on a un isomorphisme
$\alpha_{12}:R^{+(1)}f_{!}\rightarrow R^{+(2)}f_{!}$ tel que pour
toute autre pseudo-compactification $(p_{3}, j_{3},f_{3})$, on ait
$\alpha_{23}\circ \alpha_{12}=\alpha_{13}$.
\end{enonce*}

\begin{enumerate}
\renewcommand{\theenumi}{\alph{enumi}}
\renewcommand{\labelenumi}{\theenumi)}
\item On définit un système simplicial double
$[[X_{1},X_{2}|X]]$ en posant
\[
[[X_{1},X_{2}|X]]_{[n]\times [m]}=[X_{1}|X]_{n}\times
[X_{2}|X]_{n}
\]
De même on pose
\[
[[X_{1}|X]]_{[n]\times[m]}=[X_{1}|S]_{n}{\quoi},
\]
les morphismes $[[X_{1}|X]]_{[n]\times [m]}\rightarrow
[[X_{1}|X]]_{[n]\times [m']}$ induits par les morphismes
$[m']\rightarrow [m]$ étant tous égaux à l'identité. On définit de la
même manière les systèmes simpliciaux doubles
$[[T_{1},T_{2}|S]]$ et $[[T_{1}|S]]$, de sorte que
l'on a un diagramme commutatif
\[
\xymatrix@C=1em{
  & & [[T_1,T_2|S]] \ar[rrrr]\ar'[dr][rrrddd]_(.2){u_{f_1f_2}}|(.38)\hole& &   & &  [[T_1|S]]\ar[lddd]^{u_{f_1}}\\
  [[X_1,X_2|X]]
  \ar[rrrr]\ar[rru]^{j_{12}}\ar[rrrddd]_{u_{p_1 p_2}} & &  & & [[X_1|X]]\ar[rru]_{j_1}\ar[lddd]_(.75){u_{p_1}}  \\
  & & & & &   &\\
  & & & & & C^{\Delta\times\Delta}_S\\
  &&& C^{\Delta\times\Delta}_X \ar[rru]
}
\]
où\pageoriginale $j_{12}$ est une immersion degré par degré. Le
procédé de calcul exposé dans \sisi{VI b (\refnoexpose{Vbis.subsec2.3})}{\hyperrefext[Vbis][Vbis.subsec2.3]{2.3}} permet de construire un
morphisme
\[
R^{+(1)} f_{!}\xrightarrow{\alpha_{1}}R^{+}\bar{u_{f_{1}f_{2}}}\sisi{}{_*}\circ
j_{12!} \circ L^{+}\bar{u_{p_{1}p_{2}}}^{\sisi{+}{\ast}}=R^{{+}(1,2)}f_{!}
\]
qui s'obtient à partir d'un morphisme de triple complexes de faisceaux
sur $S$.

\item $\alpha_{1}$ \emph{est un isomorphisme}:

Pour démontrer ceci, on peut fixer le premier indice $n$ provenant de
$\bigl[[T_{1},T_{2}|S]\bigr]_{[n]\times[m]}$ dans le morphisme de triple
complexes précédent et regarder le morphisme de double complexes ainsi
obtenu. On vérifie que ce morphisme n'est autre que l'isomorphisme
défini dans le lemme \hyperref[XVII.constru7.0.5]{7.0.5} pour le diagramme
\[
\xymatrix{
**[r][X_1|X]_n \times X_2  \ar[r] \ar[d] & **[r][T_1|S]_n \times
T_2 \ar[d]\\
**[r][X_1|X]_n             \ar[r] & **[r][T_1|S]_n
}
\]

\item On construit de la même manière un isomorphisme
\[
R^{+(2)}f_{!}\xrightarrow{\alpha_{2}}R^{+}\bar{u_{f_{1}f_{2}}}_{\sisi{+}{*}}\circ
j_{12!} \circ L^{+}\bar{u_{p_{1}p_{2}}}\sisi{}{^*}
\]
et l'on pose $\alpha_{12}=\alpha^{-1}_{2}\circ \alpha_{1}$.
\end{enumerate}

Pour vérifier la relation $\alpha_{23}\circ \alpha_{12}=\alpha_{13}$,
on utilise un objet simplicial triple: on laisse au lecteur le soin de
vérifier que l'on obtient un diagramme commutatif:\pageoriginale
$$
\xymatrix{
  R^{+(2,3)}{f_!} \ar[rd]& R^{+(2)}f_! \ar[r] \ar[l]& R^{+(1,2)}f_! \ar[ld]\\
                &         R^{+(1,2,3)}f_!  &                  \\
R^{+(\sisi{2}{3})}f_! \ar[uu]\sisi{}{\ar[rd]}     &                         &  R^{+(1)}f_! \ar[uu]\sisi{}{\ar[ld]}\\
                &    R^{+(1,3)}f_! \ar[uu]        &
}
$$

\begin{enonce*}{Construction 7.1.4}\etiquette[7.1.4]{XVII.constru7.1.4}
Soient $Z\xrightarrow{g}Y$ et $Y\xrightarrow{f}X$ deux morphismes de
$\Sch_{!R}$. Pour tout triple $((p_{1},j_{1},f_{1}),
(p_{2},j_{2},g_{2}), (p_{3},j_{3},(fg)_{3}))$, où
$(p_{1},j_{1},f_{1})$ est une pseudo-compactification de $f$,
$(p_{2},j_{2},g_{2})$ une pseudo-compactification de $g$ et
$(p_{3},j_{3},(fg)_{3})$ une pseudo-compactification de $fg$, on a un
isomorphisme $\sisi{R^{+(3)}fg_{!}\isomfrom  R^{+(2)}g_{!}\circ
R^{+(1)}f_{!}}{R^{+(1)}f_!\circ R^{+(2)}g_!\isomto R^{+(3)}(fg)_!}$ compatible avec les isomorphismes $\alpha_{ij}$ définis
précédemment.
\end{enonce*}

Nous procéderons en plusieurs étapes:

\setcounter{subsubsection}{4}
\paragraph{}\etiquette[7.1.4.1]{XVII.par7.1.4.1}
Soit $(p_{1},j_{1},f_{1})$ (\resp $(p_{2},j_{2},f_{2})$) une
pseudo-compactification de $f:Y\rightarrow X$ (\resp de
$g:Z\rightarrow Y$):
$$\vcenter{\etiquette[7.1.4.1.1]{XVII.eq7.1.4.1.1}
\xymatrix{
  &     & Z\ar[d]^{g}  & Z_2\ar[l]_{p_2}\ar@{^{(}->}[r]^{j_2}  & T_2 \ar[lld]^{g_2}\\
T_1  \ar[rrd]_{f_1}& Y_1 \ar[r]^{p_1} \ar@{_{(}->}[l]_{j_1}& Y \ar[d]^{f} &      &    \\
    &     & X        &      &
}}\leqno{(7.1.4.1.1)}
$$

Le morphisme $Z'=\fprod{Z}{Y_{1}}{Y}\rightarrow Y_{1}$ déduit de $g$
par changement de base est muni d'une pseudo-compactification:
$$\vcenter{\etiquette[7.1.4.1.2]{XVII.eq7.1.4.1.2}
\xymatrix{
  T'_2 = \fprod{T_2}{Y_1}{Y}   \ar[rdd]_{g'_2}& \fprod{Z_2}{Y_1}{Y} =Z'_2 \ar[d]^{p'_2}
\ar@{_{(}->}[l]_{j'_2}\\
&  \fprod{Z}{Y_1}{Y} = Z' \ar[d]^{g'}\\
&    Y_1
}}\leqno{(7.1.4.1.2)}
$$


On\pageoriginale choisit maintenant une \sisi{}{pseudo-}compactification de
\sisi{$j_{1}\cdot g'_{2}$.
$$\etiquette[7.1.4.1.3]{XVII.eq7.1.4.1.3}
\vcenter{
\xymatrix{
T_{12}        \ar[d]_{g''_2} & \fprod{T_2}{Y_1}{Y}=T'_2
\ar@{_{(}->}[l]_(.6){j'_1} \ar[d]^{g'_2}\\
T_1            &   Y_1 \ar@{_{(}->}[l]
}}\leqno(7.1.4.1.3)
$$%
}{$j_1\circ g'_2$\nde{dans le texte original, on a pris une compactification, qui est possible grâce au théorème de compactification. Le but de l'appendice étant d'éviter ce théorème (\cf \nderef{XVII.ndeapp}), l'éditeur s'est permis de la remplacer par une pseudo-compactification et de légèrement modifier \hyperref[XVII.par7.1.4.2]{7.1.4.2}. Une autre solution serait de se borner aux pseudo-compactifications $(p_1,j_1,f_1)$ dont $f_1$ est \emph{projectif}.}
$$\etiquette[7.1.4.1.3]{XVII.eq7.1.4.1.3}
\vcenter{
\xymatrix{
T_{12}        \ar[dd]_{g''_2} &T_2''\ar[d]\ar[d]^p\ar@{_{(}->}[l]_{j'_1}\\
& \fprod{T_2}{Y_1}{Y}=T'_2\ar[d]^{g'_2}\\
T_1            &   Y_1 \ar@{_{(}->}[l]_{j_1}
}}\leqno(7.1.4.1.3)
$$}%
qui définit une pseudo-compactification du composé $fg$: on notera
$R^{+(12)}\sisi{\int g}{(fg)}_{!}$ le foncteur correspondant.

\paragraph{}\etiquette[7.1.4.2]{XVII.par7.1.4.2}
\emph{Les données précédentes permettent de définir un
isomorphisme} $R^{+(12)}\sisi{\int g}{(fg)}_{!}\isomfrom  R^{+(1)}f_{!}\sisi{\cdot}{\circ} R^{+(2)}g_{!}$:

On dispose d'un diagramme commutatif, où $Z_2''=\fprod{T''_2}{Z'_2}{T'_2}$ et toutes les flèches se
déduisent naturellement de \eqhyperref[XVII.eq7.1.4.1.1]{7.1.4.1.1}, \eqhyperref[XVII.eq7.1.4.1.2]{7.1.4.1.2} et
\eqhyperref[XVII.eq7.1.4.1.3]{7.1.4.1.3}.
$$\etiquette[7.1.4.2.1]{XVII.eq7.1.4.2.1}
\vcenter{
\xymatrix{
[T_{12}|X] \ar[rd]_{g''_2} &   &[T\sisi{'}{''}_2|Y] \ar[rd]_{g'_2\sisi{}{p}} \ar[ll]_{j'_1}&
[Z\sisi{'}{''}_{2}|Z] \ar[r]^{r}\ar[l]_{j\sisi{'}{''}_2} \ar[d]_{l}&
C^{\Delta}_Z \ar[d]^{g}\\
            &    [T_1|X] \ar[rrrd]_{u_{f_1}} &  & [Y_1|Y]\ar[r]_{u_{p_1}}\ar[ll]_{j_1} & C_Y^{\Delta} \ar[d]^{f}\\
            &       &  &   &  C_X^{\Delta}
}}\leqno(7.1.4.2.1)
$$

On\pageoriginale va définir une flèche naturelle:
\[
\beta:R^{+(1)}f_{!}\circ R^{+(2)}g_{!}\longrightarrow R^{+(12)}\sisi{fg}{(fg)}_{!}
\]
dont on va montrer que c'est un isomorphisme.

Pour cela nous utiliserons deux lemmes:

\begin{enonce*}{Lemme 7.1.4.2.2}\etiquette[7.1.4.2.2]{XVII.lem7.1.4.2.2}
Soit un diagramme commutatif:
\[
\xymatrix{
X_1    \ar[d]_{p_1}     \ar[r]^{j_1} & T_1 \ar[d]^{h'}\\
X    \ar[d]_{f}       \ar[r]_{j}           & T \ar[ld]^{h}\\
    S       &
}
\]
où $h$, $h'$ sont propres, $p_{1}$ est propre et surjectif et $j$,
$j_{1}$ sont les immersions. Alors la flèche canonique
\[
R^{+}h\sisi{}{_*}\circ j_{!}\longrightarrow R^{+(1)}\sisi{(hh')}{f}_{!}
\]
est un isomorphisme.
\end{enonce*}

En effet, d'après (\hyperref[XVII.constru7.0.5]{7.0.5}) la flèche
\[
j_{!}\longrightarrow R^{+}\bar{u}_{h'\sisi{}{*}}\circ j_{1!}\circ
L^{+}\bar{u}_{p_{1}}\sisi{}{^*}
\]
est un isomorphisme. D'où le résultat en appliquant $R^{+}h\sisi{}{_*}$.

\begin{enonce*}{Lemme 7.1.4.2.3}\etiquette[7.1.4.2.3]{XVII.lem7.1.4.2.3}
Avec les notations du diagramme \emph{\eqhyperref[XVII.eq7.1.4.2.1]{7.1.4.2.1}}, les flèches
canoniques
\begin{gather*}
R^{+(1)}f_{!}\circ R^{+}\bar{u}_{p_{1}\sisi{}{*}}\longrightarrow
R^{+}\bar{u}_{f_{1}\sisi{}{*}}\circ j_{1!}\\
R^{+(2)}g_{!}\circ R^{+}\bar{r}\sisi{}{_*}\longrightarrow
R^{+}\bar{u}_{p_{1}\sisi{}{*}}\circ R^{+}\sisi{g'_{2}}{(g_2'p)_*}\circ j\sisi{'}{''}_{2!}
\end{gather*}
sont des isomorphismes.
\end{enonce*}

Nous n'établirons que le premier isomorphisme, laissant au lecteur le
soin d'en\pageoriginale déduire le deuxième.

Avec des notations évidentes\sisi{, on pose $C=\fprod{[Y_{1}|Y]}{X_{1}}{X}$
et $\bar{C}=\fprod{[T_{1}|S]}{T_{1}}{S}$, de sorte que l'on a un
diagramme commutatif:
\[
\xymatrix{
                                & \bar{C} \ar[rdd]|\hole \ar[rr]  & & T_1 \ar[ldd]^{f_1}\\
C                 \ar[rr]  \ar[ru]^{i}  \ar[d]^{q}    &    &  Y_1  \ar[ru]_{j_1} \ar[d]^{p_1}&      \\
[Y_1|Y] \ar[rr]_{u_{p_1}}&  &      Y         &
}
\]
où $i$ est une immersion, $q$ est propre et surjectif. On construit
alors un diagramme
\[
\xymatrix{
\bigl[\bigl[C \underset{q}{\bigm|}[Y_1|Y]\bigr]\bigr] \ar[r]^{i} \ar[d]^{u_q}&
[[\bar{C}|C^{\Delta}_Y]] \ar[ldd]^{k}\\
|X_1 | X |                             \ar[d]         &\\
Y                       &                        &
}
\]
et il résulte de (\hyperref[XVII.lem7.1.4.2.2]{7.1.4.2.2}) que le foncteur $R\bar{u}j_{1}\circ
j_{1}$ s'identifie au foncteur $R^{+}\bar{k}\circ i_{!}\circ
L^{+}\bar{u}_{q}$ et, grâce à (\hyperref[XVII.lem7.0.3]{7.0.3}) et au théorème du changement de
base ce dernier foncteur s'identifie à $R^{+(1)}f_{!}\circ
R^{+}\bar{u}_{p_{1}}$. }{ et les notations de \hyperref[XVII.constru7.1.3]{7.1.3} a), on a un diagramme commuatif
\[\xymatrix{&[Y_1|Y]\ar[rr]^{j_1}\ar[dd]^(.7){\bar{u}_{p_1}}|\hole &&[T_1|X]\ar[dd]^{\bar{u}_{f_1}}\\
[[Y_1,Y_1|Y]]\ar[ur]^{k}\ar[dd]_{q}\ar[rr]^(.7)i && [[T_1,T_1|X]]\ar[dd]\ar[ur]\\
&Y\ar[rr]^(.3){f}|\hole && X\\
[Y_1|Y]\ar[ur]^{\bar{u}_{p_1}}\ar[rr]^{j_1} &&[T_1|X]\ar[ur]_{\bar{u}_{f_1}} }\]
où $q$ et $k$ sont les deux projections de $[[Y_1,Y_1|Y]]$ à $[Y_1|Y]$. Grâce à la descente cohomologique, $R^+\bar{u}_{f_1 *} \circ j_{1!}$ s'identifie à $R^+\bar{u}_{f_1 *} \circ j_{1!} \circ R^+ q_* L^+ q^*$, qui, d'après \hyperref[XVII.lem7.0.3]{7.0.3} appliqué aux faces antérieure et supérieure du cube ci-dessus, s'identifie à $R^+\bar{u}_{f_1*}\circ j_{1!}\circ R^+k_* \circ L^+q^*$. En vertu du théorème de changement de base, ce dernier foncteur s'identifie à $R^+\bar{u}_{f_1 *} \circ j_{1!} \circ L^+\bar{u}_{p_1*} \circ R^+\bar{u}_{p_1 *}$, qui est par définition $R^{+(1)}f_{!}\circ
R^{+}\bar{u}_{p_{1}*}$.\nde{la démonstration dans le texte original est légèrement différente %peu compréhensible
et utilise \hyperref[XVII.lem7.1.4.2.2]{7.1.4.2.2}.}

}%
Le lemme (\hyperref[XVII.lem7.1.4.2.3]{7.1.4.2.3}) étant établi, on dispose
d'une suite d'isomorphismes
\begin{align*}
R^{+(\sisi{3}{12})} \sisi{fg}{(fg)}_{!} &= R^{+}\bar{u}_{f_{1}\sisi{}{*}}\circ R^{+}g''_{2}\circ
j'_{1!}\circ j\sisi{'}{''}_{2!}\circ L^{+}\bar{r}^{\sisi{+}{*}}\\
&\isomfrom  R^{+}\bar{u}_{f_{1}\sisi{+}{*}}\circ j_{1!}\circ R^{+}
\sisi{g'_{2}}{(g'_2p)_*}\circ j'_{2!}\circ L^{+}\bar{r}^{\sisi{+}{*}}\\
&\isomfrom  R^{+(1)}f_{!}\circ R^{+}\bar{u}_{p_{1}\sisi{+}{*}}\circ
R^{+}\sisi{g'_{2+}}{(g'_2p)_*}\circ j'_{2!}\circ L^{+}\bar{r}\sisi{}{^*}\\
&\isomfrom  R^{+(1)}f_{!}\circ R^{+(2)}g_{!}\circ
R^{+}\bar{r}\sisi{}{_*}\circ L^{+}\bar{r}^{\sisi{+}{*}}\\
&\isomfrom  R^{+(1)}f_{!}\circ R^{+(2)}g_{!}
\end{align*}
dont\pageoriginale le composé donne $\beta$.

\paragraph{}\etiquette[7.1.4.3]{XVII.par7.1.4.3}
Si $(p_{3},j_{3},(fg)_{3})$ est une pseudo-compactification de $fg$,
on définit un isomorphisme
\[
R^{+(1)}f_{!}\circ
R^{+(2)}g_{!}\xrightarrow{\gamma_{123}}R^{+(3)}fg_{!}
\]
en composant $\alpha_{(12)3}$ avec $\beta$.

\paragraph{}\etiquette[7.1.4.4]{XVII.par7.1.4.4}
Soient maintenant $(p_{1'},j_{1'},f_{1'})$ et $(p_{2'},j_{2'},g_{2'})$
deux autres pseudo-compactifications de $f$ et de $g$ et un diagramme
\sisi{(\hyperref[XVII.par7.1.4.3]{7.1.4.3})}{\eqhyperref[XVII.eq7.1.4.1.3]{7.1.4.1.3}}; on vérifie que le diagramme suivant est commutatif:
\[
\xymatrix{
R^{+(12)}\sisi{fg}{(fg)}_! \ar[d]_{\alpha_{(12)(1'2')}} & R^{+(1)}f_! \circ
R^{+(2)}g_! \ar[d]^{\alpha_{11'}\ast \alpha_{22'}} \ar[l]_(.6){\beta}\\
R^{+(1'2')}\sisi{fg}{(fg)}_!  & R^{+(1')}f_! \circ R^{+(2')}g_! \ar[l]^(.6){\beta'}
}
\]
en considérant un diagramme du type \sisi{(\hyperref[XVII.par7.1.4.4]{7.1.4.4})}{\eqhyperref[XVII.eq7.1.4.2.1]{7.1.4.2.1}} où les objets
simpliciaux sont remplacés par des objets simpliciaux doubles.

On en conclut que le diagramme suivant est commutatif:
\[
\xymatrix{
R^{+(3)}\sisi{fg}{(fg)}_! \ar[d]_{\alpha_{33'}}& R^{+(1)}f_! \circ R^{+(2)}g_!
\ar[l]_(.6){\gamma_{123}} \ar[d]^{\alpha_{11'}\ast
  \alpha_{22'}}\\
R^{+(3')}\sisi{fg}{(fg)}_!  & R^{+(1')}f_! \circ R^{+(2')}g_! \ar[l]_(.6){\gamma_{1'2'3'}}
}
\]
pour toute pseudo-compactification $(p_{3'},j_{3'},(fg)_{3'})$ de
$fg$, ce qui achève la démonstration de \hyperref[XVII.constru7.1.4]{7.1.4}.

\subsubsection{}\etiquette[7.1.5]{XVII.subsubsec7.1.5}
En choisissant pour tout morphisme $f$ de $\Sch_{!R}$ une
pseudo-compactification, de telle manière qu'à un morphisme identique
soit associée la pseudo-compactification triviale, on définit pour
tout morphisme $f:X\rightarrow Y$ un foncteur triangulé
\[
R^{+}f_{!}:D^{+}(X,\mathcal{A}_{X})\longrightarrow
D^{+}(Y,\mathcal{A}_{Y})
\]
et\pageoriginale tout couple $(f,g)$ de morphisme, un isomorphisme
\[
R^{+}fg_{!}\xrightarrow{\substack{c_{fg}\\\sim}}
R^{+}f_{!}\circ R^{+}g_{!}
\]
De plus si $f$ est une identité, $R^{+}f_{!}$ est le foncteur
identique et $c_{fg}$ est l'identité.

\begin{enonce*}{Proposition 7.1.6}\etiquette[7.1.6]{XVII.prop7.1.6}
La correspondance $f\rightarrow R^{+}f_{!}$ définit un pseudo-foncteur
normalisé de $\Sch_{!R}$ dans la $2$-catégorie des catégories. En
d'autres termes, les catégories $D^{+}(X,\mathcal{A}_{X})$ pour $X$
objet de $\Sch_{!R}$ peuvent être organisées en catégorie cofibrée
sur $\Sch_{!R}$ en posant, pour tout morphisme $f:X\rightarrow Y$
de $\Sch_{!R}$:
\[
\Hom_{f}(K,L)=\Hom(R^{+}f_{!}K,L)
\]
avec $K\in D^{+}(X,\mathcal{A}_{X})$ et $L\in D^{+}(Y,\mathcal{A}_{Y})$.
\end{enonce*}

Il s'agit de voir que l'identification
$R^{+}fg_{!}\xrightarrow{\substack{c_{fg}\\\sim}}
R^{+}f_{!}\circ R^{+}g_{!}$ est compatible avec les composés triples. Soit
donc $h:T\rightarrow Z$, $g:Z\rightarrow Y$ et $f:Y\rightarrow X$
trois morphismes de $\Sch_{!R}$: en vertu de ce qui précède, il
suffit, pour vérifier cette compatibilité, de la vérifier après un
choix arbitraire de pseudo-compactifications pour $h$, $g$ et $f$.

On dispose d'un diagramme:
$$\vcenter{\etiquette[7.1.6.1]{XVII.eq7.1.6.1}
  \xymatrix{
 T_5 \ar[rd]& T_4 \ar[rd] \ar@{_{(}->}[l]& T_3 \ar[rd]\ar@{_{(}->}[l]
 & T'\ar@{_{(}->}[l] \ar[r]^{p_3} \ar[d]& T \ar[d]^{h} \\
        & T_3 \ar[rd]  & T_2 \ar@{_{(}->}[l] \ar[rd]  & Z'\ar[r]^{p_2}\ar@{_{(}->}[l] \ar[d]& Z
  \ar[d]^{g}\\
        &    & T_1 \ar[rrd]& Y_1\ar@{_{(}->}[l] \ar[r]^{p_1}& Y \ar[d]^{f}\\
        &    &     &     & X
}}\leqno{(7.1.6.1)}
$$
où toutes les flèches obliques sont propres et où $p_{1}$, $p_{2}$ et
$p_{3}$ sont propres et surjectifs.

Un diagramme du type (\hyperref[XVII.par7.1.4.4]{7.1.4.4}) construit sur
\eqhyperref[XVII.eq7.1.6.1]{7.1.6.1} permet alors de conclure.

\begin{enonce*}{Proposition 7.1.7}\etiquette[7.1.7]{XVII.prop7.1.7}
Si\pageoriginale $f:X\rightarrow S$ est un morphisme compactifiable,
le foncteur $Rf_{!}$ coïncide avec celui défini dans l'exposé \emph{XVII}.
\end{enonce*}

En effet, il existe un diagramme commutatif:
\[
\xymatrix{
X_1         \ar[r]^{j_1} \ar[d]_{p_1}& T_1 \ar[d]^{h'}\\
X           \ar[r]_{j}  \ar[d]_{f} & T \ar[ld]^{h}\\
S
}
\]
auquel on applique (\hyperref[XVII.lem7.1.4.2.2]{7.1.4.2.2}).

\begin{enonce*}{Définition 7.1.8}\etiquette[7.1.8]{XVII.def7.1.8}
Si $f:X\rightarrow Y$ est un morphisme de $\Sch_{!}$, le foncteur
$R^{+}f_{!}$ que l'on vient de définir s'appellera le foncteur \emph{image
directe à supports propres} (au sens des catégories dérivées). Le
foncteur $H^{q}\circ R^{+}f_{!}$, noté $R^{q}f_{!}$, s'appellera le
$q$\textsuperscript{ième} \emph{foncteur image directe à supports propres.}
\end{enonce*}

Avec ces notations, on laisse au lecteur le soin de vérifier la

\begin{enonce*}{Construction 7.1.9}\etiquette[7.1.9]{XVII.constru7.1.9}
Soit $(\pi_{1}, j_{1}, f_{1})$ une pseudo-compactification de
$f:X\rightarrow Y$ et désignons par $\pi^{n}_{1}$ le morphisme
canonique $[X_{1}|X]_{n}\rightarrow X$. Alors, pour tout complexe
$K^{\bigdot}$ de $D^{+}(X,\mathcal{A}_{X})$, il existe deux suites
spectrales birégulières:
\begin{align*}
E^{pq}_{1} &= R^{q}f_{!}(K^{p})\Longrightarrow R^{p+q}
f_{!}(K^{\bigdot})\tag{7.1.9.1}\\
E^{pq}_{1} &= R^{q}(f\circ
\pi^{p}_{1})_{!}(\pi^{p\sisi{+}{*}}_{1}(K^{\bigdot}))\Longrightarrow
R^{p+q}f_{!}(K^{\bigdot})\tag{7.1.9.2}
\end{align*}\etiquette[7.1.9.1]{XVII.eq7.1.9.1}
\etiquette[7.1.9.2]{XVII.eq7.1.9.2}%L'espace vertical suivant ce paragraphe est trop large.
\end{enonce*}

(Prendre par exemple la résolution flasque canonique de $j_{1!}\circ
L^{+}\bar{u}_{p_{1}}^{\sisi{+}{*}}(K^{\bigdot})$ (\cf XVII (\hyperref[XVII.subsec4.2]{4.2})), puis son
image directe par $u_{f_{1}}$).

\subsection[Quelques propriétés de $R^{+}f_{!}$]{Quelques propriétés de
  \mathversion{bold}$R^{+}f_{!}$}\etiquette[7.2]{XVII.subsec7.2}
\pageoriginale

\begin{enonce*}{Proposition 7.2.1}\etiquette[7.2.1]{XVII.prop7.2.1}
Soit $f:X\rightarrow Y$ un morphisme de $\Sch_{!R}$ de dimension
relative $\leq d$. Alors le foncteur $R^{+}f_{!}$ est de dimension
cohomologique $\leq 2d$.
\end{enonce*}

Il s'agit de montrer que pour tout faisceau $F$ de
$\mathcal{A}_{X}$-modules sur $X$, on a $R^{i}f_{!}(F)=0$ pour
$i>2d$. (On en déduit alors, par un argument standard, que pour tout
complexe $K^{\bigdot}$ de $D^{+}(X,\mathcal{A}_{X})$ tel que
$H^{i}(K^{\bigdot})=0$ pour $i>k$, on a $R^{i}f_{!}(K^{\bigdot})=0$
pour $i>k+2d$).

Soit $\mathcal{U}=(U_{i})_{i\in I}$ un recouvrement fini de $X$ par
des ouverts affines. On définit (\hyperref[XVII.prop6.2.9]{6.2.9}) une résolution gauche
$K^{\bigdot}_{\mathcal{U}}\rightarrow F$ de $F$ en posant:
\[
K^{n}_{\mathcal{U}}=\bigoplus_{|P|=\sisi{}{-}n+1} j_{P!}j^{*}_{P}(F)
\]
où pour chaque partie $P$ de $I$, $j_{P}$ est l'inclusion de
$\cap_{i\in P} U_{i}$ dans $X$.

On déduit de \eqhyperref[XVII.eq7.1.9.1]{7.1.9.1} une suite spectrale
\[
E^{pq}_{1}=\bigoplus_{\substack{|P|=\sisi{}{-}p+1\\ P\sisi{\in}{\subset} I\sisi{}{\text{ non vide}}}} R^{p}(fj_{P})_{!}
j^{\ast}_{P}(F)\Longrightarrow R^{p+q}f_{!}(F)
\]
et le fait que $fj_{P}$ soit compactifiable, cas justiciable de
l'exposé XVII, montre que $E^{pq}_{1}=0$ pour $p>0$ et $q>2d$; on en
déduit alors que $R^{n}f_{!}(F)=0$ pour $n>2d$, ce qu'il fallait
démontrer.

\begin{enonce*}{Proposition 7.2.2}\etiquette[7.2.2]{XVII.prop7.2.2}
Soient $f:X\rightarrow Y$ un morphisme de $\Sch_{!R}$ et
$K^{\bigdot}$ un objet de $D^{b}(X,\mathcal{A}_{X})$ de \tTor-dimension
finie. Alors $R^{+}f_{!}(K^{\bigdot})$ est de \tTor-dimension finie.
\end{enonce*}

Comme précédemment, soit $\mathcal{U}=(U_{i})_{i\in I}$ un
recouvrement fini de $X$ par des ouverts affines. La fonctorialité de
$K^{n}_{\mathcal{U}}$ permet de fabriquer une résolution gauche bornée
de $K^{\bigdot}$ par des complexes de \tTor-dimension finie, car cette
dernière propriété se vérifie point par point (\hyperref[XVII.def4.1.9]{4.1.9}): nous
noterons $K^{\bigdot\bigdot}_{\mathcal{U}}$ le double complexe ainsi
obtenu, le \emph{deuxième} indice correspondant à celui de
$K^{\bigdot}$. Avec les conventions de (\sisi{VI b.(\refnoexpose{Vbis.subsubsec2.3.2})}{\hyperrefext[Vbis][Vbis.subsubsec2.3.2]{2.3.2}}), on dispose pour
tout entier $n$ d'une suite exacte, où l'opération $\sigma_{\leq n}$
est prise dans $C^{+}(C^{+}(\Mod(X,\mathcal{A}_{X})))$:
\[
\sisi{0\longrightarrow (\sigma_{\leq
n}(K^{\bigdot\bigdot}_{\mathcal{U}}))_{s}\longrightarrow (\sigma_{\leq
n+1}(K^{\bigdot\bigdot}_{\mathcal{U}}))_{s}\longrightarrow (-1)^{n+1}
K^{n+1,\bigdot}_{\mathcal{U}}\longrightarrow 0}{0\to (-1)^{n+1}K^{n+1,\bigdot}_{\mathcal{U}}\to(\sigma_{\leq
n+1}(K^{\bigdot\bigdot}_{\mathcal{U}}))_{s}\to(\sigma_{\leq
n}(K^{\bigdot\bigdot}_{\mathcal{U}}))_{s}\to 0}
\]
d'où\pageoriginale un triangle distingué dans $D^{+}(X,\mathcal{A}_{X})$:
\[
\xymatrix@C=1em{
  & K^{n+1,\bigdot}_{\Uu} \sisi{\ar[ld]}{\ar[rd]}&    \\
  (\sigma_{\leq n}(K^{\bigdot\bigdot}_{\Uu}))_s \sisi{\ar[rr]}{\ar[ur]}&  &(\sigma_{\leq
  n+1}(K^{\bigdot\bigdot}_{\Uu}))_s \sisi{\ar[lu]}{\ar[ll]}\\
}
\]

Utilisant le fait que $R^{+}f_{!}(K^{n+1,.}_{\mathcal{U}})$ est de
\tTor-dimension finie pour tout $\sim$ (\hyperref[XVII.thm5.2.10]{5.2.10}), on en déduit que
$R^{+}f_{!}((\sigma_{\leq
0}(K^{\bigdot\bigdot}_{\mathcal{U}}))_{s})=R^{+}f_{!}(K^{\bigdot})$
est de \tTor-dimension finie.

\begin{enonce*}[remark]{Remarque 7.2.3}\etiquette[7.2.3]{XVII.rem7.2.3}
L'argument précédent ne permet pas d'affirmer que
$R^{+}f_{!}(K^{\bigdot})$ est de \tTor-dimension $\leq k$ dès que
$K^{\bigdot}$ est de \tTor-dimension $\leq k$ (\cf \sisi{(3.7)}{\hyperref[XVII.subsubsec7.3.7]{7.3.7}}).
\end{enonce*}

\begin{enonce*}{Proposition 7.2.4}\etiquette[7.2.4]{XVII.prop7.2.4}
Soit un diagramme cartésien dans $\Sch_{!R}$:
\[
\xymatrix{
X'   \ar[r]^{g'}\ar[d]_{f'}& X \ar[d]^{f}\\
S'  \ar[r]_{g} & S
}
\]
Alors pour tout complexe $K^{\bigdot}$ de $D^{+}(X,\mathcal{A}_{X})$,
le morphisme canonique
\[
g^{\ast}\circ
R^{+}f_{!}(K^{\bigdot})\xrightarrow{\mathrm{ch}_{K^{\bigdot}}}R^{+}f_{!}\sisi{}{'}\circ
{g'}^{\ast}(K^{\bigdot})
\]
est un isomorphisme.
\end{enonce*}

Soit $(\pi_{1}, j_{1}, f_{1})$ une pseudo-compactification de $f$: on
en déduit par changement de base une pseudo-compactification
$(\pi'_{1}, j'_{1}, f'_{1})$ de $f'$ et un diagramme commutatif:
\[
\xymatrix@R=2ex{
                 && T'_1 \ar'[ldddd]^(.8){f'_1}|(.55)\hole [lldddddddd]|(.25)\hole \ar[rrrr] &&&& T_1 \ar[lldddddddd]^{f_1}\\
                 &\\
X'_1 \ar[rruu]^{j'_1}\ar[ddd]_{\pi'_1} \ar[rrrr] &&&& X_1 \ar[rruu]_{j_1} \ar[ddd]^{\pi_1}&\\
                     & & &\\
                     & *{} & &\\
X'   \ar[rrrr]_{g'} \ar[ddd]_{f'} &&&& X  \ar[ddd]_{f}  &\\
                     & & &\\
                     & & &\\
S'   \ar[rrrr]_{g} &&&&  S         &
}
\]
qui\pageoriginale permet de définir la flèche $\mathrm{ch}_{K}$.

On dispose alors, d'après \eqhyperref[XVII.eq7.1.9.2]{7.1.9.2}, d'un morphisme de suites
spectrales birégulières:
\[
\xymatrix{
g^{\ast}\circ R^{q} (f \circ \pi^{p}_1)_! (\pi^{p\sisi{+}{\ast}}_1
(K^{\bigdot})) \ar[d]^{u^{pq}_1} \ar@{=>}[r]&  g^{\ast} \circ R^{p+q}f_{!}(K^{\bigdot}) \ar[d]^{\mathrm{ch}_K}\\
R^{q}(f'\circ \pi'^{p}_1)_! (\pi'^{p}_1 (g'^{\ast}(K^{\bigdot})))  \ar@{=>}[r] &
  R^{+}f'_! \circ g'^{\ast}(K\sisi{}{^{\bigdot}})
}
\]
et $\mathrm{ch}_{K^{\bigdot}}$ est un isomorphisme puisqu'il en est de même
pour les $u^{pq}_{1}$ d'après (\hyperref[XVII.thm5.2.6]{5.2.6}).

\subsection{Définition de $Rf_{!}$}\etiquette[7.3]{XVII.subsec7.3}

Nous nous proposons maintenant d'étendre « de façon
raisonnable » le foncteur $R^{+}f_{!}$ en un foncteur
$Rf_{!}:D(X,\mathcal{A}_{X})\rightarrow D(Y,\mathcal{A}_{Y})$, pour
tout morphisme $f$ de $\Sch_{!R}$.

\subsubsection{}\etiquette[7.3.1]{XVII.subsubsec7.3.1}
Soit $(p_{1}, j_{1}, f_{1})$ une pseudo-compactification de $f$ et $F$
un $\mathcal{A}_{X}$-module: prenant la résolution flasque canonique
de $j_{1!}\circ L\bar{u}_{p_{1}}^{\ast}(F)$ (\cf \hyperref[XVII.subsec4.2]{4.2}), puis son
image directe par $f_{1}$, puis le complexe simple associé au complexe
double ainsi construit, on obtient un complexe de
$\mathcal{A}_{Y}$-modules:
$$\etiquette[7.3.1.1]{XVII.eq7.3.1.1}
0\longrightarrow
\Xx^{0}(F)\xrightarrow{S^{1}(F)}\Xx^{1}(F)\ldots
\Xx^{n}(F)\xrightarrow{S^{n}(F)}
\Xx^{n+1}(F)\ldots.\leqno(7.3.1.1)
$$
fonctoriel\pageoriginale en $F$ et isomorphe dans
$D^{+}(Y,\mathcal{A}_{Y})$ à $R^{+}f_{!}(F)$.

D'après \hyperref[XVII.prop7.2.1]{7.2.1} $H^{i}(\Xx^{\bigdot}(F))=0$ pour $i>2d$, où $d$
est un entier tel que la dimension relative de $f$ soit $\leq d$, de
sorte que le complexe
$$\etiquette[7.3.1.2]{XVII.eq7.3.1.2}
0\to \Xx^{0}(F)\to
\Xx^{1}(F)\to \ldots \sisi{}{\to}
\Xx^{2d-1}(F)\to \Ker S^{2d}(F)\to
0\leqno(7.3.1.2)
$$
est canoniquement isomorphe dans $D^{+}(Y,\mathcal{A}_{Y})$ au
complexe \eqhyperref[XVII.eq7.3.1.1]{7.3.1.1}. De plus:

\begin{enonce*}{Lemme 7.3.2}\etiquette[7.3.2]{XVII.lem7.3.2}
Les foncteurs $\Xx^{0}$,
$\Xx^{1},\ldots,\Xx^{2d-1}$, $\Ker S^{2d}$ sont exacts
en $F$.
\end{enonce*}

L'exactitude des foncteurs $\Xx^{n}$ résulte de l'exactitude
de la résolution flasque canonique et du fait que l'image directe par
un morphisme d'une suite exacte courte de faisceaux flasques est
encore une suite exacte courte.

Pour voir l'exactitude de $\Ker S^{2d}$, on considère, pour une suite
exacte courte $0\rightarrow F'\rightarrow F\rightarrow F''\rightarrow
0$ de faisceaux sur $X$, le diagramme
\[
\xymatrix{
\Ker S^{2d}(F')  \ar[r] \ar[d]& \Xx^{2d+1}(F') \ar[r] \ar[d]&
\Xx^{2d+2}(F') \ar[d] \ldots\\
\ker S^{2d}(F)   \ar[r] \ar[d]& \Xx^{2d+1}{(F)} \ar[r] \ar[d]&
\Xx^{2d+2}(F)\ar[d] \ldots \\
\ker S^{2d}(F'')   \ar[r] & \Xx^{2d+1}{(F'')} \ar[r] &
\Xx^{2d+2}{(F'')} \ldots \\
}
\]
et la suite exacte d'homologie associée à la suite exacte
\[
0\longrightarrow \sigma_{\geq 2d+1}
(\Xx^{\bigdot}(F'))\longrightarrow \sigma_{\geq 2d+1}
(\Xx^{\bigdot}(F))\longrightarrow \sigma_{\geq
2d+1}(\Xx^{\bigdot}(F''))\to0
\]
permet de conclure.

On dispose maintenant du lemme général suivant:

\begin{enonce*}{Lemme 7.3.3}\etiquette[7.3.3]{XVII.lem7.3.3}
Soient $A$ et $B$ deux catégories abéliennes et $\mathrm{Je}^{\bigdot}$ un
objet de $C^{b}(\Ffex(A,B))$ où $\Ffex(A,B)$ désigne la
catégorie des foncteurs exacts de $A$ dans $B$. Alors, en associant à
tout complexe $K^{\bigdot}$ de $C(A)$ le complexe simple associé au
double complexe $\mathrm{Je}^{\bigdot}(K^{\bigdot})$ on obtient un foncteur
triangulé $K(A)\rightarrow K(B)$ qui préserve les quasi-isomorphismes.
\end{enonce*}

On\pageoriginale écrit le foncteur en question comme un composé
\[
C(A)\longrightarrow C^{b}(C(B))\xrightarrow{(\quad)_{s}}C(B)
\]
où le foncteur $(\quad)_{s}$ est le foncteur complexe simple associé et
on utilise (\sisi{VI b (\refnoexpose{Vbis.lem2.3.2.2})}{\hyperrefext[Vbis][Vbis.lem2.3.2.2]{2.3.2.2}}. Remarque): les détails sont laissés au
lecteur.

\begin{enonce*}[remark]{Exercice 7.3.4}\etiquette[7.3.4]{XVII.exer7.3.4}
Montrer que pour tout complexe $K^{\bigdot}$ de
$D^{+}(X,\mathcal{A}_{X})$, la flèche canonique
\[
\tau_{\leq 2d} \Xx^{\bigdot}(K^{\bigdot})\longrightarrow
\Xx^{\bigdot}(K^{\bigdot})
\]
induit un quasi-isomorphisme sur les complexes simples associés.
\end{enonce*}

On remarquera pour cela que la flèche canonique
\[
\varinjlim_{k\geq 2d} (\tau_{\leq k}
\Xx^{\bigdot}(K^{\bigdot}))_{s}\longrightarrow
(\Xx^{\bigdot}(K^{\bigdot}))_{s}
\]
est un isomorphisme, puis que les faisceaux de cohomologie d'un
complexe commutent aux limites inductives filtrantes.

\setcounter{subsubsection}{4}
\subsubsection{}\etiquette[7.3.5]{XVII.subsubsec7.3.5}
Soit $R^{(1)}f_{!}:D(X,\mathcal{A}_{X})\rightarrow
D(Y,\mathcal{A}_{Y})$ le foncteur défini par \eqhyperref[XVII.eq7.3.1.2]{7.3.1.2} grâce à
\sisi{\hyperref[XVII.rem7.2.3]{7.2.3}}{\hyperref[XVII.lem7.3.3]{7.3.3}}. En vertu de \sisi{\hyperref[XVII.prop7.2.4]{7.2.4}}{\hyperref[XVII.exer7.3.4]{7.3.4}} on peut le calculer directement à partir de
\eqhyperref[XVII.eq7.3.1.1]{7.3.1.1} et on peut vérifier qu'il « ne dépend pas » de la
pseudo-compactification choisie. Il faudrait alors vérifier les
énoncés \hyperref[XVII.constru7.1.3]{7.1.3}, \hyperref[XVII.constru7.1.4]{7.1.4} et
\hyperref[XVII.prop7.1.6]{7.1.6} en supprimant les exposants $+$, ce que
le rédacteur n'a pas eu le courage de faire.

\begin{enonce*}{Définition 7.3.6}\etiquette[7.3.6]{XVII.def7.3.6}
Le foncteur $Rf_{!}:D(X,\mathcal{A}_{X})\rightarrow
D(Y,\mathcal{A}_{Y})$ construit en \hyperref[XVII.subsubsec7.3.5]{7.3.5} s'appelle encore
\emph{foncteur image directe à supports propres.}
\end{enonce*}

\setcounter{subsubsection}{6}
\subsubsection{}\etiquette[7.3.7]{XVII.subsubsec7.3.7}
On peut alors montrer que le foncteur $Rf_{!}$ commute aux changements
de base et que les énoncés (\hyperref[XVII.prop5.2.9]{5.2.9}) et
(\hyperref[XVII.thm5.2.10]{5.2.10}) restent vrais
pour tout morphisme de $\Sch_{!R}$, ce qui justifie les
constructions que l'on vient de faire.

\begin{thebibliography}{\sisi{TF}{99}}
\sisi{\bibitem[1]{SGA4-17:DP}}{\bibitem{SGA4-17:DP}} A. Dold\pageoriginale  and D. Puppe. Homologie nicht\sisi{ }{-}additiver
Funktoren. Anwendungen. Ann. Inst. Fourier \sisi{XI 1961}{\textbf{11} (1961), 201--312}.

\sisi{\bibitem[1]{SGA4-17:GZ}}{\bibitem{SGA4-17:GZ}} Gabriel-Zisman. Calculus of fraction and homotopy
theory, Ergebnisse, Bd 35, Springer 1967.

\sisi{\bibitem[1]{SGA4-17:Gi}}{\bibitem{SGA4-17:Gi}} J. Giraud. Méthode de la descente, Mémoire nᵒ 2 de la Société Mathématique de France (1964).

\sisi{\bibitem[TF]{SGA4-17:TF}}{\bibitem{SGA4-17:TF}} R. Godement. Théorie des faisceaux, Hermann,
1958. Act. Scient. Ind. nᵒ 1252, (Paris)\sisi{ (cité [TF])}{}.

\sisi{\bibitem[1]{SGA4-17:Gr}}{\bibitem{SGA4-17:Gr}} A. Grothendieck. TDTE IV: les schémas de
Hilbert. Séminaire Bourbaki 221. Mai 1961. Repris dans FGA.

\sisi{\bibitem[RD]{SGA4-17:RD}}{\bibitem{SGA4-17:RD}} R. Hartshorne. Residues and Duality. Lecture Notes
in Mathematics nᵒ 20. Springer 1966.

\sisi{\bibitem[1]{SGA4-17:La}}{\bibitem{SGA4-17:La}} D. Lazard. Thèse --- Paris 1968.

\sisi{}{\bibitem{SGA4-17:Mu} D. Mumford, \emph{Geometric Invariant Theory}, Ergebnisse der Mathematik und ihrer Grenzgebiete 34, Springer-Verlag, 1965.}

\sisi{}{\bibitem{SGA4-17:Na} M. Nagata, « A generalization of the imbedding problem of an abstract variety in a complete variety », \emph{J. Math. Kyoto Univ.} \textbf{3} (1963), 89--102}

\sisi{\bibitem[1]{SGA4-17:Ro}}{\bibitem{SGA4-17:Ro}} N. Roby. Lois polynômes et lois formelles en théorie
des modules. Ann. E.N.S. \textbf{80} \sisi{1963, p.}{(1963), }213--\sisi{2}{3}48.

\sisi{\bibitem[CD]{SGA4-17:CD}}{\bibitem{SGA4-17:CD}} J.L. Verdier. Catégories dérivées (état 0),
multicopié par l'I.H.E.S.

\sisi{\bibitem[2]{SGA4-17:V2}}{\bibitem{SGA4-17:V2}} J.L. Verdier. Catégories dérivées (à paraître dans
Grundlehren der Math. Wiss, Springer).

\sisi{\bibitem[3]{SGA4-17:V3}}{\bibitem{SGA4-17:V3}} J.L. Verdier. Dualité dans la cohomologie des
espaces localement compacts. Séminaire Bourbaki 300. Nov 1965.
\end{thebibliography}


%0.2 no ref. 1957, 1329, 5575, 5594, 5821
%2.1.2.2 no ref. 1386
%2.1.2.3 no ref  1412
%2.1.12.1 no ref 1641
%3.2.4.3 no ref  2680
%5.1.8.2 no ref 3925
%XII 4.1 no ref 4690, 6830
%6.3.4.3 no ref 7561, 7876, 8060,8409
%6.3.11 no ref 7736, 7833
%6.2.17 no ref 8042
%6.3.17.1 no ref 8078
%6.3.2.3 no ref 8271
%4.3.0 no ref 8569
%3.7 no ref 9229


\ifx\danslelivre\undefined
\end{document}
\fi
